\documentclass[12pt,oneside]{book}

% ============================================================
% DERECHOS REALES — MANUAL DE ESTUDIO INTEGRADO
% Documento autónomo: no depende de plantillas ni archivos externos.
% ============================================================

\newcommand{\universidad}{Universidad de Buenos Aires (UBA)}
\newcommand{\facultad}{Facultad de Derecho}
\newcommand{\carrera}{Derecho Civil}
\newcommand{\materia}{Derechos Reales}
\newcommand{\titulo}{Derechos Reales}
\newcommand{\subtitulo}{Manual de estudio integrado}
\newcommand{\catedra}{Cátedra Prof. Farina}
\newcommand{\autor}{Isaac Edgar Camacho Ocampo}
\newcommand{\anio}{2026}

\usepackage[utf8]{inputenc}
\usepackage[T1]{fontenc}
\usepackage[spanish,es-nodecimaldot]{babel}
\usepackage[
top=2.5cm,
bottom=2.5cm,
left=3cm,
right=2.5cm,
headheight=15pt
]{geometry}
\usepackage{amsmath}
\usepackage{amssymb}
\usepackage{mathtools}
\usepackage{graphicx}
\usepackage{float}
\usepackage{caption}
\usepackage{subcaption}
\usepackage{booktabs}
\usepackage{array}
\usepackage{longtable}
\usepackage{tabularx}
\usepackage{enumitem}
\usepackage{xcolor}
\usepackage[most]{tcolorbox}
\usepackage{fancyhdr}
\usepackage{ragged2e}
\usepackage[
colorlinks=true,
linkcolor=blue!50!black,
citecolor=green!40!black,
urlcolor=blue!60!black,
pdftitle={Derechos Reales. Manual de estudio integrado},
pdfauthor={Isaac Edgar Camacho Ocampo},
pdfsubject={Derechos Reales},
bookmarks=true
]{hyperref}

\setcounter{tocdepth}{2}
\setcounter{secnumdepth}{2}
\setlength{\parindent}{0pt}
\setlength{\parskip}{0.5em}
\setlist[itemize]{topsep=0.3em, itemsep=0.2em, parsep=0pt}
\setlist[enumerate]{topsep=0.3em, itemsep=0.2em, parsep=0pt}
\clubpenalty=10000
\widowpenalty=10000
\raggedbottom
\renewcommand{\arraystretch}{1.25}
\captionsetup{font=small, labelfont=bf}
\newcolumntype{L}{>{\raggedright\arraybackslash}X}

\newcommand{\abs}[1]{\left|#1\right|}
\newcommand{\norm}[1]{\left\lVert#1\right\rVert}
\newcommand{\vect}[1]{\mathbf{#1}}
\newcommand{\deriv}[2]{\frac{d#1}{d#2}}
\newcommand{\pderiv}[2]{\frac{\partial#1}{\partial#2}}

% ------------------------------------------------------------
% CAJAS ACADÉMICAS (uso semántico)
% ------------------------------------------------------------
\newtcolorbox{definicion}[1]{
    enhanced, breakable,
    title={Definición: #1}, fonttitle=\bfseries\small,
    colback=blue!3, colframe=blue!50!black,
    boxrule=0.7pt, arc=1.5mm, left=6pt, right=6pt, top=4pt, bottom=4pt
}
\newtcolorbox{conceptoclave}[1]{
    enhanced, breakable,
    title={Concepto clave: #1}, fonttitle=\bfseries\small,
    colback=teal!4, colframe=teal!50!black,
    boxrule=0.7pt, arc=1.5mm, left=6pt, right=6pt, top=4pt, bottom=4pt
}
\newtcolorbox{importante}{
    enhanced, breakable,
    title={Importante}, fonttitle=\bfseries\small,
    colback=yellow!8, colframe=orange!70!black,
    boxrule=0.7pt, arc=1.5mm, left=6pt, right=6pt, top=4pt, bottom=4pt
}
\newtcolorbox{ejemplo}[1]{
    enhanced, breakable,
    title={Ejemplo: #1}, fonttitle=\bfseries\small,
    colback=green!4, colframe=green!40!black,
    boxrule=0.7pt, arc=1.5mm, left=6pt, right=6pt, top=4pt, bottom=4pt
}
\newtcolorbox{regla}[1]{
    enhanced, breakable,
    title={Regla / texto normativo: #1}, fonttitle=\bfseries\small,
    colback=purple!3, colframe=purple!50!black,
    boxrule=0.7pt, arc=1.5mm, left=6pt, right=6pt, top=4pt, bottom=4pt
}
\newtcolorbox{ejercicio}[1]{
    enhanced, breakable,
    title={Ejercicio: #1}, fonttitle=\bfseries\small,
    colback=gray!5, colframe=gray!60!black,
    boxrule=0.7pt, arc=1.5mm, left=6pt, right=6pt, top=4pt, bottom=4pt
}
\newtcolorbox{nota}{
    enhanced, breakable,
    title={Nota}, fonttitle=\bfseries\small,
    colback=white, colframe=gray!55,
    boxrule=0.5pt, arc=1.5mm, left=6pt, right=6pt, top=4pt, bottom=4pt
}
\newtcolorbox{advertencia}{
    enhanced, breakable,
    title={Advertencia}, fonttitle=\bfseries\small,
    colback=red!4, colframe=red!55!black,
    boxrule=0.7pt, arc=1.5mm, left=6pt, right=6pt, top=4pt, bottom=4pt
}

% ------------------------------------------------------------
% CUADRO CORNELL (tabla real; admite corte de página)
%   \cornell{título}{filas \cq{pregunta}{respuesta}...}{resumen}
% ------------------------------------------------------------
\newcommand{\cq}[2]{\textbf{#1} & #2 \\[0.45em]}
\newcommand{\cornell}[3]{%
\par\bigskip\noindent{\large\bfseries Cuadro Cornell: #1}\par\nopagebreak\smallskip
\begingroup
\small
\renewcommand{\arraystretch}{1.15}
\begin{longtable}{@{}>{\raggedright\arraybackslash}p{0.26\textwidth}>{\raggedright\arraybackslash}p{0.66\textwidth}@{}}
\toprule
\textbf{Preguntas / palabras clave} & \textbf{Notas / respuestas} \\
\midrule
\endfirsthead
\toprule
\textbf{Preguntas / palabras clave} & \textbf{Notas / respuestas} \\
\midrule
\endhead
#2
\midrule
\multicolumn{2}{@{}p{\dimexpr0.92\textwidth+2\tabcolsep\relax}@{}}{\vspace{0.2em}\textbf{Resumen.} #3\vspace{0.2em}} \\
\bottomrule
\end{longtable}
\endgroup
}

% ------------------------------------------------------------
% ENCABEZADOS Y PIES
% ------------------------------------------------------------
\pagestyle{fancy}
\fancyhf{}
\fancyhead[L]{\small\nouppercase{\leftmark}}
\fancyhead[R]{\small\materia}
\fancyfoot[C]{\thepage}
\renewcommand{\headrulewidth}{0.4pt}
\renewcommand{\footrulewidth}{0pt}
\fancypagestyle{plain}{
    \fancyhf{}
    \fancyfoot[C]{\thepage}
    \renewcommand{\headrulewidth}{0pt}
}

\begin{document}
\frontmatter

\begin{titlepage}
\thispagestyle{empty}
\begin{center}
\vspace*{1cm}
{\large \textsc{\universidad}}\\[0.2cm]
{\large \textsc{\facultad}}\\[0.2cm]
{\large \carrera}

\vspace{3cm}
{\Huge \textbf{\materia}}

\vspace{1cm}
{\LARGE \subtitulo}

\vspace{0.4cm}
{\normalsize \catedra}

\vspace{1.2cm}
\rule{0.70\textwidth}{0.5pt}

\vspace{0.4cm}
{\large Teoría general, posesión, dominio, condominio y propiedad horizontal,\\ derechos sobre cosa ajena, prescripción, acciones y publicidad registral}

\vspace{0.4cm}
\rule{0.70\textwidth}{0.5pt}

\vfill
{\large \autor}

\vspace{0.3cm}
{\large \anio}
\end{center}
\vfill
\begin{flushleft}
\small\textit{Este es un modesto aporte a los estudiantes de «Derechos Reales»; no pretende sustituir el material oficial de la materia.}
\end{flushleft}
\end{titlepage}

\tableofcontents

\mainmatter

\chapter*{Introducción general}
\addcontentsline{toc}{chapter}{Introducción general}
\markboth{Introducción general}{Introducción general}

Este manual reúne, en un único desarrollo, el estudio de los derechos reales en el Código Civil y Comercial de la Nación (CCyC o CCCN) y su contraste con el régimen anterior. Los derechos reales organizan el modo en que una sociedad reconoce y protege la relación de poder de las personas con las cosas: quién puede usarlas, gozarlas y disponer de ellas, frente a quiénes puede hacer valer ese poder y con qué límites.

El recorrido va de lo general a lo particular. Primero se fijan los fundamentos (qué distingue a un derecho real de un derecho personal, cuáles son sus principios y cuáles son los catorce derechos reales que reconoce la ley); luego se estudia cómo se adquieren, se hacen conocer y se extinguen, y cuáles son las relaciones de poder de hecho que los acompañan. Con esas bases se recorren los derechos reales sobre cosa propia (dominio, condominio, propiedad horizontal y las nuevas formas de propiedad) y, a continuación, los derechos reales sobre cosa ajena, de disfrute y de garantía. El final del libro estudia cómo el tiempo permite adquirir derechos, cómo se defienden la posesión y los derechos reales, cómo se hacen conocer los derechos inmobiliarios y cómo se protege la vivienda frente a los acreedores.

\begin{center}
\begin{tabularx}{\textwidth}{@{}>{\raggedright\arraybackslash}p{0.30\textwidth}>{\raggedright\arraybackslash}X@{}}
\toprule
\textbf{Parte} & \textbf{Pregunta que organiza la lectura} \\
\midrule
\textbf{I. Los derechos reales en el sistema jurídico} & ¿Qué es un derecho real, cómo se distingue de un derecho personal y cuáles reconoce la ley? \\
\textbf{II. Adquisición, publicidad, transmisión y extinción} & ¿Cómo se adquiere, se hace conocer, se transmite y se extingue un derecho real? \\
\textbf{III. Relaciones de poder: posesión y tenencia} & ¿Qué relaciones de hecho puede tener una persona con una cosa y qué efectos producen? \\
\textbf{IV. El dominio} & ¿Qué facultades reúne el derecho real más extenso, cómo se adquiere y qué límites tiene? \\
\textbf{V. Condominio y medianería} & ¿Qué ocurre cuando varias personas son titulares de una misma cosa o cuando dos inmuebles colindantes comparten un muro? \\
\textbf{VI. Propiedad horizontal} & ¿Cómo coexisten el dominio exclusivo sobre una unidad y el condominio sobre lo común? \\
\textbf{VII. Nuevas formas de propiedad} & ¿Cómo se organizan los conjuntos inmobiliarios, el tiempo compartido y los cementerios privados? \\
\textbf{VIII. Derechos reales de disfrute sobre cosa ajena} & ¿Cómo se desmembra el dominio para que otro use y goce de una cosa, y qué conserva el propietario? \\
\textbf{IX. Derechos reales de garantía} & ¿Cómo se asegura el cumplimiento de una obligación mediante un privilegio sobre una cosa determinada? \\
\textbf{X. Tiempo, posesión y defensa de los derechos reales} & ¿Cómo el tiempo transforma una relación de hecho en derecho real, y cómo se defienden la posesión y el derecho? \\
\textbf{XI. Publicidad registral y protección de la vivienda} & ¿Cómo se conocen los derechos reales inmobiliarios, y cómo se protege la vivienda de los acreedores? \\
\bottomrule
\end{tabularx}
\end{center}

\section*{Relevancia profesional}

Los contenidos de este manual se presentan de manera constante en la práctica del derecho. Quien ejerce la abogacía, la actividad notarial, judicial o registral necesita estos conocimientos para, entre otras tareas:

\begin{itemize}
    \item asesorar en compraventas de inmuebles y unidades funcionales, y en emprendimientos organizados como conjuntos inmobiliarios, tiempo compartido o cementerios privados;
    \item redactar escrituras públicas de constitución de usufructo, superficie, servidumbre, hipoteca u otros derechos reales, y de afectación de la vivienda;
    \item distinguir la extinción de la cancelación de una garantía, para evitar que un cliente pierda un bien por una hipoteca ya pagada pero no cancelada;
    \item elegir la acción posesoria o real adecuada y producir la prueba que cada una exige;
    \item asesorar sobre la prescripción adquisitiva, sobre la publicidad registral y sus plazos críticos, y sobre la protección de la vivienda.
\end{itemize}

\begin{nota}
Los capítulos incorporan cuadros Cornell (algunos más de uno): en la columna izquierda se formulan preguntas o palabras clave cuya respuesta se encuentra en el texto del capítulo, y la última fila contiene un resumen. Sirven para recuperar activamente lo estudiado.
\end{nota}

\begin{nota}
Siempre que las notas originales citan artículos de forma inconsistente entre sí, el texto conserva ambas referencias y lo señala mediante un comentario \texttt{TODO} en el código fuente, sin resolverlo con información externa. Esas marcas señalan también imprecisiones o contenidos incompletos presentes en las fuentes y requieren verificación. Las abreviaturas utilizadas son CCyC o CCCN (Código Civil y Comercial de la Nación) y CN (Constitución Nacional).
\end{nota}

% ============================================================
\part{Los derechos reales en el sistema jurídico}
% ============================================================

\chapter[Derechos personales y derechos reales]{Derechos personales y derechos reales}\label{cap:dp-dr}
% ------------------------------------------------------------

El punto de partida es una distinción que organiza todo el derecho patrimonial: la que separa los derechos personales de los derechos reales. Comprenderla con precisión permite entender por qué los derechos reales reciben una regulación propia, de orden público, y por qué el resto de este manual dedica tanto espacio a sus principios comunes.

\section{La teoría clásica o del vínculo}

El Código Civil y Comercial vigente adopta la distinción entre \textbf{derechos reales} y \textbf{derechos personales}, conocida como \textbf{teoría clásica o del vínculo}, que es la teoría aplicable en la mayoría de los sistemas occidentales de derecho. Los derechos personales ---también llamados \textbf{derechos obligacionales}--- y los derechos reales ---también llamados \textbf{derechos de las cosas}--- presentan una profunda distinción y reciben una regulación específica cada uno.

La teoría los distingue según el vínculo que establecen:

\begin{itemize}
    \item \textbf{derechos personales:} vínculo entre sujetos (entre personas);
    \item \textbf{derechos reales:} vínculo directo entre una persona y una cosa.
\end{itemize}

Ambas categorías integran el \textbf{patrimonio} de una persona. Según el apunte, el Código adopta esta teoría porque organiza el sistema patrimonial occidental moderno y permite dar respuestas claras acerca de quién puede exigir qué y sobre qué objeto.

\section{El derecho personal u obligacional}

\begin{definicion}{Obligación (art.~724 CCyC)}
«La obligación es una relación jurídica en virtud de la cual el acreedor tiene el derecho de exigir del deudor una prestación destinada a satisfacer un interés lícito y ante el incumplimiento obtener forzadamente la satisfacción de dicho interés.»
\end{definicion}

\subsection{Elementos}

El derecho personal tiene tres elementos esenciales:

\begin{itemize}
    \item \textbf{sujeto activo o acreedor:} quien tiene el derecho de exigir la prestación;
    \item \textbf{sujeto pasivo o deudor:} quien está obligado a cumplirla;
    \item \textbf{prestación:} la conducta debida, que puede ser de dar, de hacer o de no hacer (omisión).
\end{itemize}

El vínculo se representa como una \textbf{ola o vínculo circular} que encierra dentro de sí al acreedor, al deudor y a la prestación. Es un vínculo \textbf{cerrado}, lo que explica su carácter relativo.

\subsection{Caracteres}

\begin{itemize}
    \item \textbf{Relativos:} el vínculo existe solo entre las partes; el acreedor solo puede exigir el cumplimiento a su propio deudor.
    \item \textbf{Autonomía de la voluntad:} los particulares pueden crear, modificar y extinguir obligaciones; la ley es supletoria.
    \item \textbf{Objeto:} una conducta humana (dar, hacer, no hacer).
    \item \textbf{Fuente:} nacen de hechos o actos jurídicos, esencialmente el contrato, y también de la obligación de no dañar (responsabilidad civil).
    \item \textbf{Extinción:} se extinguen por cumplimiento, por prescripción liberatoria u otras formas de extinción de las obligaciones. Son siempre temporales.
    \item \textbf{No se adquieren} por el mero transcurso del tiempo.
\end{itemize}

\begin{ejemplo}{obligación alimentaria, crédito y préstamo}
En la \textbf{obligación alimentaria}, el hijo (acreedor alimentario) solo puede exigir la prestación a quienes la ley identifica como obligados (padres, abuelos, etc.), y no a un tercero. Del mismo modo, el acreedor de un \textbf{crédito} (por ejemplo, el banco) solo puede reclamar el pago a su deudor.

Si se presta una suma de dinero a un amigo, nace una obligación: el prestamista es el acreedor, el amigo es el deudor y la prestación es dar, es decir, devolver el dinero. Solo puede reclamarse al amigo, y no a un tercero (su pareja o su vecino): por eso el derecho personal es \textbf{relativo}. Si el deudor no paga, el acreedor puede acudir a la justicia para obligarlo; pero cuando cobra, el vínculo se termina. Y si el acreedor no reclama durante años, el derecho se extingue por prescripción: el tiempo juega en su contra.
\end{ejemplo}

\section{El derecho real}

\begin{definicion}{Derecho real (art.~1882 CCyC)}
«El derecho real es el poder jurídico de estructura legal que se ejerce directamente sobre su objeto en forma autónoma y que le atribuye a su titular las facultades de persecución y preferencia y las demás previstas en este código.»
\end{definicion}

El derecho real implica un vínculo \textbf{directo e inmediato} entre el \textbf{titular} y la \textbf{cosa}; estos son sus dos elementos. Ese vínculo es tan intenso que permite al titular:

\begin{itemize}
    \item \textbf{disponer de la cosa u obtener algún beneficio de ella} en forma directa, sin intermediarios;
    \item \textbf{perseguir la cosa} en poder de quien se encuentre;
    \item \textbf{ser preferido} a cualquier otro que pudiera tener sobre ella algún derecho real o personal de fecha posterior.
\end{itemize}

\subsection{Caracteres}

\begin{itemize}
    \item \textbf{Absolutos o \textit{erga omnes}:} son oponibles a toda la comunidad, que debe respetar el derecho del titular. El sujeto pasivo es, por tanto, toda la comunidad, con una obligación general de no perturbar al titular.
    \item \textbf{Estructura legal y orden público:} son creados y regulados exclusivamente por la ley (\textit{numerus clausus}, art.~1884). Los particulares no pueden crear ni modificar derechos reales, y todo acto contrario es nulo.
    \item \textbf{Objeto:} la cosa, es decir, un bien material. Excepcionalmente, bienes inmateriales taxativamente previstos.
    \item \textbf{Nacimiento:} título suficiente y modo (tradición), ley o prescripción adquisitiva.
    \item \textbf{Duración:} pueden ser perpetuos (como el dominio) o temporales.
    \item \textbf{Facultades:} persecución (\textit{ius persequendi}) y preferencia.
\end{itemize}

\begin{ejemplo}{la compra de un departamento}
Quien compra un departamento, firma la escritura y recibe la entrega (tradición) adquiere el dominio, que es un derecho real. Su relación no es con una persona específica, sino directamente con la cosa: es dueño y no necesita permiso de nadie para usarla, disfrutarla o venderla. Todos deben respetar esa propiedad, no solo el vendedor; por eso el derecho es \textit{erga omnes}. Si la cosa fuera robada u ocupada ilegalmente y transmitida a un tercero, el titular puede perseguirla: el derecho real «sigue» a la cosa sin importar en qué manos se encuentre (\textit{ius persequendi}).
\end{ejemplo}

\section{Cuadro comparativo}

\begin{table}[H]
\centering
\small
\begin{tabularx}{\textwidth}{>{\raggedright\arraybackslash}p{0.19\textwidth}>{\raggedright\arraybackslash}X>{\raggedright\arraybackslash}X}
\toprule
\textbf{Aspecto} & \textbf{Derecho personal} & \textbf{Derecho real} \\
\midrule
Vínculo & Entre sujetos (acreedor y deudor) & Entre titular y cosa \\
Oponibilidad & Relativo: solo entre las partes & Absoluto: \textit{erga omnes} \\
Sujeto pasivo & El deudor determinado & Toda la comunidad \\
Objeto & Conducta humana (dar, hacer, no hacer) & La cosa (bien material) \\
Creación y regulación & Autonomía de la voluntad; la ley es supletoria & Exclusivamente la ley; normas de orden público \\
Nacimiento & Hechos o actos jurídicos: contrato, daño & Título y modo (tradición), ley, prescripción adquisitiva \\
Duración & Siempre temporal & Perpetuo o temporal \\
Prescripción & Liberatoria: extingue el derecho & Adquisitiva: permite adquirirlo \\
Facultades de persecución y preferencia & No hay persecución & Sí \\
Ejemplo & Contrato de compraventa (obligación de entregar) & Dominio sobre un inmueble \\
\bottomrule
\end{tabularx}
\caption{Derechos personales y derechos reales}
\end{table}

\begin{conceptoclave}{la diferencia esencial}
El derecho personal es un vínculo entre personas: relativo y temporal. El derecho real es un vínculo entre una persona y una cosa: absoluto, \textit{erga omnes}, regulado solo por la ley y defendido mediante la persecución de la cosa donde quiera que esté.
\end{conceptoclave}

\cornell{derechos personales y derechos reales}{
\cq{¿Qué teoría adopta el CCyC y qué criterio usa?}{La teoría clásica o del vínculo: los derechos personales vinculan sujetos; los reales vinculan a una persona con una cosa. Ambos integran el patrimonio.}
\cq{Art.~724: ¿qué es una obligación y cuáles son sus elementos?}{Relación jurídica por la que el acreedor puede exigir al deudor una prestación que satisface un interés lícito y, ante el incumplimiento, obtener forzadamente su satisfacción. Elementos: acreedor, deudor y prestación (dar, hacer, no hacer).}
\cq{¿Por qué el derecho personal es relativo?}{Porque el vínculo es cerrado: existe solo entre las partes y el acreedor solo puede exigir a su deudor.}
\cq{Art.~1882: ¿qué es un derecho real y qué facultades confiere?}{Poder jurídico de estructura legal ejercido directamente sobre su objeto, en forma autónoma. Confiere persecución y preferencia, además de las previstas en el Código.}
\cq{¿Por qué el derecho real es absoluto?}{Porque es oponible \textit{erga omnes}: toda la comunidad debe respetar al titular.}
\cq{Compare la prescripción en ambas esferas.}{En los personales es liberatoria (extingue el derecho); en los reales, adquisitiva (permite adquirirlo). Los personales son siempre temporales; los reales pueden ser perpetuos.}
}{El Código distingue los derechos personales (vínculo entre acreedor y deudor, relativo, temporal, regido por la autonomía de la voluntad) de los reales (vínculo directo con la cosa, absoluto, de estructura legal y orden público, con persecución y preferencia). Ambos integran el patrimonio.}

\chapter[Los derechos reales en el Código Civil y Comercial]{Los derechos reales en el Código Civil y Comercial}\label{cap:ccyc}
% ------------------------------------------------------------

Una vez delimitado el concepto, corresponde ubicar a los derechos reales dentro del Código y precisar sobre qué objeto recaen. Para ello son centrales los artículos 15, 16 y 17.

\section{Estructura metodológica del Código}

El CCyC tiene una \textbf{metodología novedosa} respecto del Código de Vélez Sársfield. Su estructura es la siguiente.

\begin{table}[H]
\centering
\small
\begin{tabularx}{\textwidth}{>{\raggedright\arraybackslash}p{0.22\textwidth}>{\raggedright\arraybackslash}p{0.27\textwidth}>{\raggedright\arraybackslash}X}
\toprule
\textbf{Ubicación} & \textbf{Denominación} & \textbf{Contenido relevante} \\
\midrule
Título Preliminar & Principios rectores & Buena fe, abuso del derecho, abuso de posición dominante, preeminencia del orden público, tratados de derechos humanos como fuente y guía de interpretación; incluye el Capítulo 4, sobre bienes y cosas \\
Libro Primero & Parte General & Persona, capacidad \\
Libro Segundo & Relaciones de Familia & Derecho de familia \\
Libro Tercero & Derechos Personales & Obligaciones, contratos \\
Libro Cuarto & Derechos Reales & Teoría general y derechos reales en particular (arts.~1882 a 2276) \\
Libro Quinto & Transmisión de derechos por causa de muerte & Sucesiones \\
Libro Sexto & Disposiciones comunes & A derechos reales y personales \\
\bottomrule
\end{tabularx}
\caption{Estructura del Código Civil y Comercial}
\end{table}
% TODO: verificar consistencia entre fuentes originales (una versión denomina «Título Sexto» y otra «Libro VI» a las disposiciones comunes).

\begin{importante}
El estudio de los derechos reales \textbf{excede el Libro Cuarto}: deben integrarse el Título Preliminar, las disposiciones comunes del Libro Sexto y el resto del ordenamiento.
\end{importante}

\section{Patrimonio, bienes y cosas: artículos 15, 16 y 17}

\begin{regla}{art.~15 CCyC, titularidad}
«Las personas son titulares de los derechos individuales sobre los bienes que integran su patrimonio conforme con lo que se establece en este Código.»
\end{regla}
% TODO: verificar referencia presente en las fuentes originales (una versión transcribe el art.~15 sin la expresión final «conforme con lo que se establece en este Código»).

El patrimonio está integrado por bienes, y estos comprenden tanto los derechos personales como los derechos reales y las cosas materiales. Ambas categorías conviven en el patrimonio. El artículo 15 expresa, además, el sistema de propiedad privada (véase el capítulo~\ref{cap:legalidad}).

\begin{regla}{art.~16 CCyC, bienes y cosas}
«Los derechos pueden recaer sobre bienes susceptibles de valor económico. Cuando son materiales se llaman cosas. Las disposiciones sobre las cosas se aplican a la energía y a las fuerzas naturales susceptibles de ser puestas al servicio del hombre.»
\end{regla}

Esta norma es esencial para determinar el \textbf{objeto} de los derechos reales. Los bienes son los derechos susceptibles de valor económico; cuando son materiales (tangibles) se llaman cosas. Toda cosa es un bien, pero no todo bien es una cosa: por ejemplo, un crédito es un bien pero no una cosa. La energía y las fuerzas naturales se asimilan a las cosas.

\begin{regla}{art.~17 CCyC, derechos sobre el cuerpo humano}
Los derechos sobre el cuerpo humano o sus partes no tienen valor comercial, sino afectivo, científico, terapéutico, humanitario o social. Su disposición está sujeta a lo dispuesto en las leyes especiales (por ejemplo, órganos y genes).
\end{regla}

Consecuencia directa: el cuerpo humano y sus partes \textbf{no son objeto de derechos reales}, porque el artículo 16 exige que el bien tenga valor económico para ser considerado cosa.

\section{Los principios comunes: artículos 1882 a 1891}

Los artículos 1882 a 1891 regulan los principios comunes a todos los derechos reales. Esta parte general es la base teórica de toda la materia: muchos de sus principios no se repiten al estudiar cada derecho real en particular. Por eso, cada vez que se aborda un derecho real específico conviene «exponer en acto» lo aprendido sobre las disposiciones comunes, aplicándolas al derecho de que se trate (así se hace, por ejemplo, con el dominio en el capítulo~\ref{cap:dominio}).

\cornell{los derechos reales en el Código}{
\cq{¿Dónde se regulan los derechos reales y basta con el Libro Cuarto?}{En el Libro Cuarto (arts.~1882 a 2276), pero su estudio lo excede: se integran el Título Preliminar, las disposiciones comunes y el resto del ordenamiento.}
\cq{¿Qué contiene el Título Preliminar?}{Los principios rectores de todo el Código: buena fe, abuso del derecho, abuso de posición dominante, orden público y tratados de derechos humanos. Incluye el capítulo sobre bienes y cosas.}
\cq{¿Qué integra el patrimonio (art.~15)?}{Los bienes: derechos personales, derechos reales y cosas materiales. Las personas son titulares de derechos individuales sobre ellos.}
\cq{¿Cuándo un bien es una cosa (art.~16)? ¿Todo bien es cosa?}{Cuando es material. Toda cosa es un bien, pero no todo bien es cosa (el crédito es un bien, no una cosa). La energía y las fuerzas naturales se asimilan a las cosas.}
\cq{¿Por qué el cuerpo humano no es objeto de derechos reales (art.~17)?}{Porque sus partes tienen valor afectivo, científico, terapéutico, humanitario o social, pero no valor comercial ni económico, y se rigen por leyes especiales.}
}{Los derechos reales se regulan en el Libro Cuarto, cuyos arts.~1882 a 1891 contienen los principios comunes. Los arts.~15, 16 y 17 fijan la titularidad sobre los bienes del patrimonio, el concepto de cosa como bien material con valor económico y la exclusión del cuerpo humano como objeto de derechos reales.}

\chapter[Legalidad, orden público y fundamento constitucional]{Legalidad, orden público y fundamento constitucional}\label{cap:legalidad}
% ------------------------------------------------------------

Los caracteres de los derechos reales no son una enumeración arbitraria: responden a una idea común. Como el derecho real debe ser respetado por toda la comunidad, solo la ley puede crearlo y regularlo, y su ejercicio queda sometido a los intereses de la sociedad. Este capítulo desarrolla esa idea y los principios que la expresan.

\section{\textit{Numerus clausus} y ley federal}

\begin{regla}{art.~1884 CCyC, estructura de los derechos reales}
«La regulación de los derechos reales en cuanto a sus elementos, contenido, adquisición, constitución, modificación, transmisión, duración y extinción está establecida solo por la ley. Es nula la configuración de un derecho real no previsto en la ley o la modificación de su estructura.»
\end{regla}

Este principio de \textit{numerus clausus} (número cerrado) significa que los particulares \textbf{no pueden crear nuevos derechos reales}: solo la ley los crea y regula. Los derechos reales son \textbf{catorce y solo catorce}, y hoy se encuentran todos codificados. La sanción por apartarse de la regulación legal es la \textbf{nulidad} del acto. El derecho real es, además, un señorío de la voluntad que se ejerce en forma autónoma e independiente de otra voluntad.

La competencia para crearlos es exclusivamente federal: se trata de \textbf{ley federal}, porque en materia de derechos reales solo la Nación es competente para legislar.

\begin{ejemplo}{la lista cerrada}
Los derechos reales son como una lista cerrada escrita por el Congreso, o como los planetas del sistema solar: solo existen los que la ley nombra. Si dos personas firman un contrato en el que dicen «creamos un nuevo derecho real que se llama semi-dominio» (o «semi-usufructo», o «super-dominio compartido especial»), ese acuerdo es nulo en cuanto al derecho real: no existe. En cambio, en los derechos personales las partes pueden crear contratos atípicos: allí hay libertad. En los derechos reales, no, porque afectan a toda la sociedad.
\end{ejemplo}

\section{Normas de orden público}

Los particulares deben ceñirse estrictamente a la regulación legal de cada derecho real, porque las normas que los rigen son de \textbf{orden público}: expresan la voluntad de la Nación y no pueden ser modificadas por los particulares.

\begin{regla}{art.~12 CCyC}
Las convenciones particulares no pueden dejar sin efecto las leyes en cuya observancia está interesado el orden público.
\end{regla}

La razón de fondo es que el derecho real de cualquier país responde a un \textbf{esquema de sociedad único para toda la Nación}. Tiene que ver con tres cuestiones: el modo en que una nación organiza su sistema de propiedad y su sistema de defensa de las relaciones de poder respecto de las cosas; el modo en que define el alcance de las facultades que está dispuesta a reconocer a cada titular; y las limitaciones que está dispuesta a aceptar en función de los intereses de la comunidad.

\section{Constitucionalización e internacionalización del derecho privado}

El CCyC evidencia un cambio de paradigma respecto de la legislación anterior, esencialmente porque incorpora los tratados internacionales de derechos humanos como ley vigente. Por eso puede hablarse de \textbf{constitucionalización} e \textbf{internacionalización del derecho privado}. Los artículos 1 y 2 son fundamentos del Código y, a la vez, una regla de interpretación: todo el sistema se mira bajo esa perspectiva.

En ese marco conviven en tensión varias normas.

\begin{table}[H]
\centering
\small
\begin{tabularx}{\textwidth}{>{\raggedright\arraybackslash}p{0.15\textwidth}>{\raggedright\arraybackslash}X>{\raggedright\arraybackslash}p{0.27\textwidth}}
\toprule
\textbf{Artículo} & \textbf{Contenido} & \textbf{Efecto} \\
\midrule
Art.~15 & Las personas son titulares de los derechos individuales sobre los bienes que integran su patrimonio. & Consagra el sistema de propiedad privada. \\
Art.~14 & La ley no ampara el ejercicio abusivo de los derechos individuales cuando pueda afectar al ambiente y a los derechos de incidencia colectiva. & Limita los derechos individuales frente al interés colectivo. \\
Art.~240 & Establece un límite a los derechos individuales cuando resulte necesario compatibilizarlos con los derechos de incidencia colectiva. & Preeminencia del interés colectivo sobre el individual. \\
\bottomrule
\end{tabularx}
\caption{Propiedad privada y sus límites}
\end{table}

El derecho al ambiente y el derecho a respirar son ejemplos de derechos de incidencia colectiva que pueden entrar en tensión con el ejercicio de los derechos reales. En el sistema jurídico existe una clara \textbf{preeminencia del interés colectivo por sobre el interés individual}; por eso, en ocasiones (por ejemplo, en una situación de contingencia), resulta necesario modelar el alcance del ejercicio de un derecho real en función de los intereses de la comunidad.

El derecho real se vincula así con grandes temas del derecho: el ejercicio del derecho de propiedad, los límites al dominio, el reconocimiento de la propiedad como derecho humano y la función social. Son temas susceptibles de enfrentar ideológicamente a distintos sectores de la doctrina nacional e internacional (se desarrollan en el capítulo~\ref{cap:limites}).

\section{\textit{Nemo plus iuris} y convalidación}

\begin{definicion}{\textit{nemo plus iuris} (art.~399 CCyC)}
Nadie puede transmitir a otro un derecho mejor o más extenso que el que tiene, y nadie puede adquirir sobre un objeto un derecho mejor o más extenso que el que tenía aquel de quien lo adquiere.
\end{definicion}

Es un principio de origen romano que limita la transmisión de derechos y es el límite infranqueable de la transmisión. Se aplica a los derechos reales y a los personales y, si bien admite excepciones previstas en la ley, constituye la regla general. Quien transmite un derecho que no tiene realiza un acto que nace nulo.

\begin{regla}{art.~1885 CCyC, convalidación}
«Si alguien constituye o transmite un derecho real que no tiene, y lo adquiere posteriormente, la constitución o transmisión queda convalidada.»
\end{regla}
% TODO: verificar consistencia entre fuentes originales (el art.~1885 se transcribe con redacciones distintas y, en el desarrollo del dominio, se lo describe con un alcance diferente: ver capítulo sobre adquisición del dominio).

La convalidación es la \textbf{excepción} al \textit{nemo plus iuris}: si después de la transmisión nula el transmitente adquiere el derecho, el acto queda automáticamente saneado, como si hubiera sido válido desde el principio. El apunte señala que existe para proteger a quien adquirió de buena fe y dar seguridad jurídica al sistema.

\begin{ejemplo}{la venta de la cosa ajena}
Quien tiene una cosa en su poder pero no es propietario, y la transmite, no hace adquirir la propiedad a quien la recibe: no puede transmitir un derecho mejor que el que tenía. El acto nace nulo. Pero si después el transmitente adquiere la propiedad (por ejemplo, porque el dueño se la vende o la hereda), el acto queda convalidado retroactivamente.
\end{ejemplo}

\section{Persecución y preferencia}

Son las dos facultades características del derecho real.

\begin{definicion}{\textit{ius persequendi} (persecución)}
El titular puede \textbf{perseguir la cosa} en poder de quien se encuentre, sin importar si cambió de titular: el derecho la sigue. Se la representa como una flecha que persigue la cosa sin importar en qué manos esté.
\end{definicion}

\begin{definicion}{\textit{ius preferendi} (preferencia)}
El titular es \textbf{preferido} a cualquier otro que pudiera tener sobre la cosa algún derecho real o personal de fecha posterior.
\end{definicion}

Es una diferencia clave con los derechos personales, donde no hay persecución. Ambas facultades derivan de la relación directa entre titular y cosa y de su carácter absoluto.

\cornell{legalidad y orden público}{
\cq{¿Qué consagra el art.~1884 y cuál es la sanción?}{El \textit{numerus clausus}: solo la ley crea y regula los derechos reales (elementos, contenido, adquisición, constitución, modificación, transmisión, duración y extinción). La sanción es la nulidad.}
\cq{¿Por qué las normas de derechos reales son de orden público y de ley federal?}{Porque expresan un esquema de sociedad único para toda la Nación (sistema de propiedad, defensa de las relaciones de poder y límites aceptados por la comunidad); solo la Nación legisla. Art.~12: las convenciones no pueden dejar sin efecto leyes de orden público.}
\cq{¿Qué articulan los arts.~15, 14 y 240?}{La propiedad privada (art.~15) y sus límites frente al ambiente y a los derechos de incidencia colectiva (arts.~14 y 240), con preeminencia del interés colectivo.}
\cq{Relacione \textit{nemo plus iuris} y convalidación.}{Regla general: nadie transmite un derecho mejor que el que tiene (acto nulo). Excepción (art.~1885): si el transmitente adquiere luego el derecho, el acto queda convalidado.}
\cq{¿Qué son la persecución y la preferencia?}{Persecución: seguir la cosa en poder de quien se encuentre. Preferencia: prevalecer sobre derechos reales o personales de fecha posterior. Son propias del derecho real.}
}{Solo la ley federal crea y regula los derechos reales (\textit{numerus clausus}), bajo normas de orden público y con sanción de nulidad. Ese carácter responde a un esquema de sociedad único y a la preeminencia del interés colectivo. El \textit{nemo plus iuris} limita la transmisión; la convalidación la sana si el transmitente adquiere después el derecho. La persecución y la preferencia son las facultades características del derecho real.}

\chapter[El catálogo y la clasificación de los derechos reales]{El catálogo y la clasificación de los derechos reales}\label{cap:catalogo}
% ------------------------------------------------------------

Si solo la ley crea derechos reales, cabe conocer cuáles son. Este capítulo presenta la lista cerrada del artículo 1887 y las clasificaciones del artículo 1888, que serán la guía para ordenar los derechos que se estudian en las partes siguientes.

\section{Los catorce derechos reales (art.~1887)}

El artículo 1887 enumera taxativamente los derechos reales. Solo existen estos y no hay más.

\begin{table}[H]
\centering
\small
\begin{tabularx}{\textwidth}{>{\centering\arraybackslash}p{0.07\textwidth}>{\raggedright\arraybackslash}X>{\raggedright\arraybackslash}X}
\toprule
\textbf{N.º} & \textbf{Derecho real} & \textbf{Recae sobre cosa} \\
\midrule
1 & Dominio & Propia \\
2 & Condominio & Propia (varios titulares) \\
3 & Propiedad horizontal & Propia \\
4 & Conjuntos inmobiliarios & Propia \\
5 & Tiempo compartido & Propia \\
6 & Cementerio privado & Propia \\
7 & Superficie & Propia (con propiedad superficiaria) o ajena (sin propiedad superficiaria) \\
8 & Usufructo & Ajena \\
9 & Uso & Ajena \\
10 & Habitación & Ajena \\
11 & Servidumbre & Ajena \\
12 & Hipoteca & Ajena (garantía) \\
13 & Anticresis & Ajena (garantía) \\
14 & Prenda & Ajena (garantía) \\
\bottomrule
\end{tabularx}
\caption{Los derechos reales del art.~1887}
\end{table}
% TODO: verificar consistencia entre fuentes originales (una versión numera la superficie dos veces, según exista o no propiedad superficiaria, lo que daría 15 filas para 14 derechos; aquí se la cuenta una sola vez).

\begin{nota}
Regla mnemotécnica propuesta en los apuntes para recordar la lista: \textit{Do-Con-PH-CI-TC-CP-Su-Us-U-H-S-Hi-An-P}.
\end{nota}

La enumeración se ordena en este manual así: el dominio y sus límites (Parte~IV); el condominio y la medianería (Parte~V); la propiedad horizontal (Parte~VI); los conjuntos inmobiliarios, el tiempo compartido y el cementerio privado (Parte~VII); los derechos de disfrute sobre cosa ajena, es decir, superficie, usufructo, uso, habitación y servidumbre (Parte~VIII); y los derechos reales de garantía, esto es, hipoteca, anticresis y prenda (Parte~IX). Las Partes~X y~XI estudian el tiempo como fuente de adquisición, la defensa de la posesión y de los derechos reales, y su publicidad registral.

\section{Clasificación (art.~1888)}

\subsection{Sobre cosa propia y sobre cosa ajena}

Los derechos sobre \textbf{cosa propia} recaen sobre cosas del propio titular: dominio, condominio, propiedad horizontal, conjuntos inmobiliarios, tiempo compartido, cementerio privado y superficie (con propiedad superficiaria). Los derechos sobre \textbf{cosa ajena} recaen sobre cosas de otra persona y constituyen una \textbf{carga o gravamen} real para el propietario: usufructo, uso, habitación, servidumbre, superficie (sin propiedad superficiaria), hipoteca, anticresis y prenda.

\begin{ejemplo}{dominio e hipoteca}
Una persona es titular de dominio (derecho real sobre cosa propia) y constituye una hipoteca a favor de un banco para garantizar un crédito. Ha creado un nuevo derecho real: el banco tiene ahora un derecho real sobre esa cosa (hipoteca: derecho real sobre cosa ajena, de garantía y accesorio). Coexisten el dominio del titular y la hipoteca del banco.
\end{ejemplo}

\subsection{Principales y accesorios}

\begin{itemize}
    \item \textbf{Principales:} los derechos sobre cosa propia y los derechos de disfrute sobre cosa ajena (usufructo, servidumbre, uso, habitación, superficie). Son independientes: existen por sí solos.
    \item \textbf{Accesorios:} los derechos reales de garantía (hipoteca, prenda, anticresis). Garantizan una obligación principal, dependen de un crédito al que acceden y lo siguen.
\end{itemize}

\subsection{Registrables y no registrables}

Según su adquisición, transmisión o extinción requieran o no inscripción en un registro público. Son \textbf{registrables} los inmuebles (Registro de la Propiedad Inmueble) y los automotores (Registro de Automotores); son \textbf{no registrables} las cosas muebles en general.

\subsection{Ejercicio por posesión o sin ella}

Algunos derechos reales se ejercen mediante la \textbf{posesión} (poder de hecho sobre la cosa) y otros por otros modos; las fuentes mencionan para estos últimos la \textbf{cuasiposesión}. Este tema se desarrolla en la Parte~III.
% TODO: verificar consistencia entre fuentes originales (una versión contrapone «posesión» y «cuasiposesión»; otra, «con posesión» y «sin posesión»).

\section{Casos de aplicación}

\begin{ejercicio}{creación de un derecho real inexistente}
\textbf{Enunciado.} Pablo firma un contrato con Marta por el cual se crea un «derecho real de semi-usufructo» que no existe en ninguna ley. ¿Qué sucede?

\textbf{Resolución.} El acto es nulo (art.~1884 CCyC, \textit{numerus clausus}). Solo la ley puede crear derechos reales; los particulares no pueden crear uno que no esté en el artículo 1887.
\end{ejercicio}

\begin{ejercicio}{venta de un inmueble ajeno}
\textbf{Enunciado.} Juan vende a María un inmueble que pertenece a Carlos. Después, Juan hereda ese inmueble de Carlos. ¿Qué pasa con la venta?

\textbf{Resolución.} Originalmente el acto era nulo (\textit{nemo plus iuris}: Juan no podía transmitir lo que no tenía). Al heredar el inmueble, la transmisión queda convalidada (art.~1885 CCyC) y María es ahora legítima propietaria.
\end{ejercicio}

\begin{ejercicio}{hipoteca y venta del inmueble hipotecado}
\textbf{Enunciado.} El banco tiene una hipoteca sobre la casa de Roberto, quien la vende a Sofía. ¿Puede el banco ejecutar la hipoteca frente a Sofía?

\textbf{Resolución.} Sí. La hipoteca es un derecho real que otorga la persecución: el banco puede perseguir la cosa en manos de cualquier persona que la tenga, incluida Sofía, que adquirió la casa con la carga hipotecaria.
\end{ejercicio}

El apéndice~\ref{ap:ejercicios} reúne ejercicios adicionales de evaluación sobre toda la teoría general.

\cornell{catálogo y clasificación}{
\cq{¿Cuántos derechos reales hay y cuál es su fuente?}{Catorce (art.~1887), todos creados y regulados por la ley: dominio, condominio, propiedad horizontal, conjuntos inmobiliarios, tiempo compartido, cementerio privado, superficie, usufructo, uso, habitación, servidumbre, hipoteca, anticresis y prenda.}
\cq{Distinga cosa propia y cosa ajena (art.~1888).}{Propia: el derecho recae sobre cosa del propio titular. Ajena: recae sobre cosa de otro y es carga o gravamen para su propietario. La superficie es propia si hay propiedad superficiaria y ajena si no la hay.}
\cq{¿Qué diferencia a los principales de los accesorios?}{Los principales son independientes (cosa propia y disfrute de cosa ajena). Los accesorios, de garantía, dependen de un crédito y lo siguen.}
\cq{¿Qué son los derechos registrables?}{Los que recaen sobre cosas cuya adquisición, transmisión o extinción requiere inscripción en un registro público (inmuebles y automotores).}
\cq{Caso: el banco hipotecario y el comprador del inmueble.}{El banco puede perseguir la cosa en manos del comprador, porque la hipoteca es un derecho real con persecución.}
}{El art.~1887 fija una lista cerrada de catorce derechos reales y el art.~1888 los clasifica según recaigan sobre cosa propia o ajena, sean principales o accesorios, registrables o no registrables, y se ejerzan o no por posesión.}

% ============================================================
% ============================================================


% ============================================================
\part{Adquisición, publicidad, transmisión y extinción}
% ============================================================

\chapter[Adquisición de los derechos reales: título y modo]{Adquisición de los derechos reales: título y modo}\label{cap:adquisicion}
% ------------------------------------------------------------

Conocidos los principios y el catálogo, corresponde preguntar cómo se llega a ser titular de un derecho real. El sistema reconoce distintas formas de adquisición; la más importante para los actos entre vivos exige la concurrencia de dos elementos: el título suficiente y el modo suficiente.

\section{Adquisición originaria y derivada}

\begin{table}[H]
\centering
\small
\begin{tabularx}{\textwidth}{>{\raggedright\arraybackslash}p{0.19\textwidth}>{\raggedright\arraybackslash}X>{\raggedright\arraybackslash}p{0.28\textwidth}}
\toprule
\textbf{Tipo} & \textbf{Concepto} & \textbf{Ejemplos} \\
\midrule
Originaria & El sistema reconoce la posibilidad de adquirir el derecho real en forma libre y absoluta, sin que derive de un vínculo jurídico anterior. & Cosas sin dueño (\textit{res nullius}); prescripción adquisitiva (usucapión); adquisición legal. \\
Derivada, entre vivos & El derecho real es transmitido por un enajenante a un adquirente: uno transmite y el otro adquiere. Puede ser a título oneroso o gratuito. & Compraventa, permuta, donación, dación en pago, cesión. \\
Derivada, por causa de muerte & La transmisión opera por el fallecimiento del titular: el heredero ocupa el lugar del causante y adquiere todo lo que este tenía. & Sucesión hereditaria. \\
\bottomrule
\end{tabularx}
\caption{Formas de adquisición de los derechos reales}
\end{table}

\begin{regla}{art.~1896 CCyC}
El juez no puede constituir derechos reales ni imponer su constitución.
\end{regla}

Es una expresión del principio de legalidad: el derecho real se adquiere por los modos previstos en la ley, no por decisión judicial. En toda adquisición derivada rige además el \textit{nemo plus iuris} (art.~399; véase el capítulo~\ref{cap:legalidad}), que impide adquirir un derecho mejor o más extenso que el que tenía el transmitente.

\section{La teoría del título y modo}

\begin{regla}{art.~1892 CCyC}
El derecho real se adquiere con la concurrencia del título suficiente más el modo suficiente.
\end{regla}

La adquisición derivada de derechos reales por actos entre vivos requiere necesariamente ambos elementos. Uno sin el otro no es suficiente.

\subsection{El título suficiente}

\begin{definicion}{título suficiente}
Es el acto jurídico revestido de las condiciones de fondo y de forma establecidas por la ley, que tiene por finalidad transmitir o constituir un derecho real. El acto jurídico causal de la transmisión constituye el título suficiente; su causa puede ser una compraventa, una cesión, una permuta, una dación en pago o una donación.
\end{definicion}

El título no debe identificarse en su sentido instrumental (el papel), sino en su \textbf{sentido abstracto o intangible}. Para ilustrarlo, el apunte recurre a una analogía: así como no es posible decir «te doy mi amor, dejámelo en el buzón», ante la falta de título suficiente no se resuelve afirmando que el título fue «dejado en el buzón». Entregar un instrumento que no plasma un acto jurídico causal equivale a no tener título suficiente.

Los artículos 259 y 260 aportan los conceptos de base:

\begin{itemize}
    \item \textbf{Art.~259, acto jurídico:} el acto voluntario lícito que tiene por fin inmediato la adquisición, modificación o extinción de relaciones o situaciones jurídicas.
    \item \textbf{Art.~260, acto voluntario:} el ejecutado con discernimiento, intención y libertad, que se manifiesta por un hecho exterior.
    \item \textbf{Art.~1017:} en materia de inmuebles, el acto jurídico causal debe plasmarse en \textbf{escritura pública}.
\end{itemize}

\subsection{El modo suficiente: la tradición}

El modo es el elemento material propio de los derechos reales que se ejercen por la posesión.

\begin{definicion}{tradición}
Es la entrega material de la cosa: una persona entrega la cosa a otra, que la recibe. El poder de hecho que tenía el sujeto se transmite, con la entrega, a quien la adquiere.
\end{definicion}

Quien entrega es el \textbf{tradente} y quien recibe es el \textbf{accipiens}. La tradición es \textbf{constitutiva} en todos los casos en que se adquieren derechos reales que se ejercen por la posesión: sin tradición no se adquirió el derecho real.

\begin{regla}{art.~1924 CCyC}
La tradición no puede suplirse por fórmulas escritas.
\end{regla}

Cuando una escritura pública dice que «en este acto las partes se entregan la tradición», se pretende suplantar un acto material por una fórmula escrita en el papel, lo cual no puede hacerse: la posesión no se transmite con palabras escritas. Las únicas excepciones son la \textit{traditio brevi manu} y el constituto posesorio (capítulo~\ref{cap:brevi}).

\section{Requisitos de la tradición traslativa de dominio}

No toda entrega material implica transferencia de dominio. Para que la tradición sea traslativa deben concurrir los siguientes requisitos.

\begin{table}[H]
\centering
\small
\begin{tabularx}{\textwidth}{>{\raggedright\arraybackslash}p{0.25\textwidth}>{\raggedright\arraybackslash}X>{\raggedright\arraybackslash}p{0.21\textwidth}}
\toprule
\textbf{Requisito} & \textbf{Descripción} & \textbf{Fundamento} \\
\midrule
Legitimación del transmitente & Debe ser realizada por el propietario: solo el dueño puede transmitir el dominio. & Art.~399, \textit{nemo plus iuris} \\
Acto causal (título suficiente) & Debe existir una causa válida: compraventa, donación, dación en pago, permuta, etc. & Art.~1892 \\
Capacidad de ambas partes & Tanto quien transmite como quien adquiere deben ser capaces. & Art.~1892 \\
Escritura pública (inmuebles) & El acto causal debe plasmarse en escritura pública. & Art.~1017 \\
Tradición material efectiva & La entrega física de la cosa, salvo los dos casos de excepción. & Art.~1924 \\
\bottomrule
\end{tabularx}
\caption{Requisitos de la tradición traslativa de dominio}
\end{table}

Si faltan los requisitos de legitimación, título y capacidad, lo que se transmite es meramente la relación de poder que el tradente tenía con la cosa, y no el derecho (véase el capítulo~\ref{cap:dinamica}).

\section{La inscripción registral}

El efecto de la inscripción registral depende de la clase de bien.

\begin{table}[H]
\centering
\small
\begin{tabularx}{\textwidth}{>{\raggedright\arraybackslash}p{0.17\textwidth}>{\raggedright\arraybackslash}p{0.19\textwidth}>{\raggedright\arraybackslash}X}
\toprule
\textbf{Bien} & \textbf{Efecto} & \textbf{Consecuencia práctica} \\
\midrule
Inmuebles & Declarativo & El derecho nace con el título más el modo (escritura más tradición). La inscripción solo lo hace oponible \textit{erga omnes}, es decir, conocido por toda la comunidad. \\
Automotores & Constitutivo & Sin inscripción (transferencia registral) no se es dueño del automotor: el derecho nace con la inscripción. \\
\bottomrule
\end{tabularx}
\caption{Efectos de la inscripción registral}
\end{table}

\cornell{adquisición de los derechos reales}{
\cq{Distinga adquisición originaria y derivada.}{Originaria: no deriva de un vínculo anterior (cosas sin dueño, prescripción adquisitiva, adquisición legal). Derivada: un enajenante transmite a un adquirente, entre vivos o por causa de muerte.}
\cq{¿Qué dispone el art.~1896 y qué principio expresa?}{El juez no puede constituir derechos reales ni imponer su constitución: principio de legalidad.}
\cq{¿Qué exige el art.~1892?}{Título suficiente más modo suficiente para la adquisición derivada entre vivos.}
\cq{¿Por qué el título se entiende en sentido abstracto?}{Porque es el acto jurídico causal que plasma la transmisión, no el instrumento: entregar un papel que no plasma un acto causal es no tener título.}
\cq{¿Por qué la tradición no puede suplirse por una cláusula de escritura (art.~1924)?}{Porque es un acto material; la posesión no se transmite con palabras escritas.}
\cq{Compare la inscripción registral en inmuebles y automotores.}{Inmuebles: declarativa (oponibilidad). Automotores: constitutiva (sin inscripción no hay dominio).}
}{El derecho real se adquiere por vía originaria o derivada. En la derivada entre vivos rige la teoría del título y modo (art.~1892): acto jurídico causal válido más tradición material, que no puede suplirse por fórmulas escritas. La inscripción es declarativa en inmuebles y constitutiva en automotores.}

\chapter[Excepciones a la tradición material]{Excepciones a la tradición material}\label{cap:brevi}
% ------------------------------------------------------------

El ordenamiento admite solo dos casos en que la adquisición opera sin entrega material efectiva: la \textit{traditio brevi manu} y el constituto posesorio. Fuera de ellos, el principio de legalidad no permite apartarse del requisito de tradición. Ambos casos se explican por una misma idea: evitar un dispendio material inútil cuando la cosa ya está en poder de quien debe recibirla, o seguirá estándolo.

\section{\textit{Traditio brevi manu} (arts.~1892 y 1933)}

\begin{definicion}{\textit{traditio brevi manu}}
Es la excepción que permite que quien ya tenía la cosa como \textbf{tenedor} (por ejemplo, como locatario) pase a tenerla como \textbf{titular del derecho real} (propietario) por un acto jurídico causal, sin necesidad de tradición material, haciéndolo constar en la escritura. La relación de poder \textbf{asciende de categoría}: de tenencia a dominio.
\end{definicion}

\begin{ejemplo}{la inquilina que compra su departamento}
María alquiló durante muchos años un departamento. Se le presentó la posibilidad de comprarlo, coincidiendo con el interés del propietario en venderlo, y firmó la escritura pública de compraventa (título suficiente). Sin la \textit{traditio brevi manu}, debería contratar una mudadora, sacar sus muebles, salir en calidad de inquilina, esperar que el propietario ingresara y que luego se lo entregara como propietaria: un dispendio inútil. La ley admite que, haciéndolo constar en la escritura, María pase de tenedora (locataria) a propietaria sin ningún movimiento material.
\end{ejemplo}

Existe una variante: si el comprador es otro adquirente y María necesita seguir alquilando, puede continuar en el inmueble como inquilina del nuevo propietario. María, que tenía la cosa a nombre del propietario original, pasa a tenerla a nombre del adquirente sin movimiento material, con debida constancia notarial y fecha cierta.

\section{Constituto posesorio (arts.~1892 y 1923)}

\begin{definicion}{constituto posesorio}
Es el caso inverso: quien tenía la cosa como \textbf{propietario} la transmite a un tercero y pasa a tenerla como \textbf{tenedor} (locatario, comodatario, etc.), sin necesidad de tradición material y con la debida constancia notarial. La relación de poder \textbf{desciende de categoría}.
\end{definicion}

\begin{ejemplo}{el propietario endeudado de 2001}
En el año 2001, un señor con muchas deudas decide vender su departamento para pagarlas y propone al adquirente alquilárselo. Lo que cambia es la calidad en que se encuentra respecto de la cosa: antes era propietario, ahora es locatario. Sin el constituto posesorio, debería entregar materialmente la posesión al adquirente y luego reingresar como inquilino. La ley evita ese disparate: el propietario que vende permanece en el inmueble, pero como locatario del nuevo dueño, con debida constancia notarial y fecha cierta.
\end{ejemplo}

\section{Comparación}

\begin{table}[H]
\centering
\small
\begin{tabularx}{\textwidth}{>{\raggedright\arraybackslash}p{0.22\textwidth}>{\raggedright\arraybackslash}X>{\raggedright\arraybackslash}X}
\toprule
 & \textbf{\textit{Traditio brevi manu}} & \textbf{Constituto posesorio} \\
\midrule
Artículos & 1892 y 1933 & 1892 y 1923 \\
Situación previa & Tenedor (locatario, comodatario) & Propietario (titular del dominio) \\
Situación posterior & Propietario & Tenedor (locatario, comodatario) \\
Movimiento & Ascendente ($\uparrow$) & Descendente ($\downarrow$) \\
Ejemplo & La inquilina que compra su departamento & El propietario de 2001 que vende y queda como inquilino \\
Tradición material & No requerida & No requerida \\
Constancia & En la escritura & Notarial, con fecha cierta \\
\bottomrule
\end{tabularx}
\caption{Excepciones a la tradición material}
\end{table}
% TODO: verificar referencia presente en las fuentes originales (el art.~1933 se cita para la \textit{traditio brevi manu} en esta parte y para la obligación de restituir en el capítulo sobre efectos de las relaciones de poder; el art.~1923 se cita para el constituto posesorio y para la tradición como modo bilateral).

\cornell{excepciones a la tradición}{
\cq{¿Cuáles son los únicos casos de adquisición sin entrega material?}{La \textit{traditio brevi manu} y el constituto posesorio; fuera de ellos rige la tradición material.}
\cq{¿Qué es la \textit{traditio brevi manu}?}{Quien tenía la cosa como tenedor pasa a tenerla como propietario por un acto causal, sin entrega material, haciéndolo constar en la escritura. Es un ascenso de categoría.}
\cq{¿Qué es el constituto posesorio?}{Quien tenía la cosa como propietario la transmite y pasa a tenerla como tenedor, sin entrega material y con constancia notarial y fecha cierta. Es un descenso de categoría.}
\cq{¿Qué problema práctico resuelve cada una?}{Evitar una entrega material inútil: salir y reingresar al inmueble con otra calidad (María; el propietario de 2001).}
}{La tradición debe ser material y efectiva, salvo dos excepciones legales: la \textit{traditio brevi manu} (tenedor que pasa a propietario) y el constituto posesorio (propietario que pasa a tenedor), ambas con constancia en instrumento y sin movimiento material.}

\chapter[Publicidad, adquisición legal, transmisión y extinción]{Publicidad, adquisición legal, transmisión y extinción}\label{cap:publicidad}
% ------------------------------------------------------------

\section{Publicidad (art.~1893)}

Los derechos reales deben ser respetados por todos y, si deben ser respetados por todos, \textbf{deben ser conocidos por todos}. De allí el principio de publicidad, regulado en el artículo 1893, que se vincula con la oponibilidad a terceros de las relaciones de poder.

\begin{table}[H]
\centering
\small
\begin{tabularx}{\textwidth}{>{\raggedright\arraybackslash}p{0.25\textwidth}>{\raggedright\arraybackslash}X}
\toprule
\textbf{Tipo de publicidad} & \textbf{Descripción} \\
\midrule
Posesoria & Desde la existencia misma del hombre en sociedad, el vínculo más claro con una cosa es tenerla: la posesión es en sí misma publicidad. Opera para toda clase de cosas. \\
Registral & La inscripción en un registro público permite que los terceros conozcan quién es el titular y hace oponible el derecho. Opera cuando se trata de cosas registrables. Su efecto es declarativo en inmuebles y constitutivo en automotores. \\
\bottomrule
\end{tabularx}
\caption{Tipos de publicidad de los derechos reales}
\end{table}

\section{Adquisición legal (art.~1894)}

Existen casos en que la ley misma dispone que el derecho se adquiere, sin que concurran título y modo en sentido estricto. Son casos de adquisición legal u originaria:

\begin{enumerate}
    \item el derecho real de \textbf{habitación del cónyuge supérstite}: al fallecer uno de los cónyuges, la ley reconoce al sobreviviente el derecho real de habitación sobre el inmueble asiento del hogar conyugal, gratuito y vitalicio, siempre que se cumplan los requisitos legales;
    \item el derecho real de \textbf{habitación del conviviente supérstite}: el Código incorporó al conviviente (unión convivencial) entre los beneficiarios de este derecho, manifestación del cambio de paradigma del nuevo cuerpo normativo, a veces llamada «el conviviente impensado»;
    \item los \textbf{condominios con indivisión forzosa} (capítulo~\ref{cap:medianeria});
    \item los derechos de los adquirentes de buena fe y a título oneroso, cuyo desarrollo se presenta en el capítulo~\ref{cap:muebles}.
\end{enumerate}
El régimen de la habitación del cónyuge y del conviviente supérstites se desarrolla en el capítulo~\ref{cap:habsuperstite}.
% TODO: verificar consistencia entre fuentes originales (una versión enumera tres casos de adquisición legal en el art.~1894 y otra agrega un cuarto; el régimen de cosas muebles se cita luego como «arts.~1894 a 1895»).

\section{Transmisión}

Todos los derechos, reales o personales, son en principio transmisibles. La transmisión de derechos reales se rige por la teoría del título y modo: quien transmite lo hace con título suficiente y modo, y quien adquiere los recibe de ese modo. El \textit{nemo plus iuris} es el límite infranqueable: nadie puede transmitir más derechos de los que tiene. Los casos particulares se analizan al tratar cada derecho real.

\section{Extinción (art.~1907)}

El derecho real se extingue en forma relativa o absoluta.

\begin{table}[H]
\centering
\small
\begin{tabularx}{\textwidth}{>{\raggedright\arraybackslash}p{0.17\textwidth}>{\raggedright\arraybackslash}X>{\raggedright\arraybackslash}p{0.16\textwidth}>{\raggedright\arraybackslash}p{0.22\textwidth}}
\toprule
\textbf{Tipo} & \textbf{Descripción} & \textbf{¿Continúa?} & \textbf{Ejemplo} \\
\midrule
Relativa & Se extingue para el titular actual. & Sí, en otro sujeto. & Venta, donación, dación en pago. \\
Absoluta & El derecho desaparece por completo; como el derecho real es un vínculo entre titular y cosa, destruida la cosa ese vínculo se extingue. & No. & Destrucción de la cosa. \\
Consolidación & Se reúnen en una misma persona el derecho real principal y el que lo desmembraba. & Se consolida en uno solo. & El usufructuario adquiere el dominio pleno; el nudo propietario adquiere el usufructo. \\
\bottomrule
\end{tabularx}
\caption{Extinción de los derechos reales}
\end{table}

La consolidación es análoga a lo que en los derechos personales se denomina confusión: el derecho menor queda absorbido.

\section{El derecho real como relación de poder}

El derecho real es una \textbf{relación de poder del titular respecto de la cosa}. Ese esquema, esa línea o flecha, es lo que permite perseguir la cosa, como ocurre en las hipotecas, aun cuando hubiera cambiado de titular. La relación es defendida y protegida por el sistema jurídico y debe ser respetada por toda la sociedad; de allí la necesidad de la publicidad, posesoria o registral, a pesar de los límites que la Nación decida imponer en función de los intereses de la comunidad.

Desde el nacimiento y a lo largo de toda su existencia, los seres humanos tienen contactos permanentes con las cosas. Ese contacto directo e inmediato tiene una protección jurídica que debe ser respetada, porque es también una manera de respetar la paz social.

Sin embargo, no siempre los vínculos de poder están en cabeza del titular del derecho real: no siempre las personas viven en las tierras de las que son propietarias. Existen diversas relaciones de poder y diversas calidades en que un sujeto puede vincularse con una cosa. A ellas se dedica la Parte~III.

\begin{conceptoclave}{síntesis de la adquisición y extinción}
La adquisición derivada por actos entre vivos exige título suficiente y modo suficiente (art.~1892), sin que la tradición pueda suplirse por fórmulas escritas, salvo la \textit{traditio brevi manu} y el constituto posesorio. Para su oponibilidad, los derechos reales se publicitan mediante la posesión o la inscripción registral. Se extinguen en forma relativa (continúan en otro titular) o absoluta (típicamente, por destrucción de la cosa), siempre en el marco de la preeminencia del interés colectivo que proclama el art.~240.
\end{conceptoclave}

\cornell{publicidad, adquisición legal y extinción}{
\cq{¿Por qué los derechos reales requieren publicidad?}{Porque deben ser respetados por todos y, por tanto, conocidos por todos (art.~1893). Puede ser posesoria o registral.}
\cq{¿Qué casos de adquisición legal enumera el art.~1894?}{Habitación del cónyuge supérstite, habitación del conviviente supérstite y condominios con indivisión forzosa; según una fuente, también los adquirentes de buena fe a título oneroso.}
\cq{Distinga extinción relativa y absoluta (art.~1907).}{Relativa: el derecho se extingue para un titular pero continúa en otro. Absoluta: desaparece, típicamente por destrucción de la cosa.}
\cq{¿Qué es la consolidación?}{La reunión en una persona del derecho real principal y del que lo desmembraba (usufructuario que adquiere el dominio); análoga a la confusión.}
\cq{¿Por qué el derecho real es una «relación de poder»?}{Porque es un vínculo directo del titular con la cosa, protegido por el sistema y exigible a todos; permite perseguir la cosa y exige publicidad.}
}{La publicidad (posesoria o registral) hace conocido y oponible el derecho real. La ley adquiere por sí ciertos derechos (habitación del cónyuge y del conviviente, condominios con indivisión forzosa). La extinción puede ser relativa, absoluta o por consolidación. El derecho real es una relación de poder protegida por el sistema jurídico.}

% ============================================================
% ============================================================


% ============================================================
\part{Relaciones de poder: posesión y tenencia}
% ============================================================

\chapter[Las relaciones entre la persona y la cosa]{Las relaciones entre la persona y la cosa}\label{cap:poder}
% ------------------------------------------------------------

La tradición, la publicidad posesoria y la defensa de la posesión presuponen una pregunta previa: qué relaciones puede tener una persona con una cosa. Esas relaciones son \textbf{taxativas}: se denominan y regulan de manera específica y no pueden crearse otras distintas, pues serían nulas, en consonancia con el \textit{numerus clausus} del artículo 1884.

\section{Las relaciones de hecho reconocidas}

Las únicas relaciones de hecho que el ordenamiento reconoce son la \textbf{yuxtaposición local}, la \textbf{tenencia} y la \textbf{posesión}. Para la teoría posesoria interesan, según los fundamentos del anteproyecto, solo la posesión y la tenencia, con el agregado del \textbf{servidor de la posesión} al solo fin de la defensa extrajudicial. Las relaciones de poder se regulan en los artículos 1908 a 1940.

\subsection{La yuxtaposición local}

\begin{definicion}{yuxtaposición local}
Mero \textbf{contacto físico casual} entre una persona y una cosa.
\end{definicion}

Existe algo de voluntad (sin ella no habría contacto), pero no está dirigida hacia esa cosa en particular y no hacia otra.

\begin{ejemplo}{el colectivo}
Una persona sube a un colectivo, se apoya en una barra, se sienta y, cuando sube alguien que lo necesita más, cede el asiento. Hay contacto con las cosas, pero no está direccionado por el sujeto a esa cosa y no a otra.
\end{ejemplo}

\subsection{Los servidores de la posesión}

Quien utiliza una cosa en virtud de una relación de dependencia, hospedaje u hospitalidad se llama \textbf{servidor de la posesión} (art.~1911).

\begin{table}[H]
\centering
\small
\begin{tabularx}{\textwidth}{>{\raggedright\arraybackslash}p{0.17\textwidth}>{\raggedright\arraybackslash}X>{\raggedright\arraybackslash}X}
\toprule
\textbf{Vínculo} & \textbf{Ejemplo} & \textbf{Consecuencia jurídica} \\
\midrule
Dependencia & El obrero y su máquina de trabajo, el empleado y su computadora: la cuida y la usa todos los días, pero si lo despiden no se la lleva. & Quien sirve a la posesión no tiene \textit{animus} ni \textit{corpus} propio: posee en nombre del empleador. \\
Hospedaje & El huésped de un hotel al que asignan la habitación 22: si no hay agua, llama al conserje; con otra habitación, esa sería «su habitación». & El huésped no tiene posesión ni tenencia, solo un contacto derivado del servicio. \\
Hospitalidad & El amigo que visita una casa y se queda unos días: usa la cama, el baño, la heladera, y al irse no se lleva esas cosas. & El vínculo dura mientras dura la hospitalidad; no genera posesión ni tenencia. \\
\bottomrule
\end{tabularx}
\caption{Servidores de la posesión}
\end{table}

\section{Posesión, tenencia y servidor de la posesión}

\begin{definicion}{posesión (art.~1909 CCyC)}
Hay posesión cuando una persona, por sí o por medio de otra, ejerce un poder de hecho sobre una cosa, comportándose como titular de un derecho real, lo sea o no.
\end{definicion}

Tiene dos elementos: el \textbf{\textit{corpus}} (contacto material con la cosa) y el \textbf{\textit{animus}} o \textit{animus domini} (ánimo de dueño): el poseedor se comporta como si la cosa fuera propia, sin reconocer en otro un señorío superior, lo sea o no el dueño.

\begin{definicion}{tenencia (art.~1910 CCyC)}
Hay tenencia cuando una persona, por sí o por medio de otra, ejerce un poder de hecho sobre una cosa y se comporta como representante del poseedor.
\end{definicion}

Tiene un solo elemento, el \textit{corpus}: no hay \textit{animus}, porque el tenedor reconoce en otro la propiedad. El caso típico es el inquilino, que tiene relación directa con la cosa y numerosos derechos sobre ella, pero reconoce en otro la propiedad: paga el alquiler y, si le reclaman por desperfectos, responde que la propiedad no es suya.

\begin{table}[H]
\centering
\small
\begin{tabularx}{\textwidth}{>{\raggedright\arraybackslash}p{0.2\textwidth}>{\raggedright\arraybackslash}X>{\raggedright\arraybackslash}X>{\raggedright\arraybackslash}X}
\toprule
 & \textbf{Poseedor} & \textbf{Tenedor} & \textbf{Servidor} \\
\midrule
Elementos & \textit{Corpus} y \textit{animus} & Solo \textit{corpus} & Ni \textit{corpus} ni \textit{animus} propios \\
Actitud ante la cosa & Se comporta como titular, lo sea o no & Reconoce en otro la propiedad & Usa la cosa por dependencia, hospedaje u hospitalidad \\
Ejemplo & Fermín, que se lleva el teléfono & Juan, que lo tiene hasta que vuelva su dueño & El empleado y su computadora \\
\bottomrule
\end{tabularx}
\caption{Poseedor, tenedor y servidor de la posesión}
\end{table}
% TODO: verificar referencia presente en las fuentes originales (el art.~1911 se cita en una parte para el servidor de la posesión y, en otra, para la presunción de posesión).

\begin{ejemplo}{el teléfono: Juan y Fermín}
Una persona (A) espera un llamado y debe ausentarse un momento. Le pide a Juan que le tenga el teléfono y que, si llaman, avise que ya vuelve. Llega Carolina y le pide que se lo preste porque perdió el suyo; Juan responde que no puede, porque el teléfono no es suyo sino de A, y cuando A regresa, Juan espontáneamente le extiende el teléfono. \textbf{Juan es tenedor:} reconoce en otro un señorío superior.

En otra versión, A deja el teléfono sobre el escritorio y Fermín lo toma, lo guarda en su mochila, sale del aula, habla en el bar y lo lleva a su casa. Se comporta como si fuera dueño: no reconoció la propiedad de A, pues de haberla reconocido no se lo habría llevado. \textbf{Fermín es poseedor:} tiene \textit{corpus} y \textit{animus} y se comporta como propietario, lo sea o no.
\end{ejemplo}

\cornell{relaciones de poder}{
\cq{¿Cuáles son las relaciones de hecho reconocidas y por qué son taxativas?}{Yuxtaposición local, tenencia y posesión. Son taxativas por el \textit{numerus clausus}: no pueden crearse otras bajo pena de nulidad.}
\cq{¿Qué es la yuxtaposición local?}{Mero contacto físico casual con una cosa, sin voluntad dirigida a esa cosa en particular (el pasajero del colectivo).}
\cq{¿Qué es el servidor de la posesión?}{Quien usa la cosa por dependencia, hospedaje u hospitalidad. No tiene posesión ni tenencia; solo importa para la defensa extrajudicial.}
\cq{Defina posesión y tenencia por sus elementos.}{Posesión (art.~1909): \textit{corpus} más \textit{animus domini}. Tenencia (art.~1910): solo \textit{corpus}; el tenedor reconoce en otro la propiedad y se comporta como representante del poseedor.}
\cq{Aplique al teléfono: ¿quién es tenedor y quién poseedor?}{Juan es tenedor, porque reconoce el señorío de A. Fermín es poseedor, porque se lleva el teléfono comportándose como dueño.}
}{Las relaciones de poder reconocidas son taxativas. La posesión exige \textit{corpus} y \textit{animus domini}; la tenencia, solo \textit{corpus}, reconociendo en otro la propiedad; el servidor de la posesión usa la cosa por dependencia, hospedaje u hospitalidad y solo interesa para la defensa extrajudicial.}

\chapter[Presunciones, clasificación y defensa de la posesión]{Presunciones, clasificación y defensa de la posesión}\label{cap:clasificacion}
% ------------------------------------------------------------

\section{Presunciones y reglas posesorias}

La ley resuelve la imposibilidad práctica de conocer el verdadero vínculo que cada persona tiene con las cosas que usa mediante un régimen de presunciones. Quien tiene una relación de hecho con una cosa y se comporta como titular del derecho real se presume poseedor, lo sea o no. Esto hace posible la vida en sociedad: sería imposible verificar cada relación de poder que se observa cotidianamente, por lo que, si no existe un interés legítimo en conocer el verdadero contenido jurídico de una relación, se respetan las múltiples relaciones de hecho observadas.

\begin{regla}{presunciones de posesión, legitimidad y buena fe}
\begin{itemize}
    \item \textbf{Art.~1911 (posesión):} toda relación de poder se presume posesión, excepto prueba en contrario.
    \item \textbf{Art.~1916 (legitimidad):} las relaciones de poder se presumen legítimas, a menos que exista prueba en contrario.
    \item \textbf{Art.~1919 (buena fe):} la posesión se presume de buena fe, a menos que exista prueba en contrario o que el título sea de mala fe.
\end{itemize}
\end{regla}

Otras reglas en materia posesoria:

\begin{itemize}
    \item \textbf{Innecesariedad del título (art.~1917):} el poseedor no tiene obligación de producir su título, salvo que deba exhibirlo como obligación inherente a su posesión.
    \item \textbf{Fecha y extensión (art.~1914):} si media título, se presume que la relación de poder comienza desde la fecha del título y tiene la extensión que en él se indica.
    \item \textbf{Relatividad de los vicios (art.~1921):} los vicios de la posesión son relativos respecto de aquel contra quien se ejercen; aun si la posesión es viciosa, solo puede alegar el vicio quien lo sufrió.
    \item \textbf{Inmutabilidad de la causa:} nadie puede cambiar por sí mismo la causa de su posesión. Quien comenzó como tenedor continúa siéndolo; quien era poseedor de mala fe continúa siéndolo.
\end{itemize}

\section{Clasificación de la posesión}

\begin{table}[H]
\centering
\small
\begin{tabularx}{\textwidth}{>{\raggedright\arraybackslash}p{0.2\textwidth}>{\raggedright\arraybackslash}X>{\raggedright\arraybackslash}X}
\toprule
\textbf{Tipo} & \textbf{Definición} & \textbf{Ejemplo} \\
\midrule
Legítima & Ejercicio de un derecho real adquirido conforme a la ley: título suficiente, modo suficiente y capacidad de ambas partes. & El propietario que compró con escritura pública e hizo la tradición. \\
Ilegítima de buena fe & Sin título suficiente, pero el poseedor cree que tiene derecho a poseer (error excusable). & Quien compra por boleto de compraventa (falta la escritura); quien adquiere un libro usado sin saber que era robado. \\
Ilegítima de mala fe simple & El poseedor duda de la legitimidad de su adquisición, pero igualmente adquiere. & Quien compra un reloj de alta gama a precio ínfimo en un lugar inusual, sin preguntarse por qué se lo venden tan barato. \\
Ilegítima de mala fe viciosa, muebles & Posesión adquirida mediante vicios específicos de las cosas muebles. & Hurto, estafa, abuso de confianza. \\
Ilegítima de mala fe viciosa, inmuebles & Posesión adquirida mediante vicios específicos de los inmuebles. & Clandestinidad, violencia, abuso de confianza. \\
\bottomrule
\end{tabularx}
\caption{Clasificación de la posesión}
\end{table}

El vicio común a muebles e inmuebles es el \textbf{abuso de confianza}.

\begin{ejemplo}{abuso de confianza en un inmueble}
Un propietario presta a un amigo un inmueble desocupado que va a vender pronto. Cuando le informa que lo vendió y que debe entregar la posesión, el amigo responde que él es poseedor, que tiene \textit{animus} y que allí se queda. Hay una intervención del título basada en un abuso de confianza (véase el capítulo~\ref{cap:dinamica}).
\end{ejemplo}

\begin{importante}
Lo que se califica es el \textbf{vínculo}, no la persona. El derecho califica de legítima o ilegítima la relación que la persona tiene con la cosa, no su rostro, su nombre ni su historia: una persona honesta puede tener una posesión ilegítima respecto de una cosa, y un delincuente puede ser poseedor legítimo del inmueble que adquirió con título suficiente y tradición.
\end{importante}

\section{La defensa de la posesión}

La posesión se protege aunque sea ilegítima o viciosa, por razones de orden público:

\begin{enumerate}
    \item para proteger a la persona del poseedor, el anterior y el nuevo;
    \item para proteger la paz social, evitando la violencia privada: quien tiene derecho a la posesión debe recuperarla por vía legal, no por la fuerza;
    \item incluso el verdadero propietario que toma la posesión por la fuerza puede ser demandado por el poseedor despojado.
\end{enumerate}

Las acciones concretas con que se la defiende se estudian en el capítulo~\ref{cap:posesorias}.

La carga de la prueba recae siempre en quien controvierte la relación de poder y no en el poseedor: su posesión se presume legítima y de buena fe.

\cornell{presunciones, clasificación y defensa}{
\cq{¿Qué presunciones rigen y para qué sirven?}{De posesión (art.~1911), de legitimidad (1916) y de buena fe (1919). Permiten la vida social sin verificar cada relación de poder; la prueba contraria corresponde a quien controvierte.}
\cq{¿Qué significa que los vicios sean relativos (art.~1921)?}{Solo puede alegar el vicio quien lo sufrió.}
\cq{¿Qué quiere decir que nadie puede cambiar por sí la causa de su posesión?}{Que el tenedor sigue siendo tenedor y el poseedor de mala fe sigue siéndolo, aunque pase el tiempo.}
\cq{Clasifique la posesión.}{Legítima; ilegítima de buena fe; de mala fe simple (duda); de mala fe viciosa (hurto, estafa, abuso de confianza en muebles; clandestinidad, violencia, abuso de confianza en inmuebles).}
\cq{¿Qué se califica: la persona o el vínculo?}{El vínculo con la cosa.}
\cq{¿Por qué se defiende incluso la posesión viciosa?}{Para proteger a la persona y la paz social, evitando la violencia privada.}
}{Las relaciones de poder se presumen posesión, legítima y de buena fe. La posesión se clasifica según su legitimidad y vicios; se califica el vínculo y no a la persona. Se protege incluso la posesión ilegítima para preservar la paz social.}

\chapter[Dinámica de las relaciones de poder]{Dinámica de las relaciones de poder}\label{cap:dinamica}
% ------------------------------------------------------------

Este capítulo sigue la vida de una relación de poder: cómo se adquiere, qué efectos produce, cómo puede alterarse su causa, cómo se acumula en el tiempo y cómo se extingue.

\section{Adquisición}

\begin{regla}{art.~1922 CCyC, adquisición de poder}
Para adquirir una relación de poder sobre una cosa, esta debe ser establecida voluntariamente: a) por sujeto capaz, excepto las personas menores de edad, para quienes es suficiente que tengan diez años; b) por medio de contacto con la cosa, de la posibilidad física de establecerlo, o cuando ella ingresa en el ámbito de custodia del adquirente.
\end{regla}

Además, la cosa debe estar libre de toda otra relación excluyente (\textbf{posesión vacua}, art.~1926). Los modos de adquisición son:

\begin{itemize}
    \item \textbf{unilateral:} el apoderamiento;
    \item \textbf{bilateral:} la tradición, por excelencia (art.~1923): una persona (el tradente) entrega la cosa a otra, que ocupa la posición que tenía el tradente;
    \item \textbf{originaria:} la prescripción adquisitiva, entre otros modos.
\end{itemize}

\begin{table}[H]
\centering
\small
\begin{tabularx}{\textwidth}{>{\raggedright\arraybackslash}X>{\raggedright\arraybackslash}X}
\toprule
\textbf{Tradición traslativa de dominio} & \textbf{Tradición posesoria (mera)} \\
\midrule
Requiere: (1) que la realice el legitimado, es decir, el titular del derecho; (2) título suficiente (acto jurídico causal); (3) capacidad de ambas partes. & Si no se dan esos requisitos, se transmite meramente la relación de poder que el tradente tiene con la cosa. Nadie puede transmitir un derecho mejor o más extenso que el que tenía. \\
\bottomrule
\end{tabularx}
\caption{Tradición traslativa y tradición posesoria}
\end{table}

Si un poseedor no titular transmite la posesión, solo transmite el hecho de la posesión; pero si el propietario (legitimado, capaz y con título suficiente) transmite, lo que transmite es el derecho de poseer.

\section{Efectos}

\subsection{Derechos inherentes a la posesión}

Son los que tiene todo aquel que se encuentra en poder de una cosa, sea poseedor o tenedor; por ejemplo, los límites al dominio.

\begin{ejemplo}{ventana en pared medianera}
Un vecino, como propietario, tiene amplias facultades, pero también límites: no puede abrir una ventana de vista en la pared medianera, porque quien está del otro lado tiene derecho a que no se construya. Obligación y derecho son las dos caras de una misma moneda.
\end{ejemplo}

\subsection{Obligaciones del poseedor y del tenedor}

\begin{regla}{art.~1933 CCyC}
El poseedor o tenedor debe restituir la cosa a quien tenga derecho de reclamarla, aunque no se la haya dado a él.
\end{regla}

Quien tiene una relación de poder con la cosa y no tiene derecho real tiene la \textbf{obligación de restituirla} y el \textbf{derecho a que le sean restituidas las mejoras} realizadas en su conservación.

\subsection{La \textit{laudatio auctoris}}

El tenedor tiene la \textbf{obligación legal de identificar al poseedor} (\textit{laudatio auctoris}) cuando recibe una demanda dirigida al propietario, bajo apercibimiento de daños y perjuicios.

\begin{ejemplo}{el inquilino que no avisa}
El inmueble está ocupado por un inquilino, que es tenedor. Recibe una demanda de reivindicación y sostiene que, como no es el titular, la demanda no lo afecta: no la contesta ni avisa al propietario. La conducta es gravísima, porque priva al propietario de la posibilidad de defenderse en juicio.
\end{ejemplo}

\section{Intervención del título}

La regla general es que nadie puede cambiar por sí mismo la especie de su relación de poder: ni el tiempo, ni la voluntad interna, ni el pensamiento modifican el título.

\begin{itemize}
    \item \textbf{Intervención unilateral:} el tenedor, por actos exteriores, manifiesta su voluntad de privar al poseedor de la cosa.
    \item \textbf{Intervención bilateral:} existe acuerdo entre poseedor y tenedor (art.~1915).
\end{itemize}

\begin{ejemplo}{intervención unilateral}
El propietario presta su propiedad a un amigo para que viva allí hasta la venta. Cuando se produce la venta, le pide que desocupe; el amigo responde que supuso otra cosa y que ahora posee, porque posee y su posesión se presume legítima. Ha intervertido unilateralmente el título. Del mismo modo, quien recibió el teléfono en tenencia afirma ante quien se lo entregó que nadie se lo dejó y que es suyo: hay un acto exterior. No hay conversión natural por el paso del tiempo ni por haber pensado de otro modo.
\end{ejemplo}

\section{Unión de posesiones}

\begin{regla}{art.~1901 CCyC}
El heredero continúa la posesión de su causante. El sucesor particular puede unir su posesión a la de sus antecesores, siempre que derive inmediatamente de las otras. En la prescripción breve las posesiones unidas deben ser de buena fe y estar unidas por un vínculo jurídico.
\end{regla}

La unión de posesiones es una \textbf{suma temporal} del tiempo de dos posesiones de distintos poseedores, útil para alcanzar el número de años requerido para la \textbf{prescripción adquisitiva}.

\begin{ejemplo}{unión de posesiones}
Un poseedor que posee desde hace diez años cede su posesión por escritura pública; el nuevo poseedor acumula ese tiempo al propio para alcanzar los veinte años de la prescripción larga.
\end{ejemplo}

La unión de posesiones, su aplicación a la prescripción larga y a la breve y la diferencia entre el heredero y el sucesor particular se desarrollan en el capítulo~\ref{cap:prescripcion}.

\section{Actos posesorios}

\begin{regla}{art.~1928 CCyC}
Constituyen actos posesorios sobre la cosa inmueble: su cultura, percepción de frutos, amojonamiento o impresión de signos materiales, mejora, exclusión de terceros y, en general, su apoderamiento por cualquier modo que se obtenga.
\end{regla}

Son actos propios de quien ejerce el \textit{animus domini} y que un tenedor difícilmente realizaría, como construir, hacer mejoras o plantaciones en un inmueble ajeno que va a dejar al cabo del período locativo, salvo que se trate de un acto conservatorio. Quien invoca la posesión debe acreditar que tuvo la disposición física del inmueble y que esa relación era la propia de un poseedor, mediante actos de intervención material sobre la cosa.

\section{Extinción}

\begin{regla}{art.~1931 CCyC}
La posesión y la tenencia se extinguen cuando el sujeto que la ejerce: a) se extingue la cosa; b) otro priva al titular de la cosa; c) se ha imposibilitado el ejercicio de la posesión en forma perdurable; d) por el abandono.
\end{regla}

\begin{ejemplo}{abandono y pérdida}
Un colchón en la pila de residuos se presume abandonado, pues nadie guarda un colchón así: se extingue la relación de poder por abandono. Si un teléfono cae al mar mientras su poseedor saca una foto, es imposible recuperarlo y la relación de poder se extingue.
\end{ejemplo}

\cornell{dinámica de las relaciones de poder}{
\cq{¿Qué se exige para adquirir una relación de poder (arts.~1922 y 1926)?}{Voluntad de un sujeto capaz (diez años para menores), contacto con la cosa, posibilidad física de establecerlo o ingreso en el ámbito de custodia, y que la cosa esté libre de otra relación excluyente (posesión vacua).}
\cq{Distinga tradición traslativa de dominio y tradición posesoria.}{La traslativa requiere legitimado, título suficiente y capacidad. Sin ellos se transmite solo la relación de poder del tradente.}
\cq{¿Qué obligaciones tienen el poseedor y el tenedor sin derecho real (art.~1933)?}{Restituir la cosa a quien pueda reclamarla y derecho al reembolso de mejoras de conservación. El tenedor debe además identificar al poseedor (\textit{laudatio auctoris}).}
\cq{Distinga intervención unilateral y bilateral del título.}{Unilateral: el tenedor, con actos exteriores, pretende privar al poseedor. Bilateral: acuerdo entre ambos (art.~1915). El tiempo o el pensamiento no cambian el título.}
\cq{¿Para qué sirve la unión de posesiones (art.~1901)?}{Sumar el tiempo de posesiones sucesivas para la prescripción adquisitiva; en la breve, deben ser de buena fe y unidas por vínculo jurídico.}
\cq{¿Cómo se extinguen la posesión y la tenencia (art.~1931)?}{Por extinción de la cosa, privación por otro, imposibilidad perdurable de ejercicio y abandono.}
}{La relación de poder se adquiere voluntariamente por un sujeto capaz mediante contacto o custodia, sobre una cosa libre de otra relación excluyente. Genera derechos inherentes y obligaciones de restitución, no cambia de causa por voluntad propia salvo intervención del título, puede sumarse a la de los antecesores y se extingue por pérdida de la cosa, privación, imposibilidad o abandono.}

\chapter[Adquisición legal sobre cosas muebles: art.~1895]{Adquisición legal de derechos reales sobre cosas muebles}\label{cap:muebles}
% ------------------------------------------------------------

Este régimen, de gran interés en materia de orden público, resuelve un conflicto entre dos intereses legítimos: el del propietario que quiere conservar su cosa y el del tercero que la adquirió confiando en la apariencia.

\section{Marco: clasificación de las cosas}

Las cosas son muebles o inmuebles. Conforme al art.~227, son muebles las que pueden desplazarse, por sí mismas o por una fuerza externa. Dentro de las muebles se distinguen las registrables y las no registrables. El régimen estudiado aquí corresponde a las muebles \textbf{no registrables}.

\section{El art.~1895}

\begin{regla}{art.~1895 CCyC}
La posesión de buena fe del subadquirente de cosas muebles no registrables que no sean hurtadas o perdidas es suficiente para adquirir los derechos reales principales, excepto que el verdadero propietario pruebe que la adquisición fue gratuita.
\end{regla}

Requisitos: \textbf{posesión de buena fe} del subadquirente y que la cosa \textbf{no sea hurtada ni perdida}. Con solo ellos, por decisión legal, el poseedor adquiere la propiedad y la defiende frente a todos.

\begin{importante}
El \textbf{título oneroso no es requisito} de la adquisición: el poseedor adquiere la propiedad la haya recibido a título oneroso o gratuito. El título oneroso solo importa para repeler la acción reivindicatoria del anterior propietario, y la carga de probar la gratuidad recae sobre quien reivindica.
\end{importante}

La norma se aplica cuando el propietario se desprende voluntariamente de la cosa, aunque su voluntad hubiera estado viciada. Quien recibió la cosa fue tenedor y, defraudando la confianza, la transmitió a un tercero, el \textbf{subadquirente}.

\begin{ejemplo}{el libro de Derechos Reales}
Un propietario presta un libro de Derechos Reales a una persona, quien, defraudando su confianza, lo vende o regala a Juan. Hay tradición y Juan tiene la posesión, de buena fe; la cosa no fue robada ni perdida, porque el propietario la entregó voluntariamente. Cuando el propietario se encuentra con Juan e identifica el libro por una dedicatoria, pretende reivindicarlo, y Juan se lo niega.

Cada parte tiene su verdad, pero no es posible multiplicar los libros; hay que decidir. El artículo 1895, en consonancia con la mayoría de los ordenamientos modernos, \textbf{sacrifica al propietario y protege al subadquirente}. La única posibilidad del propietario es probar que la adquisición de Juan fue gratuita, lo que presenta dificultades fácticas y jurídicas.
\end{ejemplo}

\section{Seguridad estática y seguridad dinámica}

\begin{table}[H]
\centering
\small
\begin{tabularx}{\textwidth}{>{\raggedright\arraybackslash}X>{\raggedright\arraybackslash}X}
\toprule
\textbf{Seguridad estática} & \textbf{Seguridad dinámica} \\
\midrule
Protege al propietario actual, a quien ya tiene el derecho real y desea conservarlo (en el ejemplo, el dueño del libro). & Protege la circulación de los bienes y la seguridad de las transacciones, a quien adquiere confiando en la apariencia (Juan, el subadquirente de buena fe). \\
Favorece la estabilidad del derecho ya constituido. & Favorece el tráfico jurídico y la circulación rápida de las cosas muebles, que no dejan rastro instrumental. \\
El art.~1895 la sacrifica. & El art.~1895 la privilegia, elección casi uniforme en los sistemas occidentales modernos. \\
\bottomrule
\end{tabularx}
\caption{Seguridad estática y dinámica en las cosas muebles}
\end{table}

El Código adopta así la máxima «\textbf{posesión vale título}»: adquisición legal de derechos reales sobre muebles por subadquirentes de buena fe cuando la cosa no ha sido hurtada ni perdida. La misma tensión se presenta en los inmuebles, con otra solución, en el capítulo~\ref{cap:adq-dominio}.

\cornell{adquisición legal sobre muebles}{
\cq{¿Qué requisitos exige el art.~1895?}{Posesión de buena fe del subadquirente y que la cosa no sea hurtada ni perdida; se aplica a muebles no registrables cuando el propietario se desprendió voluntariamente.}
\cq{¿Es requisito el título oneroso?}{No para adquirir. Solo importa para repeler la reivindicación: el propietario puede recuperar si prueba que la adquisición fue gratuita.}
\cq{¿Qué hace el art.~1895 con los dos intereses en conflicto?}{Sacrifica al propietario (seguridad estática) y protege al subadquirente (seguridad dinámica).}
\cq{¿Qué significa «posesión vale título»?}{El subadquirente de buena fe de una cosa no hurtada ni perdida adquiere el derecho real por la sola posesión.}
}{El art.~1895 privilegia la seguridad de las transacciones: el subadquirente de buena fe de una cosa mueble no registrable, no hurtada ni perdida, adquiere el derecho real principal, salvo que el propietario pruebe que adquirió gratuitamente.}

% ============================================================
% ============================================================


% ============================================================
\part{El dominio}
% ============================================================

\chapter[Concepto, caracteres y clasificación del dominio]{Concepto, caracteres y clasificación del dominio}\label{cap:dominio}
% ------------------------------------------------------------

Quien es dueño de una casa puede, como propietario, hacer casi todo con ella: venderla, alquilarla, reformarla. Esa propiedad puede ser plena o estar afectada por una hipoteca o un usufructo. Cabe preguntarse además cómo se llegó a ser propietario, por compra con escritura o por herencia. Y, aun siendo dueño, no es posible hacer ruido hasta las tres de la mañana, tapar la ventana del vecino ni bloquear el río. Ese es el recorrido de esta parte: concepto y caracteres del dominio, su clasificación, su adquisición y, finalmente, sus límites.

\section{Las disposiciones comunes aplicadas al dominio}

El CCyC utiliza, además de los principios rectores aplicables a todo el ordenamiento, partes generales y disposiciones comunes que no se repiten en la regulación de cada derecho real. Por eso, cada vez que se aborda un derecho real particular, es metodológicamente fundamental «exponer en acto» lo aprendido sobre las disposiciones comunes, lo que implica aplicar la nueva metodología del Código en la actividad cotidiana. Aplicado al dominio, el ejercicio consiste en afirmar que:

\begin{itemize}
    \item es un \textbf{derecho real} y, por tanto, un poder jurídico, de \textbf{estructura legal};
    \item se ejerce directamente sobre un objeto y, muy excepcionalmente, sobre un derecho cuando la ley específicamente lo indique;
    \item sus normas son de \textbf{orden público} y no pueden ser modificadas por los particulares, bajo pena de nulidad;
    \item es un derecho real sobre \textbf{cosa propia}, que se ejerce por la \textbf{posesión};
    \item se adquiere con \textbf{título suficiente y modo}, y puede adquirirse en forma \textbf{originaria o derivada}.
\end{itemize}

\section{Dominio y propiedad}

Para determinar si «dominio» y «propiedad» son sinónimos, y qué protege la CN cuando garantiza el derecho de propiedad, hay que distinguir dos acepciones.

\begin{table}[H]
\centering
\small
\begin{tabularx}{\textwidth}{>{\raggedright\arraybackslash}p{0.23\textwidth}>{\raggedright\arraybackslash}X>{\raggedright\arraybackslash}p{0.17\textwidth}}
\toprule
\textbf{Concepto} & \textbf{Alcance} & \textbf{Referencia} \\
\midrule
Propiedad en sentido amplio & Abarca derechos personales y reales: los bienes materiales e inmateriales que integran el patrimonio. & CN, arts.~14 y 17 \\
Propiedad en sentido restringido (técnico) & Es el derecho real de dominio: el más extenso y completo que puede recaer sobre una cosa determinada. & CCyC, art.~1941 \\
\bottomrule
\end{tabularx}
\caption{Propiedad en sentido amplio y restringido}
\end{table}

El art.~14 de la CN garantiza a todos los habitantes usar y disponer de la propiedad, y el art.~17 declara que la propiedad es inviolable. La CN se refiere a la propiedad en sentido amplio: protege tanto derechos personales como reales.

\section{Regulación normativa}

El dominio se regula en tres niveles: el Título Tercero del Libro Cuarto (arts.~1941 a 1982), las disposiciones generales en materia de derechos reales (arts.~1882 a 1907) y los principios generales del ordenamiento (arts.~10, 11, 12 y 14).

\section{Concepto}

\begin{definicion}{dominio perfecto (art.~1941 CCyC)}
El dominio es el derecho real en virtud del cual una cosa se encuentra sometida a la voluntad y a la acción de una persona. El dominio pleno o perfecto es el que reúne las facultades de usar, gozar y disponer material y jurídicamente de una cosa, dentro de los límites previstos por la ley. Se presume perfecto hasta que se pruebe lo contrario.
\end{definicion}

Es el derecho más extenso que puede tener una persona sobre una cosa. Los demás derechos reales no son sino fragmentos del dominio: \textbf{desmembraciones} que el titular, que reúne todas las facultades, va constituyendo al separar alguna y conferirla a otro, mediante derechos de disfrute o de garantía. Cuando se habla de «dominio» sin más, se hace referencia al dominio perfecto, pues el art.~1941 lo presume perfecto.

\section{Caracteres del dominio perfecto}

El dominio perfecto es absoluto, exclusivo y perpetuo.

\subsection{Absoluto}

Otorga todas las facultades (usar, gozar y disponer, material y jurídicamente) reunidas en una misma persona. No significa que el titular pueda hacer «lo que quiere» sin restricción: el ejercicio no debe ser abusivo y está limitado por normas de orden público, en interés de la comunidad o de la vecindad (capítulo~\ref{cap:limites}).

\subsection{Exclusivo}

\textbf{Art.~1943:} el dominio es exclusivo y no puede tener más de un titular. Sin exclusividad no hay dominio: si dos personas tuvieran un derecho real sobre la misma cosa, sería \textbf{condominio}, un derecho rotundamente diferente (Parte~V). \textbf{Art.~1944 (facultad excluyente):} el propietario puede excluir a cualquiera del uso y goce de la cosa, remover los objetos puestos en ella y encerrar los inmuebles con muros, cercos o fosos.

\subsection{Perpetuo}

No tiene límites en el tiempo y subsiste con independencia de que el titular lo ejerza o no; se perpetúa incluso más allá de la vida humana, razón por la cual existe el sistema de sucesiones. Quien poseyera la cosa por el tiempo y con las calidades legales adquirirá el dominio por prescripción adquisitiva.

\section{Facultades del titular}

\begin{itemize}
    \item \textbf{Disponer jurídicamente:} incluye \textbf{enajenar} (transmitir) a título oneroso (venta, permuta) o gratuito (donación, dación en pago); «enajenar» no es sinónimo de vender.
    \item \textbf{Constituir derechos reales} sobre la cosa: usufructo, hipoteca, servidumbre.
    \item \textbf{Administrar:} por ejemplo, dar la cosa en comodato, gratuito u oneroso.
    \item \textbf{Abandonar:} si recae sobre inmuebles, debe hacerse por escritura pública.
    \item \textbf{Disponer materialmente:} construir, destruir (con los límites legales) y transformar.
    \item \textbf{Usar y gozar:} obtener los frutos, civiles (locaciones) o naturales (cosechas).
\end{itemize}

\begin{ejemplo}{el usufructo}
El titular de dominio tiene todas las facultades. Al constituir un usufructo, «saca» de lo que tiene el uso y el goce y los pone en el usufructuario, que será titular de un derecho real sin ser propietario, con la fortaleza de los derechos reales: orden público, publicidad, contenido patrimonial y protección legal. Tiene todo lo que tiene el titular de dominio, menos la propiedad.
\end{ejemplo}

\section{Extensión (art.~1945)}

El dominio de una cosa comprende los objetos que forman un todo con ella o son sus accesorios. El dominio de un inmueble se extiende al subsuelo y al espacio aéreo, y se presume que pertenecen al dueño todas las construcciones, plantaciones y forestaciones existentes (salvo propiedad horizontal o derecho de superficie). Como decían los romanos, el dueño lo es «desde el cielo hasta el infierno».

\begin{ejemplo}{accesorios y construcciones}
Quien compra un inmueble adquiere las cañerías de agua y gas incrustadas en sus paredes; quien compra un vehículo adquiere las ruedas, aunque separadas sean cosas muebles. Todo lo construido en el terreno pertenece al dueño del terreno.
\end{ejemplo}

\section{Clasificación (art.~1946)}

\begin{table}[H]
\centering
\small
\begin{tabularx}{\textwidth}{>{\raggedright\arraybackslash}p{0.17\textwidth}>{\raggedright\arraybackslash}X>{\raggedright\arraybackslash}p{0.15\textwidth}>{\raggedright\arraybackslash}p{0.18\textwidth}}
\toprule
\textbf{Tipo} & \textbf{Característica} & \textbf{Afecta el carácter} & \textbf{Normativa} \\
\midrule
Perfecto & Sin gravamen, sin plazo, sin condición & Ninguno & Art.~1941 \\
Revocable & Sometido a plazo o condición resolutoria (máximo diez años) & Perpetuo & Arts.~1946 y 1965 \\
Fiduciario & Ejercido durante un plazo para luego transmitir & Perpetuo & Arts.~1701 y ss. \\
Desmembrado & Gravado con un derecho real a favor de un tercero & Absoluto & Normas propias de cada derecho \\
\bottomrule
\end{tabularx}
\caption{Tipos de dominio}
\end{table}

El art.~1946 establece que los tipos de dominio imperfecto son solo tres y que son de orden público.

\begin{itemize}
    \item \textbf{Revocable (art.~1965):} sometido a condición o plazo resolutorios, con máximo de diez años. Cumplida la condición o vencido el plazo, el dueño debe transmitir la cosa a quien se la transmitió originariamente. Si a los diez años no se revoca, queda convertido en dominio perfecto.
    \item \textbf{Fiduciario (arts.~1701 y ss.):} regulado en el Libro Tercero junto con el contrato de fideicomiso; lo ejerce el adquirente durante un tiempo, al cabo del cual debe transmitir la propiedad. El \textbf{fideicomiso} es el contrato; el \textbf{dominio fiduciario} es el derecho real.
    \item \textbf{Desmembrado:} gravado con un derecho real a favor de un tercero (usufructo, hipoteca); se rige por las normas de cada derecho constituido y afecta el carácter absoluto.
\end{itemize}

\cornell{concepto y caracteres del dominio}{
\cq{¿Dominio y propiedad son sinónimos?}{No. La propiedad en sentido amplio (CN, arts.~14 y 17) abarca derechos personales y reales; en sentido restringido es el dominio (art.~1941).}
\cq{Defina el dominio perfecto (art.~1941).}{Derecho real por el que una cosa está sometida a la voluntad y acción de una persona, con facultades de usar, gozar y disponer material y jurídicamente, dentro de los límites legales; se presume perfecto.}
\cq{¿Cuáles son sus caracteres y qué significa cada uno?}{Absoluto (reúne todas las facultades, sin ser ilimitado), exclusivo (un solo titular; si no, es condominio; art.~1943) y perpetuo (sin límite temporal).}
\cq{¿Qué comprende la extensión del dominio (art.~1945)?}{Accesorios, subsuelo, espacio aéreo y construcciones, plantaciones y forestaciones, salvo propiedad horizontal o superficie.}
\cq{Distinga dominio revocable, fiduciario y desmembrado.}{Revocable: plazo o condición resolutoria de hasta diez años. Fiduciario: se ejerce por un plazo para luego transmitir. Desmembrado: gravado con derecho real de un tercero.}
}{El dominio es el derecho real más extenso, que reúne las facultades de usar, gozar y disponer; es absoluto, exclusivo y perpetuo, se extiende a accesorios, subsuelo, espacio aéreo y lo construido, y se clasifica en perfecto y tres tipos de imperfecto.}

\chapter[Adquisición del dominio]{Adquisición del dominio}\label{cap:adq-dominio}
% ------------------------------------------------------------

\section{Modos de adquisición}

Los modos de adquirir el dominio son los hechos y actos jurídicos que la ley prevé para adjudicar en propiedad una cosa mueble o inmueble a una persona. Los hechos ilícitos (hurto, robo) no son modos reconocidos. Los artículos 1947 a 1963 regulan los distintos modos, con casuística específica: animales de caza, enjambres, pesca, acrecimiento de inmuebles por hechos de la naturaleza (aluvión, avulsión), plantaciones, etc.

\begin{table}[H]
\centering
\small
\begin{tabularx}{\textwidth}{>{\raggedright\arraybackslash}p{0.15\textwidth}>{\raggedright\arraybackslash}p{0.16\textwidth}>{\raggedright\arraybackslash}X>{\raggedright\arraybackslash}p{0.2\textwidth}}
\toprule
\textbf{Criterio} & \textbf{Subtipo} & \textbf{Descripción} & \textbf{Ejemplo} \\
\midrule
Por el acto & Entre vivos & Transmisión en vida de ambas partes & Compraventa, donación \\
Por el acto & Mortis causa & Transmisión por fallecimiento & Herencia \\
Por el objeto & Título universal & Se adquieren todos los derechos del causante & Herencia universal \\
Por el objeto & Título singular & Se adquiere un bien determinado & Legado de cosa cierta \\
Por la causa & Originario & La ley permite tomar la cosa por sí mismo & Apropiación, prescripción \\
Por la causa & Derivado & Se transmite de un titular anterior & Tradición (título y modo) \\
\bottomrule
\end{tabularx}
\caption{Clasificación de los modos de adquisición del dominio}
\end{table}

\textbf{Apropiación (art.~1947).} Son susceptibles de apropiación las cosas muebles sin dueño o abandonadas por su dueño. Los animales domésticos y domesticados no son apropiables.

\textbf{Modos derivados.} Remiten a la teoría del título y modo (capítulo~\ref{cap:adquisicion}); el modo derivado por excelencia es la tradición, y uno sin el otro elemento no basta. Si un poseedor no titular transmite la posesión, solo transmite el hecho de la posesión; si lo hace el propietario con capacidad y título suficiente, transmite el derecho de poseer.

\section{Protección del subadquirente de buena fe y a título oneroso}

Este tema refleja la fortaleza del derecho real: ante varias personas con derechos sobre una misma cosa, determina a priori hacia quién se inclina la protección de la ley.

\subsection{Regla y excepciones}

\textbf{Art.~399:} nadie puede transmitir a otro un derecho mejor ni más extenso que el que tiene, salvo las excepciones expresamente establecidas. La expresión «salvo» abre la puerta a excepciones; una es la \textbf{convalidación} (art.~1885), que en este desarrollo se describe así: si alguien transmite una cosa de la que no era propietario y luego la adquiere antes de que se declare la nulidad del acto, la transmisión se convalida. % TODO: verificar consistencia entre fuentes originales (esta descripción del art.~1885, referida a muebles no registrables y a la nulidad declarada, difiere de la redacción transcripta en el capítulo sobre legalidad).

\subsection{El subadquirente de inmuebles (art.~392)}

\begin{regla}{art.~392 CCyC}
Todos los derechos reales o personales transmitidos a terceros sobre un inmueble o muebles registrables por una persona que ha resultado adquirente en virtud de un acto nulo quedan sin ningún valor, pudiendo serle opuestos por quienes tengan interés legítimo, excepto respecto de los subadquirentes de buena fe y a título oneroso. Los subadquirentes no pueden ampararse en su buena fe si el acto se ha realizado sin intervención del titular del derecho.
\end{regla}

\begin{ejemplo}{Diego, Miriam y Juan}
Diego (enajenante) vende a Miriam (adquirente). Miriam vende luego a Juan (subadquirente). Después se declara nulo el acto entre Diego y Miriam (por ejemplo, por incapacidad de Diego). ¿Qué sucede con los derechos de Juan?
\end{ejemplo}

\begin{table}[H]
\centering
\small
\begin{tabularx}{\textwidth}{>{\raggedright\arraybackslash}p{0.2\textwidth}>{\raggedright\arraybackslash}X>{\raggedright\arraybackslash}X}
\toprule
\textbf{Seguridad} & \textbf{A quién protege} & \textbf{En el ejemplo} \\
\midrule
Estática & Al propietario original, quien tiene el derecho real & Protege a Diego: el inmueble volvería a él \\
Dinámica & Al adquirente en el tráfico, quien confió en la apariencia & Protege a Juan: consolida su adquisición \\
Opción del art.~392 & Seguridad dinámica, para el subadquirente de buena fe y título oneroso & Juan hace suyo el inmueble \\
\bottomrule
\end{tabularx}
\caption{Seguridad estática y dinámica en los inmuebles}
\end{table}

Para que Juan quede protegido se requiere: haber adquirido de \textbf{buena fe}, a \textbf{título oneroso}, y que el enajenante (Diego) haya \textbf{participado} de alguna manera en el acto después declarado nulo. En ningún caso hay protección para Miriam, el adquirente intermedio, que siempre sufre los efectos de la nulidad.

\begin{importante}
\textbf{Excepción a la excepción:} si hubo sustitución de personalidad (alguien se hizo pasar por Diego), la protección \textbf{no} alcanza a Juan, aunque haya actuado de buena fe y a título oneroso, porque el enajenante debe haber participado en el acto nulo.
\end{importante}

\begin{nota}
En el ejercicio profesional se actúa sabiendo que, en infinidad de situaciones, no es posible atender la justicia de cada uno. Ante la tensión entre derechos todos atendibles, la norma se ve forzada a decidir a priori sobre quién recae la protección, con la mirada puesta en el orden público y el interés de la comunidad. La misma elección en materia de cosas muebles se estudió en el capítulo~\ref{cap:muebles}.
\end{nota}

\cornell{adquisición del dominio}{
\cq{¿Qué son los modos de adquisición y cuáles no se reconocen?}{Hechos y actos jurídicos previstos por la ley para adjudicar una cosa en propiedad; los hechos ilícitos (hurto, robo) no lo son.}
\cq{Clasifique los modos de adquisición.}{Por el acto (entre vivos o mortis causa), por el objeto (título universal o singular) y por la causa (originario o derivado).}
\cq{¿Qué es la apropiación (art.~1947)?}{Modo originario sobre cosas muebles sin dueño o abandonadas; los animales domésticos y domesticados no son apropiables.}
\cq{En el caso Diego–Miriam–Juan, ¿qué requisitos protegen a Juan (art.~392)?}{Buena fe, título oneroso y participación del enajenante en el acto luego anulado. No se protege a Juan si hubo sustitución de personalidad; Miriam nunca está protegida.}
\cq{¿Qué tipo de seguridad privilegia el art.~392?}{La dinámica: la del subadquirente que confió en la apariencia, en detrimento del propietario original.}
}{El dominio se adquiere por modos originarios (apropiación, prescripción) o derivados (título y modo). El subadquirente de inmuebles de buena fe y a título oneroso queda protegido frente a la nulidad del acto anterior si el titular participó en él, con prioridad de la seguridad dinámica.}

\chapter[Límites al dominio]{Límites al dominio}\label{cap:limites}
% ------------------------------------------------------------

\section{¿Es ilimitado el dominio?}

Si el dominio es el derecho más extenso, cabe preguntarse si es ilimitado. El ejercicio del dominio, en Argentina y en la gran mayoría de los sistemas occidentales, reconoce infinidad de límites, tanto en interés de los vecinos como de la comunidad. Esos límites se han ido ampliando y continúan modelándose en armonía con los derechos humanos, nuevo eje rector de coherencia del ordenamiento nacional.

\section{Evolución histórica}

El art.~2513 del Código de Vélez Sarsfield otorgaba al titular la facultad de desnaturalizar la cosa y degradarla; sin embargo, en una nota al art.~2506, el propio Vélez citaba una ley de Partidas según la cual el dominio debe ejercerse «según Dios y según fuero». La facultad de destruir la cosa fue suprimida por la Ley~17.711, que incorporó la noción de \textbf{abuso del derecho}.

\begin{table}[H]
\centering
\small
\begin{tabularx}{\textwidth}{>{\raggedright\arraybackslash}p{0.22\textwidth}>{\raggedright\arraybackslash}X>{\raggedright\arraybackslash}X}
\toprule
\textbf{Momento} & \textbf{Regulación} & \textbf{Criterio de abuso} \\
\midrule
Ley~17.711 & Introdujo la noción de abuso del derecho & Cuando no se tenían en cuenta las razones del legislador al dictar la norma \\
CCyC (art.~10) & Modela todo el ordenamiento & Cuando el ejercicio contraría los principios en que se asienta el sistema normativo \\
\bottomrule
\end{tabularx}
\caption{Evolución del abuso del derecho}
\end{table}

\section{La función social de la propiedad}

El art.~21 del Pacto de San José de Costa Rica reconoce el derecho de toda persona al uso y goce de sus bienes y la facultad de la ley de subordinarlos al interés social. El Pacto fue incorporado a la Constitución mediante el art.~75, inc.~22, y el CCyC lo refleja en su articulado.

La expresión «función social de la propiedad» fue suprimida del art.~15 tal como lo votó el Congreso, aunque figuraba en el proyecto original; se sostiene que la supresión no modifica ni aminora la función social, que rige en el ordenamiento nacional y forma parte expresamente de catorce constituciones provinciales. Al tiempo que limita el ejercicio absoluto del derecho del titular, también lo protege del ejercicio ilimitado o abusivo de los demás.

\begin{importante}
El dominio en Argentina es un derecho conformado e integrado por la función social de la propiedad. No hay contradicción entre ser propietario y tener límites.
\end{importante}

\section{Clasificación de los límites (art.~1970)}

Las limitaciones al dominio privado en interés público se rigen por el derecho administrativo; el uso del dominio sobre inmuebles debe ejercerse conforme a las normas administrativas de cada jurisdicción, y los límites del Capítulo sobre relaciones de vecindad rigen en subsidio.

\begin{table}[H]
\centering
\small
\begin{tabularx}{\textwidth}{>{\raggedright\arraybackslash}p{0.15\textwidth}>{\raggedright\arraybackslash}p{0.19\textwidth}>{\raggedright\arraybackslash}X>{\raggedright\arraybackslash}p{0.19\textwidth}}
\toprule
\textbf{Tipo} & \textbf{Fundamento} & \textbf{Regulación} & \textbf{Indemnización} \\
\midrule
Interés público & Beneficio de la comunidad & Derecho administrativo (ilimitadas en número y clase) & No, en principio \\
Interés de los vecinos & Convivencia y tutela de la vecindad & CCyC, arts.~1970 a 1982 & No, en principio, salvo daño efectivo \\
\bottomrule
\end{tabularx}
\caption{Tipos de restricciones al dominio}
\end{table}

Son ejemplos de interés público las normas de higiene, tranquilidad, horarios de actividades, velocidades máximas y mínimas y prohibición de estacionar.

\section{Restricciones en interés de los vecinos}

No son taxativas. El CCyC regula expresamente: cláusulas de inenajenabilidad (art.~1972), inmisiones inmateriales (art.~1973), camino de sirga (art.~1974), aguas (arts.~1975 a 1977), luces y vistas (arts.~1978 y 1979), árboles, arbustos y plantas (art.~1980) y la obligación de soportar instalaciones provisorias y el paso de personas en obras linderas. Siempre primarán las normas administrativas del lugar: el derecho privado cede ante el orden público.

\subsection{Cláusulas de inenajenabilidad (art.~1972)}

Son nulas si obligan al propietario a no enajenar a título oneroso o a no constituir derechos reales; son válidas solo cuando se obliga a no enajenar a persona determinada. Los donantes pueden prohibir enajenar por un máximo de diez años y los condóminos pueden obligarse a no dividir por un máximo de diez años. Si el dueño no desea enajenar, no enajena, pues el dominio es perpetuo; lo que la ley anula es el compromiso futuro, reflejo de la fortaleza del derecho real y de la seguridad jurídica.

\subsection{Inmisiones inmateriales (art.~1973)}

Las molestias que ocasionan el humo, calor, olores, luminosidad, ruidos, vibraciones o inmisiones similares por actividades en inmuebles vecinos no deben exceder la \textbf{normal tolerancia}, teniendo en cuenta las condiciones del lugar y aunque medie autorización administrativa. Los jueces pueden disponer la cesación o disminución de las molestias hasta lo aceptable, y/o la indemnización de los daños, ponderando especialmente el respeto debido al uso regular de la propiedad, la prioridad en el uso, el interés general y las exigencias de la producción.

\begin{table}[H]
\centering
\small
\begin{tabularx}{\textwidth}{>{\raggedright\arraybackslash}p{0.2\textwidth}>{\raggedright\arraybackslash}X>{\raggedright\arraybackslash}X}
\toprule
\textbf{Criterio} & \textbf{Descripción} & \textbf{Ejemplo} \\
\midrule
Uso regular de la propiedad & Cómo se usa la propiedad que genera la inmisión & Una fundición realiza un uso regular de su propiedad industrial \\
Prioridad en el uso & Quién estaba primero en el lugar & Si la escuela estaba antes, quien compra al lado sabe a qué se expone \\
Interés general & Si la actividad sirve a la comunidad & Una escuela o un hospital tienen un interés público que pesa en la ponderación \\
Exigencias de la producción & Si la actividad no puede realizarse sin la inmisión & Una fundición no puede funcionar sin generar calor \\
\bottomrule
\end{tabularx}
\caption{Criterios de la normal tolerancia}
\end{table}

\begin{ejemplo}{una fundición y una escuela}
Una fundición de hierro generaba calor que se percibía desde el inmueble contiguo. El juez ponderó la prioridad en el uso (la fundición estaba primero), el interés de la producción, las exigencias de la actividad y el uso regular de la propiedad; impuso medidas para amortiguar el calor (aislantes y materiales), pero no ordenó el cese de la actividad. Del mismo modo, quien adquiere un inmueble al lado de una escuela primaria no puede reclamar por el ruido del recreo: pesan la prioridad en el uso y el interés general, y los valores inmobiliarios reflejan la zona.
\end{ejemplo}

\subsection{Camino de sirga (art.~1974)}

El dueño de un inmueble colindante con las orillas de los cauces o situado en las márgenes de los canales debe dejar libre una franja de quince metros, a la que puede acceder para posibilitar su uso. Originalmente era mayor porque se usaba para jalar embarcaciones por la orilla. Es una restricción legal de interés público (facilitar la navegación) que no genera indemnización; el propietario no puede construir, plantar ni realizar obras que afecten la navegación o el tránsito de personas o animales, y quien se sienta perjudicado está legitimado para pedir la remoción. El Estado puede reglamentar el uso pero no construir allí: tendría que expropiar, fundado en una ley especial. Tampoco puede el titular obstaculizar el curso de las aguas ni modificar su velocidad o dirección (salvo obras defensivas), bajo pena de remoción e indemnización.

\subsection{Aguas (art.~1976)}

El titular del fundo inferior debe recibir las aguas, arena y piedras que descienden naturalmente del fundo superior y no puede realizar obras que alteren ese escurrimiento. Si lo hace, el titular del fundo superior puede pedir la remoción e incluso la indemnización.

\subsection{Luces y vistas (arts.~1978 y 1979)}

\begin{table}[H]
\centering
\small
\begin{tabularx}{\textwidth}{>{\raggedright\arraybackslash}p{0.22\textwidth}>{\raggedright\arraybackslash}p{0.2\textwidth}>{\raggedright\arraybackslash}X>{\raggedright\arraybackslash}p{0.1\textwidth}}
\toprule
\textbf{Tipo} & \textbf{Función} & \textbf{Distancias exigidas} & \textbf{Art.} \\
\midrule
Vistas frontales & Permiten ver el inmueble vecino & Mínimo tres metros del lindero & 1978 \\
Vistas laterales u oblicuas & Permiten ver en ángulo & Mínimo sesenta centímetros del lindero & 1978 \\
Luces & Solo permiten entrada de luz y aire, no la vista & Superiores a 1,80 m desde el piso, salvo fijas o no transparentes & 1979 \\
\bottomrule
\end{tabularx}
\caption{Luces y vistas}
\end{table}

\begin{ejemplo}{vista y elevación del vecino}
En principio no puede impedirse que el lindero ejerza su derecho a elevar su construcción aunque prive de vista y luz. Se menciona un caso en la Patagonia en que la elevación de un vecino privó directamente la vista del Nahuel Huapi.
\end{ejemplo}

\subsection{Árboles, arbustos y plantas (art.~1980)}

El Código introduce un criterio subjetivo: ya no se fijan distancias mínimas de plantación, sino que nadie puede tener árboles, arbustos u otras plantas que causen molestias que excedan la normal tolerancia. El vecino perjudicado puede exigir que sean retirados, cortados o podados, y puede cortar él mismo las \textbf{raíces} que penetren en su inmueble, pero \textbf{no} las ramas, que solo puede cortar el propietario o a su pedido. Por ejemplo, las raíces del gomero pueden levantar los pisos y veredas de la casa vecina. Los apuntes mencionan casos de extrema violencia, incluso homicidios, por conflictos de árboles limítrofes.

\section{Regla general final}

\begin{importante}
En todos los casos de límites al dominio primarán las normas administrativas del lugar (alturas de ventanas, luces, distancias de árboles, etc.), por preeminencia del orden público sobre el derecho privado. Los límites no son taxativos: el orden público es dinámico porque son dinámicos los intereses de la comunidad, y los límites pueden modificarse en la medida en que respondan a un interés de orden público.
\end{importante}

\cornell{límites al dominio}{
\cq{¿Por qué el dominio absoluto no es ilimitado?}{Porque su ejercicio no debe ser abusivo y reconoce límites de interés público y vecinal, que se amplían en armonía con los derechos humanos.}
\cq{¿Cómo evolucionó el abuso del derecho?}{Vélez admitía desnaturalizar la cosa; la Ley~17.711 incorporó el abuso del derecho (razones del legislador); el CCyC (art.~10) lo vincula con los principios del sistema.}
\cq{¿Qué lugar ocupa la función social de la propiedad?}{Integra el dominio: limita su ejercicio absoluto y lo protege del abuso de otros; no hay contradicción entre ser propietario y tener límites.}
\cq{Distinga límites de interés público y de vecindad (art.~1970).}{Los primeros se rigen por el derecho administrativo (ilimitados en número y clase); los segundos, por los arts.~1970 a 1982, no taxativos. Ambos, en principio, sin indemnización.}
\cq{¿Cómo se determina la normal tolerancia (art.~1973)?}{Ponderando uso regular de la propiedad, prioridad en el uso, interés general y exigencias de la producción.}
\cq{Compare luces y vistas, y el régimen de raíces y ramas.}{Vistas: 3~m (frontales) y 60~cm (laterales); luces: más de 1,80~m. El vecino puede cortar raíces, no ramas.}
}{El dominio absoluto está sometido a límites de interés público (derecho administrativo) y vecinal (arts.~1970 a 1982), fundados en la función social de la propiedad y en el abuso del derecho. Las normas administrativas primarán siempre y los límites no son taxativos.}

% ============================================================
% ============================================================


% ============================================================
\part{Condominio y medianería}
% ============================================================

\chapter[Concepto y caracteres del condominio]{Concepto y caracteres del condominio}\label{cap:condominio}
% ------------------------------------------------------------

El condominio es el primer derecho real en que aparece «el otro»: la pluralidad de titulares sobre una misma cosa. De esa circunstancia nacen sus caracteres, sus facultades restringidas y la abundancia de normas pensadas para el conflicto.

\section{El condominio en el sistema de derechos reales}

Es uno de los catorce derechos reales del art.~1887 y comparte los rasgos comunes: es un derecho real sobre cosa propia, absoluto y de contenido patrimonial, regulado por normas de orden público, oponible \textit{erga omnes} a partir de la publicidad (posesoria o registral), se ejerce por la posesión y se adquiere por título suficiente y modo.

\section{Relevancia práctica}

Es uno de los derechos reales de mayor aplicación cotidiana: aparece en herencias donde varios herederos quedan como titulares de un bien, en parejas que compran un inmueble o un vehículo, en familias que comparten campos y en empresas que cotitularizan activos. Permite asesorar a los condóminos antes del conflicto, redactar convenios de uso y goce, iniciar o defender acciones de partición, resolver conflictos sobre uso exclusivo y valor locativo y proteger a los acreedores que quieran ejecutar una parte indivisa. La vida humana genera condominio por herencia, por afecto o por negocio.

\section{Un derecho autónomo y esencialmente inestable}

El condominio \textbf{no es un dominio «de dos»}: es un derecho autónomo con normas propias. La diferencia fundamental es la \textbf{exclusividad}: en el dominio el titular es único dueño; en el condominio aparece el otro y, con él, la posibilidad del conflicto, por lo que su regulación está plagada de normas para cuando el conflicto surge.

\begin{regla}{Vélez Sársfield}
«El condominio es esencialmente inestable.»
\end{regla}

Las relaciones humanas son inestables y la finitud es propia de los seres humanos. Quienes adquieren en condominio se miran a los ojos al hacerlo, pero con el tiempo pueden no tener el mismo rostro, la misma historia, los mismos afectos ni los mismos intereses, y por la muerte de alguno sus herederos ocuparán su lugar. Para fines de inversión existen figuras más adecuadas, como las sociedades o los fideicomisos.

\begin{ejemplo}{Juan y Carolina}
Juan y Carolina, pareja con un proyecto de vida en común, compraron un departamento en la Ciudad Autónoma de Buenos Aires, donde residen, y luego un automóvil que usan de manera compartida. Este caso se retoma en los capítulos siguientes para analizar los conflictos posibles.
\end{ejemplo}

\section{Concepto (art.~1983)}

\begin{definicion}{condominio (art.~1983 CCyC)}
Condominio es el derecho real de propiedad sobre una cosa que pertenece en común a varias personas y que corresponde a cada una por una parte indivisa.
\end{definicion}

Las partes de los condóminos se \textbf{presumen iguales}, salvo que la ley o el título dispongan una proporción diferente: si uno aportó el 70\,\% y el otro el 30\,\%, ese será el porcentaje de cada uno; si el título nada dice, 50\,\% y 50\,\%.

\section{Caracteres}

\begin{table}[H]
\centering
\small
\begin{tabularx}{\textwidth}{>{\raggedright\arraybackslash}p{0.2\textwidth}>{\raggedright\arraybackslash}X>{\raggedright\arraybackslash}X}
\toprule
\textbf{Carácter} & \textbf{Descripción} & \textbf{Ejemplo (Juan y Carolina)} \\
\midrule
Unidad de objeto & Recae sobre una sola cosa, mueble o inmueble & El auto o el departamento es una sola cosa \\
Pluralidad de sujetos & Dos o más titulares & Juan y Carolina (o sus herederos) \\
Partes ideales (indivisas) & Ninguno tiene una parte material determinada & Ninguno es dueño de las ruedas delanteras ni de las traseras \\
\bottomrule
\end{tabularx}
\caption{Caracteres del condominio}
\end{table}

Si los condóminos no quisieran continuar juntos, no podría resolverse que cada uno se lleve dos ruedas: el derecho no recae sobre las ruedas, ni sobre el baño o la cocina del departamento, sino sobre una \textbf{parte indivisa}. Cada uno puede decir que es su auto, pero lo es con otro.

\begin{ejemplo}{fallecimiento de un condómino}
Si Carolina fallece y deja tres hijos, el condominio pasa a ser de cuatro personas: Juan, que conserva su mitad, y los tres herederos de Carolina, que comparten la otra mitad en partes iguales por sucesión \textit{mortis causa}.
\end{ejemplo}

\cornell{concepto y caracteres del condominio}{
\cq{¿Por qué el condominio no es un «dominio de dos»?}{Porque es un derecho autónomo: falta la exclusividad. El otro condómino introduce el conflicto, por lo que la regulación se orienta a resolverlo.}
\cq{¿Qué significa que sea «esencialmente inestable»?}{Que las relaciones humanas y la vida de los titulares cambian (conflictos, muerte, herederos), lo que pone en riesgo la comunidad.}
\cq{Defina el condominio (art.~1983) y la presunción de partes.}{Derecho real de propiedad sobre una cosa común a varias personas, por parte indivisa. Las partes se presumen iguales salvo proporción distinta en la ley o el título.}
\cq{¿Cuáles son sus caracteres?}{Unidad de objeto, pluralidad de sujetos y partes ideales o indivisas, no materiales.}
\cq{¿Qué ocurre si fallece un condómino?}{Sus herederos ocupan su lugar y comparten su parte indivisa.}
}{El condominio es un derecho real autónomo, de unidad de objeto, pluralidad de sujetos y partes indivisas presumidas iguales, esencialmente inestable por la presencia del otro titular y por la finitud humana.}

\chapter[Facultades, gastos y administración del condominio]{Facultades, gastos y administración del condominio}\label{cap:facultades-cond}
% ------------------------------------------------------------

Las facultades de los condóminos dependen de aquello sobre lo que recaen: son \textbf{amplísimas sobre la parte indivisa} y \textbf{restringidas sobre la cosa en común}, precisamente porque está el otro.

\section{Facultades sobre la parte indivisa (art.~1989)}

\begin{regla}{art.~1989 CCyC}
El condómino puede enajenar y gravar su parte indivisa sin el consentimiento de los restantes condóminos. Los acreedores pueden embargarla y ejecutarla sin esperar el resultado de la partición, que les es inoponible.
\end{regla}

El condómino puede \textbf{enajenar} (vender, donar o transferir), \textbf{gravar} (por ejemplo, hipotecar) su parte, y sus acreedores pueden embargarla y llevarla a subasta.

\section{Facultades sobre la cosa en común}

El condómino puede \textbf{usar y gozar} la cosa, pero de forma \textbf{conjunta} con los demás y \textbf{sin alterar su destino}; puede usarla en forma exclusiva siempre que no excluya al otro ni cambie el destino.

\subsection{Convenio de uso y goce exclusivo (art.~1987)}

\begin{regla}{art.~1987 CCyC}
Los condóminos pueden convenir el uso y goce alternado de la cosa común o que se ejercite de manera exclusiva y excluyente sobre determinadas partes materiales.
\end{regla}

Es una de las novedades más importantes del Código de 2014: permite pacificar una cuestión interna que a veces provocaba conflictos serios, sobre todo tras el fallecimiento de alguno (personas que durante veinte o treinta años convivieron sin dificultad y, al morir uno, el otro se dice tan dueño de la casa de atrás como de la de adelante).

\begin{ejemplo}{dos casas sobre un terreno}
Sobre un terreno en condominio hay una casa adelante y otra atrás. Registralmente es un solo inmueble, pero los condóminos pueden convenir que uno use exclusivamente la de adelante y el otro la de atrás. Sin convenio, ambos son titulares de una parte indivisa de todo: un poco del terreno, un poco de cada casa.
\end{ejemplo}

\subsection{Uso exclusivo no pactado y valor locativo (art.~1988)}

\begin{regla}{art.~1988 CCyC}
El condómino que usa la cosa común en forma exclusiva está obligado, excepto pacto en contrario, a satisfacer una indemnización, desde que el otro condómino lo solicita.
\end{regla}

Si Juan y Carolina se separan, ella se va y Juan queda viviendo solo en el departamento, ¿puede ella reclamar alquileres por todos esos años? El \textbf{valor locativo} es una compensación como si se pagara un alquiler, pero \textbf{no es un alquiler}: el condómino usa la cosa ejerciendo su facultad de usar y gozar, y si nada se le pide se presume que lo hace con el consentimiento del otro. Solo procede \textbf{desde que el otro condómino lo notificó fehacientemente}: si Carolina notificó a los tres días, procede desde entonces; si notificó dos o tres años después, solo desde esa fecha. No puede reclamar compensaciones anteriores a la notificación.

\section{Disposición, gastos y deudas}

Sobre la cosa en su totalidad, ningún condómino puede disponer materialmente de una parte determinada sin autorización de los demás, ni realizar actos de disposición jurídica, como alquilarla, sin el consentimiento de la mayoría. Si son tres condóminos y dos dicen que sí y uno que no, este no puede oponerse.

Los \textbf{gastos comunes} se afrontan en proporción al porcentaje de participación (si Juan tiene el 70\,\% pagará el 70\,\%; si el título nada dice, por mitades). Si un condómino contrata con un tercero (por ejemplo, un servicio de pintura), él es el obligado ante el tercero, pero puede recuperar internamente de los demás su parte; esto genera debates sobre si los gastos fueron de conservación o de mejora, necesarios o no.

\section{Administración y asamblea (arts.~1993 y 1994)}

\begin{regla}{art.~1994 CCyC}
Cualquier condómino puede convocar a una asamblea de condóminos, citando a los demás en forma fehaciente, indicando la finalidad y con antelación suficiente. La resolución de la mayoría absoluta de los condóminos computada según el valor de las partes indivisas es obligatoria para todos.
\end{regla}

Presenta dificultades prácticas en condominios de solo dos titulares, donde no resulta claro cómo convocar ni dónde alcanzar acuerdo; si hay empate, el art.~1994 establece que \textbf{decide la suerte}. El régimen es útil en condominios con varios titulares, pero no es de aplicación amplia.

\cornell{facultades y administración del condominio}{
\cq{¿Qué puede hacer el condómino con su parte indivisa (art.~1989)?}{Enajenarla y gravarla sin consentimiento de los demás; sus acreedores pueden embargarla y ejecutarla sin esperar la partición.}
\cq{¿Y sobre la cosa en común?}{Usar y gozar conjuntamente y sin alterar el destino; no disponer materialmente de una parte ni alquilar sin consentimiento de la mayoría.}
\cq{¿Qué permite el art.~1987?}{Convenir uso alternado o exclusivo y excluyente de partes materiales, pacificando conflictos internos.}
\cq{¿Cuándo procede el valor locativo (art.~1988) y qué naturaleza tiene?}{Desde la notificación fehaciente del condómino que lo reclama; no es un alquiler, sino una compensación por uso exclusivo. Antes se presume consentimiento.}
\cq{¿Cómo se reparten los gastos y cómo decide la asamblea (art.~1994)?}{Gastos en proporción a la parte; la asamblea decide por mayoría absoluta según el valor de las partes, y en caso de empate, la suerte.}
}{Las facultades del condómino son amplísimas sobre su parte indivisa y restringidas sobre la cosa común. El uso y goce puede pactarse; el uso exclusivo no pactado da lugar a valor locativo desde la notificación; los gastos se reparten según las partes y la asamblea decide por mayoría computada por valor.}

\chapter[Partición e indivisión en el condominio]{Partición e indivisión en el condominio}\label{cap:particion}
% ------------------------------------------------------------

\section{Clasificación}

\begin{table}[H]
\centering
\small
\begin{tabularx}{\textwidth}{>{\raggedright\arraybackslash}p{0.18\textwidth}>{\raggedright\arraybackslash}X>{\raggedright\arraybackslash}X>{\raggedright\arraybackslash}X}
\toprule
\textbf{Tipo} & \textbf{Definición} & \textbf{Consecuencia} & \textbf{Plazo máximo} \\
\midrule
Sin indivisión forzosa & No hay causa legal que impida dividir & Cualquier condómino puede pedir la partición en cualquier momento & Sin plazo: acción imprescriptible \\
Con indivisión forzosa & Hay una causa legal que restringe la partición & Partición temporalmente impedida según la fuente & Diez años (convención y testamento); cinco años (sentencia) \\
\bottomrule
\end{tabularx}
\caption{Tipos de condominio según la partición}
\end{table}

La indivisión forzosa surge únicamente de la ley, aunque con distintas fuentes. Se distingue la \textbf{indivisión temporaria}, de plazo determinado, y la \textbf{perdurable}, vinculada a la naturaleza del objeto, como en la medianería (capítulo~\ref{cap:medianeria}).

\section{Condominio sin indivisión forzosa}

\begin{regla}{arts.~1996 y 1999 CCyC}
Todo condómino puede, en cualquier tiempo, pedir la partición de la cosa. La acción es imprescriptible (art.~1996). El derecho a pedir la partición es irrenunciable por tiempo indeterminado (art.~1999).
\end{regla}

Es imprescriptible, al igual que el dominio es perpetuo, e irrenunciable por tiempo indeterminado: no puede pactarse definitivamente que nunca se pedirá.

\begin{ejemplo}{la cláusula prohibitiva de la partición}
Una pareja que compra «el departamento de sus sueños» pide incluir una condición: que nunca puedan dividirlo. Quien asesora no puede garantizarlo, porque el derecho a pedir la partición es irrenunciable por tiempo indeterminado y, si Carolina muere, es dudoso para qué querrían sus herederos tener el departamento con Juan.
\end{ejemplo}

\section{Condominio con indivisión forzosa temporaria: fuentes}

\begin{table}[H]
\centering
\small
\begin{tabularx}{\textwidth}{>{\raggedright\arraybackslash}p{0.18\textwidth}>{\raggedright\arraybackslash}p{0.12\textwidth}>{\raggedright\arraybackslash}X>{\raggedright\arraybackslash}p{0.2\textwidth}}
\toprule
\textbf{Fuente} & \textbf{Art.} & \textbf{Regla} & \textbf{Plazo} \\
\midrule
Convención & 2000 & Los condóminos pueden convenir suspender la partición; si no fijan plazo, o es incierto o superior, se considera celebrada por diez años & Hasta diez años, renovable por igual término \\
Sentencia judicial & 2001 & A petición de parte, cuando la partición es nociva para el interés de los condóminos o de alguno, el juez puede disponerla con carácter diferido & Hasta cinco años, renovable judicialmente por otros cinco \\
Acto de última voluntad & 2330 & El testador puede imponer a los herederos, aun mayores de edad, la indivisión de la herencia o de algún bien & Hasta diez años \\
\bottomrule
\end{tabularx}
\caption{Fuentes de la indivisión temporaria}
\end{table}

\begin{itemize}
    \item \textbf{Convención.} Si Carolina muere, sus herederos deben cumplir el pacto, que también los obliga.
    \item \textbf{Sentencia.} Puede \textbf{demorar} una partición que no tenía restricción, si se demuestra perjudicial para todos, o \textbf{anticipar} una partición restringida, a pedido de quienes acrediten circunstancias graves.
    \item \textbf{Última voluntad.} Es muy frecuente en la práctica («les dejo todo a mis hijos, pero que nunca lo separen»), pero el plazo máximo es de diez años y no puede ser indefinido.
\end{itemize}

\section{Partición nociva y derechos humanos}

La partición nociva existía ya en el Código derogado. El ejemplo clásico es el campo sembrado: si uno pide la partición en plena cosecha, el otro puede oponerse judicialmente porque sería antieconómica y perjudicial para todos.

\begin{importante}
Ya existen fallos que demoran la partición con fundamento en los derechos humanos, los derechos del niño, de los ancianos y de los enfermos: algunos jueces la demoran cinco años y, tras nueva evaluación, otros cinco, ya no por perjuicio económico común sino por el interés del condómino anciano o enfermo que vive en el inmueble.
\end{importante}

Esto no significa que la partición no pueda pedirse: es un derecho rotundo. Ante un caso concreto y cierta prueba, el juez puede demorarla o anticiparla.

\cornell{partición e indivisión}{
\cq{¿Qué caracteres tiene el derecho a pedir la partición (arts.~1996 y 1999)?}{Es imprescriptible e irrenunciable por tiempo indeterminado.}
\cq{¿Puede pactarse que nunca se divida?}{No: solo se puede suspender la partición por hasta diez años (art.~2000), renovables.}
\cq{Enumere las fuentes de indivisión temporaria con sus plazos.}{Convención (hasta diez años, art.~2000), sentencia (hasta cinco años renovables por otros cinco, art.~2001) y testamento (hasta diez años, art.~2330).}
\cq{¿Qué es una partición nociva?}{La que perjudica el interés de los condóminos o de alguno; ejemplo: pedirla en plena cosecha.}
\cq{¿Qué rol cumplen hoy los derechos humanos?}{Fundan fallos que demoran la partición en interés de condóminos ancianos, enfermos o niños, sin negar el derecho a pedirla.}
}{La partición puede pedirse siempre salvo indivisión forzosa, que solo la ley establece, con fuentes convencional, judicial o testamentaria de plazo máximo diez o cinco años. La partición nociva y los derechos humanos permiten demorarla o anticiparla según el caso.}

\chapter[Medianería: estructura, fundamentos y clasificación de los muros]{Medianería: estructura, fundamentos y clasificación de los muros}\label{cap:medianeria}
% ------------------------------------------------------------

La medianería es la regulación legal de muros, cercos y fosos dentro del condominio con indivisión forzosa de origen legal. Para quien ejerce en la práctica inmobiliaria, notarial o registral, comprender este régimen permite asesorar a propietarios sobre derechos y obligaciones en construcciones colindantes, resolver conflictos sobre muros y cercos, redactar escrituras que reflejen la verdadera naturaleza del inmueble, orientar trámites registrales y evitar litigios costosos. Todo acto de compraventa, construcción o modificación de un inmueble toca este régimen. Cada muro es una especie de acuerdo silencioso entre vecinos: ambos necesitan encerrarse (derecho) y ambos deben contribuir (obligación). El estudio avanza de lo más simple (qué tipos de muros hay) a lo más complejo (quién paga qué y cuándo).

\section{Estructura del condominio con indivisión forzosa}

Es un condominio de \textbf{origen legal} con dos facetas, reguladas en el Título Cuarto del Libro Cuarto.

\begin{definicion}{condominio sobre accesorios indispensables (arts.~2004 y 2005)}
Condominio sobre accesorios indispensables para el uso común de dos inmuebles colindantes, como pasillos de acceso y pasadizos comunes. Cada condómino puede usarlos conforme a su destino, sin afectar al otro, y no puede pedir la partición.
\end{definicion}

\begin{definicion}{condominio sobre muros, cercos y fosos (arts.~2006 a 2030)}
Condominio sobre muros de cerramiento entre inmuebles colindantes. Es la faceta compleja, por su vinculación con la medianería.
\end{definicion}

\section{Un límite al derecho de propiedad}

La medianería otorga a cada propietario de inmuebles colindantes urbanos el \textbf{derecho y la obligación recíprocos} de construir un muro de cerramiento. En este régimen se advierte la potencia del derecho real: el muro \textbf{no requiere la conformidad del vecino}; quien construye puede hacerlo aunque el vecino se oponga o no contribuya, y la ley autoriza a edificar \textbf{la mitad del muro sobre el terreno del vecino}. Quien construye adelanta el cumplimiento de una obligación legal que también pesaba sobre el vecino y por eso tiene derecho a reclamar. Con las normativas actuales, los muros de cerramiento miden habitualmente \textbf{tres metros de alto por treinta centímetros de espesor} («pared de treinta»).

\section{Fundamentos}

\begin{itemize}
    \item \textbf{Constitucionales.} CN, art.~18 (inviolabilidad del domicilio) y art.~19 («Las acciones privadas de los hombres que de ningún modo ofendan al orden y a la moral pública, ni perjudiquen a un tercero, están solo reservadas a Dios, y exentas de la autoridad de los magistrados»): protegen la privacidad e intimidad personal y familiar, que dependen de la posibilidad de encerrarse.
    \item \textbf{Facultad de exclusión y encerramiento (art.~1944).} Si cada propietario tiene esa facultad, su vecino también; la faceta obligacional surge de la normativa del condominio con indivisión forzosa.
    \item \textbf{Económicos.} Permite economizar el suelo, escaso y no siempre cultivable; el ahorro se multiplica por la enorme cantidad de inmuebles urbanos.
    \item \textbf{Técnicos.} Solidez: los inmuebles se «abrazan», se apoyan unos en otros y se fortalecen mutuamente.
\end{itemize}

\section{Clasificación de los muros}

Siguiendo la estructura didáctica de Marina Mariani de Vidal, se formulan dos preguntas: \textbf{dónde está} el muro (criterio físico) y \textbf{de quién es} (criterio jurídico).

\subsection{Clasificación física}

\begin{table}[H]
\centering
\small
\begin{tabularx}{\textwidth}{>{\raggedright\arraybackslash}p{0.2\textwidth}>{\raggedright\arraybackslash}p{0.12\textwidth}>{\raggedright\arraybackslash}X}
\toprule
\textbf{Muro} & \textbf{Art.~2006} & \textbf{Definición} \\
\midrule
Lindero, separativo o divisorio & inc.~a & Sinónimos: demarca el inmueble y lo delimita del colindante. \\
Incaballado & inc.~b & Lindero que se asienta parcialmente en cada inmueble; lo típico es quince centímetros en cada terreno («pared de 30»: 15 y 15). \\
Contiguo & inc.~c & Lindero que se asienta totalmente en uno de los inmuebles, de modo que el filo coincide con el límite separativo. \\
De elevación & inc.~d & Lindero que excede la altura del muro de cerramiento ordinario (tres metros), sea incaballado o contiguo, en la altura que supera esa medida. \\
Enterrado & inc.~h & Ubicado debajo del nivel del suelo, sin servir de cimiento a una construcción en superficie. \\
\bottomrule
\end{tabularx}
\caption{Clasificación física de los muros}
\end{table}

Por ejemplo, un muro de cinco metros es medianero hasta los tres y de elevación en los dos metros que superan esa altura.

\subsection{Clasificación jurídica: el muro medianero}

\begin{definicion}{muro medianero}
Es el lindero común, perteneciente en condominio a ambos colindantes. Surge desde la construcción.
\end{definicion}

Desde que uno o ambos vecinos construyen el muro, sea incaballado o contiguo, el muro \textbf{pertenece a los dos}, aunque no haya sido construido en el terreno de ambos; por eso la ley permite que Juan lo construya avanzando quince centímetros sobre el terreno de su vecina María. El condominio es de \textbf{fuente legal}: nace medianero desde la construcción, hasta los tres metros de altura.

\begin{table}[H]
\centering
\small
\begin{tabularx}{\textwidth}{>{\raggedright\arraybackslash}p{0.28\textwidth}>{\raggedright\arraybackslash}X}
\toprule
\textbf{Característica} & \textbf{Descripción} \\
\midrule
Origen legal & El condominio surge por ley desde la construcción. \\
Titularidad & Pertenece a ambos condóminos por igual (50\,\% y 50\,\%). \\
Derechos del no constructor & Usar el muro, anclar construcciones, reclamar contribución. \\
\bottomrule
\end{tabularx}
\caption{Características del muro medianero}
\end{table}

El muro puede ser construido por ambos vecinos, por uno solo, con comunidad de gastos o pagado por uno solo. % TODO: contenido incompleto o ambiguo en las notas originales (la transcripción del pasaje sobre la construcción compartida o unilateral es confusa).
Si lo pagó uno solo, el vecino que no lo construyó es igualmente titular, porque el muro es medianero desde la construcción. Como el muro construido por Juan pertenece a María en un 50\,\%, ella, aunque no haya contribuido, puede \textbf{anclar} construcciones al muro, \textbf{elevar} construcciones por sobre los tres metros, \textbf{prolongar} cimientos hacia abajo y \textbf{destruir y reconstruir} el muro si es necesario para su obra.

\cornell{medianería: estructura, fundamentos y muros}{
\cq{¿Cuáles son las dos facetas del condominio con indivisión forzosa?}{Accesorios indispensables (arts.~2004 y 2005, como pasillos) y muros, cercos y fosos (arts.~2006 a 2030).}
\cq{¿Por qué la medianería muestra la potencia del derecho real?}{Porque el muro puede construirse sin conformidad del vecino, incluso sobre la mitad de su terreno, al adelantar una obligación legal recíproca.}
\cq{¿Cuáles son sus fundamentos?}{Constitucionales (CN arts.~18 y 19: privacidad e intimidad), facultad de encerramiento (art.~1944), económicos (ahorro de suelo) y técnicos (solidez).}
\cq{Distinga muro incaballado, contiguo, de elevación y enterrado.}{Incaballado: asentado parcialmente en cada inmueble. Contiguo: totalmente en uno. De elevación: la altura que excede tres metros. Enterrado: bajo el nivel del suelo sin servir de cimiento.}
\cq{¿De quién es el muro y desde cuándo?}{Es medianero, condominio legal por mitades, desde que se construye (hasta tres metros); el vecino que no pagó es igualmente titular con derechos plenos.}
}{La medianería impone a los propietarios colindantes el derecho y la obligación recíprocos de cerramiento, fundada en la privacidad, la economía de suelo y la solidez. Los muros se clasifican por su ubicación y por su titularidad; el muro nace medianero desde su construcción.}

\chapter[Régimen del muro medianero, presunciones y medianería rural]{Régimen del muro medianero, presunciones y medianería rural}\label{cap:regimen-medianero}
% ------------------------------------------------------------

\section{Acción de cobro de medianería}

\begin{definicion}{acción de medianería (art.~2019)}
Desde la construcción del muro nace el derecho a reclamar el cobro de la medianería, haya utilizado o no el vecino el muro.
\end{definicion}

La mora es \textbf{automática} (art.~886): no requiere interpelación ni demanda previa. La acción prescribe a los \textbf{cinco años} desde la construcción (art.~2560); pasado ese plazo se extingue el derecho a reclamar, incluso si el vecino comienza a usar el muro después. La defensa es la prescripción liberatoria.

\begin{ejemplo}{terreno baldío y uso posterior}
Se construye y paga un muro; no vive nadie al lado, pasa el tiempo y nadie lo utiliza. A los seis años, el propietario comienza a construir y usa el muro: la acción de cobro ya prescribió. Puede usar los cimientos necesarios para tres metros y la pared de treinta, porque el muro es suyo en condominio.
\end{ejemplo}

\section{Muro de elevación y muro enterrado}

En un mismo muro conviven dos regímenes: hasta tres metros, \textbf{medianero}; desde tres metros hacia arriba, \textbf{privativo} del constructor. Para la parte que excede tres metros, la obligación nace cuando el vecino efectivamente la usa, apoyando en el segmento elevado. El muro enterrado sigue un régimen similar: la obligación de pagar nace desde la utilización específica (art.~2019).

\section{Imposibilidad de liberarse}

No es posible liberarse de la obligación de pagar, ni de los gastos de conservación o reconstrucción, mediante abdicación o declinación de la parte. Como el muro ya nació medianero, para liberarse habría que transmitir el 50\,\% al vecino, lo que requiere nuevos planos y escritura pública: un procedimiento costoso.

\section{Presunciones legales (arts.~2010 a 2013)}

Cuando hay construcciones antiguas y es difícil determinar la naturaleza del muro, la ley establece presunciones que admiten prueba en contrario, pero son fuertes.

\begin{definicion}{presunción de medianería (art.~2010)}
A menos que se pruebe lo contrario, un muro lindero entre dos edificios que supera los tres metros se presume medianero hasta esa altura, y la parte que los supera se presume privativa del dueño del edificio más alto.
\end{definicion}

No se aplican cuando el muro separa patios, huertos o jardines (y no un edificio), ni cuando hay documentos o signos que identifican la verdadera naturaleza o permiten destruir la presunción.

\section{Medianería rural}

En predios rurales la división no se hace con muros sino con \textbf{cercas vivas} (árboles alineados), \textbf{zanjas}, cercos de alambre u otras estructuras (arts.~2032 y 2033); el titular tiene la facultad de levantar o excavar el cerramiento, como se ve al costado de una ruta. Los \textbf{árboles y arbustos} plantados en la línea divisoria pueden ser medianeros, contiguos o incaballados, según su ubicación respecto de los muros, cercos y fosos linderos, tanto en predios rurales como urbanos (arts.~2034 a 2036). Si causan perjuicio a un condómino, puede exigir su remoción en cualquier tiempo, e incluso su reemplazo. % TODO: contenido incompleto o ambiguo en las notas originales (el apunte no precisa el alcance del reemplazo).

\section{Síntesis}

\begin{conceptoclave}{medianería}
La medianería es una institución del derecho real que limita el derecho de propiedad en favor de la vida urbana moderna. Nace por ley desde la construcción de un muro entre inmuebles colindantes. Sus fundamentos son constitucionales (arts.~18 y 19), económicos (ahorro de terreno) y técnicos (solidez). Ambos propietarios son condóminos con derechos y obligaciones recíprocos, con independencia de quién financió la obra. La ley presume medianería hasta tres metros y prevé regímenes especiales para los muros de elevación y enterrados; las zonas rurales tienen regímenes propios. La acción de cobro prescribe a los cinco años.
\end{conceptoclave}

Para resolver un caso conviene distinguir: (1) dónde está el muro (ubicación física); (2) de quién es (naturaleza jurídica); y (3) qué régimen le corresponde según su tipo.

\cornell{régimen del muro medianero}{
\cq{¿Cuándo nace la acción de cobro y cuándo prescribe?}{Nace desde la construcción, aunque el vecino no use el muro (art.~2019); la mora es automática y prescribe a los cinco años, sin renacer por uso posterior.}
\cq{¿Qué régimen tiene el muro de elevación?}{Medianero hasta tres metros y privativo del constructor por encima; la obligación sobre el excedente nace cuando el vecino lo usa.}
\cq{¿Puede un vecino liberarse declinando su parte?}{No: debería transmitir su 50\,\% por escritura pública y con nuevos planos.}
\cq{¿Qué presume el art.~2010 y cuándo no se aplica?}{Medianería hasta tres metros y privatividad del dueño del edificio más alto por encima. No se aplica entre patios, huertos o jardines ni con documentos o signos en contrario.}
\cq{¿Cómo se divide un predio rural?}{Con cercas vivas, zanjas o alambrados (arts.~2032 y 2033); los árboles divisorios pueden ser medianeros (arts.~2034 a 2036) y su remoción puede exigirse si perjudican.}
}{El cobro de la medianería nace con la construcción, con mora automática y prescripción a los cinco años. La ley presume medianería hasta tres metros, con regímenes diferenciados para elevación y muros enterrados. No es posible liberarse de las obligaciones sin transmitir la parte, y la medianería rural se ordena con cercas vivas, zanjas y alambrados.}

% ============================================================
% ============================================================


% ============================================================
\part{Propiedad horizontal}
% ============================================================

\chapter[Concepto y naturaleza de la propiedad horizontal]{Concepto y naturaleza de la propiedad horizontal}\label{cap:ph}
% ------------------------------------------------------------

La propiedad horizontal, regulada en los artículos 2037 a 2069, es uno de los regímenes más presentes en la práctica profesional: la mayoría de la población urbana vive o trabaja en inmuebles sometidos a él. Conocerlo permite asesorar en la compraventa de unidades funcionales (verificando deudas de expensas y reglamentos), redactar o revisar reglamentos y contratos de administración, iniciar juicios ejecutivos por expensas, defender o representar consorcistas en conflictos de asamblea, asesorar sobre prehorizontalidad y articularlo con el usufructo, el embargo, la subasta y la protección de la vivienda.

\section{Del dominio y el condominio a la propiedad horizontal}

El \textbf{dominio} tiene un titular exclusivo y le otorga un poder directo e inmediato sobre la cosa, con las facultades de usar, gozar y disponer (\textit{ius utendi, ius fruendi, ius abutendi}); sobre un inmueble, el dominio se extiende a todo lo plantado y construido.

\begin{ejemplo}{Juan y sus hijos}
Juan adquiere un terreno, lo construye y luego edifica sobre esa casa una planta alta con dos departamentos, con la ilusión de que cada uno sea usado por un hijo. Con el tiempo, decide enajenarles los departamentos. Materialmente se observan tres departamentos, pero jurídicamente existe \textbf{un solo inmueble}, el de Juan; si sus acreedores lo ejecutan, ejecutarán el inmueble con todo lo plantado y construido. Juan puede disponer de toda la cosa o de una parte indivisa: si dona a cada hijo el 25\,\%, nace un \textbf{condominio}, con unidad de objeto, pluralidad de sujetos y partes indivisas, sin que se haya transmitido un departamento determinado a cada hijo.
\end{ejemplo}

En el ámbito de los derechos reales, \textbf{enajenar} es transmitir el derecho de dominio, sea por compraventa, donación o dación en pago.

\begin{table}[H]
\centering
\small
\begin{tabularx}{\textwidth}{>{\raggedright\arraybackslash}X>{\raggedright\arraybackslash}X}
\toprule
\textbf{Dominio} & \textbf{Condominio} \\
\midrule
Titular único & Pluralidad de sujetos \\
Poder directo e inmediato & Partes indivisas o ideales \\
Se extiende a lo construido y plantado & No se corresponde con una parte material determinada \\
Un solo inmueble jurídico & Un solo inmueble jurídico (misma unidad de objeto) \\
\bottomrule
\end{tabularx}
\caption{Dominio y condominio}
\end{table}

\section{El derecho real de propiedad horizontal}

La cuestión es cómo cada departamento puede pertenecer en propiedad a cada titular. La propiedad horizontal permite constituir \textbf{unidades independientes dentro de un mismo inmueble}: un \textbf{dominio exclusivo sobre partes privativas} inescindiblemente unido a un \textbf{condominio sobre las partes comunes}. No es simplemente dominio más condominio: es un derecho único de estructura dual.

\begin{definicion}{propiedad horizontal (art.~2037 CCyC)}
La propiedad horizontal es el derecho real que se ejerce sobre un inmueble propio que otorga a su titular facultades de uso, goce y disposición material y jurídica que se ejercen sobre partes privativas y sobre partes comunes de un edificio, de conformidad con lo que establece este Título y el respectivo reglamento de propiedad horizontal.
\end{definicion}

Un mismo inmueble puede ser objeto de dominio, de condominio o estar afectado a propiedad horizontal.

\subsection{Origen}

Fue creada por la \textbf{Ley~13.512} en 1948 y cumplió una gran función social. Antes, el Código de Vélez Sársfield prohibía específicamente la división horizontal de los inmuebles. Su contenido fue incorporado y actualizado en los arts.~2037 a 2069 del CCyC (Ley~26.994, vigente desde agosto de 2015).

\subsection{Inmuebles afectables y trascendencia}

No solo pueden afectarse propiedades en altura, sino también las de una sola planta: con frente y fondo, «en tira» y dúplex. El sistema es, según los apuntes, una de las mayores creaciones en materia de derechos reales del siglo XX y superó largamente sus objetivos originales; a 70 años de su creación % TODO: la cifra de años figura así en las notas originales.
se demuestra su gran trascendencia social. Prueba de ello es que los \textbf{conjuntos inmobiliarios} son una regulación especial de la propiedad horizontal, por lo que se los llama \textbf{propiedades horizontales especiales} (capítulo~\ref{cap:conjuntos}).

\cornell{concepto de la propiedad horizontal}{
\cq{¿Qué es la propiedad horizontal?}{Derecho real autónomo (arts.~2037 a 2069): dominio exclusivo sobre partes privativas inescindiblemente unido a condominio sobre partes comunes.}
\cq{¿En qué se diferencia del dominio y del condominio?}{El dominio tiene un titular y recae sobre todo el inmueble; el condominio, pluralidad de sujetos con partes indivisas; la propiedad horizontal permite dominio exclusivo sobre cada unidad.}
\cq{¿Por qué el condominio no permite que cada hijo tenga su departamento?}{Porque las partes son indivisas e ideales: no corresponden a una parte material determinada y sigue habiendo un solo inmueble.}
\cq{¿De dónde proviene y qué inmuebles puede abarcar?}{De la Ley~13.512 (1948), incorporada al CCyC; vale para edificios en altura y también para propiedades de una planta, en tira y dúplex.}
\cq{¿Qué son las propiedades horizontales especiales?}{Los conjuntos inmobiliarios, regulación especial de la propiedad horizontal.}
}{La propiedad horizontal articula un dominio exclusivo sobre la unidad con un condominio inescindible sobre lo común. Creada por la Ley~13.512 y recogida por el CCyC, permite unidades independientes dentro de un mismo inmueble.}

\chapter[Afectación, unidad funcional y partes comunes]{Afectación, unidad funcional y partes comunes}\label{cap:afectacion}
% ------------------------------------------------------------

\section{Requisitos para afectar un inmueble}

El inmueble debe ser material y jurídicamente apto. Cada unidad funcional debe tener salida independiente a la vía pública, acceso común de todos los propietarios a los servicios de agua y tanques, y cumplir las normativas locales de construcción. En inmuebles ya construidos pueden requerirse obras, con costos adicionales y, a veces, imposibilidad material por falta de espacio.

\begin{table}[H]
\centering
\small
\begin{tabularx}{\textwidth}{>{\raggedright\arraybackslash}X>{\raggedright\arraybackslash}X}
\toprule
\textbf{Requisitos administrativos} & \textbf{Requisitos jurídicos} \\
\midrule
Plano de construcción habitual & Confección del reglamento de copropiedad \\
Plano de mensura (esencial) & Escritura pública del reglamento \\
Aprobación por la municipalidad local & Inscripción en el Registro de la Propiedad Inmueble (constitutiva) \\
\bottomrule
\end{tabularx}
\caption{Requisitos para someter un inmueble a propiedad horizontal}
\end{table}

El \textbf{plano de mensura} es esencial: sin él no existiría la propiedad horizontal. Divide e individualiza los sectores privativos y comunes y fija el porcentual que cada parte privativa tiene sobre la cosa común. La \textbf{inscripción} del reglamento es \textbf{constitutiva}: hace nacer el estado de propiedad horizontal.

\section{La unidad funcional}

El objeto del derecho es la \textbf{unidad funcional}, sobre la que recae el dominio exclusivo (parte privativa), inescindiblemente unida a la parte común. Pueden ser departamentos, locales o cocheras, siempre con salida independiente a la vía pública, directa o por un pasaje común.

\section{Partes comunes (arts.~2040 y 2041)}

Son las destinadas a permitir el acceso a las partes privativas (como el ascensor) y a garantizar el funcionamiento del consorcio. Ningún propietario tiene derecho exclusivo sobre ellas y siempre se determinan en el reglamento. Algunas pueden estar destinadas al uso exclusivo de un propietario, como una azotea que es común pero está destinada al uso exclusivo de una unidad.

\begin{regla}{art.~2041 CCyC, cosas y partes necesariamente comunes}
No pueden recibir otro destino en el reglamento, bajo pena de nulidad: el terreno (siempre común); los muros exteriores y los que separan cada unidad funcional; el piso de cada unidad (techo del vecino inferior); el techo de cada unidad (piso del vecino superior).
\end{regla}

\section{La teoría del cubo de aire}

\begin{ejemplo}{el octavo piso}
Quien vive en el octavo piso tiene un piso que es común (techo del vecino del séptimo) y un techo que es común (piso del vecino del noveno). ¿Cuál es entonces el objeto de su derecho?
\end{ejemplo}

El objeto es el \textbf{cubo de aire}: el volumen interior que el plano de mensura identificó como su unidad funcional (en el ejemplo, la unidad «8.º C»), que la arquitectura y la ingeniería hicieron materialmente posible. El Código lo denomina «volumen». El titular es dueño, además, de un \textbf{porcentual indiviso} del terreno inescindiblemente unido a la unidad; de otro modo no tendría ni terreno, ni paredes laterales, ni piso, porque todo ello es común. % TODO: contenido incompleto o ambiguo en las notas originales (la referencia a la terraza del piso 12 y al piso ocho pisos más abajo no es clara).

\cornell{afectación y estructura jurídica}{
\cq{¿Qué requisitos administrativos y jurídicos exige la afectación?}{Administrativos: plano de construcción y plano de mensura aprobado. Jurídicos: reglamento en escritura pública e inscripción constitutiva en el Registro de la Propiedad.}
\cq{¿Por qué es esencial el plano de mensura?}{Porque sin él no existe la propiedad horizontal: individualiza partes privativas y comunes y fija el porcentual.}
\cq{¿Cuál es el objeto del derecho?}{La unidad funcional (departamento, local, cochera) con dominio exclusivo, unida inescindiblemente a la parte común.}
\cq{¿Qué es el cubo de aire?}{El volumen interior identificado como unidad funcional; el titular no es dueño de paredes, pisos ni techos, que son comunes.}
\cq{¿Qué partes son necesariamente comunes (art.~2041)?}{Terreno, muros exteriores y separativos de unidades, y piso y techo de cada unidad; no pueden recibir otro destino bajo pena de nulidad.}
}{La afectación requiere plano de mensura aprobado y reglamento en escritura inscripto con efecto constitutivo. El objeto es la unidad funcional como cubo de aire con porcentual indiviso del terreno; terreno, muros, pisos y techos son necesariamente comunes.}

\chapter[El consorcio y las expensas]{El consorcio de propietarios y las expensas}\label{cap:consorcio}
% ------------------------------------------------------------

\section{El consorcio como persona jurídica (art.~2044)}

Cuando se transmite al menos una unidad, es decir, cuando hay al menos dos propietarios diferentes, nace el \textbf{consorcio}.

\begin{definicion}{consorcio (art.~2044 CCyC)}
El conjunto de los propietarios de las unidades funcionales constituye la persona jurídica consorcio. Tiene personería jurídica y se rige por las normas del Capítulo y las del reglamento de propiedad horizontal.
\end{definicion}

Tiene nombre y domicilio, capacidad de adquirir derechos y contraer obligaciones, patrimonio formado por las expensas, una voluntad común que se forma en las asambleas y un órgano ejecutivo, el \textbf{administrador}, representante legal.

\subsection{Responsabilidad de los propietarios (art.~143)}

\begin{ejemplo}{falla del ascensor}
Falla el ascensor y un tercero sufre un daño: demanda al consorcio. Responde la persona jurídica, con su patrimonio de expensas.
\end{ejemplo}

Se cita la postura de Marina Mariani de Vidal, con la que las fuentes manifiestan coincidir: ser propietario no puede ser un velo de indemnidad para no responder por los daños que puede ocasionar la cosa, en virtud de la responsabilidad objetiva. La responsabilidad de los propietarios es siempre \textbf{subsidiaria e ilimitada}.

\section{Facultades y obligaciones de los propietarios}

\begin{table}[H]
\centering
\small
\begin{tabularx}{\textwidth}{>{\raggedright\arraybackslash}X>{\raggedright\arraybackslash}X}
\toprule
\textbf{Facultades} & \textbf{Obligaciones y prohibiciones} \\
\midrule
Usar y gozar de la unidad (facultades materiales) & Cumplir el reglamento de copropiedad \\
Enajenar la unidad funcional & Sostener económicamente al consorcio (pagar expensas) \\
Constituir derechos reales & Permitir el acceso a la unidad para reparaciones \\
Realizar actos de administración (alquilar) & No destinar la unidad a un fin diferente al admitido por el reglamento \\
Facultades siempre unidas a la parte proporcional de la cosa común & No realizar actividades riesgosas o peligrosas que afecten al edificio \\
\bottomrule
\end{tabularx}
\caption{Facultades y obligaciones de los propietarios}
\end{table}

El art.~2069 faculta al consorcio a accionar para que cesen las infracciones, por la vía procesal más rápida. Por ejemplo, ante una pérdida en un departamento cuyo titular no permite el acceso, con riesgo de inundar pisos inferiores, el consorcio tiene acción, incluso para pedir el desalojo si está ocupado por terceros.

\section{Las expensas}

\subsection{Concepto y fundamento (art.~2048)}

Para Mariani de Vidal, las expensas son la savia que alimenta al consorcio: sin el sostenimiento de los gastos por todos los propietarios, el sistema sucumbe. Por eso la normativa de orden público es rigurosa y el crédito por expensas goza de grandes privilegios.

\begin{ejemplo}{falta de pago de expensas}
Si la mayoría no paga, el consorcio no puede pagar la luz de los servicios comunes: dejan de funcionar el ascensor, la luz del edificio y el agua. Los apuntes proponen los casos de Juan (piso 14, con secuelas de un accidente) y de María (embarazo a término, piso 17, cuyo médico debería subir por la escalera con un respirador eléctrico, sin luz), que muestran por qué el pago es esencial.
\end{ejemplo}

\begin{regla}{art.~2048 CCyC}
Cada propietario debe atender los gastos de conservación y reparación de su propia unidad funcional. Asimismo, debe pagar las expensas comunes ordinarias de administración y las extraordinarias de reparación o mejora que resuelva la asamblea, así como las que, estando determinadas en el reglamento, se devenguen en forma proporcional a su parte indivisa.
\end{regla}

\subsection{Tipos}

\begin{table}[H]
\centering
\small
\begin{tabularx}{\textwidth}{>{\raggedright\arraybackslash}X>{\raggedright\arraybackslash}X}
\toprule
\textbf{Ordinarias} & \textbf{Extraordinarias} \\
\midrule
Regulares y periódicas & No siempre existen \\
Sueldo del encargado, luz, seguridad, seguros, mantenimiento de ascensores & Cambio del sistema de ascensores, portero eléctrico, pintura del edificio \\
No requieren aprobación previa de la asamblea & Deben ser decididas por la asamblea en todos los casos \\
Las liquida el administrador mensualmente & Exceden las facultades ordinarias del administrador \\
\bottomrule
\end{tabularx}
\caption{Expensas ordinarias y extraordinarias}
\end{table}

\subsection{Obligados al pago (arts.~2049 y 2059)}

Están obligados los \textbf{propietarios} (obligación esencial e irrenunciable), los \textbf{poseedores} por cualquier título, los \textbf{usufructuarios} (el nudo propietario también es responsable), el \textbf{adquirente} (adquiere también la deuda existente) y el \textbf{enajenante} (continúa respondiendo por las deudas al tiempo de la enajenación).

\begin{importante}
Los \textbf{locatarios} se obligan ante el consorcio solo si se obligaron a pagar las expensas, en una relación de derecho personal con sus locadores; si no pagan, es motivo de desalojo. El consorcio \textbf{no tiene acción directa} contra ellos.
\end{importante}

\subsection{La expensa como obligación \textit{propter rem}}

La doctrina no acuerda si es \textit{propter rem} en sentido propio; al menos sería una \textit{propter rem sui generis}. La característica de la obligación \textit{propter rem} es que \textbf{sigue a la cosa}, pero en las expensas el enajenante no queda desobligado.

\begin{table}[H]
\centering
\small
\begin{tabularx}{\textwidth}{>{\raggedright\arraybackslash}X>{\raggedright\arraybackslash}X}
\toprule
\textbf{Adquirente} & \textbf{Enajenante} \\
\midrule
Por deudas devengadas \textbf{antes} de la adquisición: responde solo con la cosa & Continúa respondiendo con todo su patrimonio por las deudas anteriores \\
Por deudas devengadas \textbf{durante} su titularidad: responde con la cosa y con todo su patrimonio & Ya no tiene la cosa: responde exclusivamente con su patrimonio \\
\bottomrule
\end{tabularx}
\caption{Responsabilidad del adquirente y del enajenante}
\end{table}

\begin{ejemplo}{deuda al enajenar}
Una unidad (octavo C) tiene una deuda de expensas de \$10.000. Su titular la enajena y la adquiere Juan. Juan responde desde ese momento por los \$10.000 y el enajenante también. El consorcio tiene acción contra ambos; quien pague repetirá contra quien corresponda.
\end{ejemplo}

\begin{importante}
\textbf{Inoponibilidad.} Todo convenio entre enajenante y adquirente sobre las deudas de expensas es inoponible al consorcio (art.~2049), por ejemplo, una cláusula de la escritura por la que el adquirente se hace cargo de la deuda a cambio de un menor precio. Es un instrumento de derecho personal que produce efectos solo entre las partes.
\end{importante}

\subsection{El plenario «Servicios Eficientes c/ Sabrá»}

Plenario de la Cámara Civil (1999), obligatorio para todas las salas por diez años. A partir de él, que siguió aplicándose aun después de cesar su obligatoriedad, el adquirente en subasta, que se libera de impuestos y tasas, no se libera de las deudas de expensas cuando el monto obtenido en la subasta no alcanza a solventarlas. % TODO: contenido incompleto o ambiguo en las notas originales (el pasaje original contiene una frase confusa sobre el precio obtenido).
Al cumplirse diez años, en 2009, el Dr. Jorge Alterini, entonces director de la revista La Ley, solicitó al docente un estudio publicado en esa revista, que constató que la doctrina era aplicada por todas las salas de la Cámara Civil y había llegado a la justicia comercial.

\subsection{Privilegios y protecciones}

\begin{itemize}
    \item \textbf{Prescripción bienal:} el art.~2562 fija dos años para ejecutar el cobro de expensas.
    \item \textbf{Privilegio especial} en ejecuciones individuales: frente a otros acreedores que subastan el inmueble.
    \item No se les opone la afectación del inmueble al \textbf{régimen de protección de la vivienda} (anterior bien de familia; véase el capítulo~\ref{cap:vivienda}).
    \item Se autorizan \textbf{intereses punitorios} además de moratorios.
    \item \textbf{Procedimiento ejecutivo:} el certificado de deuda del administrador (suscripto también por el consejo de propietarios, si existe) es título ejecutivo que habilita embargo y subasta.
\end{itemize}

\cornell{consorcio y expensas}{
\cq{¿Qué es el consorcio y cuándo nace?}{La persona jurídica formada por los propietarios (art.~2044), con nombre, domicilio, patrimonio de expensas, voluntad por asamblea y representante (administrador). Nace con al menos dos propietarios.}
\cq{¿Responden los propietarios por las deudas del consorcio?}{Responde primero el consorcio; la responsabilidad de los propietarios es subsidiaria e ilimitada (postura de Mariani de Vidal).}
\cq{Distinga expensas ordinarias y extraordinarias.}{Ordinarias: regulares, sin aprobación previa, liquidadas mensualmente. Extraordinarias: no siempre existen y las decide la asamblea.}
\cq{¿Quiénes deben pagar expensas y qué pasa con adquirente y enajenante?}{Propietarios, poseedores, usufructuarios, adquirente y enajenante; locatarios solo frente a su locador. Los convenios entre enajenante y adquirente son inoponibles al consorcio.}
\cq{¿Qué dijo el plenario «Servicios Eficientes c/ Sabrá»?}{El adquirente en subasta responde por las expensas en la medida en que el producido no las cubra; la doctrina se mantuvo y llegó a la justicia comercial.}
\cq{¿Qué privilegios tiene el crédito por expensas?}{Prescripción bienal, privilegio especial, inoponibilidad del régimen de vivienda, intereses punitorios y vía ejecutiva con certificado de deuda.}
}{El consorcio es una persona jurídica sostenida por las expensas, cuya obligación recae sobre propietarios, poseedores y aun sobre quien ya enajenó. La expensa es una obligación \textit{propter rem sui generis} muy tutelada, y la doctrina del plenario protege el cobro frente a la subasta.}

\chapter[Gobierno del consorcio, reglamento, extinción y prehorizontalidad]{Gobierno del consorcio, reglamento, extinción y prehorizontalidad}\label{cap:gobierno-ph}
% ------------------------------------------------------------

\section{El administrador (art.~2067)}

Es el \textbf{representante legal} del consorcio y un \textbf{mandatario}: administra su patrimonio, paga impuestos y sueldos (el encargado, el ayudante), convoca a la asamblea, redacta el orden del día (la identificación clara, concreta y fehaciente de los puntos a tratar), ejecuta sus decisiones y recauda y administra las expensas ordinarias y las extraordinarias decididas. Debe rendir cuentas y tiene la responsabilidad propia de todo mandatario.

\section{Las asambleas (arts.~2058 a 2063)}

La asamblea es el órgano deliberativo y la única vía para resolver asuntos de interés común: forma la voluntad de la persona jurídica. Su decisión obliga a presentes, ausentes y a quienes votaron en contra, siempre que se cumplan las mayorías y se trate el orden del día.

\begin{table}[H]
\centering
\small
\begin{tabularx}{\textwidth}{>{\raggedright\arraybackslash}X>{\raggedright\arraybackslash}X}
\toprule
\textbf{Ordinaria} & \textbf{Extraordinaria} \\
\midrule
Se reúne al menos una vez al año & Se convoca cuando lo exige la situación \\
Trata los temas regulares y la rendición de cuentas & Puede convocarla el administrador, propietarios con más del 5\,\% o autoconvocarse \\
Evalúa presupuestos de gastos o reparaciones & Para autoconvocarse se necesitan los 2/3 de la totalidad de propietarios \\
 & Trata temas urgentes o que no admiten demora \\
\bottomrule
\end{tabularx}
\caption{Asambleas ordinarias y extraordinarias}
\end{table}

\begin{table}[H]
\centering
\small
\begin{tabularx}{\textwidth}{>{\raggedright\arraybackslash}p{0.19\textwidth}>{\raggedright\arraybackslash}X>{\raggedright\arraybackslash}X}
\toprule
\textbf{Mayoría} & \textbf{Descripción} & \textbf{Aplicación} \\
\midrule
Absoluta & Doble cómputo: más del 50\,\% del porcentual y más de la mitad de los votos por unidad funcional & Obras nuevas o mejoras importantes y costosas \\
Unanimidad & 100\,\% de los propietarios & Obras que afecten la estructura en beneficio de uno o pocos; que modifiquen los porcentuales de contribución y propiedad; hipotecar el terreno; cambiar el destino de partes comunes \\
Dos tercios & 2/3 de la totalidad de propietarios & Modificar el reglamento \\
De valor & Más de la mitad del valor del inmueble & Destrucción parcial o total: demolición, venta o reconstrucción \\
\bottomrule
\end{tabularx}
\caption{Tipos de mayorías}
\end{table}

\section{El consejo de copropietarios (art.~2064)}

Órgano \textbf{facultativo} de fiscalización que controla y acompaña la actividad del administrador; cuando existe, suscribe el certificado de deuda.

\section{El reglamento (art.~2056)}

Es un \textbf{contrato de adhesión}: cada nuevo propietario adhiere al contrato que dio vida al sistema. Se redacta en escritura pública, rige derechos y obligaciones e integra el título suficiente de la unidad funcional.

\begin{regla}{art.~2056 CCyC, cláusulas obligatorias}
Son obligatorias para su validez, entre otras: la determinación del terreno; las unidades funcionales y las complementarias; la ubicación de las partes privativas; las partes comunes; el porcentual para el pago de las expensas; las facultades de las asambleas; las mayorías requeridas para las distintas decisiones.
\end{regla}

\section{Extinción del sistema}

La propiedad horizontal recae sobre cosas y carece de suelo propio: si el edificio se destruye, subsiste un \textbf{condominio sobre el terreno}. Ante destrucción o grave deterioro, la \textbf{mayoría de valor} resuelve demolición, reconstrucción o venta. Esa reconstrucción permite al propietario que votó en contra enajenar su parte; si no hay oferentes, el consorcio puede llegar a adquirirla.

\section{Prehorizontalidad (arts.~2070 y 2071)}

Se aplica a los contratos de compraventa de las constructoras sobre unidades a construir y luego someter a propiedad horizontal. Protege a los adquirentes por boleto, titulares de derechos personales, que no pueden adquirir el derecho real porque el inmueble aún no está afectado. Se impone al propietario (generalmente la constructora) la obligación de contratar un \textbf{seguro} en favor de los adquirentes. \textbf{No se aplica} cuando la propiedad horizontal resulta de una partición de condominio, en inmuebles del dominio privado del Estado ni en construcciones con financiamiento público o fideicomisos de organismos oficiales.

\cornell{gobierno del consorcio y prehorizontalidad}{
\cq{¿Qué papel cumple el administrador?}{Representante legal y mandatario: recauda y liquida expensas, paga sueldos e impuestos, convoca asambleas, redacta el orden del día, ejecuta decisiones y rinde cuentas (art.~2067).}
\cq{Distinga asamblea ordinaria y extraordinaria.}{Ordinaria: al menos anual, temas regulares y cuentas. Extraordinaria: cuando la situación lo exige; convocable por el administrador, por propietarios con más del 5\,\% o autoconvocada con 2/3.}
\cq{Mencione las mayorías y un caso de cada una.}{Absoluta (doble cómputo): obras importantes. Unanimidad: cambiar porcentuales o destino de partes comunes. Dos tercios: modificar el reglamento. De valor: demolición, venta o reconstrucción.}
\cq{¿Qué naturaleza tiene el reglamento?}{Contrato de adhesión en escritura pública, parte del título suficiente, con cláusulas obligatorias (art.~2056).}
\cq{¿Cómo se extingue el sistema?}{Por destrucción: queda un condominio sobre el terreno; la mayoría de valor decide demolición, venta o reconstrucción.}
\cq{¿Qué protege la prehorizontalidad y cuándo no se aplica?}{A los adquirentes por boleto antes de la afectación, mediante seguro; no se aplica a particiones de condominio, inmuebles del dominio privado del Estado ni financiamiento público.}
}{La voluntad del consorcio se forma en asambleas con mayorías según la trascendencia del acto y la ejecuta el administrador, bajo fiscalización facultativa del consejo. El reglamento es un contrato de adhesión; ante la destrucción subsiste un condominio sobre el terreno, y la prehorizontalidad protege a quienes compran antes de que nazca el derecho real.}

% ============================================================
% ============================================================


% ============================================================
\part{Nuevas formas de propiedad: conjuntos inmobiliarios, tiempo compartido y cementerios privados}
% ============================================================

\chapter[Conjuntos inmobiliarios]{Conjuntos inmobiliarios}\label{cap:conjuntos}
% ------------------------------------------------------------

Quien ejerza como abogado, escribano o asesor inmobiliario se enfrentará con operaciones sobre inmuebles en countries, barrios cerrados, parques industriales o clubes náuticos. Saber bajo qué régimen está organizado cada emprendimiento, derechos reales o personales, determina qué derechos tiene el cliente, qué riesgos enfrenta, si puede hipotecar el bien o afectarlo a la protección de la vivienda y qué ocurre si el desarrollista quiebra.

\section{Marco histórico}

Bajo la denominación de \textbf{conjuntos inmobiliarios}, el CCyC regula formas de organización habituales durante décadas bajo los nombres de barrios cerrados, countries y clubes de campo: nuevas manifestaciones de propiedad inmobiliaria con amplio desarrollo incluso antes de que existiera una ley nacional.

\subsection{Organización bajo derechos personales}

La mayoría se organizaba en el ámbito de los \textbf{derechos personales}: la \textbf{tierra madre} pertenecía en dominio a una persona jurídica (por ejemplo, una sociedad anónima) con dominio absoluto, exclusivo y perpetuo. Lo que adquiría el comprador \textbf{no era un derecho real, sino una acción}, que, por la autonomía de la voluntad, le conferiría derechos sobre un lote.

\subsection{Los riesgos}

Muchos emprendimientos funcionaron sin dificultad, pero otros revelaron la fragilidad del sistema: la tierra madre fue subastada en ejecuciones individuales o la persona jurídica quebró. Cuando se subasta la tierra, se subasta con lo plantado y lo construido, y los compradores perdían todo.

\section{Regulación provincial y debate doctrinal}

Ante la ausencia de una ley nacional, la Provincia de Buenos Aires dictó una norma, inconstitucional o susceptible de serlo, que exigía afectar los nuevos emprendimientos a propiedad horizontal. Sus fundamentos de inconstitucionalidad eran dos: las provincias carecen de competencia para legislar en propiedad, reservada al Congreso, y en la propiedad horizontal clásica el \textbf{terreno es bien común}, incompatible con la naturaleza de un country, pues quien adquiere en un barrio cerrado difícilmente quiera compartir con todos el terreno donde construirá. La Provincia sabía esto, pero no encontraba otro modo de dar seguridad; la sanción de una ley nacional era imprescindible. Marina Mariani de Vidal caracterizó a los inmuebles de countries organizados bajo propiedad horizontal como «grandes castillos construidos sobre arena».

Gran parte de la doctrina sostuvo que debía legislarse en el ámbito de los \textbf{derechos reales}, mediante un derecho real autónomo o una propiedad horizontal especial que comparta lo común y respete la propiedad individual del terreno. El proyecto del Código incluía en el art.~1887 a los conjuntos inmobiliarios como derechos reales, pero su art.~2075 admitía organizarlos también como derechos personales o mixtos, a elección del desarrollista. La ponencia presentada en el Quinto Congreso de Derecho Privado rechazó esa posibilidad: al urbanizador le conviene el ámbito personal, más flexible y menos oneroso porque no exige plano aprobado, escritura pública de cada propietario ni inscripción registral para la oponibilidad.

El adquirente atiende a la ubicación, el precio, el tipo de construcción, los servicios, la distancia a las lagunas, el tiempo de traslado al trabajo y la cercanía de un colegio; en esas cuestiones «el derecho debe guardar un respetuoso silencio». Pero difícilmente comprenda las implicancias de estar bajo derechos personales o reales, y esa función protectora corresponde al derecho: cautelar el orden público y el interés social es función del derecho real, que considera la finitud humana (de allí las sucesiones), la oponibilidad \textit{erga omnes} y la extensión del derecho más allá de la vida del adquirente.

\begin{table}[H]
\centering
\small
\begin{tabularx}{\textwidth}{>{\raggedright\arraybackslash}X}
\toprule
\textbf{Ventajas exclusivas del titular de un derecho real} \\
\midrule
Afectar el inmueble al régimen de protección de la vivienda familiar \\
Derecho de habitación vitalicio y gratuito del cónyuge o conviviente supérstite \\
Persecución y preferencia \\
Constituir nuevos derechos reales, como usufructo e hipoteca \\
Oponibilidad \textit{erga omnes} mediante inscripción registral \\
\bottomrule
\end{tabularx}
\caption{Ventajas del derecho real frente al derecho personal}
\end{table}

\section{Regulación en el CCyC}

Se los incorpora en el art.~1887 como derechos reales y, en el tránsito a la sanción definitiva, se modificó el art.~2075 eliminando la posibilidad de constituirlos como derechos personales. A partir de la sanción (31 de agosto de 2015), todo conjunto inmobiliario debe constituirse exclusivamente como derecho real. % TODO: verificar referencia presente en las fuentes originales (los apuntes transcriben el art.~1887 con un texto que parece una paráfrasis).

\begin{regla}{art.~2075, tercer párrafo}
Todos los conjuntos inmobiliarios preexistentes a la sanción del nuevo ordenamiento que se hubieran establecido como derechos personales, o en los que coexistan derechos personales y derechos reales, se deben adecuar a las previsiones normativas que regulan este derecho real.
\end{regla}

\begin{regla}{art.~2073 CCyC, tipos}
Son conjuntos inmobiliarios los clubes de campo, barrios cerrados o privados, parques industriales, empresariales o náuticos, y todo otro emprendimiento urbanístico independientemente del destino de vivienda permanente o temporaria, laboral, comercial o empresarial que tenga, comprendidos asimismo aquellos que contemplan usos mixtos.
\end{regla}

No es el destino lo que convierte al emprendimiento en conjunto inmobiliario, sino los elementos del art.~2074.

\begin{regla}{art.~2074 CCyC, elementos característicos}
Sin que su enumeración sea taxativa: cerramiento, partes comunes y privativas, estado de indivisión forzosa y perpetua de las partes, lugares y bienes comunes, reglamento que establece órganos de funcionamiento, limitaciones y restricciones a los derechos particulares y régimen disciplinario, obligación de contribuir con los gastos y cargas comunes y entidad con personería jurídica que agrupe a los propietarios de las unidades privativas.
\end{regla}

\subsection{La adecuación de los preexistentes}

El art.~2075 ordena adecuarlos, pero \textbf{no fija plazo} para los organizados como derechos personales o mixtos \textbf{ni sanción} para quienes no lo hagan; coexistirán durante años conjuntos anteriores a 2015 como derechos personales y posteriores como derechos reales. La adecuación requiere planos de mensura aprobados, reglamento, escritura pública de afectación e inscripción registral de cada unidad; es más conveniente y segura, pero de \textbf{elevadísimos costos}, lo que explica la ausencia de plazo y sanción.

\section{Estructura y funcionamiento}

Los conjuntos son una \textbf{propiedad horizontal especial}. En la propiedad horizontal lo esencialmente común es el terreno; en el conjunto, conforme al art.~2077, cada propietario tiene la \textbf{propiedad exclusiva del lote}: así como en aquella el derecho recae sobre un «cubo de aire», aquí recae sobre la tierra.

\begin{regla}{art.~2077 CCyC, unidad funcional}
El derecho de propiedad del propietario en un conjunto inmobiliario se ejerce sobre unidades funcionales privativas, cuya superficie se destina a uso privado. La unidad funcional puede encontrarse construida o en proceso de construcción; debe contar con independencia funcional vinculada a su destino y salida a la vía pública por vías o elementos comunes o en forma directa.
\end{regla}

\begin{table}[H]
\centering
\small
\begin{tabularx}{\textwidth}{>{\raggedright\arraybackslash}X>{\raggedright\arraybackslash}X}
\toprule
\textbf{Propiedad horizontal clásica} & \textbf{Conjunto inmobiliario} \\
\midrule
El terreno es siempre común a todos los copropietarios & Cada propietario tiene propiedad exclusiva de su lote \\
Se requiere edificio ya construido para afectar & Puede estar construido o en proceso de construcción \\
Típicamente, edificios en altura o plantas divididas & Típicamente, lotes individuales con cerramiento perimetral \\
La unidad es un «cubo de aire» & La unidad incluye la tierra (el lote) \\
Destino predominantemente residencial o comercial en altura & Cualquier destino: residencial, industrial, náutico, deportivo \\
Arts.~2037 a 2072 & Arts.~2073 a 2086 (propiedad horizontal especial) \\
\bottomrule
\end{tabularx}
\caption{Propiedad horizontal clásica y conjunto inmobiliario}
\end{table}

\subsection{Partes comunes y privativas}

Son necesariamente comunes las partes del terreno destinadas a vías de circulación, acceso y comunicación; las destinadas a actividades deportivas, recreativas y sociales; y lo que el reglamento determine (art.~2079, que se transcribe en lo esencial). Son privativas las unidades funcionales, construidas o en construcción, con independencia funcional y salida a la vía pública, que por la naturaleza del emprendimiento suele ser indirecta, a través del cerramiento perimetral.

\subsection{Persona jurídica, reglamento y normas urbanísticas}

El conjunto conforma una \textbf{persona jurídica}, como el consorcio, y se constituye con un \textbf{reglamento} que define partes propias y comunes, funcionamiento y asambleas, régimen disciplinario, gastos y su cálculo y el porcentual de contribución. Se aplican además las \textbf{normas urbanísticas} de cada jurisdicción: son normas administrativas, y no puede constituirse un conjunto en cualquier lugar; localización, límites perimetrales y aspectos edilicios son regulados por las normas provinciales y municipales.

\subsection{Facultades, obligaciones y límites del reglamento}

El propietario debe ejercer su derecho conforme a la normativa, el reglamento y las normas de convivencia, respetar los valores constructivos, paisajísticos, arquitectónicos y ecológicos y pagar las expensas en la proporción del reglamento, que puede fijar otras contribuciones por servicios (por ejemplo, polo). Existen restricciones que no se observan en la propiedad horizontal clásica: régimen de invitados, reglas sobre uso de la pileta y canchas, canon adicional por ciertas instalaciones y restricciones arquitectónicas, paisajísticas y ecológicas.

\begin{importante}
El reglamento puede \textbf{restringir} el uso de servicios y espacios comunes, pero \textbf{nunca prohibir} que un propietario venda o transmita su unidad, porque limitaría la transmisión inmobiliaria. Sí puede establecer un \textbf{derecho de preferencia} a favor del consorcio o de los demás propietarios (igualar la oferta antes que un tercero).
\end{importante}

\section{Conclusiones}

Todo conjunto constituido desde el 31 de agosto de 2015 solo puede organizarse como derecho real; los preexistentes personales o mixtos deben adecuarse aunque no haya plazo ni sanción; la adecuación es costosa y compleja, por lo que coexistirán ambas organizaciones durante años; y el profesional debe identificar el régimen de cada emprendimiento para asesorar correctamente.

\cornell{conjuntos inmobiliarios}{
\cq{¿Qué problema tenían los emprendimientos organizados como derechos personales?}{El comprador solo tenía una acción de una sociedad titular de la tierra madre; si la tierra se subastaba o la sociedad quebraba, perdía todo, incluso lo construido.}
\cq{¿Por qué la propiedad horizontal clásica no servía y era inconstitucional la norma bonaerense?}{Porque obligaba a compartir el terreno, contradictorio con un country, y las provincias carecen de competencia para legislar en propiedad.}
\cq{¿Qué dispone el CCyC sobre su constitución y los preexistentes (arts.~1887 y 2075)?}{Desde el 31/8/2015 solo como derecho real; los anteriores deben adecuarse, sin plazo ni sanción, lo que se explica por los elevados costos.}
\cq{¿Qué elementos los caracterizan (arts.~2073 y 2074)?}{No el destino sino el cerramiento, partes comunes y privativas interdependientes, indivisión forzosa, reglamento, contribución a gastos y personería jurídica.}
\cq{Compare con la propiedad horizontal clásica.}{En el conjunto el lote (tierra) es privativo; en la clásica el terreno es común y la unidad es un cubo de aire. El conjunto puede estar en construcción y tener cualquier destino.}
\cq{¿Qué límite tiene el reglamento sobre la transmisión?}{No puede prohibir la venta, pero sí establecer derecho de preferencia.}
}{Los conjuntos inmobiliarios son una propiedad horizontal especial con lote exclusivo, partes comunes, reglamento y persona jurídica. Desde 2015 solo pueden constituirse como derecho real y los anteriores deben adecuarse, aunque sin plazo ni sanción.}

\chapter[Tiempo compartido]{Tiempo compartido}\label{cap:tc}
% ------------------------------------------------------------

Quien adquiere una semana en un complejo turístico, o una parcela en un cementerio privado, confía en que su adquisición perdurará. El Código intentó dar seguridad jurídica a esas decisiones convirtiéndolas en derechos reales, regulación que todavía genera dudas. El tiempo compartido (arts.~2087 a 2102) se enumera en el art.~1887, inc.~e). Quien ejerce la abogacía enfrenta clientes que lo adquieren sin comprender sus derechos: conocer la figura permite asesorar sobre riesgos y garantías, redactar o revisar contratos de afectación, detectar irregularidades en la comercialización y actuar ante ejecuciones hipotecarias que afecten a los usuarios.

\section{Concepto y objeto}

\begin{definicion}{tiempo compartido (art.~2087 CCyC)}
Hay tiempo compartido si uno o más bienes están afectados al uso periódico y por turno para alojamiento, hospedaje, comercio, turismo, industria u otros fines, que ejercen sobre ellos un derecho que los habilita a su uso y goce. Puede recaer sobre inmuebles o muebles.
\end{definicion}

Se lo asocia con el turismo por ser la modalidad conocida, pero el Código pretende abrir paso a nuevos negocios; por ejemplo, un tomógrafo u otro equipo costoso de medicina nuclear usado por turnos, por un lugar o por otro, durante una semana o período determinado.

\subsection{La modalidad turística en el ámbito de los derechos personales}

La persona adquiere, por un contrato de derecho personal, una \textbf{unidad de tiempo}, generalmente de siete días (o catorce) al año en un complejo vacacional, sin pagar servicios o con un costo muy inferior. Existen sistemas de intercambio dentro del país o en el exterior, como \textbf{Interval}. Quienes lo usan habitualmente no han tenido dificultades, pero como lo lleva adelante una persona jurídica, los inmuebles pueden ser subastados o la persona quebrar, comprometiendo el derecho adquirido por muchos años o a perpetuidad.

\begin{importante}
La regulación como derecho real busca proteger a los adquirentes de los avatares de las personas jurídicas: \textbf{quiebra, inhibición y subasta} de los inmuebles afectados.
\end{importante}

\section{Afectación e inscripción}

Cuando recae sobre inmuebles, el emprendedor debe afectar el inmueble mediante \textbf{escritura pública} con \textbf{inscripción registral} para su oponibilidad \textit{erga omnes} (art.~2089): quien pretenda comprar el inmueble para su propio proyecto, al pedir el informe registral, verá que está afectado a tiempo compartido. Con carácter previo a los anuncios y la comercialización, debe realizarse la afectación en un \textbf{registro de prestadores} y en el Registro de la Propiedad conforme a la ley especial (art.~2092). Según las notas, no existe aún ley especial más que la que regula los tiempos compartidos turísticos, por lo que serían los únicos comercializables; esta ausencia es uno de los mayores obstáculos para la plena operatividad del derecho real.

\section{Sujetos intervinientes}

\begin{table}[H]
\centering
\small
\begin{tabularx}{\textwidth}{>{\raggedright\arraybackslash}p{0.22\textwidth}>{\raggedright\arraybackslash}X}
\toprule
\textbf{Figura} & \textbf{Rol} \\
\midrule
Propietario & Titular del inmueble; puede afectarlo (o el condómino junto a los demás). \\
Emprendedor & Organiza y gestiona el emprendimiento; puede afectar con autorización del propietario, como mandatario. \\
Comercializador & Ofrece y vende las unidades de tiempo al público. \\
Administrador & Gestiona el funcionamiento; puede ser el propio emprendedor o un tercero. \\
Usuarios & Adquirentes del uso periódico: el sujeto que la normativa busca proteger. \\
\bottomrule
\end{tabularx}
\caption{Figuras del tiempo compartido (art.~2100, según las notas)}
\end{table}

Solo puede afectar el inmueble el \textit{dominus}: el propietario o el condómino que, con las partes ideales de todos, conforma el derecho al todo; también el emprendedor con autorización, como mandatario. En definitiva, el titular del dominio, por medio de un administrador.

\section{Interrogantes sobre la naturaleza del derecho}

Según las notas, el art.~2011 establece que al derecho del adquirente «se le aplican las normas de los derechos reales». % TODO: verificar referencia presente en las fuentes originales (el art.~2011 figura en otro capítulo de los apuntes como regla de medianería; la numeración citada aquí puede ser errónea).
De allí surgen interrogantes: ¿se le aplican las normas de los derechos reales, o es un derecho real? ¿Cómo se aplica la norma de un derecho real a un derecho personal, o viceversa? ¿Por qué ha de ser necesariamente uno u otro? Y sobre la titularidad: el propietario tiene un derecho real indiscutido, pero los usuarios, ¿tienen un derecho real sobre cosa propia?

\begin{importante}
\textbf{Interrogante central sin resolver:} el Código dice que al derecho del adquirente se le aplican las normas de los derechos reales, pero no queda claro si es un derecho real en sí mismo. Generará un extenso trabajo de la doctrina y la jurisprudencia.
\end{importante}

\section{Protección de los usuarios}

Al afectar, el inmueble debe estar \textbf{libre de gravámenes} y el propietario, emprendedor, administrador y comercializador \textbf{no deben estar inhibidos} para disponer. El Código admite hipotecas posteriores a la afectación, que \textbf{no son oponibles} a los usuarios (art.~2093): los acreedores del propietario no afectan el derecho de los usuarios.

Las notas señalan dificultades prácticas: si la hipoteca no se paga, el acreedor hipotecario tiene un derecho real con preferencia y puede subastar el inmueble. Es habitual que el titular hipoteque el inmueble precisamente para realizar el emprendimiento, y un cliente puede adquirir una unidad en un complejo importante de la zona norte cuyo inmueble estaba hipotecado.

Los usuarios tienen derechos y deberes: ejercer el uso en la fecha y del modo que correspondan y pagar las expensas necesarias; estas tienen \textbf{cobro ejecutivo} (arts.~2096 y 2097).

\section{Relación de consumo}

El Código incorpora el tiempo compartido como una \textbf{relación de consumo}, algo que ya se había hecho antes por la gran problemática que generaba (art.~2100, según las notas): las relaciones entre propietario, emprendedor, comercializador, administrador y usuarios se rigen por las normas de consumo. % TODO: verificar referencia presente en las fuentes originales (el art.~2100 se cita tanto para la enumeración de sujetos como para la relación de consumo).
Todos deben cumplir las normas nacionales, provinciales o municipales: la regulación del derecho real es nacional, pero el funcionamiento es competencia de las provincias.

\cornell{tiempo compartido}{
\cq{¿Qué es el tiempo compartido (art.~2087)?}{Uso periódico y por turno de uno o más bienes, inmuebles o muebles, para alojamiento, turismo, comercio, industria u otros fines.}
\cq{¿Qué riesgos tenía en derechos personales y qué busca el derecho real?}{Subasta o quiebra de la persona jurídica comprometía el derecho; el derecho real protege contra quiebra, inhibición y subasta.}
\cq{¿Cómo se afecta el inmueble?}{Por escritura pública con inscripción registral (art.~2089), previa a la comercialización y con registro de prestadores (art.~2092).}
\cq{¿Quiénes intervienen?}{Propietario, emprendedor, comercializador, administrador y usuarios.}
\cq{¿Qué interrogante central plantea el Código?}{Si el derecho del adquirente es un derecho real o solo se le aplican las normas de los derechos reales.}
\cq{¿Cómo se protege al usuario frente a hipotecas posteriores?}{Son inoponibles al usuario (art.~2093), aunque subsisten dificultades prácticas si el acreedor ejecuta.}
}{El tiempo compartido es un uso periódico y por turno de bienes afectados mediante escritura inscripta, protegido frente a acreedores posteriores y sometido a normas de consumo. Su naturaleza real sigue siendo discutida y falta la ley especial de registro.}

\chapter[Cementerios privados]{Cementerios privados}\label{cap:cp}
% ------------------------------------------------------------

Los cementerios privados (arts.~2103 a 2113) son el segundo derecho real de esta parte que plantea, como el tiempo compartido, interrogantes sobre su naturaleza.

\section{Antecedentes}

La organización de los cementerios estuvo ligada a lo religioso: se ubicaban en torno a las iglesias. A partir de \textbf{1821} se produce la \textbf{secularización} y surgen los \textbf{cementerios públicos}, del dominio público municipal, destinados a dar sepultura; el derecho de los particulares sobre las sepulturas es una \textbf{concesión de uso} administrativa. Con el aumento de la población, se saturaron y hace unos cincuenta años surgieron los \textbf{cementerios privados}. Un marco provincial de ejemplo: la Ordenanza General~221 de 1978 y la Ley~9.000/94 autorizan cementerios privados con fines especulativos (de lucro), destinados a la inhumación de restos humanos; el poder de policía mortuoria y el servicio de inhumaciones y cremación es ejercido por las municipalidades.

\section{Concepto y estructura}

\begin{definicion}{cementerio privado (art.~2113 CCyC)}
Se consideran cementerios privados los inmuebles de propiedad privada afectados a la inhumación de restos humanos. Requieren habilitación municipal, cerramiento, y tendrán partes comunes y partes privativas sujetas a indivisión forzosa.
\end{definicion}

El Código pretende regularlos como derecho real, el que mayor seguridad confiere. Aunque la ley no lo aclara, pueden afectar el inmueble el titular de dominio y todos los condóminos que representen el cien por ciento de la propiedad.

\begin{table}[H]
\centering
\small
\begin{tabularx}{\textwidth}{>{\raggedright\arraybackslash}X>{\raggedright\arraybackslash}X}
\toprule
\textbf{Partes comunes (indivisión forzosa)} & \textbf{Partes exclusivas (dominio privativo)} \\
\midrule
Calles y circulaciones internas & Parcelas individuales de sepultura \\
Servicios centrales & Sepulcros y construcciones individuales \\
Capillas y lugares de culto & Derecho a inhumar, exhumar y trasladar restos \\
Instalaciones generales del predio & \\
\bottomrule
\end{tabularx}
\caption{Partes comunes y exclusivas del cementerio privado}
\end{table}

\section{Reglamento y derechos del titular}

\begin{regla}{art.~2105 CCyC, reglamento de afectación y uso}
Debe contener: descripción del inmueble; identificación de las partes y lugares comunes y de los servicios comunes; identificación de las parcelas de dominio exclusivo; modo de pago del canon locativo y los gastos; normas de acceso y circulación de titulares y visitantes; normas administrativas sobre inhumaciones, exhumaciones y cremaciones. Debe existir un registro de inhumados y de sepulturas.
\end{regla}

El \textbf{poder de policía mortuoria} está siempre en cabeza del Estado. El titular de la parcela puede realizar inhumaciones, exhumaciones, reducciones y traslados y construir sepulcros, conforme a las normas del cementerio y el poder de policía mortuoria. El administrador debe garantizar el funcionamiento de las instalaciones y servicios comunes para que quien adquirió pueda usar los servicios contratados.

\section{Inembargabilidad de las parcelas (art.~2110)}

\begin{regla}{art.~2110 CCyC}
Las parcelas exclusivas son inembargables e inejecutables, salvo por saldo del precio de compra, precio de construcción de sepulcros, expensas, tasas, impuestos y contribuciones; en esos casos son embargables y ejecutables.
\end{regla}

\begin{importante}
La inembargabilidad es la diferencia fundamental entre el cementerio privado como derecho real y como derecho personal: protege el destino del bien frente a las deudas del propietario del cementerio, de gran trascendencia para quien adquiere con miras al descanso eterno.
\end{importante}

\section{Relación de consumo y aplicación de normas de derechos reales}

\begin{regla}{arts.~2111 y 2112 CCyC}
La relación entre el propietario y el administrador del cementerio y los usuarios es una relación de consumo sujeta a las normas del derecho del consumidor (art.~2111). Al derecho de sepultura se le aplican las normas de los derechos reales (art.~2112).
\end{regla}

Se plantea el mismo interrogante que en el tiempo compartido: sobre el titular de dominio no hay duda de que tiene un derecho real sobre cosa propia, pero si el derecho de sepultura es un derecho real, ¿por qué habría que aplicarle las normas de los derechos reales?

\begin{importante}
En ambos casos el Código usa la misma fórmula, «se le aplican las normas de los derechos reales», sin afirmar que sea un derecho real en sí mismo. Esta ambigüedad es el principal interrogante que deberá resolver la doctrina y la jurisprudencia. Se trata de regulaciones muy novedosas que requerirán tiempo.
\end{importante}

\cornell{cementerios privados}{
\cq{¿Cómo evolucionó la organización de los cementerios?}{Del vínculo religioso a la secularización de 1821 (cementerios públicos, concesión de uso) y, por saturación, a los cementerios privados con fines de lucro.}
\cq{¿Qué es un cementerio privado (art.~2113)?}{Inmueble privado afectado a inhumación, con habilitación municipal, cerramiento y partes comunes y privativas en indivisión forzosa.}
\cq{¿Qué partes son comunes y cuáles exclusivas?}{Comunes: calles, servicios centrales, capillas e instalaciones. Exclusivas: parcelas y sepulcros, con el derecho de inhumar, exhumar y trasladar.}
\cq{¿Qué contiene el reglamento (art.~2105)?}{Descripción del inmueble, partes y servicios comunes, parcelas exclusivas, canon y gastos, normas de acceso y administrativas, y registro de inhumados y sepulturas.}
\cq{¿Por qué es relevante la inembargabilidad (art.~2110)?}{Protege las parcelas de las deudas del propietario del cementerio, salvo saldo de precio, construcción de sepulcros, expensas, tasas e impuestos.}
\cq{¿Qué ambigüedad comparten tiempo compartido y cementerios?}{Ambos usan la fórmula «se aplican las normas de los derechos reales» sin afirmar que sean derechos reales.}
}{El cementerio privado es un inmueble afectado a la inhumación, con partes comunes en indivisión forzosa y parcelas exclusivas inembargables, bajo poder de policía mortuoria estatal y normas de consumo. Como el tiempo compartido, plantea el interrogante de si es un derecho real o solo se le aplican sus normas.}

\begin{conceptoclave}{consideración conjunta}
Tanto en el tiempo compartido como en los cementerios privados hay muchos interrogantes que no pueden resolverse en esta etapa y que serán despejados con un largo trabajo de la doctrina y la jurisprudencia: son regulaciones novedosas que buscan dar seguridad jurídica convirtiendo en derechos reales adquisiciones antes organizadas en el ámbito personal.
\end{conceptoclave}


% ============================================================
\part{Derechos reales de disfrute sobre cosa ajena}
% ============================================================

\chapter[Marco clasificatorio y dominio desmembrado]{Marco clasificatorio y dominio desmembrado}\label{cap:marco}
% ------------------------------------------------------------

El estudio de los derechos reales de disfrute sobre cosa ajena exige partir de la clasificación de los derechos reales y de las facultades del titular de dominio. A partir de allí se comprende cómo el dominio puede desmembrarse y cómo cada derecho que se estudia en los capítulos siguientes deja al propietario con un dominio imperfecto.

\section{La clasificación del art.~1888 aplicada a la cosa ajena}

La clasificación del art.~1888 se presentó en el capítulo~\ref{cap:catalogo}; aquí se retoma la distinción que organiza los capítulos de esta parte y de la siguiente.


El artículo 1888 distingue entre \textbf{derechos reales sobre cosas total o parcialmente propias} y \textbf{derechos reales sobre cosa ajena}. Dentro de estos últimos se distinguen los \textbf{derechos reales de disfrute} y los \textbf{derechos reales de garantía}.

Así como los derechos sobre cosa propia son derechos principales, también lo son los derechos de disfrute sobre cosa ajena. En cambio, los derechos de garantía son siempre, y en todos los casos, derechos \textbf{accesorios} de la obligación que garantizan.

\subsection{Extensión del derecho según su tipo}

El derecho del titular de dominio recae sobre toda la cosa. En cambio:

\begin{itemize}
    \item en el \textbf{condominio}, el derecho real recae sobre una parte indivisa;
    \item en la \textbf{propiedad horizontal}, el objeto del derecho es una parte propia privativa, inescindiblemente unida a un porcentaje en la cosa común;
    \item lo mismo ocurre en la propiedad horizontal especial de los conjuntos inmobiliarios.
\end{itemize}

En los derechos reales sobre cosa ajena, el derecho recae siempre sobre una cosa que pertenece en propiedad a otra persona.

\section{Facultades del titular de dominio}

El titular de dominio tiene la facultad de \textbf{usar} la cosa, \textbf{gozar} de sus frutos y utilidades, y \textbf{disponer} material y jurídicamente de ella. En el marco de esas facultades puede constituir derechos reales; administrar la cosa, por ejemplo celebrando un contrato de locación; donarla, enajenarla, darla en permuta y darla en pago. También puede, conservando la propiedad, constituir derechos de disfrute a favor de otro: es el caso del \textbf{dominio desmembrado}.

\section{Desmembración del dominio}

Cuando el titular de dominio constituye un usufructo a favor de una persona (el \textbf{usufructuario}), desmembra el dominio pleno. Desplaza íntegramente al usufructuario toda la facultad de usar y gozar la cosa, y pasa a denominarse \textbf{nudo propietario}. En una expresión figurada, sería como si dijera: ``Te lo he dado todo; me he quedado con el corazón vacío; me he quedado con la propiedad desnuda''.

El nudo propietario conserva, sin embargo, la \textbf{facultad de disponer}. Puede enajenar la cosa, pero la transmite juntamente con aquello que tiene (la nuda propiedad), de modo que quien adquiere lo hace como adquirente de la nuda propiedad. El usufructo, a su vez, continúa hasta su extinción sin ser afectado por el cambio de titularidad registral: en los inmuebles se mantiene vigente incluso cuando la nuda propiedad se transmite a otra persona. Extinguido el usufructo, el nudo propietario recupera la facultad de uso y goce y vuelve a conformarse un dominio pleno.

\section{Una constante: la imperfección del dominio}

Cada derecho real sobre cosa ajena que se estudia a continuación produce una imperfección o desmembración del dominio, con manifestaciones propias:

\begin{table}[H]
\centering
\begin{tabularx}{\textwidth}{>{\raggedright\arraybackslash}p{0.22\textwidth}>{\raggedright\arraybackslash}X>{\raggedright\arraybackslash}p{0.14\textwidth}}
\toprule
\textbf{Derecho} & \textbf{Cómo se manifiesta la imperfección del dominio} & \textbf{Capítulo} \\
\midrule
Usufructo & El nudo propietario conserva solo la facultad de disponer y no puede afectar el uso y goce del usufructuario. & \ref{cap:usufructo} \\
Superficie & El propietario conserva la disposición material y jurídica, pero tiene la prohibición de turbar el derecho del superficiario y enajena con la superficie. & \ref{cap:superficie} \\
Servidumbre & El dominio del fundo sirviente pasa a ser imperfecto: la servidumbre es una carga que lo grava. & \ref{cap:servidumbre} \\
Garantías reales & Al constituir la garantía el titular imperfecciona su dominio; conserva sus facultades pero no puede disminuir el valor de la garantía. & \ref{cap:garantias-constituyente} \\
\bottomrule
\end{tabularx}
\end{table}

\cornell{clasificación y dominio desmembrado}{
\cq{¿Cómo clasifica el art.~1888 a los derechos reales y qué carácter tiene cada categoría?}{Distingue derechos sobre cosa total o parcialmente propia y derechos sobre cosa ajena; estos últimos son de disfrute o de garantía. Los derechos sobre cosa propia y los de disfrute son principales; los de garantía son siempre accesorios de la obligación que garantizan.}
\cq{¿Sobre qué recae el derecho en el dominio, el condominio y la propiedad horizontal?}{En el dominio, sobre toda la cosa; en el condominio, sobre una parte indivisa; en la propiedad horizontal (y en los conjuntos inmobiliarios), sobre una parte privativa inescindiblemente unida a un porcentaje en la cosa común.}
\cq{¿Qué facultades tiene el titular de dominio y cómo puede ejercer la de disposición?}{Usar, gozar y disponer material y jurídicamente. Puede constituir derechos reales, administrar (locación), donar, enajenar, permutar, dar en pago y, conservando la propiedad, constituir derechos de disfrute a favor de otro (dominio desmembrado).}
\cq{¿Qué se desplaza y qué se conserva al constituirse un usufructo?}{Se desplaza íntegramente al usufructuario la facultad de usar y gozar; el titular pasa a ser nudo propietario y conserva la facultad de disponer, que ejerce transmitiendo la nuda propiedad.}
\cq{¿Por qué el usufructo subsiste si el nudo propietario enajena?}{Porque continúa hasta su extinción, es decir, hasta el cumplimiento del plazo, sin ser afectado por el cambio de titularidad registral.}
}{Los derechos reales se clasifican (art.~1888) en derechos sobre cosa propia y sobre cosa ajena; entre estos últimos, los de disfrute son principales y los de garantía accesorios. El titular de dominio puede usar, gozar y disponer; al constituir derechos sobre cosa ajena desmembra su dominio, que queda imperfecto. En el usufructo el nudo propietario conserva solo la disposición; al extinguirse el derecho se recupera el dominio pleno.}

\chapter{Usufructo}\label{cap:usufructo}
% ------------------------------------------------------------

El usufructo es el derecho de disfrute más amplio sobre cosa ajena: transfiere al usufructuario todo el uso y goce y deja al propietario con la nuda propiedad. Este capítulo recorre su concepto, sus caracteres, su constitución, las obligaciones y derechos de las partes, su extinción y los alcances de su fortaleza.

\section{Concepto y definición legal}

El derecho real de usufructo confiere a su titular la facultad de usar y gozar la cosa que pertenece a otro, con una gran obligación: el cargo de \textbf{conservar la sustancia}, es decir, no alterar la cosa, su materialidad ni su destino. Su obligación principal consiste en restituir la cosa al titular una vez extinguido el usufructo.

\begin{definicion}{Usufructo (art.~2129 CCyC)}
``El usufructo es el derecho real de usar y gozar y disponer jurídicamente de un bien ajeno sin alterar su sustancia''.
\end{definicion}

Hay alteración de la sustancia, si se trata de una \textbf{cosa}, cuando se modifica la materia, la forma o el destino; si se trata de un \textbf{derecho}, cuando se lo menoscaba.

\subsection{Alcance de la facultad de disposición del usufructuario}

El artículo 2129 se refiere a la facultad de disponer jurídicamente. Sin embargo, la facultad de disponer de la cosa está en cabeza del nudo propietario: solo el titular del dominio puede disponer de ella. Lo que el usufructuario puede disponer, y esto resulta sumamente novedoso en materia de usufructo, es aquello que tiene: el \textbf{derecho de uso y goce temporal} de la cosa. Ese derecho puede constituirse en forma onerosa o gratuita; cuando nada se dice, se presume oneroso.

\section{Caracteres fundamentales}

\begin{itemize}
    \item \textbf{Temporalidad.} En todos los casos el usufructo es temporal. El tiempo se fija en el momento de la constitución y consta en el instrumento. Si nada se dice, el usufructo es \textbf{vitalicio}: si algo es esencialmente temporal, es la vida humana.
    \item \textbf{Intransmisibilidad a los herederos.} Este derecho temporal no es transmisible a los herederos.
    \item \textbf{Carácter patrimonial.} Es un derecho real de contenido patrimonial que integra el patrimonio del usufructuario.
\end{itemize}

\section{Constitución del usufructo}

Constituir un usufructo significa afectar el dominio.

\subsection{Legitimados}

\begin{definicion}{Legitimados a constituir (art.~2131 CCyC)}
``Solo pueden constituir usufructo: (1) el titular de dominio; (2) el titular del derecho real de propiedad horizontal sobre la parte privativa; (3) el condominio respecto de la parte ideal sobre la que recae su derecho real; (4) el superficiario cuando existe propiedad superficiaria''.
\end{definicion}

Los conjuntos inmobiliarios no están identificados explícitamente en el artículo 2131. Fue tema de debate si existe fundamento para que el titular de un conjunto inmobiliario, que es una propiedad horizontal especial, no pueda constituirlo y sí pueda hacerlo el titular de una propiedad horizontal clásica. Se estima rotundamente que \textbf{no hay diferencia}: puede constituirlo tanto el propietario de un departamento como el de un inmueble de un conjunto inmobiliario. Considerarlo de otra manera resultaría violatorio del \textit{numerus clausus} interpretado, porque el conjunto inmobiliario es esencialmente una propiedad horizontal especial y se le aplican supletoriamente las normas propias de esta.

\subsection{Sujetos beneficiarios y duración}

El usufructo puede constituirse a favor de una persona física o jurídica, y también a favor de varias personas en forma simultánea.
% TODO: contenido incompleto o ambiguo en las notas originales (el apunte agrega que "el fallecimiento de una no ofrece a las otras"; la frase aparece así y no se completó)

\begin{itemize}
    \item \textbf{Persona física:} será vitalicio o por el plazo por el cual se haya constituido.
    \item \textbf{Persona jurídica:} si nada se dice, el plazo máximo es de \textbf{50 años}.
\end{itemize}

\subsection{Modos y modalidades}

\begin{importante}
\textbf{Art.~2133 CCyC:} ``En ningún caso puede imponerse judicialmente la constitución de un usufructo. Siempre será su constitución consensual, ya sea por contrato oneroso o gratuito''.
\end{importante}

Además del consenso, puede constituirse por acto de última voluntad, es decir, por testamento. Puede transmitirse el uso y goce con reserva de la nuda propiedad; donarse la propiedad con reserva de usufructo para el donante; o transmitirse a unos la nuda propiedad y a otro el usufructo.

El usufructo puede constituirse en forma pura y simple, condicional, sujeta a plazo resolutorio o con algún cargo. No puede, en ningún caso, sujetarse a condiciones o plazos suspensivos, salvo en la constitución testamentaria, cuando la condición o el plazo suspensivo se cumpla antes del fallecimiento del testador.

\section{Obligaciones y derechos del usufructuario}

\subsection{Obligaciones anteriores al uso y goce}

Además del cargo de conservar la sustancia y la obligación de restituir, el usufructuario tiene dos obligaciones previas a entrar en el uso y goce: el inventario y la fianza.

\begin{itemize}
    \item \textbf{Inventario.} Consiste en inventariar el estado de la cosa y es una garantía tanto para el nudo propietario como para el usufructuario. Si ambas partes son capaces, puede realizarse por instrumento privado; si interviene un incapaz, debe realizarse por \textbf{escritura pública}.
\end{itemize}

\begin{definicion}{Fianza (art.~2139 CCyC)}
``La fianza es la garantía de que, una vez extinguido el usufructo, el usufructuario restituirá la cosa en el estado de conservación en que debe restituir, en todos los casos es optativa''.
\end{definicion}

\subsection{Derechos}

El usufructuario tiene un amplio uso y goce de la cosa. Puede percibir los frutos, incluidos los frutos civiles como los alquileres; constituir derechos reales y personales; y transmitir el derecho que tiene a otro. Todos los derechos constituidos, o la propia transmisión, se extinguen con el usufructo en el tiempo previsto originariamente: un plazo temporal definido o la vida, si se trata de una persona humana y el usufructo fue vitalicio.

\subsection{Gastos, mejoras y expensas}

El usufructuario debe hacer frente a los gastos que la cosa demande para su conservación.
% TODO: contenido incompleto o ambiguo en las notas originales (la frase sobre el derecho del propietario y la obligación del usufructuario aparece truncada)
Las mejoras facultativas \textbf{no son indemnizables}. Se ha incorporado en forma definida la obligación de pagar las expensas cuando se trate del usufructo de un inmueble afectado al régimen de propiedad horizontal.

\subsection{Ejecutabilidad}

El usufructo puede embargarse y ejecutarse porque es un derecho real patrimonial. En el caso de un inmueble, puede ser adquirido en subasta; el adquirente deberá dar garantías de la restitución y adquirirá el usufructo hasta su extinción.

\section{Obligaciones del nudo propietario}

\begin{importante}
\textbf{Art.~2151 CCyC:} ``La obligación principal del nudo propietario es no afectar de ninguna manera el uso y goce del usufructuario. Si lo hiciera, si lo afectara de alguna manera, el usufructuario podría llegar a exigir el cese del usufructo e inclusive dar lugar a la reducción de precio o indemnización''.
\end{importante}

El nudo propietario tiene también derechos que a su vez son obligaciones: conserva la facultad de disponer de la cosa material y jurídicamente, pero con una gran limitación: \textbf{no puede afectar el usufructo}.

\section{Extinción del usufructo}

El usufructo se extingue por:

\begin{itemize}
    \item \textbf{Cumplimiento del plazo.} Si es vitalicio también hay plazo, solo que incierto, pues durará lo que dure la vida del usufructuario. En las personas jurídicas se extingue por la extinción de la persona jurídica, por el cumplimiento del plazo o, si no se pactó uno menor, a los 50 años.
% TODO: contenido incompleto o ambiguo en las notas originales (el apunte dice "la vida de los usuarios")
    \item \textbf{No uso durante diez años.} Es una forma de extinción propia del usufructo.
    \item \textbf{Uso abusivo.} Si el usufructuario hiciera un uso abusivo y alterara la sustancia, comprobado judicialmente en un proceso específico, a solicitud del nudo propietario podrá decretarse la extinción por falta de conservación de la sustancia o uso abusivo de la cosa.
\end{itemize}

Cualquiera sea la causa de extinción, el usufructuario debe restituir la cosa y se extinguen todos los derechos reales y personales que hubiese constituido sobre el usufructo.

\section{Fortaleza del usufructo}

El usufructo es un derecho muy fuerte, sobre todo si se constituye sobre un inmueble de forma vitalicia. El titular adquiere sobre la cosa un poder jurídico de estructura legal, \textbf{oponible \textit{erga omnes}}, ejercido por la posesión, que le confiere las facultades de \textbf{persecución y preferencia}.
% TODO: contenido incompleto o ambiguo en las notas originales (el apunte dice "El usuario adquiere sobre la cosa")
A partir de la constitución, el propietario tiene un \textbf{dominio imperfecto}: si se embarga y ejecuta su dominio, el adquirente en subasta deberá esperar la extinción del usufructo para recuperar el uso y goce.

\section{Uso indebido: el usufructo como instrumento para evitar la sucesión}

En los últimos años se ha popularizado la idea de que el usufructo sirve para evitar la sucesión futura. Es frecuente la consulta de quien quiere donar a un hijo y constituir un usufructo a su favor para vivir hasta su muerte, de modo que al fallecer el hijo no deba hacer la sucesión.

Evitar la sucesión no es la finalidad de este derecho real. Se señalan los siguientes argumentos:

\begin{itemize}
    \item La vida es finita y el momento de su finitud es incierto: no se puede saber si, donado el inmueble, la muerte del usufructuario ocurrirá primero y quien donó deba hacer su sucesión para que el inmueble vuelva a él.
% TODO: contenido incompleto o ambiguo en las notas originales (el apunte utiliza la expresión "sucesión del dolor")
    \item No se puede saber si el usufructuario hará uso del derecho de disposición del usufructo que tiene a partir de la regulación de 2015. Tampoco coincide con la realidad el imaginario de quien constituye el derecho con la esperanza de que lo utilice a lo largo de toda su vida: hoy se sabe que es un derecho del cual el usufructuario puede disponer.
% TODO: contenido incompleto o ambiguo en las notas originales (el apunte dice "el propietario puede disponer")
\end{itemize}

El usufructo es un derecho real muy potente que impone una evaluación muy rigurosa y amplia de toda la gama del derecho privado antes de asesorar a un cliente respecto de su constitución.

\cornell{usufructo}{
\cq{¿Qué es el usufructo y qué significa disponer jurídicamente en el art.~2129?}{Derecho real principal, de disfrute y temporal, de usar y gozar y disponer jurídicamente de un bien ajeno sin alterar su sustancia. La disposición no es de la cosa (que corresponde solo al nudo propietario) sino del derecho de uso y goce temporal del usufructuario, constituible en forma onerosa o gratuita (se presume oneroso).}
\cq{¿Cuándo hay alteración de la sustancia?}{En una cosa, cuando se modifica la materia, la forma o el destino; en un derecho, cuando se lo menoscaba.}
\cq{¿Es siempre temporal? ¿Se transmite a los herederos?}{Sí, siempre es temporal; si nada se dice es vitalicio (50 años máximo para personas jurídicas si nada se dice). No es transmisible a los herederos.}
\cq{¿Quién puede constituirlo y cómo?}{Los legitimados del art.~2131, con igual solución para los conjuntos inmobiliarios por aplicación supletoria de la propiedad horizontal. Siempre de forma consensual o por testamento, nunca judicialmente (art.~2133). Admite modalidad pura y simple, condicional, a plazo resolutorio o con cargo; no admite condición ni plazo suspensivo, salvo en la constitución testamentaria si se cumplen antes del fallecimiento del testador.}
\cq{¿Qué obligaciones tiene el usufructuario antes del uso y goce?}{Inventario (privado si las partes son capaces; escritura pública si interviene un incapaz) y fianza, que es optativa (art.~2139).}
\cq{¿Qué derechos tiene y qué ocurre con mejoras, gastos y expensas?}{Amplio uso y goce; percibir frutos (incluidos los civiles); constituir derechos reales y personales; transmitir su derecho; todo se extingue con el usufructo. Las mejoras facultativas no son indemnizables; paga gastos de conservación y expensas en propiedad horizontal. Su derecho es patrimonial, embargable y ejecutable (el adquirente en subasta da garantías de restitución).}
\cq{¿Cuál es la obligación principal del nudo propietario (art.~2151)?}{No afectar de ninguna manera el uso y goce del usufructuario; si lo hace, este puede exigir el cese del usufructo y reducción de precio o indemnización.}
\cq{¿Cómo se extingue y qué efectos tiene?}{Por cumplimiento del plazo (cierto o incierto), no uso durante diez años, uso abusivo comprobado judicialmente a pedido del nudo propietario y extinción de la persona jurídica usufructuaria. Se restituye la cosa y caen los derechos constituidos sobre el usufructo.}
\cq{¿Por qué es un derecho fuerte y por qué no sirve para evitar la sucesión?}{Es un poder jurídico oponible \textit{erga omnes}, ejercido por la posesión, con persecución y preferencia, que deja al propietario con dominio imperfecto frente a embargos. No evita la sucesión por el plazo incierto (puede ser necesaria la sucesión del constituyente) y porque el usufructuario puede disponer del derecho desde 2015.}
}{El usufructo (art.~2129) permite usar, gozar y disponer jurídicamente de un bien ajeno sin alterar su sustancia. Es temporal (vitalicio si nada se dice), intransmisible a herederos, oponible \textit{erga omnes} y subsiste ante el cambio de titularidad de la nuda propiedad. Se constituye por consenso o testamento, nunca judicialmente. El usufructuario debe inventariar, conservar y restituir; el nudo propietario no puede afectar el uso y goce (art.~2151). Se extingue por plazo, no uso durante diez años o uso abusivo, y no es un instrumento adecuado para evitar la sucesión.}

\chapter{Uso y habitación}\label{cap:usohab}
% ------------------------------------------------------------

Tras el usufructo, cuya facultad de uso y goce es íntegra, el análisis avanza hacia dos derechos reales de disfrute limitados: el derecho real de uso y el derecho real de habitación.

\section{Derecho real de uso}

\begin{definicion}{Derecho real de uso}
El derecho real de uso es el derecho de \textbf{usar y gozar} de una cosa que pertenece a otro en propiedad, pero \textbf{en forma limitada}.
\end{definicion}

La diferencia fundamental con el usufructo reside en el límite. En el usufructo la facultad de uso y goce es íntegra y comprende toda la facultad que tenía el titular dominial.
% TODO: contenido incompleto o ambiguo en las notas originales (la redacción original de esta frase es confusa: «uso y goce como su fruto ario de toda la facultad que tenía el titular dominial»)
En el derecho de uso aparece el \textbf{límite}: el titular puede usar y gozar de la cosa de otro, pero de manera limitada.

\subsection{El título constitutivo y el límite}

Los límites deben estar \textbf{claramente consignados en el título constitutivo}. Este requisito es tan riguroso que, aun cuando se constituya un derecho real de uso sobre un inmueble mediante escritura pública y se lo denomine como tal, si el límite no consta claramente, \textbf{la figura se convierte en usufructo}.

\begin{quote}
\textbf{Artículo 2154 CCCN.} \textit{El uso es ese derecho real que consiste en usar y gozar de la cosa ajena, ya sea en una parte material o indivisa, en la extensión y con los límites establecidos en el título, que comparte con el usufructuario la carga de no alterar su sustancia.}
\end{quote}

La carga de no alterar la sustancia se explica porque, una vez extinguido el derecho, la cosa debe restituirse en el estado en que se encontraba al momento de su constitución.

\subsection{Sujeto, frutos y límites a la disposición}

\begin{itemize}
    \item \textbf{Sujeto.} Solo puede constituirse a favor de una persona humana. No necesariamente debe constituirse para cubrir las necesidades del usuario y su familia, como lo exigía la normativa anterior, pero sí debe estar ceñido a un límite al momento de su constitución, relacionado con la persona, con el derecho que se le confiere y con el tiempo.
    \item \textbf{Frutos.} Si existieran únicamente frutos y, una vez abastecida la cantidad que corresponde al usuario y su familia, no hubiese más, los frutos del usuario son preferidos a los del titular dominial. Es un derecho sumamente fuerte en tanto derecho real, aunque su ejercicio sea limitado.
    \item \textbf{Transmisibilidad.} El titular no puede transmitirlo ni constituir sobre la cosa derechos reales. El Código no aclara lo relativo a los derechos personales.
    \item \textbf{Inembargabilidad.} Los derechos de disfrute son patrimoniales e integran el patrimonio; no obstante, si el uso está limitado a las necesidades del usuario y su familia, es inembargable.
\end{itemize}

\begin{quote}
\textbf{Artículo 2157 CCCN.} \textit{Si el uso se limita a las necesidades del usuario y su familia, el derecho es inembargable. El titular de este derecho no puede constituir sobre él derechos reales, ni puede trasmitirlo a título oneroso o gratuito.}
\end{quote}

\section{Derecho real de habitación}

\begin{definicion}{Derecho real de habitación}
El objeto del derecho real de habitación es la \textbf{facultad de morar}, es decir, de habitar en un inmueble que pertenece a otro, ya sea en todo el inmueble o en una parte material. Solo puede constituirse a favor de una \textbf{persona humana}.
\end{definicion}

\begin{nota}
Queda planteada, para reflexión, la cuestión de la inscripción registral cuando el derecho recae no sobre todo el inmueble sino sobre una parte.
% TODO: contenido incompleto o ambiguo en las notas originales (la duda se plantea pero no se resuelve)
\end{nota}

\begin{quote}
\textbf{Artículo 2159 CCCN.} \textit{Supletoriamente se le aplican las normas del usufructo al derecho real de habitación.}
\end{quote}

\begin{table}[H]
\centering
\begin{tabularx}{\textwidth}{>{\raggedright\arraybackslash}p{0.30\textwidth}>{\raggedright\arraybackslash}X}
\toprule
\textbf{Aspecto} & \textbf{Limitación del habitador} \\
\midrule
Transmisión entre vivos & No es transmisible por actos entre vivos. \\
Transmisión \textit{mortis causa} & No es transmisible por causa de muerte. \\
Derechos reales sobre la cosa & No puede constituirlos. \\
Derechos personales sobre la cosa & No puede constituirlos. \\
Ejecución por acreedores & No es ejecutable por los acreedores del habitador. \\
Gastos e impuestos & Debe abonar los impuestos, contribuciones y reparaciones necesarias para la conservación, en proporción a la parte que ocupa. \\
\bottomrule
\end{tabularx}
\end{table}

El derecho de habitación se estudia en los capítulos siguientes en una de sus aplicaciones más relevantes: la que protege al cónyuge y al conviviente supérstite (capítulo~\ref{cap:habsuperstite}).

\cornell{uso y habitación}{
\cq{¿En qué se diferencia el derecho de uso del usufructo y qué ocurre si el título no fija el límite?}{En el usufructo el uso y goce es íntegro; en el uso es limitado y los límites deben constar claramente en el título constitutivo. Si no constan, la figura se convierte en usufructo (art.~2154).}
\cq{¿Qué carga comparte el usuario con el usufructuario y por qué?}{No alterar la sustancia, porque al extinguirse el derecho la cosa debe restituirse en el estado en que estaba al constituirse.}
\cq{¿A favor de quién puede constituirse el uso y qué límite debe tener?}{Solo de una persona humana. No debe necesariamente cubrir las necesidades del usuario y su familia, pero debe ceñirse a un límite relacionado con la persona, el derecho conferido y el tiempo.}
\cq{¿Qué sucede con los frutos?}{Si, satisfechas las necesidades del usuario y su familia, no sobran, los del usuario son preferidos a los del titular dominial: es un derecho real fuerte aunque limitado.}
\cq{¿Es embargable o transmisible el uso?}{Si se limita a las necesidades del usuario y su familia es inembargable; no puede constituirse sobre él derechos reales ni transmitirse a título oneroso o gratuito (art.~2157).}
\cq{¿Qué es la habitación y qué limitaciones tiene el habitador?}{Facultad de morar en un inmueble ajeno, en todo o en parte, a favor de una persona humana; se le aplican supletoriamente las normas del usufructo (art.~2159). No transmite entre vivos ni por causa de muerte, no constituye derechos reales ni personales, no es ejecutable por sus acreedores y paga impuestos, contribuciones y reparaciones necesarias en proporción a la parte que ocupa.}
}{El uso es un derecho de usar y gozar de cosa ajena de manera limitada; el límite debe constar en el título, pues de lo contrario la figura se convierte en usufructo. Se constituye a favor de personas humanas, no es transmisible y, si se limita a las necesidades del usuario y su familia, es inembargable. La habitación otorga la facultad de morar en un inmueble ajeno, en todo o en parte, con limitaciones de transmisión, constitución de derechos y ejecución, y con obligación de afrontar gastos en proporción a la parte ocupada.}

\chapter{Habitación del cónyuge y del conviviente supérstites}\label{cap:habsuperstite}
% ============================================================

El derecho de habitación estudiado en el capítulo anterior tiene una aplicación especialmente relevante: la protección de la vivienda del cónyuge y del conviviente que sobreviven a quien fue titular del inmueble. Este capítulo recorre el problema que la ley procura resolver, su antecedente en el Código de Vélez, el régimen vigente para el cónyuge (art.~2383), el régimen del conviviente (art.~527) y el fundamento común de ambos.

\section{Naturaleza y ubicación en el Código}

El derecho real de habitación del cónyuge supérstite es un derecho real sobre cosa ajena, de los denominados de disfrute, con una particularidad: \textbf{no se encuentra regulado en el Libro Cuarto} del Código Civil y Comercial, donde se tratan conjuntamente los demás derechos reales, sino en el \textbf{Libro Quinto}, dedicado a la transmisión de los derechos por causa de muerte. El artículo que lo regula es el \textbf{2383}. Su ubicación en el libro de las sucesiones se explica porque este derecho se inserta en el conjunto de normas que pretenden enfrentar las consecuencias de la finitud humana.

\begin{nota}
Quien ejerce la abogacía, en el ámbito independiente, notarial, judicial o de asesoramiento, se encontrará con situaciones en las que un cónyuge o conviviente supérstite necesite ser informado y protegido en relación con la vivienda familiar. El material señala que el derecho de habitación era prácticamente desconocido incluso para los propios profesionales.
\end{nota}

\section{El problema y su antecedente normativo}

\begin{ejemplo}{El caso de la tía Cota (año 1972)}
La tía Cota estaba casada con el tío Pedro, quien era viudo de su primer matrimonio y tenía un hijo ya adulto de esa unión anterior. Vivían en su casa propia; el tío trabajaba y no tenían grandes lujos ni necesidades. Hacia el año 1972 el tío Pedro falleció y se inició la sucesión. La casa en la que ambos habían vivido quedó dividida en un \textbf{50\% indiviso para la tía Cota y un 50\% indiviso para el hijo del tío Pedro}.

Cuando el hijo pidió la partición, el inmueble se puso en venta. Con el 50\% indiviso cobrado, descontados los gastos de sucesión, inmobiliaria y escritura, y con una pequeña pensión que apenas le alcanzaba para vivir, la tía Cota no pudo comprar nada y \textbf{quedó en la calle}.
\end{ejemplo}

Este modo de perder la vivienda, denominado ``pedir la parte'', sigue siendo una preocupación que viudos y viudas manifiestan con temor en la consulta profesional: qué sucederá si los otros herederos piden la parte.

\subsection{El artículo 3574 bis del Código de Vélez}

Hasta el año 1974 no existía protección alguna para el cónyuge supérstite en situaciones como la descripta: cada vez que el hijo pedía la partición, el cónyuge viudo quedaba sin vivienda.

\begin{quote}
\textbf{Artículo 3574 bis CC (Ley 20.798, incorporado en 1974).} \textit{Incorporó el derecho real vitalicio y gratuito del cónyuge supérstite de habitar en el inmueble que fue asiento del hogar conyugal, siempre que fuera el único inmueble habitable del causante. Este derecho prevalecía aun ante el derecho de los herederos.}
\end{quote}

Esta norma presentaba limitaciones importantes:

\begin{itemize}
    \item el juez \textbf{no podía imponerlo de oficio}: solo se reconocía a solicitud de parte;
    \item era un derecho \textbf{prácticamente desconocido}, tanto por sus titulares como por los propios colegas, por lo que, pese a existir desde 1974, viudas y viudos continuaban quedando sin vivienda por no invocarlo;
    \item el derecho \textbf{se extinguía} si el cónyuge supérstite contraía nuevas nupcias.
\end{itemize}

\section{Régimen actual del cónyuge supérstite (art.~2383)}

\begin{quote}
\textbf{Artículo 2383 CCCN.} \textit{El cónyuge supérstite tiene derecho real de habitación vitalicio y gratuito de pleno derecho sobre el inmueble de propiedad del causante que constituyó el último hogar conyugal y que a la apertura de la sucesión no se encontraba en condominio con otras personas. Este derecho es inoponible a los acreedores del causante.}
\end{quote}

Las características centrales del régimen se resumen en el siguiente cuadro:

\begin{table}[H]
\centering
\begin{tabularx}{\textwidth}{@{}>{\raggedright\arraybackslash}p{0.25\textwidth}>{\raggedright\arraybackslash}X>{\raggedright\arraybackslash}p{0.17\textwidth}@{}}
\toprule
\textbf{Característica} & \textbf{Descripción} & \textbf{¿Novedad respecto del art.~3574 bis?} \\
\midrule
Vitalicio & Se ejerce durante toda la vida del cónyuge supérstite. & Sí \\
Gratuito & No se paga contraprestación por el uso del inmueble. & Sí \\
De pleno derecho & Se adquiere sin necesidad de solicitud judicial ni acto voluntario (art.~1894). & Sí \\
Sobre inmueble propio o ganancial & No importa la calificación del bien: puede ser propio o ganancial del causante. & Sí \\
Sin requisito de único inmueble & Ya no es necesario que sea el único inmueble habitable del causante. & Sí \\
No se extingue por nuevas nupcias & No se pierde si el cónyuge supérstite se casa nuevamente o forma una unión convivencial. & Sí \\
Inoponible a acreedores del causante & No puede hacerse valer frente a quienes el causante debía dinero. & Límite \\
Inmueble no en condominio & El inmueble debe haber sido de propiedad exclusiva del causante al momento del fallecimiento. & Requisito \\
\bottomrule
\end{tabularx}
\end{table}

\subsection{Adquisición legal (art.~1894)}

\begin{quote}
\textbf{Artículo 1894 CCCN.} \textit{Se adquieren por imperio de la ley el derecho real de habitación del cónyuge supérstite y del conviviente supérstite, y los derechos reales sobre inmuebles que surgen de la ley.}
\end{quote}

Se trata de un caso de \textbf{adquisición legal} de los derechos reales: el derecho se adquiere por el solo hecho del fallecimiento del causante, sin necesidad de ningún acto jurídico adicional. Esto supera el principal problema práctico de la normativa anterior, que exigía la solicitud judicial.

\subsection{Tensión entre el cónyuge y los herederos}

Existe una tensión entre dos intereses. Por un lado, el cónyuge supérstite necesita conservar el inmueble que constituye el asiento del hogar conyugal. Por el otro, los herederos tienen derecho a suceder a su causante, y sin duda lo tienen.

\begin{importante}
El inmueble es uno: o se concede al cónyuge supérstite el derecho real de habitación de forma vitalicia, o se define a favor del derecho de los herederos. \textbf{El derecho ha optado por el derecho del cónyuge supérstite.}
\end{importante}

\subsection{Constitucionalización del derecho privado}

El cambio de paradigmas producido con el Código Civil y Comercial de 2015 implica la denominada ``constitucionalización'' (o ``internalización'') del derecho privado: los principios y derechos reconocidos en la \textbf{Constitución Nacional} y en los \textbf{tratados de derechos humanos} penetran y modelan el derecho privado.

Bajo este paradigma, si se protege al cónyuge viudo, no se puede perder esa protección por el solo hecho de que se case nuevamente: ello sería un castigo y una contradicción con la lógica protectoria del instituto. El derecho se defiende, por lo tanto, con independencia de si el cónyuge supérstite contrae nuevas nupcias, constituye una nueva unión convivencial o permanece solo.

\section{Régimen del conviviente supérstite (art.~527)}

La segunda gran novedad en esta materia es el reconocimiento de la \textbf{unión convivencial} como forma de familia protegida por el ordenamiento. Así como en algún momento se distinguía entre hijos legítimos e ilegítimos, distinción hoy superada, actualmente se reconoce que existen ``familias'': el término familia ya no está ligado exclusivamente a la institución matrimonial. El Código reconoce la unión convivencial como una formación de familia, con o sin hijos, que puede durar muchos años y en la que uno de los convivientes puede quedar, de un día para el otro, sin dónde vivir por el fallecimiento del otro.

\begin{quote}
\textbf{Artículo 527 CCCN (efectos de la convivencia).} \textit{El conviviente supérstite que carece de vivienda propia habitable o de bienes suficientes que aseguren el acceso a ésta, puede invocar el derecho real de habitación gratuito por un plazo máximo de dos años sobre el inmueble de propiedad del causante que constituyó el último hogar familiar y que a la apertura de la sucesión no se encontraba en condominio con otras personas. Este derecho es inoponible a los acreedores del causante. Se extingue si el conviviente supérstite constituye una nueva unión convivencial o contrae matrimonio.}
\end{quote}

\subsection{Comparación con el régimen del cónyuge}

\begin{table}[H]
\centering
\begin{tabularx}{\textwidth}{@{}>{\raggedright\arraybackslash}p{0.24\textwidth}>{\raggedright\arraybackslash}X>{\raggedright\arraybackslash}X@{}}
\toprule
\textbf{Aspecto} & \textbf{Cónyuge supérstite} & \textbf{Conviviente supérstite} \\
\midrule
Duración & Vitalicio. & Plazo máximo de 2 años. \\
Extinción por nuevas nupcias o unión convivencial & No se extingue. & Se extingue. \\
Carencia de vivienda propia & No la exige. & Exige carecer de vivienda propia habitable o de bienes suficientes. \\
Inmueble en condominio & No opera si hay condominio. & No opera si hay condominio. \\
Adquisición & De pleno derecho (art.~1894). & De pleno derecho, pero con requisitos que los herederos pueden invocar. \\
Inoponibilidad a acreedores & Inoponible a acreedores del causante. & Inoponible a acreedores del causante. \\
Regulación en el CCCN & Libro Quinto, art.~2383. & Libro Segundo, art.~527. \\
\bottomrule
\end{tabularx}
\end{table}

\subsection{Debate sobre la adquisición de pleno derecho y los requisitos}

El artículo 527 exige que el conviviente supérstite carezca de vivienda propia habitable o de bienes suficientes para proveerse una. Surge entonces la cuestión interpretativa de cómo se compatibiliza esta exigencia con la adquisición legal del derecho (art.~1894).

La posición expuesta en el material sostiene que el conviviente supérstite \textbf{adquiere el derecho de pleno derecho}: desde el momento del fallecimiento cuenta con los dos años que la ley le confiere para habitar gratuitamente el inmueble. Es facultad de los herederos invocar, durante ese período, que el conviviente tiene medios propios de vida o que posee otro inmueble en el cual vivir.

\begin{importante}
El derecho del conviviente supérstite nace automáticamente con el fallecimiento, pero su subsistencia puede ser cuestionada por los herederos si acreditan que el conviviente posee recursos suficientes.
\end{importante}

\section{La vivienda como derecho humano fundamental}

Ambos institutos se encuadran en el reconocimiento del \textbf{derecho a la vivienda como derecho humano fundamental}. La ley ha tenido en cuenta este principio al regular el derecho real de habitación del cónyuge supérstite, de carácter vitalicio, y el del conviviente supérstite, limitado por un plazo de dos años, en correspondencia con la opción que cada familia realiza en el marco de su libertad: contraer matrimonio o constituir una unión convivencial.

\cornell{habitación del cónyuge y del conviviente supérstites}{
\cq{¿Qué es el derecho real de habitación del cónyuge supérstite y dónde está regulado?}{Es un derecho real vitalicio y gratuito, adquirido de pleno derecho (art.~1894), sobre el inmueble que constituyó el último hogar conyugal y era de propiedad exclusiva del causante al abrirse la sucesión (art.~2383). Se regula en el Libro Quinto (sucesiones) y no en el Libro Cuarto, porque se inserta en las normas que enfrentan las consecuencias de la finitud humana.}
\cq{¿Por qué era necesaria la protección y cómo operaba el art.~3574 bis?}{Sin protección, cuando un heredero pedía la partición, el cónyuge podía quedar sin vivienda (caso de la tía Cota). El art.~3574 bis (1974) otorgaba un derecho vitalicio y gratuito sobre el único inmueble habitable del causante que fue asiento del hogar conyugal, pero solo a pedido de parte, era desconocido y se extinguía por nuevas nupcias.}
\cq{¿Qué novedades introduce el CCCN respecto del régimen anterior?}{Adquisición de pleno derecho; no se extingue por nuevas nupcias ni nueva unión convivencial; no requiere que sea el único inmueble habitable; puede recaer sobre bienes propios o gananciales del causante.}
\cq{¿A quién protege la ley si hay conflicto con los herederos y por qué?}{Al cónyuge supérstite: el inmueble es uno y el derecho optó por la habitación vitalicia. Ello responde a la constitucionalización del derecho privado: perder la protección por casarse de nuevo sería un castigo contrario a la lógica protectoria.}
\cq{¿Qué establece el art.~527 sobre el conviviente supérstite?}{Habitación gratuita por un máximo de dos años sobre el inmueble del causante que fue último hogar familiar y no estaba en condominio, si carece de vivienda propia habitable o de bienes suficientes; es inoponible a los acreedores y se extingue si constituye nueva unión convivencial o contrae matrimonio.}
\cq{¿En qué se diferencian el derecho del cónyuge y el del conviviente?}{Duración (vitalicio frente a dos años), extinción por nuevas relaciones (solo el del conviviente) y carencia de vivienda propia (solo la exige el del conviviente, y los herederos pueden invocarla).}
\cq{¿Cuándo se adquiere el derecho del conviviente y qué pueden alegar los herederos?}{Según la posición expuesta, de pleno derecho desde el fallecimiento; los herederos pueden invocar durante los dos años que el conviviente tiene medios propios o posee otro inmueble.}
\cq{¿Qué fundamento común tienen ambos regímenes?}{La vivienda como derecho humano fundamental, con soluciones distintas según la opción de cada familia: matrimonio o unión convivencial.}
}{El cónyuge supérstite tiene un derecho real de habitación vitalicio, gratuito y de pleno derecho sobre el último hogar conyugal de propiedad exclusiva del causante (arts.~2383 y 1894), inoponible a los acreedores del causante, que superó las limitaciones del art.~3574 bis. El conviviente supérstite tiene, según el art.~527, un derecho gratuito de hasta dos años, sujeto a carencia de vivienda propia y extinguible por nueva unión o matrimonio. Ambos regímenes se fundan en la protección de la vivienda como derecho humano fundamental.}

\chapter{Superficie}\label{cap:superficie}
% ============================================================

La superficie es una de las figuras más innovadoras del Código Civil y Comercial, vigente desde 2015. Su rasgo distintivo es que \textbf{separa la propiedad del suelo de la propiedad de lo construido, plantado o forestado}. Este capítulo estudia su concepto y caracteres, su objeto, su adquisición y facultades, su duración y extinción, y la indemnización debida al superficiario.

\begin{importante}
Esta disociación genera posibilidades de negocio que antes eran jurídicamente imposibles en Argentina, y es una herramienta cotidiana en mercados inmobiliarios sofisticados de Europa y los Estados Unidos.
\end{importante}

\begin{nota}
Como analogía, el derecho de superficie se asemeja a la situación de un inquilino que coloca muebles propios en una casa ajena, con la diferencia de que, en este caso, el ``mueble'' puede ser un edificio entero.
\end{nota}

\section{Concepto, clasificación y caracteres}

El derecho real de superficie presenta una doble naturaleza según el CCCN.

% TODO: contenido incompleto o ambiguo en las notas originales (el apunte cita el Art. 1888 tanto para la clasificación como para la carga o gravamen real)
\begin{importante}
\textbf{Art.~1888 CCCN.} El derecho real de superficie es un derecho real \textbf{sobre cosa propia} cuando existe propiedad superficiaria, y es un derecho real \textbf{sobre cosa ajena} cuando recae sobre el derecho de plantar, construir o forestar.
\end{importante}

Los lineamientos generales de los derechos reales se aplican a la superficie de la siguiente manera:

\begin{itemize}
    \item es un \textbf{derecho principal} (art.~1889);
    \item \textbf{se ejerce por la posesión} (art.~1891);
    \item \textbf{constituye una carga o gravamen real}: todo derecho que el titular de dominio constituya a favor de otro respecto de la cosa constituye una carga o gravamen real (art.~1888);
    \item \textbf{se adquiere por título suficiente y modo suficiente} (art.~1892);
    \item siempre se constituye sobre un \textbf{inmueble ajeno};
    \item se constituye por \textbf{escritura pública}, porque recae sobre inmuebles y el título suficiente debe contenerse en dicho instrumento;
    \item la \textbf{inscripción registral} otorga oponibilidad \textit{erga omnes} a partir de su registración.
\end{itemize}

Es además un derecho \textbf{temporario}, que recae sobre un inmueble ajeno y confiere al titular las facultades de \textbf{uso, goce y disposición material y jurídica}. Para precisar el alcance de la disposición conviene comparar tres situaciones: el propietario tiene derecho sobre toda la cosa; el condómino, sobre su parte indivisa; el superficiario, en cambio, tiene derecho a disponer y transmitir únicamente aquello que constituye el objeto de su derecho. No dispone de toda la cosa, porque la superficie siempre se constituye sobre un inmueble ajeno.

\section{Objeto, emplazamiento y legitimados}

\subsection{El derecho real sobre un derecho}

Como regla general, el objeto de los derechos reales son las cosas. Solamente cuando la ley lo disponga, el derecho real puede recaer sobre un derecho. Esta es una gran novedad del CCCN respecto del Código de Vélez.

\begin{definicion}{Objeto del derecho de superficie (art.~2114)}
El derecho de superficie puede recaer sobre el \textbf{derecho de plantar, forestar o construir} en un terreno ajeno, ya sea en el suelo, el subsuelo o el espacio aéreo.
\end{definicion}

\subsection{Las dos modalidades del objeto}

\begin{definicion}{Modalidades del derecho de superficie (art.~2115)}
El objeto del derecho puede ser:
\begin{enumerate}[label=(\alph*)]
    \item el \textbf{derecho real de construir, plantar o forestar} en un inmueble ajeno, haciendo propio aquello que se construya o plante; o
    \item la \textbf{propiedad de aquello que ya está plantado o ya está construido} en la superficie, sin adquirir la propiedad del suelo.
\end{enumerate}
\end{definicion}

En todos los casos coexisten la propiedad del superficiario (ya sea el derecho real sobre el derecho a plantar, construir o forestar, o sobre lo plantado, construido o forestado) y el dominio del titular del inmueble. Este último presenta una imperfección, dada por la carga que implica la constitución de este derecho real.

\subsection{Emplazamiento y legitimados}

El derecho puede emplazarse en \textbf{todo el inmueble} o en una \textbf{parte material} del inmueble. Pueden constituirlo únicamente el \textbf{titular de dominio}, los \textbf{condóminos en su conjunto} y el \textbf{propietario de un inmueble sometido a propiedad horizontal}: solamente ellos pueden afectar el inmueble con el gravamen que genera la constitución de un derecho real a favor de terceros, en este caso, del superficiario.

\section{Adquisición, facultades y situación del propietario}

\subsection{Adquisición}

El derecho de superficie se adquiere con \textbf{título suficiente y modo}. El acto causal de la transmisión puede ser un contrato \textbf{oneroso o gratuito}, y puede realizarse por actos \textbf{entre vivos} o \textbf{mortis causa}. \textbf{En ningún caso se adquiere por prescripción adquisitiva.}

\subsection{Facultades del superficiario}

Se trata de un derecho real de especial fortaleza: el superficiario, aunque no ha adquirido la propiedad del terreno, puede constituir derechos reales de garantía sobre el objeto de su derecho. Sus facultades son:

\begin{itemize}
    \item \textbf{ejerce el derecho por la posesión} en todos los casos;
    \item puede constituir \textbf{derechos reales de garantía} sobre el objeto de su derecho (no sobre el inmueble del propietario);
    \item puede \textbf{hipotecarse un derecho}: la hipoteca puede recaer sobre el derecho de construir, plantar o forestar, o bien sobre la propiedad superficiaria. Esta es una particularidad remarcable;
    \item salvo pacto en contrario en el instrumento constitutivo, puede construir un inmueble y \textbf{afectarlo a propiedad horizontal};
    \item puede \textbf{transmitir} lo que tiene: el derecho de superficie.
\end{itemize}

\begin{importante}
Todos estos derechos están limitados en el tiempo, y esas limitaciones son de orden público.
\end{importante}

\subsection{Situación del propietario del suelo}

Durante la vigencia del derecho de superficie, el propietario del suelo conserva la disposición material y jurídica, porque sigue siendo el dueño. Sin embargo, su dominio es imperfecto.

\begin{definicion}{Dominio imperfecto del propietario}
Es la situación del titular de dominio cuyo derecho está limitado por la \textbf{prohibición de turbar el derecho del superficiario}. Debe respetarlo y, si dispone o enajena el inmueble, lo enajena con el derecho de superficie.
\end{definicion}

\section{Duración, destrucción del objeto y extinción}

% TODO: contenido incompleto o ambiguo en las notas originales (el índice temático del apunte menciona «plazos máximos legales», «plazo convencional» y «orden público», pero el desarrollo no detalla los plazos máximos)
Las limitaciones temporales del derecho son de orden público. En la práctica, resulta difícil que el derecho se constituya por un plazo pequeño; en general se contempla el plazo máximo legal.

\subsection{Destrucción del objeto}

\begin{importante}
\textbf{Art.~2122 CCCN.} El derecho real de superficie \textbf{no se extingue por la destrucción del objeto}. Si se destruyen las construcciones, el superficiario debe reconstruir en el plazo de \textbf{seis (6) años}; cuando el derecho recae sobre plantaciones o forestaciones, debe plantar o forestar en el plazo de \textbf{tres (3) años}. Si no cumple, opera la extinción del derecho.
\end{importante}

En consecuencia, la extinción del objeto no produce la extinción inmediata del derecho.

\subsection{Modos de extinción}

El derecho de superficie se extingue por:

\begin{enumerate}
    \item \textbf{vencimiento del plazo} convencional o legal;
    \item \textbf{renuncia expresa} del superficiario;
    \item \textbf{condición resolutoria}, cuando se cumple la condición a la que el derecho estaba sujeto;
    \item \textbf{consolidación}: las calidades de propietario y superficiario se confunden en una misma persona;
    \item \textbf{no uso durante 10 años}, cuando el derecho fue constituido para construir;
    \item \textbf{no uso durante 5 años}, cuando el objeto fue plantar o forestar.
\end{enumerate}

\subsection{Efectos de la extinción sobre los derechos de terceros}

El superficiario puede constituir derechos reales y personales a favor de terceros, siempre limitados al plazo de duración de la superficie. Los efectos de la extinción sobre esos derechos dependen del modo en que aquella se produzca:

\begin{table}[H]
\centering
\begin{tabularx}{\textwidth}{@{}>{\raggedright\arraybackslash}p{0.30\textwidth}>{\raggedright\arraybackslash}X@{}}
\toprule
\textbf{Modo de extinción} & \textbf{Efectos sobre terceros} \\
\midrule
Extinción anticipada (antes del plazo) & Los derechos reales y personales constituidos a favor de terceros se mantienen y conservan sus efectos hasta la extinción del plazo original. \\
Extinción por vencimiento del plazo & El propietario adquiere todo lo construido, plantado o forestado que exista al vencimiento, libre de los derechos reales o personales que el superficiario haya constituido. \\
\bottomrule
\end{tabularx}
\end{table}

\section{Indemnización al superficiario y normas supletorias}

\subsection{Indemnización (art.~2126)}

\begin{importante}
Producida la extinción, el titular del suelo, que es quien adquiere lo plantado, construido o forestado, \textbf{debe indemnizar al superficiario}, excepto pacto en contrario. El monto se fija:
\begin{enumerate}[label=(\alph*)]
    \item por las partes en el acto constitutivo;
    \item por acuerdos posteriores; y
    \item si no lo hubieran fijado, por \textbf{resolución judicial}, tomando en cuenta los \textbf{valores incorporados al inmueble en los últimos dos años} y \textbf{sin descontar amortizaciones}.
\end{enumerate}
\end{importante}

La aplicación práctica presenta complejidad: resulta difícil prever, con décadas de anticipación (por ejemplo, 50 o 70 años), una indemnización que contemple todo lo incorporado a lo largo del tiempo. Por esa razón se trata de un derecho que probablemente genere mucho interés en el futuro, dado que será difícil que se constituya por un plazo breve y, en general, se contempla el plazo máximo legal.

\subsection{Normas supletorias}

\begin{table}[H]
\centering
\begin{tabularx}{\textwidth}{@{}>{\raggedright\arraybackslash}p{0.38\textwidth}>{\raggedright\arraybackslash}X@{}}
\toprule
\textbf{Tipo de superficie} & \textbf{Normas supletorias} \\
\midrule
Derecho a construir, plantar o forestar (recae sobre un derecho) & Normas de las limitaciones de uso y goce del usufructo, en forma supletoria. \\
Propiedad superficiaria (construcciones, forestaciones o plantaciones ya existentes) & Normas del espíritu del instituto y, supletoriamente, las normas del dominio revocable. \\
\bottomrule
\end{tabularx}
\end{table}

\begin{nota}
Dado que no existe experiencia previa con esta figura, será necesario que transcurra el tiempo y que se abra paso a nuevas modalidades negociales que pongan en funcionamiento este derecho real.
\end{nota}

\section{Utilidad profesional de la figura}

El material señala que la comprensión de la superficie permite redactar escrituras públicas de constitución sobre inmuebles urbanos, rurales o forestales; asesorar a desarrolladores inmobiliarios que deseen construir en terrenos ajenos sin comprarlos (modelo BTS, \textit{Build-to-Suit}); estructurar fideicomisos o emprendimientos de propiedad horizontal levantados sobre suelo de terceros, algo cada vez más frecuente en la Ciudad de Buenos Aires; registrar y oponer \textit{erga omnes} el derecho ante el Registro de la Propiedad Inmueble; y calcular y negociar la indemnización al superficiario cuando el derecho se extingue (art.~2126).

\cornell{superficie}{
\cq{¿Cómo clasifica el art.~1888 a la superficie y cuáles son sus caracteres?}{Es derecho sobre cosa propia cuando existe propiedad superficiaria y sobre cosa ajena cuando recae sobre el derecho de plantar, construir o forestar. Es principal, se ejerce por la posesión, constituye carga sobre el dominio, se adquiere por título y modo, recae sobre inmueble ajeno, exige escritura pública y es oponible \textit{erga omnes} desde la inscripción.}
\cq{¿Qué novedad plantea el CCCN y cuáles son las dos modalidades del objeto (arts.~2114 y 2115)?}{Permite que el objeto de un derecho real sea un derecho. Puede ser (a) el derecho real de construir, plantar o forestar, haciendo propio lo construido o plantado, o (b) la propiedad de lo ya construido o plantado, sin adquirir el suelo.}
\cq{¿Quiénes pueden constituirla y cómo se adquiere?}{El titular de dominio, los condóminos en su conjunto y el propietario en propiedad horizontal. Se adquiere por título suficiente y modo, a título oneroso o gratuito, entre vivos o \textit{mortis causa}; nunca por prescripción adquisitiva.}
\cq{¿Qué facultades tiene el superficiario?}{Posee, constituye derechos reales de garantía sobre el objeto de su derecho, puede hipotecar el derecho de construir, plantar o forestar o la propiedad superficiaria, afectar lo construido a propiedad horizontal (salvo pacto en contrario) y transmitir su derecho.}
\cq{¿Qué imperfección tiene el dominio del propietario del suelo?}{No puede turbar el derecho del superficiario y, si enajena el inmueble, lo transmite con la carga de la superficie.}
\cq{¿La destrucción del objeto extingue el derecho (art.~2122)?}{No. El superficiario debe reconstruir en seis años o plantar o forestar en tres; si no cumple, opera la extinción.}
\cq{¿Cuáles son los modos de extinción y qué ocurre con los terceros?}{Vencimiento del plazo, renuncia expresa, condición resolutoria, consolidación y no uso (diez años para construir; cinco para plantar o forestar). Si se extingue por vencimiento, el propietario adquiere lo edificado libre de cargas; si es anticipada, los derechos de terceros subsisten hasta el plazo original.}
\cq{¿Cómo se fija la indemnización (art.~2126) y qué normas se aplican supletoriamente?}{Se indemniza al superficiario, salvo pacto en contrario, por lo fijado en el acto constitutivo, por acuerdos posteriores o por el juez (valores incorporados en los últimos dos años, sin descontar amortizaciones). Supletoriamente, usufructo para el derecho de construir, plantar o forestar; dominio revocable para la propiedad superficiaria.}
}{La superficie, regulada en los arts.~2114 a 2126, separa la propiedad del suelo de la de lo construido, plantado o forestado. Admite dos modalidades según su objeto y se adquiere por título y modo, nunca por prescripción, con escritura pública e inscripción para su oponibilidad. El superficiario posee, puede gravar y transmitir su derecho; el propietario conserva un dominio imperfecto. Se extingue por vencimiento, renuncia, condición resolutoria, consolidación o no uso; la destrucción del objeto no la extingue de inmediato. Al extinguirse, el propietario adquiere lo edificado o plantado y debe indemnizar al superficiario, salvo pacto en contrario.}

\chapter{Servidumbre}\label{cap:servidumbre}
% ============================================================

La servidumbre es el derecho real de disfrute sobre cosa ajena que, a diferencia de los anteriores, no vincula a una persona con una cosa sino a \textbf{dos inmuebles entre sí}. Este capítulo presenta su concepto, la forma en que se clasifica, cómo se constituye y quién puede hacerlo, los derechos y obligaciones de los titulares de ambos fundos, y su extinción.

\section{Concepto y objeto}

La servidumbre es un \textbf{derecho real de disfrute sobre cosa ajena} que \textbf{vincula dos fundos entre sí}. A partir de su constitución, el dominio del titular de uno de los fundos se \textbf{imperfecciona}: el dominio, que era absoluto, pasa a ser un \textbf{dominio imperfecto}, porque la constitución de la servidumbre \textbf{desmembra el dominio} de ese titular.

\begin{definicion}{Servidumbre (art.~2162 CCCN)}
La servidumbre es el derecho real que se establece entre dos inmuebles y que concede al titular del inmueble dominante determinada utilidad sobre el inmueble del fondo sirviente, que es el inmueble ajeno.
\end{definicion}

\subsection{Fundo dominante y fundo sirviente}

Para presentar la figura se utiliza un esquema en el que dos fundos (campos o potreros) se identifican como Fundo A y Fundo B, con una parte del terreno delimitada y la Ruta Nacional N.º~2 como vía pública. Se lo emplea como recurso de presentación y no porque los fundos deban ser necesariamente agrícolas. En ese esquema el \textbf{Fundo A} es el \textbf{fundo dominante} y el \textbf{Fundo B} el \textbf{fundo sirviente}.

% TODO: imagen del esquema (Fundo A, Fundo B, Ruta Nacional N.º 2) no disponible en las notas originales; se conserva su descripción textual
\begin{nota}
La imagen del esquema no se encuentra disponible en las fuentes originales; se conserva su descripción textual.
\end{nota}

Ambos fundos se vinculan a partir de la constitución del derecho real, que implica que el titular del fundo sirviente deberá prestar a favor del otro fundo cierta \textbf{utilidad}. Esa utilidad puede consistir en pasar, transitar o sacar agua, e incluso puede ser de mero recreo. En todos los casos se traduce en la obligación de \textbf{dejar que} el titular del fundo dominante haga algo, o de \textbf{abstenerse} de realizar determinada acción o actividad que el titular del dominio tendría si no existiese la servidumbre.

\begin{importante}
\textbf{En ningún caso puede tratarse de una obligación de hacer.} La carga siempre consiste en una abstención o en permitir que el otro haga.
\end{importante}

\subsection{Objeto}

El objeto de la servidumbre puede ser todo el terreno o una parte determinada de él (por ejemplo, si se tratara de transitar, la franja delimitada en el esquema). Siempre se trata de \textbf{dos inmuebles con distinta titularidad}, que tienen por objeto una parte material determinada o todo el inmueble.

\section{Clasificación de las servidumbres}

\subsection{Positivas y negativas}

El tipo de utilidad permite clasificar las servidumbres en positivas y negativas.

\begin{table}[H]
\centering
\begin{tabularx}{\textwidth}{@{}>{\raggedright\arraybackslash}X>{\raggedright\arraybackslash}X@{}}
\toprule
\textbf{Positivas} & \textbf{Negativas} \\
\midrule
La carga consiste en \textbf{soportar} el ejercicio de una acción. & El titular del fundo sirviente debe \textbf{abstenerse} de realizar alguna facultad. \\
\textit{Ejemplo:} soportar que pasen por todo el fundo o por un sector determinado de él. & \textit{Ejemplo:} servidumbre de vista; el titular sirviente no puede construir y afectar la vista del fundo dominante. \\
\bottomrule
\end{tabularx}
\end{table}

\subsection{Reales y personales}

Esta clasificación \textbf{no se refiere a la distinción entre derechos reales y derechos personales}, que son dos esferas diferentes de derechos. Las servidumbres son \textbf{siempre reales}: no existen servidumbres personales como instituto dentro de los derechos reales. Se denomina \textbf{personales} a aquellas servidumbres (siempre reales) constituidas, ya sea por voluntad de las partes o por ley, a favor de una persona determinada, teniendo en cuenta quién es el titular del fundo dominante. En cambio, la \textbf{servidumbre real} se constituye teniendo en cuenta el fundo, con independencia de quién sea su titular.

% TODO: verificar consistencia entre fuentes originales: se afirma que la servidumbre personal puede constituirse «por voluntad de las partes o por ley», pero más adelante se la describe como siempre convencional.

\subsection{Forzosas (de origen legal)}

En el caso de la servidumbre real puede tratarse de fundos encerrados que no tienen acceso a la vía pública, o de fundos que carecen de agua. En estos casos la servidumbre es de \textbf{origen legal}, porque el orden público interviene rigurosamente y la impone. El orden público no admite, por resultar desfavorable para todos y para la comunidad, un fundo que carezca de toda utilidad. El ejemplo concreto es la servidumbre de paso o de tránsito.

\begin{ejemplo}{Fundo encerrado (servidumbre legal de tránsito)}
Se tienen dos fundos. Uno de ellos no tiene salida a la vía pública (en el ejemplo, la Ruta Nacional N.º~2). Sus titulares no podrían ingresar y, si lograran hacerlo por algún medio aéreo, tampoco podrían salir: no tendría sentido. Lo mismo ocurre con los fundos que carecen de agua.

La servidumbre de origen legal impone necesariamente al titular del fundo que tiene una salida más directa a la vía pública tolerar y permitir que el titular del fundo dominante pase.
\end{ejemplo}

\begin{importante}
La mirada del orden público está puesta en la \textbf{utilidad del fundo}, no en la calidad personal del titular de dominio. Cualquiera sea el titular, aun cuando mute el titular registral por enajenación o por transmisión \textit{mortis causa} a herederos, tendrá que poder pasar.
\end{importante}

\subsection{Personales convencionales}

Se analiza ahora una servidumbre de tránsito también, pero personal. En el esquema, el inmueble A tiene salida a la vía pública por un camino lateral, un camino rural que bordea los campos y recién a muchos kilómetros sale a la ruta. No puede imponerse al titular del otro fundo que permita que A pase por su fundo solo porque le resulte más cómodo o el acceso sea más directo: el orden público solo impone la servidumbre cuando \textbf{no existe salida a la vía pública o no existe acceso al agua}.

\begin{ejemplo}{Servidumbre personal convencional}
El titular de A nunca tuvo dificultades y le resultó suficiente salir por el camino existente, con su caballo o su vehículo, hacia la ruta. Un día decide cambiar su actividad o vende, y el adquirente quiere realizar una actividad que impone un gran tránsito de camiones; ese camino le resulta poco práctico y además debe recorrer muchos kilómetros para llegar a la vía pública.

Entonces, el adquirente pacta con el titular del fundo sirviente la constitución de una servidumbre: le propone un camino determinado por el cual pueda pasar con sus camiones. Hay libertad y consenso; no hay imposición legal, es exclusivamente convencional. Las partes pactarán el precio y la duración.

Si no pudieran ponerse de acuerdo en la indemnización, A podría solicitarla judicialmente y el precio sería fijado por el juez. Nada de esto ocurre con la servidumbre personal cuando el orden público está satisfecho: la salida a la vía pública existe y la servidumbre es del interés de A y, si le cierra la ecuación económica, también del interés de B.
\end{ejemplo}
% TODO: verificar consistencia entre fuentes originales: se afirma que no hay imposición legal y que la servidumbre es exclusivamente convencional, pero también que, sin acuerdo sobre la indemnización, el juez fijaría el precio; más adelante se afirma que la servidumbre personal nunca puede constituirse judicialmente.

\subsection{Cuadro comparativo}

\begin{table}[H]
\centering
\small
\begin{tabularx}{\textwidth}{@{}>{\raggedright\arraybackslash}p{0.22\textwidth}>{\raggedright\arraybackslash}X>{\raggedright\arraybackslash}X@{}}
\toprule
\textbf{Aspecto} & \textbf{Real (forzosa)} & \textbf{Personal (convencional)} \\
\midrule
Origen & Legal; orden público & Convencional; voluntad de las partes \\
Requisito & Fundo encerrado o sin agua & Solo el interés de las partes \\
Camino & El más directo a la vía pública & El que acuerden las partes \\
Precio & Fijado por el juez si no hay acuerdo & Libremente pactado \\
Transmisión & Con el inmueble (muta el titular) & Se extingue al transmitir \\
Duración & Mientras subsiste la utilidad & Temporal; la acordada \\
Plazo máximo sin estipular & Indeterminado (mientras sea útil) & Vitalicia (persona humana); 50 años (persona jurídica) \\
\bottomrule
\end{tabularx}
\end{table}
% TODO: verificar consistencia entre fuentes originales: el cuadro equipara «real» con «forzosa» y «personal» con «convencional», mientras que el desarrollo presenta estas categorías como clasificaciones distintas.

\section{Constitución y legitimación}

\subsection{Forma}

La servidumbre siempre debe constituirse mediante \textbf{escritura pública}, porque se trata de un derecho real constituido sobre inmuebles. La escritura pública es el \textbf{título suficiente} de constitución.

\subsection{Duración de la servidumbre personal}

Como la servidumbre personal es siempre convencional en su origen, tendrá una duración y será necesariamente temporal. Si no se determina el plazo en el instrumento constitutivo, cuando se trata de persona humana se considera \textbf{vitalicia}.

\begin{regla}{Plazo de la servidumbre personal (art.~2165 CCCN)}
Se presume que la servidumbre personal es vitalicia cuando no se ha determinado un plazo en el instrumento constitutivo. En el caso de persona jurídica, el plazo máximo es de cincuenta (50) años.
\end{regla}

También en la servidumbre forzosa (tránsito, fondos encerrados, acueducto, sacar agua) el titular del fundo sirviente tiene derecho a recibir la indemnización y a no sufrir perjuicios.

\begin{importante}
En ningún caso la servidumbre personal puede ser constituida judicialmente.
\end{importante}

\subsection{Quién puede constituirla}

Puede constituir la servidumbre el \textbf{titular de dominio} y, en caso de condominio, \textbf{todos los condóminos} que totalicen el cien por ciento de la titularidad de la cosa.

\paragraph{Nudo propietario y usufructo.} Si el titular del fundo B hubiese constituido un usufructo, no puede a su vez constituir una servidumbre: el nudo propietario solo tiene el poder de disposición, y la servidumbre implica un menoscabo a ese derecho, que él mismo ha desplazado a favor del usufructuario. Cuando existe usufructo, quien tiene la facultad es el \textbf{usufructuario}, y la servidumbre estará vinculada con el tiempo que dure el usufructo, que es siempre un derecho temporal.

\paragraph{Uso, habitación y anticresis.} Si el nudo propietario hubiese constituido, en lugar de un usufructo, un derecho de uso, de habitación o de anticresis, podrá constituir la servidumbre:
\begin{itemize}
    \item con la autorización del habitador, del usuario o del acreedor anticresista; o
    \item constituyéndola de modo que comience a regir como condición necesaria cuando se hayan extinguido el uso, la habitación o la anticresis.
\end{itemize}

\section{Transmisión, derechos y obligaciones}

\subsection{Transmisión}

Las servidumbres \textbf{reales}, que vinculan los fundos entre sí, se transmiten conjuntamente con el inmueble, ya enajene el inmueble el titular del fundo dominante o el del sirviente. Como los vinculados son los fundos, quienquiera que sea el titular de uno u otro, la servidumbre se transmite con independencia de la mutación del titular.

\begin{regla}{Transmisión (art.~2172 CCCN)}
Las servidumbres reales se transmiten conjuntamente con el inmueble al que acceden, tanto la constituida a su favor como la que lo grava. La servidumbre personal se extingue con la transmisión del inmueble.
\end{regla}

Los adquirentes reciben el inmueble con la servidumbre. La \textbf{personal}, en cambio, se extingue con la transmisión, porque tanto en su constitución como en su duración se tuvo en cuenta a esa persona determinada, que es el titular de A.

\subsection{Fundo dominante}

\begin{regla}{Fundo dominante (arts.~2173 a 2175 CCCN)}
El titular del fundo dominante puede realizar en el inmueble sirviente todas las obras necesarias para el ejercicio y conservación de la servidumbre. Tiene a su cargo las mejoras necesarias, como mantener el camino en condiciones. Debe comunicar al titular del fundo sirviente cualquier perturbación, de hecho o de derecho, de que tome conocimiento, para que pueda ejercer la defensa.
\end{regla}

El titular del fundo dominante tiene la utilidad, pero debe realizar todas las mejoras necesarias y puede hacer en el inmueble sirviente todo aquello que sea necesario para que su derecho no se vea menoscabado: si se trata de un camino, debe repararlo y mantenerlo en condiciones; si se hubiese producido un hecho natural, como la caída de árboles que impidan salir, puede quitarlos; y debe comunicar al sirviente cualquier perturbación, de hecho o de derecho, para que pueda ejercer la defensa.

\subsection{Fundo sirviente}

\begin{regla}{Fundo sirviente (arts.~2180 y 2181 CCCN)}
El titular del fundo sirviente conserva la disposición jurídica y material del inmueble. Puede ejercer todas las facultades propias de su derecho de dominio, con el límite de no afectar el ejercicio de la servidumbre. Puede realizar obras, constituir otros derechos reales y mejoras, siempre que no impidan o menoscaben el derecho del titular del fundo dominante.
\end{regla}

Al titular del fundo sirviente, a quien en la servidumbre de origen legal se le impone una carga como el acueducto o el tránsito, le quedan todas las facultades materiales y jurídicas propias del dominio. El \textbf{límite} es que no puede realizar nada que afecte el derecho del titular del fundo dominante. La imperfección de su dominio desmembrado consiste en que no puede constituir otros derechos reales ni personales que de algún modo afecten el goce (la utilidad) que el titular del fundo dominante recibe de la servidumbre.

% TODO: verificar consistencia entre fuentes originales: en el cuadro Cornell de una de las versiones la prohibición de constituir derechos que afecten la utilidad figura entre las obligaciones del fundo dominante; en el desarrollo se la refiere al titular del fundo sirviente.

\subsection{Inembargabilidad}

\begin{regla}{Inembargabilidad (art.~2178 CCCN)}
La servidumbre no puede ser embargada ni vendida con independencia del inmueble al que accede, ya sea el inmueble dominante o el sirviente.
\end{regla}

A diferencia de otros derechos reales ya estudiados, el art.~2178 indica con claridad que en ningún caso la servidumbre en sí misma es ejecutable: no puede ser embargada ni vendida con independencia del fundo que grava o del fundo que beneficia.
% TODO: contenido ambiguo en las notas originales: se menciona como ejemplo de derecho ejecutable el caso «de los frutos», sin mayor precisión.

\section{Extinción}

Cuando se trata de fundos encerrados o sin agua, la servidumbre se mantiene mientras subsistan esas condiciones.

\begin{importante}
\textbf{Causales de extinción (art.~2182 CCCN).} Son causales de extinción de la servidumbre: la desaparición de toda utilidad para el fundo dominante; la falta de uso por diez (10) años; en las personales, la muerte del titular cuando no se ha fijado un plazo (o, cuando sea persona jurídica, el vencimiento del plazo máximo de cincuenta (50) años); y el vencimiento del plazo fijado en el instrumento constitutivo.
\end{importante}

\begin{itemize}
    \item \textbf{Desaparición de toda utilidad para el fundo dominante.} Si el fundo está encerrado y necesita salir a la vía pública, carecería de sentido mantener la servidumbre si luego se construyera un camino lateral: el interés del orden público desaparece y la servidumbre carece de objeto.
    \item \textbf{Falta de uso por diez años.}
    \item \textbf{Muerte del titular} en las servidumbres personales, cuando no se haya fijado un plazo: la servidumbre personal es siempre temporal, y lo temporal por excelencia es la vida humana.
    \item \textbf{Vencimiento del plazo máximo de cincuenta años} en el caso de persona jurídica, cuando no se ha fijado un plazo especial.
    \item \textbf{Vencimiento del plazo} fijado en el instrumento constitutivo.
\end{itemize}

Al extinguirse la servidumbre se extinguen todos los derechos personales o reales que pudieran estar vinculados con ella. El titular del fundo sirviente restablece su derecho: pasa a tener nuevamente un \textbf{dominio pleno} (también llamado \textbf{dominio perfecto}), porque ha desaparecido la carga que desmembraba su derecho.

\section{Síntesis}

La servidumbre es un derecho real que siempre implica la vinculación entre dos fundos. En las \textbf{servidumbres forzosas o legales} se tiene en cuenta la calidad del fundo con independencia del titular; deben constituirse por imperio de la ley, más allá de la voluntad del titular del fundo sirviente, porque el orden público ha puesto allí su mirada. Las \textbf{servidumbres personales} son siempre reales, pero se las llama personales porque en su constitución se tiene en cuenta la persona del titular; por eso son convencionales: las partes acuerdan el monto de la indemnización, la manera de pasar por el inmueble y los derechos y obligaciones que surgen de la situación consensual que las vincula durante el tiempo que hayan fijado.

\section{Utilidad profesional}

El material señala que el régimen de servidumbres se vincula con la práctica notarial y de derechos reales: redactar escrituras constitutivas (paso, acueducto, vistas) con cláusulas precisas sobre el fundo dominante, el sirviente, la utilidad concedida y el precio; asesorar al adquirente de un inmueble, que asume automáticamente la servidumbre real constituida; patrocinar reclamos judiciales de servidumbre forzosa de tránsito o acueducto cuando un fundo queda sin salida a la vía pública o sin acceso al agua, fijando la indemnización con el juez; e identificar el momento de extinción (falta de uso por 10 años, muerte del titular en las personales, desaparición de la utilidad) para cancelar la registración correspondiente.

\cornell{servidumbre}{
\cq{¿Qué es la servidumbre, qué establece el art.~2162 y por qué imperfecciona el dominio?}{Es un derecho real de disfrute sobre cosa ajena que vincula dos inmuebles de distinta titularidad: el dominante recibe una utilidad y el sirviente la presta. Su constitución desmembra el dominio del sirviente, que de absoluto pasa a imperfecto. Puede recaer sobre todo el inmueble o sobre una parte determinada.}
\cq{¿En qué se traduce la utilidad para el fundo sirviente?}{En dejar que el titular del dominante haga algo o en abstenerse de una acción que tendría si no existiese la servidumbre. Nunca es una obligación de hacer.}
\cq{¿Cuándo es positiva y cuándo negativa?}{Positiva: se soporta el ejercicio de una acción (que pasen por el fundo). Negativa: el sirviente debe abstenerse de una facultad (servidumbre de vista: no construir para no afectar la vista del dominante).}
\cq{¿Qué diferencia hay entre servidumbre real y personal, y entre forzosa y convencional?}{Las servidumbres son siempre reales como derechos. La real atiende al fundo y se transmite con el inmueble; la personal atiende a la persona del titular dominante y se extingue al transmitirse. La forzosa se impone por orden público a fundos encerrados o sin agua; la convencional nace del acuerdo sobre camino, precio y plazo.}
\cq{¿Puede imponerse una servidumbre solo porque resulta más cómoda?}{No. El orden público solo interviene si no hay salida a la vía pública o acceso al agua.}
\cq{¿Qué forma, quién puede constituirla y qué plazo tiene la personal sin plazo (art.~2165)?}{Escritura pública. Pueden hacerlo el titular de dominio o todos los condóminos; no el nudo propietario durante el usufructo (sí el usufructuario, por el tiempo del usufructo); con uso, habitación o anticresis, con autorización o con efecto diferido a su extinción. Sin plazo, es vitalicia (persona humana) o de hasta cincuenta años (persona jurídica), y nunca puede constituirse judicialmente.}
\cq{¿Cómo se transmite y qué derechos y obligaciones tienen los titulares (arts.~2172 a 2181)?}{Las reales se transmiten con el inmueble; la personal se extingue. El dominante realiza las obras necesarias, soporta las mejoras necesarias y comunica perturbaciones. El sirviente conserva las facultades del dominio con el límite de no afectar la servidumbre.}
\cq{¿Es ejecutable la servidumbre (art.~2178)?}{No: no puede ser embargada ni vendida con independencia del inmueble al que accede.}
\cq{¿Cuáles son las causales de extinción (art.~2182) y sus efectos?}{Desaparición de toda utilidad para el dominante; falta de uso por diez años; muerte del titular en las personales sin plazo; vencimiento del plazo máximo de cincuenta años en persona jurídica; vencimiento del plazo pactado. Se extinguen los derechos vinculados y el dominio sirviente recupera su plenitud.}
}{La servidumbre vincula dos fundos de distinta titularidad: el dominante recibe una utilidad y el sirviente la presta mediante un permitir o abstenerse, nunca un hacer. Se clasifica en positivas y negativas, reales y personales, forzosas y convencionales. Se constituye por escritura pública, la transmiten con el inmueble las reales y se extingue la personal, es inembargable con independencia del fundo (art.~2178) y se extingue por las causales del art.~2182, tras lo cual el dominio sirviente recupera su plenitud.}

% ============================================================
% ============================================================


% ============================================================
\part{Derechos reales de garantía}
% ============================================================

\chapter{Fundamentos históricos y conceptuales de las garantías}\label{cap:garantias-fund}
% ============================================================

Los derechos reales de garantía responden a una pregunta anterior a ellos: cómo se reduce el riesgo de que una obligación no se cumpla. Este capítulo parte de ese riesgo, recorre la evolución histórica que condujo al patrimonio del deudor como garantía común de los acreedores y presenta el privilegio como excepción a la igualdad entre ellos. Los capítulos siguientes estudian los principios comunes de las garantías reales (arts.~2184 a 2204) y, luego, la hipoteca, la anticresis y la prenda.

\begin{nota}
\textbf{Relevancia profesional.} Las garantías reales se presentan de manera constante en la práctica jurídica: la compra de una propiedad con crédito hipotecario, el financiamiento empresarial con garantía prendaria o la herencia de un bien gravado exigen explicar, redactar, ejecutar o defender una garantía. Distinguir la \textbf{extinción} de la \textbf{cancelación} puede evitar que un cliente pierda un bien por una hipoteca ya pagada pero no cancelada en el registro. Comprender el principio de \textbf{especialidad} permite redactar garantías que no sean atacadas de nulidad en juicio. Dominar la figura del \textbf{propietario no deudor} permite asesorar correctamente a quien garantiza deudas ajenas.
\end{nota}

\section{El riesgo inherente a la obligación}

El estudio de las garantías se vincula con la necesidad de \textbf{minimizar el riesgo de incumplimiento obligacional}. Existe una relación de \textbf{inherencia} entre el riesgo y la obligación, porque el incumplimiento puede originarse en infinitos factores. Pueden impedir el cumplimiento el caso fortuito, la fuerza mayor y la finitud humana, y también los deudores que, teniendo posibilidad de cumplir, deciden no hacerlo. Todos estos supuestos ratifican que la obligación, desde el momento mismo de su nacimiento, lleva implícito el riesgo del incumplimiento.

\section{Evolución histórica: de la esclavitud por deudas al patrimonio}

La garantía real por excelencia en la antigüedad fue la \textbf{esclavitud por deudas}. El hombre libre, con la categoría que tenía en ese momento histórico, era obligado a someterse a la esclavitud, por sí o a través de su familia. El esclavo era una cosa, susceptible de ser usada, comprada y vendida, de modo que su producido permitía resarcir al acreedor. La esclavitud no era un castigo, sino el medio de resarcir al acreedor.

La reforma histórica más importante fue la abolición de la esclavitud por deudas decretada por el legislador \textbf{Solón} en el siglo VI a.C., considerada tal vez la mayor reforma en materia de garantías en la historia de la humanidad. Dos o tres siglos más tarde, Roma hizo lo propio con la \textbf{Ley Poetelia Papiria}, a partir del año 326 o 313 a.C.
% TODO: contenido incompleto o ambiguo en las notas originales (el año de la Ley Poetelia Papiria se indica como 326 o 313 a.C.)

A partir de esa abolición surgió la necesidad de encontrar otro mecanismo que garantizara el cumplimiento. La respuesta fue el \textbf{patrimonio del deudor} como garantía común de los acreedores. En esa misma línea, el mundo occidental decidió prohibir la tortura, los azotes, la esclavitud por deudas y la prisión del deudor incumplidor, y desarrolló la \textbf{teoría del patrimonio}: el patrimonio del deudor responde a todas las obligaciones.

\section{El patrimonio como garantía común (arts.~242 y 743)}

\begin{quote}
\textbf{Artículo 242 CCyCN.} \textit{Todos los bienes del deudor están afectados al cumplimiento de sus obligaciones y constituyen la garantía común de sus acreedores, con excepción de aquellos bienes que el Código indique.}
\end{quote}
% TODO: verificar consistencia entre fuentes originales: el material sobre protección de la vivienda transcribe el art. 242 como «Todo el patrimonio del deudor es garantía común de todas sus obligaciones», con una redacción distinta de la citada en el material sobre garantías.

Este artículo incorpora \textbf{por primera vez} la garantía de los acreedores como norma expresa en el derecho nacional; antes se la infería del ordenamiento general. Las excepciones son los bienes esenciales para la vida de la persona y para el desarrollo de su actividad profesional o laboral. Sin firmar nada en especial, solo por contraer obligaciones, todo el patrimonio queda afectado a su cumplimiento y es susceptible de ser embargado y ejecutado, incluso el inmueble en el que se vive. Esa consecuencia explica el sistema de protección de la vivienda que se estudia en el capítulo~\ref{cap:vivienda}.

\begin{quote}
\textbf{Artículo 743 CCyCN.} \textit{Los bienes presentes y futuros del deudor constituyen la garantía común de los acreedores.}
\end{quote}

No solo responden los bienes que el deudor tiene al constituirse la obligación: si el deudor adquiere un bien en el futuro y aún no ha cumplido, ese bien también respalda la obligación. Todos los acreedores concurren en una posición igualitaria. El concepto de \textbf{acreedor} no se limita al banco: en la obligación alimentaria, por ejemplo, el acreedor es el hijo que se alimenta, y todo el patrimonio del deudor está afectado para garantizar ese crédito.

\section{Garantías personales y garantías reales}

La distinción entre garantías personales y garantías reales reproduce, en el campo de las garantías, la clásica dualidad entre derechos personales y derechos reales.

\begin{table}[H]
\centering
\begin{tabularx}{\textwidth}{@{}LL@{}}
\toprule
\textbf{Garantías personales} & \textbf{Garantías reales} \\
\midrule
Ejemplo típico: la fianza. Una persona distinta del deudor garantiza la obligación con todo su patrimonio. & Ejemplo típico: la hipoteca. Un bien determinado queda afectado al cumplimiento de una obligación. \\
\addlinespace
El patrimonio del garante puede crecer o disminuir; no hay forma de garantizar su inmutabilidad. & La cosa determinada queda sujeta al privilegio del acreedor. Este derecho es \textbf{preferente y persecutorio}. \\
\bottomrule
\end{tabularx}
\end{table}

Dentro de la clasificación de los derechos reales estudiada en el capítulo~\ref{cap:marco}, los derechos de disfrute son \textbf{principales}, al igual que los derechos sobre cosa propia; los derechos reales de garantía son, en cambio, siempre y en todos los casos \textbf{accesorios de una obligación}.

\section{Igualdad entre acreedores y privilegio}

La \textbf{igualdad entre acreedores es de orden público}. La excepción a ese principio está regulada por la ley y debe surgir de una causa legal o convencional de preferencia; en el caso de la convención, debe estar claramente identificada por la ley.

Los derechos reales de garantía (hipoteca, prenda y anticresis) constituyen un \textbf{privilegio especial} que el titular tiene sobre una cosa de propiedad del deudor. Ese privilegio lo deja a salvo de los vaivenes del patrimonio, de otras deudas y de otras obligaciones que puedan agravar o disminuir su composición.

\cornell{fundamentos de las garantías}{
\cq{¿Por qué el riesgo es inherente a la obligación y qué buscan las garantías?}{Porque el incumplimiento puede deberse a infinitos factores (caso fortuito, fuerza mayor, finitud humana, decisión del deudor de no cumplir). La obligación lleva implícito, desde su nacimiento, el riesgo del incumplimiento, y las garantías procuran minimizarlo.}
\cq{¿Cuál fue la garantía real por excelencia en la antigüedad y cómo se superó?}{La esclavitud por deudas: el hombre libre, por sí o a través de su familia, era sometido a esclavitud para resarcir al acreedor con el producido. Solón la abolió en el siglo VI a.C. y Roma lo hizo con la Ley Poetelia Papiria (326 o 313 a.C.). Surgió entonces el patrimonio del deudor como garantía común.}
\cq{¿Qué establecen los arts.~242 y 743?}{Que los bienes presentes y futuros del deudor constituyen la garantía común de los acreedores, con excepción de los bienes que el Código indique (esenciales para la vida y el trabajo). Los acreedores concurren en igualdad de condiciones.}
\cq{¿Qué diferencia hay entre garantías personales y reales?}{En la personal (fianza) una persona distinta del deudor responde con todo su patrimonio, que puede crecer o disminuir. En la real (hipoteca) un bien determinado queda afectado a la obligación y sujeto al privilegio del acreedor, derecho preferente y persecutorio.}
\cq{¿Cómo se exceptúa la igualdad entre acreedores?}{La igualdad es de orden público; solo puede exceptuarse por una causa legal o convencional de preferencia identificada por la ley. Los derechos reales de garantía son un privilegio especial sobre una cosa del deudor y son siempre accesorios de una obligación.}
}{Las garantías buscan minimizar el riesgo de incumplimiento inherente a toda obligación. Tras la abolición de la esclavitud por deudas, el patrimonio del deudor pasó a ser la garantía común de los acreedores (arts.~242 y 743), con igualdad entre ellos como regla de orden público. Los derechos reales de garantía son accesorios y otorgan un privilegio especial sobre una cosa determinada, a diferencia de las garantías personales.}

\chapter{Caracteres y principios de las garantías reales}\label{cap:garantias-carac}
% ============================================================

Este capítulo estudia el régimen común que el Código Civil y Comercial establece para todas las garantías reales: cuáles son, qué caracteres deben reunir para ser válidas y cómo operan los principios de convencionalidad, accesoriedad, especialidad e indivisibilidad. La reforma de la especialidad crediticia merece un desarrollo propio.

\section{Los tres derechos reales de garantía y su marco normativo}

Los derechos reales de garantía son la \textbf{hipoteca}, la \textbf{anticresis} y la \textbf{prenda}. Existen también garantías reales especiales, como la hipoteca aeronáutica, la hipoteca naval y la prenda de automotores, reguladas en todos los casos por leyes especiales. Se trata de derechos reales accesorios, porque dependen siempre de una obligación principal, pero siguen los principios generales del ordenamiento nacional.

\begin{importante}
\textbf{Normativa aplicable.} A los tres derechos de garantía se aplican las normas generales de todos los derechos reales (arts.~1882 a 1907) y las disposiciones comunes en materia de garantías reales (arts.~2184 a 2204). A partir de allí rigen las normas propias de cada garantía: hipoteca, arts.~2205 a 2211; anticresis, arts.~2212 a 2218; y prenda, que cuenta con disposiciones comunes a sus dos tipos (prenda de cosas y de créditos instrumentados) antes de las normas específicas.
\end{importante}

\section{Caracteres esenciales y naturales}

\begin{table}[H]
\centering
\begin{tabularx}{\textwidth}{@{}LL@{}}
\toprule
\textbf{Caracteres esenciales} & \textbf{Caracteres naturales} \\
\midrule
Su ausencia torna nula la garantía. & La ley los establece, pero las partes pueden pactar en contrario. \\
\addlinespace
1) Convencionalidad. 2) Accesoriedad. 3) Especialidad (objetiva y crediticia). & Indivisibilidad (las partes pueden pactar su divisibilidad). \\
\bottomrule
\end{tabularx}
\end{table}

% TODO: verificar consistencia entre fuentes originales: una de las fuentes incluye la indivisibilidad entre los caracteres esenciales exigidos bajo pena de nulidad (convencionalidad, accesoriedad, especialidad objetiva, especialidad crediticia e indivisibilidad), mientras que otra la presenta como carácter natural, modificable por pacto.

\section{Convencionalidad (art.~2185)}

\begin{quote}
\textbf{Artículo 2185 CCyCN.} \textit{Los derechos reales de garantía solo pueden ser constituidos por contrato celebrado por los legitimados y con las formas que indique la ley.}
\end{quote}

Los derechos reales de garantía solo pueden constituirse por \textbf{contrato}, únicamente por contrato. Otros ordenamientos jurídicos admiten garantías judiciales y garantías legales; el codificador nacional, tanto Vélez como el Código de 2015, optó por la fuente exclusivamente contractual. En materia de hipoteca, o de inmuebles como en la anticresis, por el art.~1017 deberá constituirse siempre en \textbf{escritura pública}.

\section{Accesoriedad (art.~2186)}

\begin{quote}
\textbf{Artículo 2186 CCyCN.} \textit{Los derechos reales de garantía son accesorios del crédito que aseguran. No puede existir un derecho de garantía sin crédito en el que garanticen.}
\end{quote}

\begin{importante}
La garantía no es una obligación: es el accesorio de una obligación.
\end{importante}

Existe un vínculo obligacional con tres elementos: sujeto activo, sujeto pasivo y prestación. De manera subyacente a esa relación se constituye una garantía, es decir, un vínculo directo entre el titular del derecho de garantía y una cosa determinada. Esta garantía se pone en movimiento \textbf{ante el incumplimiento} y otorga al acreedor un derecho de privilegio que recae sobre esa cosa.

\section{Especialidad}

El principio de especialidad tiene un doble aspecto.

\begin{table}[H]
\centering
\begin{tabularx}{\textwidth}{@{}LL@{}}
\toprule
\textbf{Especialidad objetiva} & \textbf{Especialidad crediticia} \\
\midrule
Determinación indubitable del objeto: qué bien específico respalda la obligación. & Determinación indubitable del monto: hasta qué suma el acreedor tiene privilegio. \\
\addlinespace
Art.~2188: la cosa o el derecho debe ser actual e individualizado en el contrato constitutivo. & Art.~2189: debe constar el monto estimado (créditos determinados) o el monto máximo (créditos indeterminados). \\
\bottomrule
\end{tabularx}
\end{table}

\begin{quote}
\textbf{Artículo 2187 CCyCN.} \textit{La garantía real puede garantizar el puro y simple crédito, el crédito a plazo, el condicional, el eventual, obligaciones de dar, de hacer o de no hacer. En todos los casos, al constituirse la garantía debe individualizarse el crédito a través de sujeto, objeto y causa, con las excepciones admitidas por la ley.}
\end{quote}

\begin{quote}
\textbf{Artículo 2188 CCyCN.} \textit{Tanto cosas como derechos pueden constituir el objeto de los derechos reales de garantía. Deben ser actuales y estar individualizados en el contrato constitutivo.}
\end{quote}

\subsection{Importancia práctica de la especialidad crediticia}

\begin{ejemplo}{Juan, Fermín y el inmueble}
Juan es propietario de un inmueble que vale 150.000 dólares. Por el art.~242, todos sus acreedores son quirografarios: si el patrimonio es insuficiente, cobran en proporción a su crédito.

Juan obtiene un crédito de 100.000 dólares y constituye a favor de Fermín una garantía hipotecaria, de modo que Fermín tiene privilegio. Si Juan paga, la garantía se extingue y nada ocurre. Si no paga, Fermín inicia la ejecución. El inmueble se remata en 110.000 dólares: Fermín cobra sus 100.000 con privilegio y los 10.000 restantes pertenecen a Juan, quedando para los demás acreedores quirografarios.
% TODO: contenido incompleto o ambiguo en las notas originales (el título del caso menciona ``el inmueble en CAVA'')

\textbf{Conclusión:} el privilegio ha reducido enormemente la garantía disponible para los acreedores quirografarios. Por eso es tan importante la especialidad crediticia: los quirografarios necesitan saber cuánto es el privilegio y cuál es el remanente.
\end{ejemplo}

\subsection{La reforma de la ley 27261}

El art.~2189 original, vigente desde el 1.º de septiembre de 2015, fue \textbf{sustituido} por la ley 27261, que rige desde el 1.º de septiembre de 2016. Esto genera tres regímenes aplicables según la fecha de constitución de la hipoteca.

\begin{table}[H]
\centering
\begin{tabularx}{\textwidth}{@{}LLL@{}}
\toprule
\textbf{Código de Vélez (hasta el 31/8/2015)} & \textbf{CCyCN original (1/9/2015 al 31/8/2016)} & \textbf{CCyCN reformado por ley 27261 (desde el 1/9/2016)} \\
\midrule
Hipoteca sobre crédito indeterminado: nula (art.~2502 del régimen anterior). & Régimen transitorio: normas del nuevo código pendientes de reforma. & Se admiten garantías sobre créditos indeterminados bajo requisitos rigurosos. Rige el texto sustituido del art.~2189. \\
\bottomrule
\end{tabularx}
\end{table}
% TODO: contenido incompleto o ambiguo en las notas originales (régimen del CCyCN original descrito de forma imprecisa; cita del art. 2502)
% TODO: verificar consistencia entre fuentes originales: el material sobre garantías atribuye la sustitución del art. 2189 a la ley 27261 (vigente desde el 1/9/2016), mientras que el material sobre hipoteca atribuye la sustitución de los arts. 2189 y 2210 a la ley 27.271 (con efectos desde el 15/9/2016).

\begin{quote}
\textbf{Artículo 2189 CCyCN (texto ley 27261, vigente desde el 1/9/2016).} \textit{En la constitución de los derechos reales de garantía debe individualizarse el crédito garantizado, indicándose los sujetos, el objeto y la causa. El monto de la garantía debe estimarse en dinero y puede no coincidir con el monto del capital del crédito. Se considera satisfecho el principio de especialidad en cuanto al crédito si la garantía se constituye en seguridad de créditos indeterminados, sea que su causa exista al tiempo de la constitución o posteriormente, siempre que el instrumento constitutivo contenga la indicación del monto máximo garantizado en todo concepto, que la garantía que se constituye es de máximo, y que los créditos garantizados han de nacer dentro del plazo de diez (10) años a contar de la constitución de la garantía.}
\end{quote}

El nuevo texto habilita dos categorías de créditos.

\begin{table}[H]
\centering
\begin{tabularx}{\textwidth}{@{}LL@{}}
\toprule
\textbf{Créditos determinados (art.~2187 y 1.er párr.\ del art.~2189)} & \textbf{Créditos indeterminados (2.º párr.\ del art.~2189, novedad)} \\
\midrule
Se individualiza la obligación: sujetos, objeto y causa. & La causa puede existir al momento de la constitución o surgir posteriormente. \\
\addlinespace
Se estima el monto en dinero; puede no coincidir con el capital del crédito. & Debe indicarse el \textbf{monto máximo garantizado} en todo concepto. \\
\addlinespace
El privilegio se extiende a capital, intereses, costas y daños posteriores al incumplimiento. & Debe indicarse que es una ``garantía de máximo'' y que los créditos nacerán dentro de los 10 años de su constitución. \\
\addlinespace
\textbf{Alcance amplio del privilegio.} & \textbf{Privilegio solo hasta el monto máximo garantizado; todo excedente es quirografario.} \\
\bottomrule
\end{tabularx}
\end{table}

\begin{importante}
La diferencia entre ``monto estimado'' y ``monto máximo'' no es meramente lingüística: expresa un concepto totalmente diferente en cuanto al alcance que tendrá el privilegio en uno y otro caso.
\end{importante}

\paragraph{Las ``hipotecas abiertas'' como antecedente.} Las hipotecas abiertas no tenían denominación específica porque no tenían recepción en el ordenamiento. El acreedor pretendía una hipoteca que garantizara todas las obligaciones del deudor. Ello generaba gran inseguridad a los quirografarios, que carecían de respuesta a preguntas como: ¿por qué monto es el privilegio?, ¿cuándo se acaba?, ¿cuál es la obligación?, ¿cómo se identifica si el deudor cumplió o no? Ante esa falta de respuestas, los jueces declaraban la nulidad de estas garantías.

\begin{quote}
\textbf{Artículo 2190 CCyCN.} \textit{La constitución de la garantía será válida aunque falte alguna especificación, si puede conocerse indubitablemente cuál es la enunciación completa de los parámetros básicos.}
\end{quote}

\section{Indivisibilidad (art.~2191)}

\begin{definicion}{Indivisibilidad}
Las hipotecas son indivisibles, pero las partes pueden pactar la divisibilidad. La garantía subsiste en cada uno de los bienes afectados a la deuda y en cada parte de ellos: todos están afectados al pago hasta su extinción.
\end{definicion}

\begin{ejemplo}{Crédito en 120 cuotas}
Si se obtiene un crédito bancario en 120 cuotas con garantía hipotecaria, el inmueble sigue gravado tanto cuando se debe la primera cuota como cuando se debe la última. La garantía subsiste hasta la cancelación total.
\end{ejemplo}

Las partes pueden pactar en el instrumento constitutivo que, cuando el deudor haya pagado 50 cuotas, la hipoteca se extinga por ese porcentaje y subsista por la diferencia, pues ello no agravia el orden público.

\begin{ejemplo}{Hipoteca sobre dos inmuebles y 240 cuotas}
En una hipoteca cuyo objeto son dos inmuebles y 240 cuotas, ambos inmuebles están gravados hasta que se cancelen las 240 cuotas; aun para la cuota 239, si se han pagado las anteriores, los dos inmuebles siguen gravados. La garantía y el privilegio no dependen del monto, sino de cómo se va pagando la obligación.
\end{ejemplo}

\cornell{caracteres y principios de las garantías}{
\cq{¿Cuáles son los derechos reales de garantía y qué normas los rigen?}{Hipoteca, anticresis y prenda; también existen la hipoteca aeronáutica, la naval y la prenda de automotores, regidas por leyes especiales. Rigen las normas generales de los derechos reales (arts.~1882 a 1907), las disposiciones comunes de las garantías (arts.~2184 a 2204) y las normas propias de cada una (hipoteca, 2205 a 2211; anticresis, 2212 a 2218).}
\cq{¿Cuáles son los caracteres esenciales y en qué se diferencian de los naturales?}{Esenciales: convencionalidad (art.~2185), accesoriedad (art.~2186) y especialidad objetiva (art.~2188) y crediticia (art.~2189); su ausencia torna nula la garantía. Los naturales están establecidos por la ley, pero pueden pactarse en contrario: es el caso de la indivisibilidad.}
\cq{¿Por qué las garantías solo pueden nacer de un contrato y por qué son accesorias?}{Porque el codificador (Vélez y el Código de 2015) optó por la fuente exclusivamente contractual; en inmuebles, escritura pública (art.~1017). Son accesorias porque no pueden existir sin el crédito que garantizan: la garantía es el accesorio de una obligación y se activa ante el incumplimiento.}
\cq{¿Qué exige la especialidad objetiva y la crediticia?}{Objetiva: determinar indubitablemente el bien, actual e individualizado en el contrato (art.~2188). Crediticia: determinar el monto hasta el cual hay privilegio, individualizando el crédito por sujeto, objeto y causa (arts.~2187 y 2189).}
\cq{¿Por qué es importante la especialidad crediticia?}{Porque el privilegio reduce la garantía disponible para los quirografarios, que necesitan saber cuánto es el privilegio y cuál es el remanente (caso de Juan y Fermín: remate en 110.000 dólares, deuda de 100.000, remanente de 10.000).}
\cq{¿Qué regímenes existen según la fecha de constitución y qué cambió con la ley 27261?}{Código de Vélez (hasta el 31/8/2015: nula la hipoteca sobre crédito indeterminado), CCyCN original (1/9/2015 a 31/8/2016) y CCyCN reformado (desde el 1/9/2016), que admite garantías sobre créditos indeterminados.}
\cq{¿Qué diferencia hay entre monto estimado y monto máximo garantizado?}{El estimado corresponde a créditos determinados y el privilegio alcanza capital, intereses, costas y daños posteriores. El máximo corresponde a créditos indeterminados (garantía de máximo, con créditos que nazcan dentro de diez años) y el privilegio llega solo hasta ese tope; el excedente es quirografario.}
\cq{¿Qué eran las hipotecas abiertas y qué dice el art.~2190?}{Hipotecas con las que el acreedor pretendía garantizar todas las obligaciones del deudor; los jueces las anulaban por la inseguridad que generaban. El art.~2190 admite la validez aunque falte alguna especificación si pueden conocerse indubitablemente los parámetros básicos.}
\cq{¿Qué es la indivisibilidad y cuándo pueden las partes modificarla?}{La garantía subsiste en cada bien afectado y en cada parte de ellos hasta la extinción total de la deuda. Las partes pueden pactar la divisibilidad (por ejemplo, extinción parcial al pagar 50 cuotas), pues no agravia el orden público.}
}{La hipoteca, la prenda y la anticresis comparten una parte general (arts.~2184 a 2204). Su validez exige convencionalidad, accesoriedad y especialidad objetiva y crediticia; la indivisibilidad se presenta como carácter natural modificable por pacto, aunque otra fuente la incluye entre los exigidos bajo pena de nulidad. La reforma del art.~2189 incorporó las garantías sobre créditos indeterminados, con privilegio acotado al monto máximo garantizado y créditos que nazcan dentro de los diez años.}

\chapter{Constituyente, propietario no deudor, extinción y cancelación}\label{cap:garantias-constituyente}
% ============================================================

Constituida la garantía, la cosa gravada sigue perteneciendo a alguien, y ese alguien puede ser el deudor o un tercero. Este capítulo estudia quién puede constituir garantías, qué puede y qué no puede hacer el propietario de la cosa gravada, cómo se protege el acreedor frente a actos que disminuyen la garantía, cuál es la situación del propietario no deudor y de qué manera se extingue y se cancela la garantía.

\section{Legitimación para constituir garantías}

Solo los titulares de derechos reales sobre cosa propia pueden constituir derechos reales de disfrute o de garantía a favor de otro titular. Es excepcional la facultad del titular del derecho de superficie de constituir hipotecas sobre el objeto de su derecho (véase el capítulo~\ref{cap:superficie}). El condómino puede constituir un gravamen hipotecario sobre su parte indivisa. Al constituir el derecho real de garantía, el titular \textbf{imperfecciona su dominio}, lo desmembra; por eso el art.~1884 indica que para el titular constituyente existe un gravamen.

\section{Facultades del constituyente (art.~2195)}

\begin{quote}
\textbf{Artículo 2195 CCyCN.} \textit{El constituyente de la garantía conserva todas las facultades inherentes a su derecho, pero no puede realizar ningún acto que disminuya el valor de la garantía. Si esto ocurre, el acreedor puede requerir la privación de plazo, o bien estimar el valor de la disminución y exigir su depósito, o bien que se otorgue otra garantía suficiente.}
\end{quote}

Se distinguen dos situaciones.

\begin{table}[H]
\centering
\begin{tabularx}{\textwidth}{@{}LL@{}}
\toprule
\textbf{Acto no consumado} & \textbf{Acto consumado} \\
\midrule
El acreedor advierte que el constituyente está por realizar un acto que puede afectar el valor de la garantía. & El deterioro o acto ya se realizó y el perjuicio está consumado. \\
\addlinespace
\textbf{Remedio:} medida de no innovar para detener la disminución del valor. & \textbf{Remedio:} caducidad de plazos (exigir el pago total aun sin mora en los servicios), depósito del valor disminuido o nueva garantía suficiente. \\
\bottomrule
\end{tabularx}
\end{table}

\section{Inoponibilidad de los actos que perjudican la garantía (art.~2196)}

\begin{quote}
\textbf{Artículo 2196 CCyCN.} \textit{En caso de ejecución, son inoponibles al acreedor los actos jurídicos celebrados en perjuicio de la garantía.}
\end{quote}

\begin{table}[H]
\centering
\begin{tabularx}{\textwidth}{@{}LL@{}}
\toprule
\textbf{Locación anterior a la constitución de la hipoteca} & \textbf{Locación posterior a la constitución de la hipoteca} \\
\midrule
El adquirente en la subasta adquiere un inmueble ya alquilado y debe respetar el plazo de locación. & La locación carece de todo efecto frente al acreedor garantizado; la subasta se realizará con el inmueble libre de ocupantes. \\
\bottomrule
\end{tabularx}
\end{table}

\section{Prohibición del pacto comisorio (art.~2198)}

\begin{quote}
\textbf{Artículo 2198 CCyCN.} \textit{Es nula la cláusula que permite al titular del derecho de garantía disponer de la cosa por otros medios que los previstos legalmente para la ejecución forzada.}
\end{quote}

\begin{definicion}{Pacto comisorio}
Cláusula por la cual el acreedor, al constituirse la garantía, establece que si no se cumple la obligación se quedará directamente con la cosa. Se tiene por no escrita.
\end{definicion}

\begin{importante}
La nulidad recae sobre la cláusula específica (el pacto comisorio) y no sobre toda la garantía. Esto contrasta con los caracteres esenciales (convencionalidad, accesoriedad, especialidad), cuya falta torna nula la garantía íntegramente.
\end{importante}

Si la obligación se incumple, el acreedor debe iniciar el proceso de ejecución judicial. El juez, con plena intervención del constituyente y del deudor, determinará el monto de la obligación, el privilegio, la subasta y todo el proceso, lo que impide la ``justicia por mano propia''.

\section{El propietario no deudor (arts.~2199 y 2200)}

En toda garantía real pueden identificarse distintos protagonistas.

\begin{table}[H]
\centering
\begin{tabularx}{\textwidth}{@{}cLL@{}}
\toprule
\textbf{Escenario} & \textbf{Actores} & \textbf{Observación} \\
\midrule
1 & Juan = deudor y constituyente; Manuel = acreedor hipotecario. & El propietario del bien gravado es el mismo deudor. Es el caso más común. \\
\addlinespace
2 & Juan = deudor; Manuel = acreedor; Juan padre = constituyente propietario no deudor. & Tercero ajeno al vínculo obligacional que constituye el gravamen para garantizar la deuda ajena. \\
\addlinespace
3 & Juan = ex deudor y constituyente; Manuel = acreedor; Joaquín = adquirente del inmueble gravado. & Adquirente de un bien gravado que no asumió expresamente la deuda; también es propietario no deudor. \\
\bottomrule
\end{tabularx}
\end{table}

\begin{quote}
\textbf{Artículo 2199 CCyCN.} \textit{Son propietarios no deudores: (a) el tercero ajeno al vínculo obligacional que constituye el gravamen a fin de garantizar la obligación del deudor; (b) el adquirente de un bien gravado que no ha asumido expresamente la deuda.}
\end{quote}

\begin{quote}
\textbf{Artículo 2200 CCyCN.} \textit{El propietario no deudor puede, ante el incumplimiento del deudor: (a) satisfacer el pago de la obligación, subrogándose en los derechos del acreedor; (b) oponer las defensas admisibles en el proceso ejecutivo; (c) si el bien es subastado, ejercer el derecho al remanente que exceda el monto del gravamen.}
\end{quote}

\textbf{Responsabilidad.} El propietario no deudor no tiene obligación de pagar: lo que responde es la cosa. Se subastará la cosa que adquirió, y responde únicamente con el objeto gravado y hasta el monto máximo expresado en el instrumento constitutivo.

\textbf{Buena fe del adquirente.} Quien adquiere un inmueble gravado no es sorprendido en su buena fe, porque los derechos reales inmobiliarios son oponibles a terceros a partir de su inscripción registral. La hipoteca constituida es oponible a terceros desde su inscripción, y la carga surge de las constancias registrales que deben solicitarse antes de realizar la operación.

\textbf{Proceso cuando el deudor no paga:}
\begin{enumerate}
    \item El acreedor intima el pago al obligado (el deudor).
    \item Si el deudor paga, la garantía se extingue.
    \item Si no paga, el juez da traslado al propietario no deudor.
    \item El propietario no deudor puede pagar voluntariamente y subrogarse en los derechos del acreedor para repetir lo pagado contra el deudor principal.
    \item Si no paga, el juez ordena la subasta. El remanente que excede el monto del gravamen pertenece al propietario no deudor, no a los acreedores quirografarios del deudor.
\end{enumerate}

\section{Extensión de la garantía y subrogación real}

Las garantías se extienden sobre el principal y todos sus accesorios. Si se constituye una hipoteca sobre un inmueble, alcanza también a sus cañerías, paredes y techos.

\begin{definicion}{Subrogación real}
La garantía se traslada de pleno derecho sobre los bienes que sustituyan al gravado. Si hay una hipoteca sobre un inmueble con seguro de incendio y el inmueble se incendia, la garantía subsiste y el privilegio recae sobre la póliza. Lo mismo ocurre ante una indemnización por expropiación.
\end{definicion}

\section{Extinción y cancelación (art.~2204)}

La distinción entre extinción y cancelación es fundamental en la práctica profesional.

\begin{table}[H]
\centering
\begin{tabularx}{\textwidth}{@{}LL@{}}
\toprule
\textbf{Extinción de la garantía} & \textbf{Cancelación de la garantía} \\
\midrule
Hecho jurídico que produce el cese del derecho real de garantía. & Acto jurídico en virtud del cual se deja sin efecto la inscripción de la garantía en el registro. \\
\addlinespace
Causas: (a) extinción de la obligación principal; (b) subasta del bien gravado con citación de los titulares; (c) confusión (acreedor = propietario constituyente). & Tipos: (a) \textbf{voluntaria}: el acreedor otorga instrumento de igual naturaleza al constitutivo y manifiesta que se extinguió la garantía; (b) \textbf{judicial}: el constituyente, ante la falta de cancelación, la solicita judicialmente. \\
\bottomrule
\end{tabularx}
\end{table}

\begin{importante}
Una hipoteca puede estar \textbf{extinguida} (porque la obligación se pagó) pero no \textbf{cancelada} (porque la inscripción registral sigue vigente). En ese caso, al solicitar un certificado, el inmueble aparece como gravado, pues falta inscribir que la garantía se extinguió. Esto genera trabajo adicional, y es habitual que las partes confundan ambos conceptos.
\end{importante}

\cornell{constituyente, propietario no deudor, extinción y cancelación}{
\cq{¿Quiénes pueden constituir garantías reales y qué efecto tiene para su dominio?}{Los titulares de derechos reales sobre cosa propia; excepcionalmente, el superficiario sobre el objeto de su derecho, y el condómino sobre su parte indivisa. Al constituirla, el titular imperfecciona su dominio (art.~1884).}
\cq{¿Qué puede y qué no puede hacer el dueño de la cosa gravada (arts.~2195 y 2196)?}{Conserva todas las facultades inherentes a su derecho, pero no puede realizar actos que disminuyan el valor de la garantía. Ante un acto no consumado, medida de no innovar; ante uno consumado, caducidad de plazos, depósito del valor disminuido o nueva garantía suficiente. En la ejecución son inoponibles al acreedor los actos celebrados en perjuicio de la garantía.}
\cq{¿Qué efectos tiene una locación anterior y una posterior a la hipoteca?}{La anterior debe ser respetada por el adquirente en subasta; la posterior carece de efecto frente al acreedor y la subasta se realiza con el inmueble libre de ocupantes.}
\cq{¿Qué es el pacto comisorio y cuáles son sus efectos (art.~2198)?}{La cláusula por la cual el acreedor se quedaría directamente con la cosa ante el incumplimiento. Es nula y se tiene por no escrita, pero el resto de la garantía conserva su validez.}
\cq{¿Quién es propietario no deudor (art.~2199) y qué puede hacer ante la ejecución (art.~2200)?}{El tercero ajeno al vínculo obligacional que constituye el gravamen, y el adquirente de un bien gravado que no asumió expresamente la deuda. Puede pagar y subrogarse, oponer las defensas admisibles en el proceso ejecutivo y reclamar el remanente si el bien se subasta. Responde solo con la cosa y hasta el monto máximo del instrumento constitutivo.}
\cq{¿Hasta dónde se extiende la garantía y qué es la subrogación real?}{Se extiende al principal y a todos sus accesorios. Por subrogación real, se traslada de pleno derecho a los bienes que sustituyan al gravado (póliza de seguro de incendio, indemnización por expropiación).}
\cq{¿Qué diferencia hay entre extinción y cancelación (art.~2204)?}{La extinción es el cese del derecho real de garantía (por extinción de la obligación principal, subasta o confusión). La cancelación es el acto que deja sin efecto la inscripción registral, voluntaria o judicial. Una hipoteca puede estar extinguida pero no cancelada y el inmueble aparecerá gravado en las constancias registrales.}
}{El constituyente conserva las facultades inherentes a su derecho pero no puede disminuir el valor de la garantía; los actos en perjuicio de ella son inoponibles al acreedor y el pacto comisorio es nulo, sin afectar la validez del resto de la garantía. El propietario no deudor responde solo con la cosa y puede pagar y subrogarse, defenderse y reclamar el remanente. La garantía alcanza los accesorios y se traslada por subrogación real. La extinción de la garantía debe distinguirse de su cancelación registral.}

\chapter{Hipoteca, anticresis y prenda}\label{cap:hip-prenda-antic}
% ============================================================

Establecido el régimen común, este capítulo estudia las tres garantías reales en particular. Para quien desea comprar un departamento, un automotor y, además, posee un campo productivo que quiere usar como garantía de un crédito, el Código ofrece tres figuras distintas según el bien y la forma de estructurar la garantía: la hipoteca, la prenda y la anticresis. Todas comparten los mismos fundamentos normativos, pero cada una tiene sus propias reglas según el objeto y la posesión.
% TODO: verificar consistencia entre fuentes originales: el ejemplo de la fuente asigna la prenda al automotor y la anticresis al campo, pero más adelante la misma fuente señala que el automotor se prenda mediante prenda con registro y que la anticresis recae sobre cosas registrables.

Las garantías reales aparecen en la práctica profesional en operaciones de compraventa de inmuebles financiadas, créditos bancarios, adquisición de automotores y maquinaria industrial, ejecuciones hipotecarias y prendarias, cancelación registral de gravámenes, asesoramiento a condóminos, superficiarios y usufructuarios, y contratos de locación con garantía de prenda de moneda extranjera.

\section{Cómo se distinguen: objeto y posesión}

Las tres garantías se diferencian esencialmente en cuanto al \textbf{objeto} sobre el que recae el derecho real, es decir, el que soporta el privilegio del acreedor y será ejecutado si el deudor no cumple.

\begin{table}[H]
\centering
\begin{tabularx}{\textwidth}{@{}>{\raggedright\arraybackslash}p{0.20\textwidth}>{\raggedright\arraybackslash}X>{\raggedright\arraybackslash}X>{\raggedright\arraybackslash}p{0.24\textwidth}@{}}
\toprule
\textbf{Garantía} & \textbf{Objeto} & \textbf{¿Quién tiene la posesión?} & \textbf{Publicidad} \\
\midrule
Hipoteca & Un inmueble; también el derecho de superficie. & El constituyente (titular del dominio); el acreedor no tiene posesión. & Registro de la Propiedad Inmueble. \\
Prenda clásica & Cosas muebles no registrables o créditos instrumentados. & Se desplaza al acreedor, quien la recupera al extinguirse la obligación. & La posesión misma. \\
Anticresis & Cosas registrables (muebles o inmuebles). & Se desplaza al acreedor o al tercero designado por las partes. & La posesión de la cosa (y la inscripción en el registro correspondiente). \\
\bottomrule
\end{tabularx}
\end{table}
% TODO: verificar consistencia entre fuentes originales: la fuente indica que en la anticresis la publicidad la confiere la posesión, y a la vez que, por tratarse de cosas registrables, en todos los casos se procede a la inscripción en el registro correspondiente.

\begin{ejemplo}{Posesión en la hipoteca}
Si una persona es propietaria de un inmueble y tramitó un crédito con garantía hipotecaria (por ejemplo, por saldo de precio en el Banco Nación), no es el banco quien tiene la posesión del departamento. El propietario está en posesión del inmueble y es el titular dominial de la cosa hipotecada.
\end{ejemplo}

\section{Hipoteca}

\subsection{Concepto}

\begin{definicion}{Hipoteca (art.~2205 CCyCN)}
``La hipoteca es un derecho real que recae sobre uno o más inmuebles individualizados que continúan en poder del constituyente y que otorgan al acreedor, ante el incumplimiento del deudor, las facultades de persecución y preferencia para cobrar con su producido el crédito garantizado.''
\end{definicion}

El artículo 2205 es claro en que el derecho real de garantía recae sobre uno o más inmuebles individualizados que \textbf{continúan en poder del constituyente}. La posesión no se desplaza al acreedor; el constituyente es siempre el propietario del inmueble. Ese propietario en ocasiones es deudor y en ocasiones no lo es: tendrá la calidad de \textbf{propietario no deudor}.

La fortaleza de este derecho real es la \textbf{persecución} y la \textbf{preferencia}. Esto permite, por ejemplo, que el titular de dominio transmita el inmueble sin que ello afecte al acreedor, quien persigue la cosa en manos de quien se encuentre, aun cuando hubiese cambiado el titular dominial. Ambas se ejercen ante el incumplimiento del deudor: el acreedor intima el pago; si el deudor paga, la hipoteca se extingue (aunque el inmueble continúe figurando como hipotecado en el registro hasta la cancelación); si no paga, se pone en marcha el procedimiento de ejecución, es decir, la subasta del inmueble.

\subsection{Legitimados (art.~2206)}

Únicamente pueden constituir hipoteca:
\begin{itemize}
    \item el \textbf{titular del dominio} (dueño pleno), que tiene el \textit{ius disponendi} que le permite enajenar y constituir derechos reales de disfrute o de garantía;
    \item los \textbf{condóminos}, respecto de su parte indivisa;
    \item quienes sean propietarios en \textbf{propiedad horizontal};
    \item los propietarios en \textbf{conjuntos inmobiliarios};
    \item el \textbf{superficiario}, que puede hipotecar su derecho.
\end{itemize}

El condómino no puede hipotecar toda la cosa, porque sus facultades respecto de ella son restringidas. Sí puede hipotecar su \textbf{parte indivisa}, que puede ser ejecutada (subastada) por el acreedor en caso de incumplimiento. El adquirente en subasta adquiriría, por ejemplo, un 33,33\,\% del inmueble y pasaría a ser condómino junto con los otros titulares.

\begin{importante}
Mientras la hipoteca sobre la parte indivisa subsista, si los condóminos realizan la partición extrajudicial, el acreedor necesariamente deberá ser parte de ese acto. De lo contrario, la partición le será \textbf{inoponible} al acreedor.
\end{importante}

\subsection{Forma y plazo de la inscripción}

Como en materia de derechos reales sobre inmuebles todo debe realizarse necesariamente en escritura pública, la hipoteca se constituye por \textbf{escritura pública}. Del mismo modo, para cancelar la inscripción registral también se requiere escritura pública; cuando esa escritura de cancelación se inscribe en el Registro de la Propiedad Inmueble, opera la cancelación.

\begin{importante}
La Ley 27.271 sustituyó los artículos 2189 y 2210 del Código Civil y Comercial. Esta reforma no se encuentra en el texto de la mayoría de las ediciones del código que se utilizan en la práctica, por lo que es fundamental tenerla en cuenta.
\end{importante}

\begin{regla}{Duración de la inscripción hipotecaria (art.~2210, según Ley 27.271)}
Los efectos del registro de la hipoteca se conservan por el término de \textbf{35 años} si antes no se renueva.
\end{regla}

\begin{table}[H]
\centering
\begin{tabularx}{\textwidth}{@{}>{\raggedright\arraybackslash}p{0.34\textwidth}>{\raggedright\arraybackslash}p{0.16\textwidth}>{\raggedright\arraybackslash}X@{}}
\toprule
\textbf{Período} & \textbf{Plazo} & \textbf{Fuente normativa} \\
\midrule
Código 2015 hasta el 15/09/2016 & 20 años & Texto original del CCyCN (art.~2210) \\
Desde el 15 de septiembre de 2016 en adelante & 35 años & Ley 27.271 (reforma del art.~2210) \\
\bottomrule
\end{tabularx}
\end{table}

\subsection{Extinción y cancelación de la hipoteca}

\begin{table}[H]
\centering
\begin{tabularx}{\textwidth}{@{}>{\raggedright\arraybackslash}p{0.17\textwidth}>{\raggedright\arraybackslash}p{0.30\textwidth}>{\raggedright\arraybackslash}X@{}}
\toprule
\textbf{Concepto} & \textbf{¿Cuándo ocurre?} & \textbf{Efecto} \\
\midrule
\textbf{Extinción} & Cuando se extingue la obligación principal (por ejemplo, pago de la última cuota del crédito). & La hipoteca se extingue por vía de consecuencia (carácter accesorio), pero el inmueble puede seguir figurando como hipotecado en el registro. \\
\textbf{Cancelación} & Cuando se otorga escritura pública de cancelación y se la inscribe en el Registro de la Propiedad Inmueble. & Recién aquí el gravamen desaparece frente a terceros. Sin cancelación, el inmueble sigue hipotecado ante el registro. \\
\bottomrule
\end{tabularx}
\end{table}

\begin{ejemplo}{Hipotecas extinguidas pero no canceladas}
En la práctica profesional hay numerosos casos en que las hipotecas se encuentran extinguidas pero no canceladas, y el propietario lo desconoce. Cuando va a vender el inmueble, allí aparece la hipoteca. Por eso es tan importante distinguir extinción de cancelación.
\end{ejemplo}

\begin{nota}
\textbf{Seguros de vida en créditos bancarios.} En los créditos bancarios en general existe un seguro para extinguir las obligaciones ante el fallecimiento del deudor. Pero este seguro no comprende la cancelación registral: extingue la obligación, pero ante terceros queda pendiente la cancelación hasta que los herederos realicen el trámite necesario para otorgar la escritura y cancelar la inscripción.
\end{nota}

La cancelación puede ser:
\begin{itemize}
    \item \textbf{voluntaria}, cuando el propio acreedor (particular o bancario) realiza la escritura de cancelación;
    \item \textbf{judicial}, cuando el acreedor no cancela voluntariamente: requiere un juicio específico de cancelación de hipoteca, en el que el constituyente demanda al acreedor;
    \item \textbf{de pleno derecho}, por mero efecto de la ley, cuando vence el plazo de inscripción sin que se haya producido la renovación (antes, 20 años; actualmente, 35 años según la Ley 27.271).
\end{itemize}

\subsection{Ejecución hipotecaria}

La ejecución de la hipoteca (subasta del inmueble ante el incumplimiento del deudor) está regulada en los Códigos de Procedimiento de cada jurisdicción, no directamente en el Código Civil y Comercial.

\cornell{hipoteca}{
\cq{¿Qué es la hipoteca y qué permiten la persecución y la preferencia (art.~2205)?}{Es el derecho real sobre uno o más inmuebles individualizados que continúan en poder del constituyente y que otorga al acreedor, ante el incumplimiento, persecución y preferencia para cobrar con su producido. El titular de dominio puede transmitir el inmueble sin afectar al acreedor, que persigue la cosa en manos de quien se encuentre.}
\cq{¿Quién conserva la posesión y cómo se da la publicidad?}{El constituyente, siempre propietario (deudor o no deudor); el acreedor no tiene posesión. La publicidad es la inscripción en el Registro de la Propiedad Inmueble.}
\cq{¿Quiénes pueden constituirla (art.~2206) y qué ocurre con la parte indivisa?}{Titular del dominio pleno, condóminos (solo su parte indivisa, que puede subastarse), propietarios en propiedad horizontal o conjuntos inmobiliarios y superficiario sobre su derecho. Si hay partición extrajudicial mientras subsiste la hipoteca, el acreedor debe ser parte; de lo contrario, la partición le es inoponible.}
\cq{¿Qué forma exige y cuál es el plazo de inscripción?}{Escritura pública, para constituirla y para cancelar la inscripción. Los efectos del registro se conservan 35 años (Ley 27.271, desde el 15/9/2016); antes, 20 años.}
\cq{¿Qué diferencia hay entre extinción y cancelación y cuáles son las formas de cancelar?}{La extinción ocurre al extinguirse la obligación principal; la cancelación requiere escritura pública inscripta, y sin ella el inmueble sigue hipotecado frente a terceros. El seguro de vida de los créditos bancarios extingue la obligación pero no cancela el registro. La cancelación puede ser voluntaria, judicial o de pleno derecho (vencimiento del plazo sin renovación).}
\cq{¿Dónde se regula la ejecución hipotecaria?}{En los Códigos de Procedimiento de cada jurisdicción.}
}{La hipoteca opera sobre inmuebles sin desplazamiento posesorio, otorga persecución y preferencia y da publicidad a través del Registro de la Propiedad Inmueble. Se constituye por escritura pública, sus efectos registrales se conservan por 35 años y debe distinguirse la extinción de la cancelación, pues la primera no implica automáticamente la segunda.}

\section{Anticresis}

\subsection{Concepto, objeto y legitimados}

\begin{definicion}{Anticresis (arts.~2212 a 2218 CCyCN)}
La anticresis es el derecho real de garantía que recae sobre cosas registrables (muebles o inmuebles). El acreedor recibe la posesión de la cosa como garantía, o la recibe el tercero designado por las partes. La publicidad la confiere la posesión, a diferencia de la hipoteca, donde la publicidad la confiere el registro.
\end{definicion}

Pueden constituir anticresis (art.~2213) el titular del dominio, los condóminos, los propietarios en propiedad horizontal, el superficiario y el \textbf{usufructuario}, que puede hacerlo porque tiene el derecho de disponer de los frutos, pero con la limitación temporal de su propio derecho: si es vitalicio, hasta su muerte; si es a término, hasta que se cumpla el plazo de 10 o 20 años según corresponda.

\subsection{Plazos máximos}

\begin{table}[H]
\centering
\begin{tabularx}{\textwidth}{@{}>{\raggedright\arraybackslash}p{0.35\textwidth}>{\raggedright\arraybackslash}X@{}}
\toprule
\textbf{Tipo de bien} & \textbf{Plazo máximo desde la constitución} \\
\midrule
Muebles registrables & 5 años \\
Inmuebles & 10 años \\
\bottomrule
\end{tabularx}
\end{table}

\begin{nota}
A diferencia de lo que ocurre con la hipoteca, no se ha prorrogado el plazo de la anticresis de 20 a 35 años (conf.~art.~2218). El artículo 2218 mantiene las duraciones de 5 y 10 años según el tipo de bien.
\end{nota}

\subsection{Derechos y deberes del acreedor anticresista}

Al tener la posesión del objeto de la garantía, el acreedor anticresista puede \textbf{usar la cosa} (usar y explotar el bien, inclusive darlo en arrendamiento) y \textbf{percibir los frutos}: si es un inmueble, los frutos son civiles (rentas, alquileres) y se imputan primero a gastos e intereses y luego a capital. Tiene, además, los deberes de \textbf{conservar} y mantener el bien, afrontar los \textbf{gastos de conservación}, \textbf{restituir} la cosa al constituyente al cumplirse el plazo (o antes si se extingue la obligación), y puede \textbf{recuperar} el costo de las mejoras necesarias; las útiles, solo si hubiesen conferido a la cosa un mayor valor. Como se trata de cosas registrables, en todos los casos se procede a la inscripción en el registro correspondiente.

\cornell{anticresis}{
\cq{¿Qué es la anticresis y en qué se diferencia de la hipoteca?}{Es el derecho real de garantía sobre cosas registrables, muebles o inmuebles. Hay desplazamiento de la posesión al acreedor (o a un tercero designado) y la publicidad la otorga la posesión y no el registro.}
\cq{¿Quiénes pueden constituirla (art.~2213)?}{Titular del dominio, condóminos, propietarios en propiedad horizontal, superficiario y usufructuario, este último limitado por la duración de su propio derecho.}
\cq{¿Cuáles son los plazos máximos (art.~2218)?}{Cinco años para muebles registrables y diez para inmuebles. No se prorrogaron de 20 a 35 años, como ocurrió con la inscripción hipotecaria.}
\cq{¿Qué derechos y deberes tiene el acreedor anticresista?}{Usa y explota la cosa (incluso arrendándola) y percibe los frutos, imputados primero a gastos e intereses y luego a capital. Debe conservar la cosa, pagar los gastos de conservación y restituirla; recupera las mejoras necesarias y las útiles que dieron mayor valor.}
}{La anticresis recae sobre cosas registrables e implica el desplazamiento posesorio al acreedor, quien puede usar la cosa y percibir sus frutos, con plazos máximos de cinco y diez años según el bien. El acreedor debe conservar la cosa y restituirla.}

\section{Prenda}

\subsection{Concepto: prenda clásica y prenda con registro}

\begin{definicion}{Prenda (art.~2219 CCyCN)}
``La prenda es el derecho real de garantía sobre cosas muebles no registrables o créditos instrumentados.''
\end{definicion}

Debe distinguirse la \textbf{prenda clásica} (la regulada por el CCyCN) de la \textbf{prenda con registro} (regulada por leyes especiales), que son dos figuras completamente distintas.

\begin{table}[H]
\centering
\begin{tabularx}{\textwidth}{@{}>{\raggedright\arraybackslash}p{0.17\textwidth}>{\raggedright\arraybackslash}X>{\raggedright\arraybackslash}X@{}}
\toprule
\textbf{Característica} & \textbf{Prenda clásica (CCyCN)} & \textbf{Prenda con registro (leyes especiales)} \\
\midrule
Objeto & Cosas muebles \textbf{no} registrables o créditos instrumentados. & Cosas muebles registrables (autos, motos, maquinaria industrial, etc.). \\
Posesión & \textbf{Siempre} se desplaza al acreedor. & El constituyente conserva la posesión y puede seguir usando el bien. \\
Publicidad & La posesión del acreedor. & La inscripción en el registro prendario o en el registro del automotor. \\
Normativa aplicable & Arts.~2219 y ss.\ del CCyCN. & Leyes especiales (Decreto-Ley 15.348/46 y normas del RNPA para automotores). \\
\bottomrule
\end{tabularx}
\end{table}

\begin{ejemplo}{Prenda con registro: automotor}
Cuando se compra un automotor y se paga en cuotas, el auto está prendado. La posesión la tiene el comprador, no la concesionaria que otorgó el préstamo ni el banco. El comprador es constituyente y a su vez tiene la posesión. La prenda registrable otorga publicidad por el registro, no por la posesión.
\end{ejemplo}

\begin{ejemplo}{Prenda con registro: industria}
Se utiliza mucho en la industria metalúrgica, donde lo que se prendan son máquinas. Esta prenda se realiza mediante su inscripción en un registro prendario. Lo importante es que el constituyente puede seguir utilizando las máquinas; de lo contrario, para obtener el crédito, el propietario debería entregar la posesión y carecería de ellas para producir.
% TODO: contenido incompleto o ambiguo en las notas originales (el apunte dice «el acreedor puede seguir utilizando las máquinas», pero el sentido del ejemplo indica que se refiere al constituyente)
\end{ejemplo}

\subsection{Legitimados y forma de constitución}

La prenda puede ser constituida por el dueño de la cosa. El contrato puede ser público o privado y en todos los casos es necesaria la \textbf{fecha cierta}. Para que se perfeccione la garantía es requisito indispensable (constitutivo) la \textbf{entrega de la cosa}, ya que el acreedor ejerce el derecho real por la posesión.

\begin{nota}
A diferencia de lo que ocurre en la hipoteca, el condómino \textbf{no} puede prendar la cosa por su parte indivisa, porque tiene derechos muy restringidos respecto de toda la cosa y no podría entregarla. En la prenda, debe ser constituida necesariamente por todos los condóminos en conjunto.
\end{nota}

\subsection{Prenda de créditos instrumentados (art.~2232)}

Existe la posibilidad de prender cualquier crédito instrumentado que pueda ser cedido. El caso típico es el pagaré.

\begin{ejemplo}{Funcionamiento de la prenda de un pagaré}
Quien es acreedor de un pagaré (alguien le debe dinero y está documentado en ese instrumento) necesita un crédito. En lugar de ceder el pagaré, lo entrega en prenda a su acreedor como garantía. Para que quede constituida la garantía no basta la mera entrega del instrumento: se requieren \textbf{dos elementos constitutivos}:
\begin{enumerate}
    \item la entrega de la posesión del instrumento al acreedor prendario;
    \item la notificación al deudor cedido (quien firmó el pagaré), informándole que ya no debe pagar al constituyente sino directamente al acreedor prendario.
\end{enumerate}
\end{ejemplo}

El acreedor prendario, a partir de la constitución de la prenda, puede realizar todas las gestiones necesarias para cobrar el crédito: intimar el pago, ejecutar el pagaré, etc. Una vez extinguida la obligación garantizada, debe restituir el instrumento y rendir cuentas al constituyente de lo que percibió, para determinar si existe algún remanente que le corresponda a este último.

\subsection{Ejecución y prenda de dinero (arts.~2228 y 2229)}

Si el deudor no cumple, se realizará el bien en subasta pública conforme a lo regulado en los artículos 2228 y 2229 del Código Civil y Comercial. Se plantea el interrogante de si puede prendarse el dinero y si la moneda extranjera puede ser objeto de prenda, teniendo en cuenta que jurídicamente es una cosa mueble.

\begin{ejemplo}{Prenda de moneda extranjera en contratos de locación}
Es muy habitual que, cuando se celebra un contrato de locación y el locatario no tiene garantía propietaria, se constituya una prenda de moneda extranjera como garantía. La moneda extranjera, jurídicamente cosa mueble, queda prendada y es restituida al locatario a la finalización del contrato.
\end{ejemplo}

\cornell{prenda}{
\cq{¿Qué es la prenda clásica y qué la distingue de la prenda con registro?}{Es la regulada por el CCyCN (art.~2219) sobre cosas muebles no registrables o créditos instrumentados; la posesión siempre se desplaza al acreedor y la publicidad es la posesión. En la prenda con registro (automotores, maquinaria industrial, regida por leyes especiales) el constituyente conserva la posesión y la publicidad la da el registro.}
\cq{¿Quién puede constituirla y con qué forma?}{El dueño de la cosa. El contrato puede ser público o privado, siempre con fecha cierta, y es requisito constitutivo la entrega de la cosa. Los condóminos no pueden prendar su parte indivisa: debe constituirse por todos en conjunto.}
\cq{¿Cómo funciona la prenda de créditos instrumentados?}{Se prendan créditos en instrumentos cedibles (por ejemplo, pagarés). Requiere entrega de la posesión del instrumento al acreedor prendario y notificación al deudor cedido para que pague a este. El acreedor puede cobrar y debe restituir el instrumento y rendir cuentas al constituyente.}
\cq{¿Cómo se ejecuta la prenda y puede prendarse el dinero?}{El bien se realiza en subasta pública conforme a los arts.~2228 y 2229. La moneda extranjera, jurídicamente cosa mueble, puede prendarse; se usa con frecuencia como garantía en locaciones y se restituye al finalizar el contrato.}
}{La prenda clásica del Código Civil y Comercial recae sobre muebles no registrables o créditos instrumentados, exige siempre la entrega de la posesión al acreedor y no debe confundirse con la prenda con registro, figura especial regulada por leyes separadas en la que el constituyente conserva la posesión.}

% ============================================================
% ============================================================


% ============================================================
\part{Tiempo, posesión y defensa de los derechos reales}
% ============================================================

\chapter{Prescripción adquisitiva}\label{cap:prescripcion}
% ============================================================

El estudio de la prescripción adquisitiva de inmuebles es, esencialmente, el estudio del \textbf{tiempo} y de sus efectos en el derecho. Este capítulo explica cómo el transcurso del tiempo, combinado con una posesión pública, continua e ininterrumpida, puede transformar a quien carece de derecho sobre una cosa en titular de un derecho real. Recorre su regulación, la prescripción larga, la breve, las causas que alteran el cómputo del plazo y el juicio mediante el cual se hace valer.

\section{El tiempo como fenómeno jurídico}

Desde una perspectiva humana, el autor Ashley sostiene que los seres humanos, desde la más remota antigüedad, han mostrado su impotencia ante el tiempo: no es posible detenerlo cuando se es feliz ni acelerarlo cuando se desea que algo concluya. San Agustín denomina a esta vivencia \textbf{tiempo existencial}: el tiempo que se percibe como efímero en los buenos momentos y como eterno en aquellos que se desea que terminen.

Los juristas no son ajenos a esta realidad. En materia de \textbf{derechos reales}, el tiempo tiene una potencia enorme. También la tiene en los derechos personales, pero allí se manifiesta esencialmente a través de la prescripción liberatoria. En los derechos reales, el tiempo es tan potente que puede convertir a quien tiene solo una \textbf{relación de hecho} con la cosa en \textbf{adquirente de un derecho real}, al cabo de cierto tiempo y bajo determinadas calidades de poder sobre la cosa.

\begin{ejemplo}{Don Ramón}
Don Ramón llegó a un terreno baldío hace 22 años, construyó su casa, pagó algunos impuestos y crió a su familia allí. El dueño original nunca apareció. ¿Puede Don Ramón escriturar?

La respuesta se construye a lo largo del capítulo: qué es la prescripción adquisitiva y sus fundamentos; cuánto tiempo y con qué calidad de posesión se requiere (prescripción larga); el caso del comprador de buena fe con título imperfecto (prescripción breve); cómo el tiempo puede pausarse o reiniciarse (suspensión e interrupción); y cómo se lleva todo ello a juicio para obtener la escritura.
\end{ejemplo}

El instituto es central en el ejercicio profesional, porque con frecuencia el tiempo y la posesión determinan quién tiene derecho sobre un inmueble. Permite asesorar a quien lleva 15 o 20 años viviendo en un terreno sin escritura, defender o atacar una acción reivindicatoria, iniciar un juicio de usucapión y producir la prueba documental necesaria, evaluar si un título tiene defectos de fondo (justo título frente a título suficiente), calcular plazos considerando suspensiones e interrupciones, y estimar el valor real de una posesión en la compraventa. La vivienda es el bien más valioso que poseen la mayoría de las familias: entender la prescripción adquisitiva permite proteger ese patrimonio, sanear títulos imperfectos y brindar seguridad jurídica.

\section{Regulación normativa}

La prescripción adquisitiva está regulada en el Código Civil y Comercial en dos lugares:

\begin{itemize}
    \item el \textbf{Libro IV}, dedicado a los derechos reales, en los arts.~1897 a 1907: es el lugar natural de la regulación, porque el efecto de la relación de poder (la posesión) es tan potente que permite la adquisición del dominio o de otros derechos reales;
    \item el \textbf{Libro V}, que contiene disposiciones comunes en materia de derechos reales y personales, en los arts.~2532 a 2565.
\end{itemize}

\begin{regla}{Art.~2565 CCyCN}
Los derechos reales principales se pueden adquirir por prescripción adquisitiva en los términos del art.~1897, a excepción del derecho real de superficie, que no puede adquirirse por prescripción adquisitiva.
\end{regla}

Los derechos reales se clasifican en principales y accesorios (estos últimos, los de garantía). Por lo tanto, pueden adquirirse por prescripción el dominio, la propiedad horizontal, los conjuntos inmobiliarios y también los derechos de disfrute sobre cosa ajena (con la excepción de la superficie, véase el capítulo~\ref{cap:superficie}). Sin embargo, como los requisitos y el tiempo son los mismos, es muy difícil que alguien inicie un juicio para adquirir, por ejemplo, un usufructo; por eso se habla usualmente de la prescripción adquisitiva \textbf{de dominio}.

\begin{definicion}{Prescripción adquisitiva (art.~1897 CCyCN)}
La prescripción para adquirir es el modo por el cual el poseedor de una cosa adquiere un derecho real sobre ella mediante la posesión durante el tiempo fijado por la ley.
\end{definicion}

\section{Prescripción adquisitiva larga}

El art.~1897 precisa dos elementos centrales: el \textbf{tiempo}, veinte años según el art.~1899, y una \textbf{posesión exigible}, que debe ser \textbf{ostensible} (denominada pública en el Código de Vélez, es decir, no clandestina) y \textbf{continua e ininterrumpida}.

\begin{regla}{Art.~1899 CCyCN}
El plazo para la prescripción adquisitiva de inmuebles es de veinte (20) años.
\end{regla}

\begin{importante}
En la prescripción larga \textbf{no hace falta justo título ni título de ningún tipo, ni buena fe}. Es un punto que suele confundirse cuando se pregunta por los requisitos.
\end{importante}

La ley exige veinte años de actos posesorios, es decir, veinte años en los que la persona haya tenido una relación de poder con la cosa en forma pública, continua e ininterrumpida. Ninguna persona es siempre buena ni siempre mala; esa calificación excede el ámbito de la ciencia jurídica. Lo que existe es una posesión que en ocasiones es legítima y en otras ilegítima, en ocasiones de buena fe y en otras de mala fe: es la calidad que inviste al poseedor en determinada circunstancia.

\subsection{Fundamentos del instituto}

Frente a la pregunta de si es posible que quien ingresó sin ningún derecho sobre la cosa llegue a ser el dueño, el instituto se apoya en los siguientes fundamentos:

\begin{table}[H]
\centering
\begin{tabularx}{\textwidth}{@{}>{\raggedright\arraybackslash}p{0.27\textwidth}>{\raggedright\arraybackslash}X@{}}
\toprule
\textbf{Fundamento} & \textbf{Descripción} \\
\midrule
Orden público & Las reglas que dan estabilidad y certeza a toda la sociedad exigen que el derecho de propiedad se defina en algún momento. \\
Seguridad jurídica & Es necesario que los derechos reales, especialmente el dominio, no queden en un estado de indefinición permanente. \\
Interés general & La sociedad se beneficia cuando un inmueble tiene un uso productivo y un responsable claro. \\
Función social de la propiedad & El poseedor que usa el inmueble para vivir, construir, pagar impuestos y contribuir con la comunidad cumple mejor la función social que el propietario ausente. \\
Saneamiento de títulos & Permite detener la cascada de nulidades e imperfecciones que podrían sucederse en la cadena de transmisiones de un inmueble. \\
\bottomrule
\end{tabularx}
\end{table}

Muchas veces un usurpador puede llegar a adquirir un inmueble a expensas del verdadero propietario. Pero para que esto ocurra debe haber existido necesariamente una \textbf{inacción del propietario} que supere el período exigido por la ley, es decir, los veinte años. A lo largo de ese tiempo el poseedor ha dado un uso al inmueble: lo utilizó para vivir, para formar una familia, construyó, pagó algunos impuestos, contribuyó con los servicios barriales, envió a sus hijos al colegio cercano, pagó cooperadoras y realizó una función de utilidad pública mucho mayor que la del propietario que no estuvo en posesión durante veinte años, no pagó impuestos, no realizó ninguna acción útil respecto de la cosa y no inició ninguna acción para recuperarla.

Es cierto que el poseedor ha ingresado ilegítimamente al inmueble, pero a partir de ese momento el titular dominial (el \textit{dominus}) cuenta con todas las acciones reales y posesorias para recuperar la posesión (véanse los capítulos~\ref{cap:posesorias} y~\ref{cap:reales}). Si transcurren veinte años sin que las haya utilizado, el orden público, en virtud de las situaciones beneficiosas para la sociedad que se han producido, tolera que el poseedor sea beneficiado con el derecho real, siempre que se cumplan determinadas condiciones.

\subsection{Unión de posesiones (art.~1901)}

La unión de posesiones permite \textbf{sumar tiempo} para alcanzar los veinte años exigidos.

\begin{regla}{Art.~1901 CCyCN}
El heredero continúa la posesión de su causante. El sucesor particular puede unir su posesión a la de sus antecesores, siempre que ambas sean de fuente legítima y estén ligadas por un vínculo jurídico.
\end{regla}

\begin{ejemplo}{Juan y la cesión de posesión}
Juan está en posesión de un inmueble. Lo adquirió mediante un boleto de compraventa y no realizó ninguna acción tendiente a la escrituración, por lo que el titular dominial sigue siendo otra persona. Pasa el tiempo y Juan tiene una posesión de, por ejemplo, seis o siete años. Toda persona puede transmitir lo que tiene: Juan no puede transmitir la cosa, porque no es suya, pero sí puede transmitir la \textbf{posesión}.

Como se trata de inmuebles, la transmisión se realiza siempre mediante una escritura pública de \textbf{cesión de posesión}, que tiene un precio y es totalmente legítima. Sin embargo, no es lo mismo, y es lo que consultan los clientes a los abogados, el precio de una posesión de cinco o seis años, a la que le restan muchos años durante los cuales puede ser reivindicada por el titular dominial, que el de una posesión a la que le falta un año para cumplir los veinte. En este último caso el riesgo es menor, aunque todavía falte un juicio de escrituración.
\end{ejemplo}

Los requisitos de la unión de posesiones son los siguientes: ambas posesiones deben derivar la una de la otra, en cadena (la que se entrega a quien la recibe); deben ser del mismo tipo; se aplica al \textbf{sucesor particular}, pues el heredero no necesita este instituto, ya que ocupa el lugar de su causante y continúa teniendo la misma posesión que él; y no hace falta buena fe ni un nexo jurídico determinado.

% TODO: verificar consistencia entre fuentes originales: la redacción del art. 1901 difiere entre el material sobre relaciones de poder (capítulo de dinámica) y el material sobre prescripción.
% TODO: verificar consistencia entre fuentes originales: el art. 1901 transcripto exige «vínculo jurídico» y «fuente legítima», mientras que la explicación posterior indica que no hace falta un nexo jurídico determinado. Se conservan ambas formulaciones tal como figuran en el apunte.

\begin{table}[H]
\centering
\begin{tabularx}{\textwidth}{@{}>{\raggedright\arraybackslash}p{0.27\textwidth}>{\raggedright\arraybackslash}X>{\raggedright\arraybackslash}X@{}}
\toprule
 & \textbf{Heredero} & \textbf{Sucesor particular} \\
\midrule
¿Necesita unión de posesiones? & No: continúa automáticamente la posesión del causante. & Sí: debe unir su posesión a la de su antecesor. \\
Vínculo jurídico & No aplica (ocupa el lugar del causante). & Posesiones del mismo tipo, derivadas la una de la otra. \\
Buena fe requerida & No aplica. & No es exigible en la prescripción larga. \\
\bottomrule
\end{tabularx}
\end{table}

\section{Prescripción adquisitiva breve}

Existen dos tipos de prescripción adquisitiva, la larga y la breve, muy distintas entre sí en requisitos, plazos y objetivo. Lo que ambas comparten es que alguien que no era titular de dominio llega a serlo por intermedio de estos institutos. La prescripción breve tiene un \textbf{plazo de diez años} y la posesión, en todos los casos, debe ser de \textbf{buena fe} y con \textbf{justo título}.

% TODO: verificar referencia presente en las fuentes originales: el art. 1898 (plazo de la prescripción breve) solo aparece citado en el cuadro comparativo del apunte; el texto no lo transcribe.

\subsection{El justo título (art.~1902)}

\begin{definicion}{Justo título (art.~1902 CCyCN)}
El justo título es el que tiene por fin transmitir un derecho real principal que se ejerce por la posesión, que está revestido de las condiciones de forma, pero que sin embargo carece de los requisitos de fondo.
\end{definicion}

El \textbf{título suficiente} es el acto jurídico revestido de las condiciones de fondo y de forma. El \textbf{justo título} tiene solo forma y carece de fondo: quien transmitió no era el propietario, o no tenía capacidad. Para que el dominio se transmita realmente, debe ser transmitido por el propietario.

\begin{ejemplo}{Sustitución de personalidad: el \textit{non domino}}
Alguien se hace pasar por el propietario con documentación falsa y realiza la transmisión de un inmueble. El adquirente ha hecho todo lo exigido: otorgó la escritura pública, realizó estudios de títulos, obtuvo certificados, y el escribano cumplió con todas las formalidades. Sin embargo, existía algo subyacente en el fondo que hace que ese título, que el adquirente cree firmemente suficiente, sea en realidad un justo título, porque faltaban las condiciones de fondo. Por eso este adquirente \textbf{no ha adquirido nada}.

Un día recibe la acción reivindicatoria del propietario, que le pide la restitución del inmueble, y recién entonces se entera. Ambos tienen una razón: al propietario se le ha cometido un delito, pues se transmitió su propiedad por quien no era el dueño (el \textit{non domino}); y el adquirente también tiene derecho, porque adquirió e hizo todo lo que jurídicamente se exige para adquirir una propiedad. Pero existe un solo inmueble.
\end{ejemplo}

Frente a este conflicto entre dos derechos sobre un único inmueble, la decisión del orden público (modificable y perfectible) se resuelve mediante el factor \textbf{tiempo}: si el titular dominial ejerce la acción reivindicatoria \textbf{antes de los diez años} de la transmisión, el derecho protege al titular de dominio; si el reclamo se realiza \textbf{después de los diez años}, el adquirente, que tenía un justo título en lugar de un título suficiente por una falencia en la capacidad o en la legitimidad de quien transmitió, tendrá por adquirido el inmueble.

\subsection{Buena fe y unión de posesiones}

A diferencia de la prescripción larga, en la breve la buena fe es \textbf{indefectible}. El adquirente no sabía: adquirió y creyó haber adquirido con un título suficiente, y no pudo vencer el obstáculo que se encontraba en la profundidad del negocio jurídico. Para la unión de posesiones se requiere que ambas posesiones sean de buena fe y que estén ligadas por un vínculo jurídico, por ejemplo la cesión de posesión; es una excepción respecto de la prescripción larga, porque la buena fe del adquirente debe estar acreditada.

\subsection{Cuadro comparativo}

\begin{table}[H]
\centering
\begin{tabularx}{\textwidth}{@{}>{\raggedright\arraybackslash}p{0.22\textwidth}>{\raggedright\arraybackslash}X>{\raggedright\arraybackslash}X@{}}
\toprule
\textbf{Característica} & \textbf{Prescripción larga} & \textbf{Prescripción breve} \\
\midrule
Plazo & 20 años (art.~1899) & 10 años (art.~1898) \\
Justo título & No es necesario & Sí, obligatorio (art.~1902) \\
Buena fe & No es necesaria & Sí, obligatoria \\
Posesión requerida & Ostensible y continua & Ostensible, continua y de buena fe \\
Origen del poseedor & Cualquiera (incluso usurpador) & Comprador con título imperfecto \\
Unión de posesiones & Cadena, sin nexo jurídico especial & Ambas de buena fe, con vínculo jurídico \\
Norma principal & Arts.~1897 y 1899 & Arts.~1898 y 1902 \\
\bottomrule
\end{tabularx}
\end{table}

\section{Suspensión e interrupción de la prescripción}

El tiempo jurídico no siempre coincide con el tiempo del almanaque. Por ejemplo, la feria judicial hace que ese mes no sea un plazo hábil judicialmente. El derecho interviene en el tiempo a través de dos institutos, la suspensión y la interrupción, que se aplican tanto a la prescripción larga como a la breve.

\subsection{La suspensión}

La suspensión \textbf{detiene el cómputo del tiempo} por el plazo que dura.

\begin{ejemplo}{Línea de tiempo de una suspensión}
Juan comienza a poseer en el año 2000. En el año 2005 se produce una causal de suspensión, que abre un paréntesis y se cierra en 2010. Ese período, encapsulado entre paréntesis, no se cuenta. Si el paréntesis comprende cinco años, en lugar de llegar a 2020 para computar veinte años será necesario llegar a 2025.
\end{ejemplo}

Las causales de suspensión están indefectiblemente pautadas por la ley y responden a razones consideradas valiosas en la puja entre los distintos derechos, por ejemplo: el matrimonio, la unión convivencial y el vínculo entre incapaces y sus tutores. Cuando existe un vínculo protegido por el Estado entre el titular de dominio y el poseedor, el tiempo que haya durado ese vínculo no se computa.

\subsection{La interrupción (art.~2544)}

En la interrupción, el plazo comienza a correr y se produce una causa que \textbf{borra todo el tiempo anterior de posesión}. El poseedor vuelve a poseer y comienza a contar de nuevo.

\begin{regla}{Art.~2544 CCyCN}
El curso de la prescripción se interrumpe: a) por petición judicial; b) por pedido de arbitraje; c) por el reconocimiento del derecho de aquel contra quien prescribe, efectuado por el poseedor.
\end{regla}

La interrupción se produce cuando el titular de dominio ha realizado una acción, como la petición judicial para recuperar la cosa: su pasividad se ve cortada. También se produce por una solicitud de arbitraje o por un reconocimiento por parte del poseedor.

\begin{table}[H]
\centering
\begin{tabularx}{\textwidth}{@{}>{\raggedright\arraybackslash}p{0.25\textwidth}>{\raggedright\arraybackslash}X>{\raggedright\arraybackslash}X@{}}
\toprule
 & \textbf{Suspensión} & \textbf{Interrupción} \\
\midrule
Efecto sobre el tiempo acumulado & Se detiene: el tiempo anterior se \textbf{conserva}. & Se borra: el tiempo anterior se \textbf{pierde}. \\
Al cesar la causa & El cómputo continúa desde donde se detuvo. & El cómputo comienza desde cero. \\
Causales típicas & Matrimonio, unión convivencial, tutela de incapaces. & Petición judicial, arbitraje, reconocimiento del poseedor. \\
Norma aplicable & Arts.~2539 a 2543 & Art.~2544 \\
Aplica a prescripción & Larga y breve & Larga y breve \\
\bottomrule
\end{tabularx}
\end{table}

\section{El juicio de usucapión}

\subsection{Cómo se hace valer la prescripción}

La adquisición del derecho real se cumple por el mero cumplimiento de los plazos y de la calidad de posesión exigida. Sin embargo, nadie lo sabe: el poseedor no puede transmitir el derecho y sigue siendo poseedor, mientras que en el Registro de la Propiedad Inmueble el inmueble continúa a nombre del titular dominial. Por eso es preciso iniciar el \textbf{juicio de prescripción adquisitiva de inmuebles}, también llamado \textbf{juicio de usucapión}.

\begin{regla}{Ley 14.159, art.~24}
El juicio de prescripción adquisitiva de inmuebles se rige por el art.~24 de la Ley 14.159, tramita por el proceso contradictorio y tiene por objeto hacer cesar el estado de incertidumbre entre la situación registral y la posesión real ejercida.
\end{regla}

El juicio busca hacer cesar la discordancia entre la relación extrarregistral real, que es la posesión (el verdadero y nuevo titular que carece de título), y la situación registral. Se persigue el reconocimiento judicial del derecho y la inscripción de la sentencia, inscribiendo el dominio por prescripción adquisitiva a nombre del poseedor.

\subsection{La prueba}

La prueba es fundamental en este juicio, y se trata de una \textbf{prueba compleja}: debe ser amplia e indubitable. Se está controvirtiendo a un titular dominial, que tenía un dominio perpetuo, por lo que el orden público está presente con toda su potencia, tanto en la regulación como en materia de prueba.

La mediación obligatoria no se ha exceptuado en la prescripción adquisitiva. Sin embargo, aunque deba convocarse la mediación, la cuestión no puede resolverse ni siquiera por acuerdo del titular de dominio. Es necesaria, por lo tanto, la realización del juicio y la producción de la prueba.

Se habla de prueba compuesta o compleja, que debe ser insospechada, clara, cabal y convincente. Lo que debe probarse es que han transcurrido los veinte años, y no basta con el cómputo cronológico: es necesario analizar que no haya existido alguna causal de suspensión o de interrupción.

\subsection{Documentos obligatorios y actos posesorios}

Deben acompañarse a la demanda:
\begin{enumerate}
    \item el \textbf{informe de dominio}, que indica quién es el titular de dominio del inmueble, a los fines de la oponibilidad a terceros;
    \item el \textbf{plano de agrimensura con fines de usucapión}, que debe estar registrado e inscripto: es tan importante que su ausencia podría llevar al juez a rechazar \textit{in limine} la demanda.
\end{enumerate}

El actor debe probar el tiempo y los actos posesorios: haber intervenido, construido, plantado, instalado cañerías, cloacas y agua, y haber pagado impuestos. No es necesario haber pagado todos los impuestos, pero la prueba de impuestos es fundamental, aunque sea de hace quince o veinte años.

\begin{importante}
Pagar todos los impuestos en el presente no sirve de ninguna manera, porque el sello fechador es el de hoy y no demuestra que el poseedor estuviera en el inmueble en ese momento. El pago de impuestos no se valora como una cuestión contributiva o tributaria de buen ciudadano que cumplió con sus deudas, sino como acreditación de que en ese momento la persona ya estaba allí, con intención y comportándose como dueño.
\end{importante}

Veinte años es mucho tiempo en la vida de las personas y también lo es en la posesión de un inmueble. Si el poseedor logra acreditar que durante veinte años ha poseído la cosa sin reconocer en otro un señorío superior (lo sea o no), comportándose como dueño, esa prueba da lugar a la prescripción adquisitiva de inmuebles.

\cornell{prescripción adquisitiva}{
\cq{¿Qué es la prescripción adquisitiva y dónde se regula?}{Es el modo por el cual el poseedor adquiere un derecho real mediante la posesión durante el tiempo fijado por la ley (art.~1897): un puente entre la relación de hecho y el derecho real. Se regula en el Libro IV (arts.~1897 a 1907) y en el Libro V (arts.~2532 a 2565). El art.~2565 habilita la adquisición de todos los derechos reales principales, salvo la superficie.}
\cq{¿Cuáles son los requisitos de la prescripción larga?}{Veinte años (art.~1899) de posesión ostensible y continua e ininterrumpida. No se exige justo título ni buena fe. Bastan los actos posesorios propios del dueño: construcción, plantación, pago de algún impuesto, conexión de servicios.}
\cq{¿Por qué puede ganar un usurpador contra el dueño?}{Por los fundamentos de orden público, seguridad jurídica, interés general, función social de la propiedad y saneamiento de títulos, frente a la inacción del propietario, que no ejerció acción real o posesoria alguna durante veinte años mientras el poseedor dio uso efectivo al inmueble.}
\cq{¿Qué es la unión de posesiones (art.~1901) y a quién se aplica?}{Permite sumar el tiempo de distintas posesiones. Se aplica al sucesor particular, no al heredero, que continúa la posesión del causante. Las posesiones deben derivar la una de la otra y ser del mismo tipo; se instrumenta por escritura de cesión de posesión.}
\cq{¿Cuáles son los requisitos de la prescripción breve y qué es el justo título (art.~1902)?}{Diez años de posesión de buena fe con justo título: un acto con forma pero sin fondo, porque el transmitente no era el propietario o carecía de capacidad. El título suficiente reúne forma y fondo. La unión de posesiones exige buena fe en ambas posesiones y un vínculo jurídico.}
\cq{¿Qué ocurre si el verdadero propietario reclama el inmueble a quien tiene justo título?}{Si ejerce la reivindicatoria antes de los diez años de la transmisión, el derecho lo protege; si lo hace después, el adquirente tiene por adquirido el inmueble.}
\cq{¿Qué diferencia hay entre suspensión e interrupción?}{La suspensión (arts.~2539 a 2543) detiene el cómputo y conserva el tiempo anterior (matrimonio, unión convivencial, tutela de incapaces). La interrupción (art.~2544) borra el tiempo acumulado y obliga a comenzar de cero (petición judicial, arbitraje, reconocimiento del poseedor). Con un paréntesis de cinco años, desde el 2000 hay que llegar al 2025.}
\cq{¿Cómo se hace valer la prescripción y qué prueba exige el juicio de usucapión?}{Mediante el juicio de usucapión (Ley 14.159, art.~24), proceso contradictorio, con informe de dominio y plano de agrimensura registrado. La prueba debe ser compleja, indubitable y convincente: tiempo sin suspensión ni interrupción y actos posesorios históricos; los impuestos pagados recientemente no sirven. La sentencia se inscribe en el Registro de la Propiedad Inmueble.}
}{La prescripción adquisitiva transforma una relación de hecho en derecho real por el transcurso del tiempo con una posesión de cierta calidad. La larga exige veinte años sin título ni buena fe; la breve, diez años con justo título y buena fe. La suspensión detiene el cómputo y la interrupción lo borra. Para que la prescripción cumplida sea oponible a terceros debe obtenerse sentencia en el juicio de usucapión e inscribirla en el Registro. Su fundamento es el orden público, la función social de la propiedad y la seguridad jurídica.}

\chapter{Acciones posesorias}\label{cap:posesorias}
% ============================================================

\begin{ejemplo}{Hilo conductor}
Un inquilino ocupa un departamento. El dueño, enojado porque el inquilino no pagó, un día entra por la fuerza y cambia la cerradura. El inquilino tiene el \textbf{hecho} de vivir allí (posesión o tenencia); el dueño tiene el \textbf{derecho} de ser propietario. La pregunta es quién gana: se trata de comprender cómo el derecho protege el hecho frente al derecho, y cuándo y cómo puede cada uno recuperar lo que le corresponde.
\end{ejemplo}

Este capítulo estudia los medios que la ley confiere para defender la relación de poder con la cosa. Parte de su concepto y fundamento, distingue la turbación del desapoderamiento, delimita el área del posesorio frente al área del petitorio, presenta las acciones específicas y la defensa extrajudicial, y se cierra con una advertencia práctica que prepara el estudio de las acciones reales (capítulo~\ref{cap:reales}).

\section{Concepto y fundamento}

Las \textbf{acciones posesorias} son los medios o herramientas que la ley confiere para defender la relación de poder. Están reguladas en los artículos 2238 a 2246 del CCyCN y tienen por objeto la protección de las relaciones de poder (la posesión y la tenencia) frente a actos de turbación o de desapoderamiento. Las relaciones de poder son tres: \textbf{la posesión, la tenencia y los servidores de la posesión}; a estos últimos se los reconoce al solo efecto de la defensa privada.

La razón de ser de estas acciones es un amplio sentido social de defensa del \textbf{orden público} y de la \textbf{paz social}. Lo que el sistema jurídico pretende evitar es la \textbf{justicia por mano propia}, que podría generar un amplio gravamen a toda la sociedad.

% TODO: verificar consistencia entre fuentes originales: aquí las relaciones de poder se enumeran como posesión, tenencia y servidores de la posesión; el capítulo sobre relaciones entre la persona y la cosa añade la yuxtaposición local.
\subsection{Todas las formas de posesión son protegidas}

La clasificación de la posesión se estudió en el capítulo~\ref{cap:clasificacion}.


\begin{table}[H]
\centering
\begin{tabularx}{\textwidth}{@{}>{\raggedright\arraybackslash}p{0.24\textwidth}>{\raggedright\arraybackslash}X>{\centering\arraybackslash}p{0.17\textwidth}@{}}
\toprule
\textbf{Tipo de posesión} & \textbf{Características} & \textbf{¿Protegida por acciones posesorias?} \\
\midrule
Legítima & Constituida conforme a la normativa; titularidad propia del derecho real. & \textbf{Sí} \\
Ilegítima de buena fe & Sin título suficiente; el poseedor cree tener derecho. & \textbf{Sí} \\
Ilegítima de mala fe simple & Sin título; el poseedor sabe que no tiene derecho. & \textbf{Sí} \\
Ilegítima viciosa & Adquirida por violencia, clandestinidad o abuso de confianza. & \textbf{Sí} \\
\bottomrule
\end{tabularx}
\end{table}

Incluso el poseedor de mala fe viciosa goza de la protección posesoria frente a terceros que actúen aun con peor derecho o por vías de hecho.

\subsection{Un título válido no da la posesión (art.~2239)}

\begin{regla}{Art.~2239 CCyCN}
Un título válido no da la posesión ni la tenencia misma, sino el derecho a requerir el poder sobre la cosa. El que no tiene sino un derecho a la posesión o a la tenencia no puede tomarla; debe demandarla por las vías legales.
\end{regla}

Este artículo cristaliza uno de los pilares de la materia: \textbf{tener el hecho de la posesión no tiene nada que ver con tener el derecho de poseer}. Se pueden dar cuatro situaciones: (a) ser titular del dominio y tener la posesión; (b) ser titular del dominio y \textbf{no} tener la posesión (porque se la robaron, la perdieron, etc.); (c) tener la posesión de hecho \textbf{sin} ser titular de ningún derecho real; (d) no tener ni la posesión ni el derecho real.

\section{Tipos de lesión: turbación y desapoderamiento}

Un poseedor o tenedor puede sufrir dos tipos de lesiones en su relación de poder con la cosa: la \textbf{turbación} y el \textbf{desapoderamiento} (despojo).

\begin{table}[H]
\centering
\begin{tabularx}{\textwidth}{@{}>{\raggedright\arraybackslash}X>{\raggedright\arraybackslash}X@{}}
\toprule
\textbf{Turbación} & \textbf{Desapoderamiento (despojo)} \\
\midrule
El poseedor o tenedor \textbf{no} es excluido absolutamente. Alguien interfiere con su relación de hecho con la cosa, pero no lo desplaza por completo. & El poseedor o tenedor es \textbf{excluido absolutamente} de la relación de hecho con la cosa. Pierde toda vinculación fáctica. \\
\addlinespace
\textit{Objetivo del turbador:} no necesariamente constituirse en poseedor, pero sí actuar como si tuviera un derecho sobre la cosa. & \textit{Objetivo del despojante:} excluir total y definitivamente al poseedor o tenedor. \\
\addlinespace
\textbf{Acción que corresponde:} acción de \textbf{mantener} (art.~2242). & \textbf{Acción que corresponde:} acción de \textbf{despojo} (art.~2241). \\
\bottomrule
\end{tabularx}
\end{table}

No toda molestia constituye turbación en sentido posesorio. No se trata, por ejemplo, de que el ruido en la puerta impida dormir. La turbación es aquella que tiene por finalidad \textbf{controvertir la relación de poder con la cosa}: debe existir una intención de ejercer un derecho sobre la cosa ajena.

\begin{ejemplo}{Turbación}
Alguien entra al terreno del poseedor y estaciona su automóvil, o descarga materiales allí. No lo desplaza totalmente, pero se comporta respecto de lo ajeno \textbf{como si fuera el dueño}. Eso es turbación. Cuando los clientes relatan estos casos, la descripción es siempre similar: la persona se comportó como si fuera el dueño.
\end{ejemplo}

\section{Área del posesorio y área del petitorio}

La relación entre las acciones posesorias y las acciones reales puede comprenderse distinguiendo dos grandes áreas de protección.

\begin{table}[H]
\centering
\begin{tabularx}{\textwidth}{@{}>{\raggedright\arraybackslash}p{0.17\textwidth}>{\raggedright\arraybackslash}X>{\raggedright\arraybackslash}X@{}}
\toprule
 & \textbf{Área del posesorio} & \textbf{Área del petitorio} \\
\midrule
Acciones & Acciones posesorias. & Acciones reales. \\
Qué se defiende & El \textbf{hecho} de la posesión. & El \textbf{derecho} real. \\
Qué se discute & Hechos, no derechos. & Derechos (el derecho real mismo). \\
Qué debe probarse & La relación de poder y la turbación o el despojo. & El derecho real (el derecho de poseer). \\
Proceso & Juicio breve: sumarísimo. & Proceso de conocimiento: amplia prueba. \\
Sentencia & Provisoria (solo hechos). & Definitiva (derechos). \\
Prescripción & Un año. & Imprescriptible (por el carácter perpetuo del dominio). \\
Legitimados & Todos los poseedores, legítimos e ilegítimos, y también los tenedores. & Solo los titulares del derecho real que se ejerce por la posesión. \\
\bottomrule
\end{tabularx}
\end{table}
% TODO: verificar consistencia entre fuentes originales: el material sobre acciones posesorias ubica las acciones reales en los arts. 2269 a 2276, mientras que el material sobre acciones reales las ubica en los arts. 2247 a 2268.

Las acciones posesorias tienen un \textbf{plazo de prescripción de un año}: si al cabo de ese año no se ha iniciado la acción, prescribe y no puede iniciarse después. El objetivo de la sentencia posesoria es el \textbf{restablecimiento al estado anterior}, y se tramita por el juicio más breve que tenga la jurisdicción. La sentencia es \textbf{provisoria}, porque solo se discutieron hechos: el titular del derecho real puede, después, iniciar la acción real correspondiente para reclamar su derecho en forma definitiva. En cambio, si el titular del derecho real inicia directamente la acción real y es vencido, la \textbf{sentencia definitiva} nada más permitirá reclamar: ha perdido su derecho definitivamente. El juego de estas acciones se da muy frecuentemente.

\subsection{Casos prácticos}

\begin{ejemplo}{Caso 1: Juan, poseedor sin derecho real, es despojado por el propietario}
Juan está en posesión de la cosa. En realidad no es titular del derecho: se comporta como si lo fuera, pero no lo es. Sin embargo, tiene un vínculo directo con la cosa y realiza actos posesorios. Cierto día llega el propietario y lo despoja violentamente.

El propietario \textbf{tiene} el derecho real (dominio), pero \textbf{no} tiene el derecho a recuperar la posesión en forma directa y por propia mano. Juan puede iniciar por juicio breve (sumarísimo) una \textbf{acción posesoria de despojo contra el titular de dominio}. Una vez reconocida su relación de hecho por el juez, el propietario deberá restituirle la posesión. La sentencia es provisoria: el propietario, para recuperar la posesión en forma definitiva, deberá iniciar la acción real reivindicatoria.
\end{ejemplo}

\begin{ejemplo}{Caso 2: el propietario puede iniciar ambas acciones}
El titular de dominio, al ser titular de un derecho real, puede iniciar las dos. En la \textbf{opción A} inicia la acción posesoria (sumarísima): si es vencido porque no logra probar su relación de hecho, puede igualmente iniciar después la acción reivindicatoria, porque la sentencia posesoria es provisoria y no hace cosa juzgada sobre el derecho. En la \textbf{opción B} inicia directamente la acción real: si es vencido, la sentencia es definitiva y nada más puede reclamar, pues ha perdido su derecho de propiedad.
\end{ejemplo}

\begin{ejemplo}{Caso 3: el inquilino expulsado por el propietario}
El propietario sostiene que el inquilino no pagó y no cumplió; corta la luz y el gas, entra violentamente a la casa y lo saca. El tenedor (inquilino) puede iniciar una acción posesoria contra el propio dueño para ser restablecido en su situación de hecho.
\end{ejemplo}

\section{Las acciones posesorias específicas}

El CCyCN \textbf{simplifica las acciones posesorias} en comparación con el Código de Vélez: regula una sola acción para la turbación (mantener) y otra para el despojo.

\subsection{Acción de mantener (art.~2242)}

\begin{regla}{Art.~2242 CCyCN}
El titular de la acción de mantener la tenencia o la posesión puede pedir que se ordene el cese de la turbación y adoptar las medidas pertinentes para evitar que vuelva a producirse.
\end{regla}

Se utiliza cuando existe \textbf{turbación}: el poseedor o tenedor no ha sido excluido totalmente, pero alguien, con intención de poseer, interfiere con su vínculo con la cosa. La sentencia buscada es que \textbf{cese la turbación} y que el juez adopte medidas preventivas para que no vuelva a ocurrir. También se aplica cuando exista una amenaza fundada de sufrir un daño y cuando se produzca una obra nueva en terreno lindero que, aunque no se realice en el propio terreno del poseedor, lo afecte de tal manera que se justifique la defensa.

\subsection{Acción de despojo (art.~2241)}

\begin{regla}{Art.~2241 CCyCN}
Acción de despojo: quien ha sido desapoderado de la tenencia o la posesión de un inmueble o de una universalidad de hecho, puede reclamar su restitución.
\end{regla}

Se utiliza cuando el poseedor o tenedor ha sido despojado totalmente. Su objetivo es \textbf{recuperar la tenencia o la posesión}, ya sea sobre toda la cosa individualmente o sobre una universalidad de hecho (por ejemplo, una manada de ganado o una biblioteca). La acción se dirige contra el despojante, los herederos del despojante, los sucesores particulares de mala fe y aun el propio dueño, cuando tomó la cosa por su propia mano sin realizar las acciones legales correspondientes. También se aplica a obras que afecten al poseedor o tenedor, cuando se realizan en terrenos vecinos pero generan consecuencias en el terreno del poseedor o tenedor.

\subsection{Los interdictos en los códigos procesales}

Los códigos procesales, tanto el nacional (CPCCN) como los provinciales, regulan \textbf{interdictos}: trámites muy rápidos para proteger determinadas turbaciones o el despojo. Se mencionan el interdicto de daño temido y el interdicto de obra nueva. La doctrina discute si los interdictos son herramientas independientes o simplemente la forma procesal en que se tramitan estas acciones posesorias; por el momento basta con saber que la discusión existe.

\section{Defensa extrajudicial o privada (art.~2240)}

La defensa extrajudicial constituye una \textbf{excepción} al principio general de la paz social y la prohibición de la justicia por mano propia.

\begin{regla}{Art.~2240 CCyCN}
Nadie puede mantener o recuperar la posesión o la tenencia de propia autoridad, excepto cuando debe protegerse y repeler una agresión con el empleo de una fuerza suficiente, en los casos en que los auxilios de la justicia llegarían demasiado tarde. El afectado debe recobrarla sin intervalo de tiempo y sin exceder los límites de la propia defensa.
\end{regla}

La figura se compara con la \textbf{legítima defensa en el derecho penal}: así como en materia penal existen casos en que la norma es impotente para evitar la reacción del afectado, en materia civil también existen situaciones que, por su virulencia e inmediatez, no pueden ser canalizadas únicamente a través de las vías judiciales.

\begin{importante}
Todos los requisitos deben darse \textbf{conjuntamente}. No alcanza con que se den algunos ni la mayoría.
\end{importante}

\begin{table}[H]
\centering
\begin{tabularx}{\textwidth}{@{}>{\raggedright\arraybackslash}p{0.25\textwidth}>{\raggedright\arraybackslash}X@{}}
\toprule
\textbf{Requisito} & \textbf{Explicación} \\
\midrule
Inmediatez & La reacción debe producirse en el mismo momento de la turbación o el despojo. No puede ser posterior. \\
Proporcionalidad & La defensa debe ser proporcional al ataque sufrido. No puede excederse en la respuesta. \\
Auxilio tardío & Los auxilios de la justicia llegarían demasiado tarde. Si pudiera llamarse a la policía o al juzgado y llegar a tiempo, no hay justificación para la defensa privada. \\
\bottomrule
\end{tabularx}
\end{table}

% TODO: contenido incompleto o ambiguo en las notas originales (la cita de la transcripción del caso y la frase final "paso a cásar" resultan ininteligibles)
\begin{ejemplo}{Proporcionalidad: el caso del ingeniero}
A un ingeniero le roban sus cosas; se detiene, sigue a los ladrones, dispara y mata. Se plantea si fue legítima defensa. La respuesta es negativa, porque el bien jurídico protegido por la ley es la vida; la actitud había sido \textbf{desproporcionada}.
\end{ejemplo}

La analogía civil es directa: en materia posesoria tampoco puede responderse en forma desproporcionada al ataque sufrido; la defensa debe ser proporcional al bien que se protege y al ataque que se recibe.

La legitimación para ejercer la defensa privada es \textbf{amplia}: todos los poseedores (con o sin derecho real), los tenedores, los titulares del derecho real y los \textbf{servidores de la posesión}, que se reconocen para la defensa privada.

% TODO: contenido incompleto o ambiguo en las notas originales (la justificación del reconocimiento de los servidores de la posesión aparece formulada de manera poco clara)
\begin{ejemplo}{Servidor de la posesión}
Alguien trabaja en la empresa en la que es empleado y tiene su máquina de escribir, o está en un banco, y alguien pretende robar: el empleado se aferra a la cosa para defenderla.
\end{ejemplo}

\section{Síntesis: de la turbación a la acción real}

\begin{importante}
\textbf{«Toda turbación es una disposición en marcha.»} Muchas veces, cuando se está tramitando una acción de mantener por turbación, en el transcurso del proceso ya se ha producido el despojo completo. Quien se atreve a comportarse sobre la cosa como si fuera el dueño, quien da el primer golpe, no necesariamente se detendrá después.
\end{importante}

Las acciones posesorias (área del posesorio) y las acciones reales (área del petitorio) conforman el panorama completo de la protección de la posesión y del derecho real. Confundir ambas esferas puede costar la recuperación del bien y, en el peor escenario, el derecho real mismo.

\cornell{acciones posesorias}{
\cq{¿Qué son las acciones posesorias y cuál es su fundamento?}{Medios que la ley (arts.~2238 a 2246) confiere para defender la relación de poder (posesión y tenencia) frente a turbación o desapoderamiento. Se fundan en el orden público y la paz social, y buscan evitar la justicia por mano propia.}
\cq{¿Se protege solo la posesión legítima?}{No: se protegen todas las formas de posesión (legítima, ilegítima de buena fe, de mala fe simple y viciosa), porque se protege la paz social.}
\cq{¿Qué establece el art.~2239 y qué situaciones pueden darse entre hecho y derecho?}{Que un título válido no da la posesión ni la tenencia, sino el derecho a requerirla por las vías legales. Pueden darse cuatro situaciones: titular con posesión; titular sin posesión; poseedor sin derecho real; ni posesión ni derecho.}
\cq{¿Qué diferencia hay entre turbación y despojo y qué acción corresponde a cada uno?}{En la turbación el poseedor o tenedor no es excluido: alguien actúa como si tuviera un derecho sobre la cosa (acción de mantener, art.~2242). En el despojo es excluido absolutamente (acción de despojo, art.~2241). Una molestia como el ruido no es turbación.}
\cq{¿Qué se discute en el posesorio y qué en el petitorio?}{En el posesorio, hechos: juicio sumarísimo, sentencia provisoria, prescripción de un año, legitimados poseedores y tenedores. En el petitorio, el derecho real: amplia prueba, sentencia definitiva, imprescriptible, legitimados los titulares.}
\cq{¿Qué ocurre si el titular de dominio inicia primero la posesoria o directamente la real y es vencido?}{Si inicia la posesoria y es vencido, puede iniciar después la reivindicatoria, porque la sentencia es provisoria. Si inicia directamente la real y es vencido, la sentencia definitiva le cierra toda posibilidad.}
\cq{¿Contra quién se dirige la acción de despojo y qué simplificó el CCyCN?}{Contra el despojante, sus herederos, los sucesores particulares de mala fe y aun el propio dueño que actuó por mano propia. El CCyCN regula una sola acción para la turbación y otra para el despojo. Los interdictos son trámites procesales rápidos cuya naturaleza discute la doctrina.}
\cq{¿Cuándo procede la defensa extrajudicial (art.~2240) y quién puede ejercerla?}{Cuando se dan conjuntamente inmediatez, proporcionalidad y auxilio tardío de la justicia; es una excepción comparable a la legítima defensa penal. Pueden ejercerla poseedores, tenedores, titulares del derecho real y servidores de la posesión.}
}{Las acciones posesorias son la primera y más veloz línea de defensa de quien, sea o no titular de un derecho real, ve perturbada o extinguida su relación de poder con una cosa. Se justifican en la paz social y en la erradicación de la justicia por mano propia; se tramitan por juicio sumarísimo, su prueba se limita al hecho posesorio y la sentencia es provisoria. La acción de mantener responde a la turbación y la de despojo al desapoderamiento; la defensa privada es una excepción de requisitos estrictos. El derecho real mismo se defiende con las acciones reales.}

\chapter{Acciones reales}\label{cap:reales}
% ============================================================

Las acciones reales son las herramientas procesales con las que se defiende a un titular cuando un derecho real (dominio, usufructo, servidumbre, hipoteca) es atacado, turbado o desconocido. Este capítulo presenta sus fundamentos, la acción reivindicatoria, las acciones negatoria, confesoria y de deslinde, y el juicio de reivindicación de inmuebles. Elegir la acción correcta es esencial: demandar por reivindicatoria cuando corresponde la negatoria implica la pérdida del proceso.

\begin{ejemplo}{Un campo con cuatro problemas distintos}
Un propietario de un campo enfrenta cuatro problemas diferentes:
\begin{itemize}
    \item un vecino le ocupa el lote por completo: corresponde la \textbf{acción reivindicatoria};
    \item otro vecino cruza su tierra como si tuviera una servidumbre: corresponde la \textbf{acción negatoria};
    \item un tercero le impide usar el camino que sí tiene derecho a usar: corresponde la \textbf{acción confesoria};
    \item un cuarto vecino no sabe dónde termina su terreno y dónde empieza el del propietario: corresponde la \textbf{acción de deslinde}.
\end{itemize}
Cada acción es una herramienta distinta para un problema distinto.
\end{ejemplo}

\section{Fundamentos}

Las acciones reales están contempladas en los arts.~2247 a 2268 del CCyCN. La \textbf{acción} es el medio, la herramienta con que cuenta el titular de un derecho para reclamar la intervención del órgano jurisdiccional a fin de que se le reconozca un derecho que ha sido lesionado o afectado. Se distinguen según la naturaleza del derecho que protegen: son \textbf{acciones reales} cuando protegen derechos reales y \textbf{acciones personales} cuando procuran la defensa de un derecho personal.

\begin{table}[H]
\centering
\begin{tabularx}{\textwidth}{@{}>{\raggedright\arraybackslash}p{0.20\textwidth}>{\raggedright\arraybackslash}X>{\raggedright\arraybackslash}X@{}}
\toprule
\textbf{Tipo de acción} & \textbf{Derecho que protege} & \textbf{Ejemplo} \\
\midrule
Acción real & Derechos reales (dominio, usufructo, hipoteca, servidumbre, etc.) & Recuperar un inmueble del que se fue despojado \\
Acción personal & Derechos personales o crediticios & Cobrar una deuda contractual \\
\bottomrule
\end{tabularx}
\end{table}

Todos los derechos reales (sobre cosa propia y sobre cosa ajena, de disfrute o de garantía) y todas las facultades reconocidas a cada titular encuentran su defensa a través de las acciones reales.

\begin{definicion}{Acciones reales (art.~2247 CCyCN)}
Las acciones reales son los medios de defender en juicio la \textbf{existencia}, \textbf{plenitud} y \textbf{libertad} de los derechos reales. En cada caso se utilizará la acción reivindicatoria, confesoria, negatoria y también la acción de deslinde.
\end{definicion}

El criterio para identificar qué acción corresponde es el \textbf{ámbito} de cada acción: qué derechos protege, quiénes son sus titulares y respecto de qué lesiones (legitimados activos) y contra quién debe iniciarse (legitimados pasivos). El art.~2248 regula la finalidad de las acciones reales, los titulares legitimados para interponerlas y contra quién hacerlo.

\begin{table}[H]
\centering
\small
\begin{tabularx}{\textwidth}{@{}>{\raggedright\arraybackslash}p{0.17\textwidth}>{\raggedright\arraybackslash}X>{\raggedright\arraybackslash}X>{\raggedright\arraybackslash}X@{}}
\toprule
\textbf{Acción} & \textbf{Lo que protege} & \textbf{Tipo de lesión} & \textbf{Objeto de la sentencia} \\
\midrule
Reivindicatoria & Existencia del derecho real & Despojo total de la posesión & Restitución de la cosa \\
Negatoria & Libertad del derecho real & Turbación: actos posesorios indebidos o más allá del límite & Hacer cesar la turbación \\
Confesoria & Plenitud del derecho real & Impedimento del ejercicio pleno de derechos inherentes a la posesión & Restablecer el pleno ejercicio \\
Deslinde & Certeza de los límites & Incertidumbre sobre límites entre fondos contiguos & Delimitar la superficie de cada inmueble \\
\bottomrule
\end{tabularx}
\end{table}

\section{Acción reivindicatoria}

La acción reivindicatoria protege la \textbf{existencia del derecho}. Protege al titular de un derecho real que se ejerce por la posesión y que ha sido despojado totalmente de ella: procede ante el despojo \textbf{total}, no ante una simple turbación, sino ante la pérdida completa del contacto con la cosa.

\subsection{Legitimados}

Están legitimados activamente el titular de los derechos reales que se ejercen por la posesión (dominio, usufructo, uso, habitación, etc.) y también el acreedor hipotecario. El reconocimiento del acreedor hipotecario es muy importante: está facultado para iniciar la acción \textit{iure propio}, es decir, por su propio derecho, porque la garantía que tiene sobre la obligación radica sobre el inmueble. Puede ocurrir que el titular de dominio despojado, que a su vez no ha podido cumplir con la obligación y cuyo inmueble posiblemente sea subastado en poco tiempo, resulte negligente en la defensa de su derecho; en ese caso la ley otorga al acreedor hipotecario la posibilidad de iniciar la reivindicatoria.

La acción se dirige contra el \textbf{poseedor} o el \textbf{tenedor}, aunque este último tenga la cosa a nombre del reivindicante: puede iniciarse contra quien tiene una relación de hecho con la cosa, ejerciendo un poder de hecho que no tiene la facultad jurídica de ejercer respecto de ella.

\begin{ejemplo}{Demanda contra Diego}
Quien se encuentra en el inmueble se presume poseedor. Se inicia entonces la acción contra Diego, porque se tomó conocimiento de que es la persona que está en posesión del inmueble. Puede ocurrir que se trate de un tenedor, aunque el actor no lo sepa. En ese caso, el tenedor que recibe la demanda tiene la \textbf{obligación de nominar}, es decir, de indicar e identificar concretamente contra quién debe dirigirse la acción. Por ejemplo: ``no soy el poseedor, estoy alquilando, soy tenedor y reconozco a tal persona''.
\end{ejemplo}

\subsection{Cosas no reivindicables (art.~2253)}

No son reivindicables los objetos inmateriales; las cosas indeterminadas o fungibles; los accesorios, si no se reivindica el principal; y los automotores inscriptos de buena fe, salvo que sean hurtados o robados (en cuyo caso tampoco son reivindicables).
% TODO: contenido incompleto o ambiguo en las notas originales (la frase sobre automotores hurtados o robados se conserva tal como figura en el apunte)

\begin{nota}
En materia de automotores se requiere otro complemento normativo para poder interpretar la regla: el Código establece un régimen puntual respecto de ellos, que no se desarrolla en las fuentes.
\end{nota}

\subsection{Efectos de la sentencia (art.~2261)}

La sentencia reivindicatoria debe reconocer el derecho: es \textbf{declarativa}, porque declara la existencia del derecho, y es también de \textbf{condena}, porque ordena restituir la cosa y, si fue solicitado, pagar los daños. Debe restituirse la cosa principal con sus accesorios; en el caso de un inmueble, debe restituirse la posesión \textbf{libre de enseres y ocupantes}; y corresponde el pago de daños y perjuicios si hubiesen sido solicitados.

\begin{importante}
La indemnización puede pedirse como accesoria de la restitución, o como vía principal sustitutiva cuando, por algún motivo durante el proceso, ya no es posible restablecer el vínculo con la cosa, pero el daño persiste.
\end{importante}

\section{Acciones negatoria, confesoria y de deslinde}

\subsection{Acción negatoria (art.~2262)}

La negatoria responde a una lesión menos grave que la reivindicatoria: no hay despojo total de la cosa, pero existe una \textbf{turbación}. El art.~2262 establece que la acción protege la \textbf{libertad del derecho real}. Procede cuando el titular de un derecho real es turbado mediante actos que manifiestan una intención de poseer o el ejercicio de derechos más allá de sus límites.

\begin{definicion}{Turbación}
No se trata de una molestia cualquiera (por ejemplo, que hagan ruido en la puerta o impidan dormir). Quien turba manifiesta una intención de poseer: hace algo que nadie haría si no fuera el dueño. El turbador actúa ``como si fuera el dueño'' o ``como si tuviera un derecho que no tiene''. Esta conducta habilita la acción negatoria.
\end{definicion}

\begin{ejemplo}{Servidumbre inexistente}
Alguien pretende tener un derecho de servidumbre que no tiene: pasa por el inmueble ajeno como si tuviera una servidumbre de paso y no la tiene.
\end{ejemplo}

\begin{ejemplo}{Ejercicio más allá de los límites del derecho}
Quien tiene una servidumbre de paso tiene derecho a pasar, pero en lugar de hacerlo por el lugar en el que está autorizado conforme a su derecho, lo hace por todo el inmueble. Ejerció el derecho más allá de sus límites, y ese ejercicio superior a los límites es el que turba el vínculo de propiedad del titular con la cosa y afecta el libre ejercicio de su derecho.
\end{ejemplo}

La turbación en muchas ocasiones es un \textbf{despojo en marcha}: lo que comenzó siendo una turbación va aumentando en virulencia. Quien comienza a descargar material en el fondo del terreno del vecino ha pasado un límite, y nada indica que no vaya a continuar. Así aparece, en materia de derechos reales, el despojo y, en materia de violencia, posiblemente la muerte.

\begin{importante}
Progresión: turbación $\rightarrow$ despojo $\rightarrow$ violencia. La negatoria actúa en la primera etapa; no debe esperarse a que la situación escale.
\end{importante}

El objeto de la sentencia es \textbf{hacer cesar la turbación}; si se trata de una servidumbre, se la reduce a sus límites. Son legitimados activos todos los titulares de derechos reales que se ejercen por la posesión, los titulares de servidumbre (cuando se impide el ejercicio) y el acreedor hipotecario, cuyo fundamento es el mismo que en la reivindicatoria: mantener el valor de la propiedad. Si se trata de una servidumbre de fondo encerrado, en caso de subasta no es lo mismo subastar un fondo con salida a la vía pública que uno sin ella, totalmente inutilizado para la vida humana, animal y comercial.

\subsection{Acción confesoria (art.~2264)}

El art.~2264 establece que la acción confesoria protege la \textbf{plenitud del derecho real}. Procede cuando el titular está impedido de ejercer plenamente los derechos inherentes a su posesión, o cuando se afectan servidumbres activas. Una guía práctica atribuida al Dr.~Paños, que se mantuvo aun bajo el Código de Vélez, indica que cuando el ámbito de la reivindicatoria (que es preciso) y el de la negatoria (que también lo es) no son aplicables, todas las demás lesiones corresponden a la confesoria: su ámbito es \textbf{residual}.

\begin{importante}
Criterio práctico: si la lesión no encaja en la reivindicatoria (despojo total) ni en la negatoria (turbación con intención de poseer o ejercicio excesivo), corresponde a la confesoria.
\end{importante}

La plenitud se afecta cuando otro realiza actos que impiden al titular ejercer libremente sus derechos y poner en acto todas sus facultades. No hay un despojo tan grave ni una turbación en el sentido de la negatoria, pero se impide el ejercicio libre de las facultades.

\begin{ejemplo}{Servidumbre activa bloqueada}
El titular de una servidumbre tiene derecho a pasar, pero no se le permite. Como titular del fondo dominante tiene derecho de paso sobre el fondo sirviente, pero se ha realizado una construcción o una tranquera que se lo impide. Esto es del ámbito de la acción confesoria.
\end{ejemplo}

\begin{ejemplo}{Restricciones y límites al dominio por vecindad}
Las restricciones y límites al dominio relativos a luces, vistas, olores y ruidos son inherentes a la posesión. El titular de dominio cuyo vecino no cumple con estas restricciones (realiza construcciones, ventanas o ruidos pese a no poder hacerlo) queda habilitado para iniciar una acción confesoria contra él para procurar su cese.
\end{ejemplo}

Son legitimados activos los titulares que ven afectados los derechos inherentes a la posesión, los titulares de servidumbres activas y los acreedores hipotecarios por derecho propio. La acción se ejerce contra quienes impiden el pleno ejercicio de los derechos reales o de las servidumbres activas.

\subsection{Comparación de las tres acciones}

\begin{table}[H]
\centering
\small
\begin{tabularx}{\textwidth}{@{}>{\raggedright\arraybackslash}X>{\raggedright\arraybackslash}p{0.19\textwidth}>{\raggedright\arraybackslash}X@{}}
\toprule
\textbf{Supuesto} & \textbf{Acción aplicable} & \textbf{Fundamento} \\
\midrule
Un vecino impide usar la servidumbre de paso (pone una tranquera) & Confesoria & Servidumbre activa impedida \\
Un vecino construye ventanas en una pared medianera violando restricciones & Confesoria & Derechos inherentes a la posesión (vecindad) \\
Un vecino cruza el fondo como si tuviera una servidumbre que no tiene & Negatoria & Turbación: actúa como si fuera el dueño \\
Un vecino ocupó el fondo completo & Reivindicatoria & Despojo total de la posesión \\
\bottomrule
\end{tabularx}
\end{table}

\subsection{Acción de deslinde}

La acción de deslinde no tiende a restablecer un derecho ni a hacer cesar un ataque, sino a hacer cesar un \textbf{estado de incertidumbre} entre fondos contiguos que no tienen demarcados sus límites. Existen dos fondos con límites confundidos, imprecisos o inciertos; quien inicia la acción (alguno de los titulares o ambos) procura delimitar, en la superficie de cada inmueble, el límite que se encuentra confundido.

\begin{importante}
En el deslinde no hay un ``agresor'' ni un ``atacante''. Hay dos fondos contiguos con límites confundidos. El objeto no es hacer cesar un ataque, sino resolver la incertidumbre sobre los límites.
\end{importante}

En todas las acciones reales, los legitimados activos son siempre titulares de derechos reales.

\section{El juicio de reivindicación de inmuebles}

El juicio de reivindicación es aquel mediante el cual se ejerce la acción reivindicatoria. Es un \textbf{juicio petitorio}: la controversia versa sobre derechos. Es posiblemente uno de los juicios de mayor dificultad en el área del derecho, comparable al juicio de simulación o a la acción pauliana en materia de derechos personales, que presentan una gran dificultad probatoria. Se enfrenta el titular del derecho real contra el poseedor, y la carga de la prueba está íntegramente en cabeza del reivindicante.

\subsection{La carga probatoria y la ``prueba diabólica''}

En igualdad de condiciones, siempre prevalece el poseedor, aunque se defienda o no, porque juegan en su favor numerosas presunciones: su relación con la cosa se presume posesión, se presume legítima y se presume de buena fe. No tiene que exhibir ni rendir título, porque posee. El reivindicante, en cambio, debe probar todo: su derecho, que ese derecho es mejor que el del poseedor y, si el poseedor exhibiera algún título, que es mejor que el que emerge de ese título.

\begin{importante}
\textbf{La prueba diabólica.} El reivindicante debe probar: (1) su derecho real; (2) que ese derecho es mejor que el del poseedor; (3) que estuvo en posesión del inmueble y fue despojado. Solo el título no es suficiente.
\end{importante}

Esta dificultad es tal que los franceses ya la llamaban \textbf{prueba diabólica}: no solo debe probarse el derecho mediante el título (lo más sencillo), sino también haber estado en posesión del inmueble que hoy se reivindica y del que se fue despojado.

\begin{table}[H]
\centering
\small
\begin{tabularx}{\textwidth}{@{}>{\raggedright\arraybackslash}p{0.26\textwidth}>{\raggedright\arraybackslash}X>{\raggedright\arraybackslash}X@{}}
\toprule
\textbf{Carga de la prueba} & \textbf{Reivindicante} & \textbf{Poseedor} \\
\midrule
Probar el título & Sí, obligatorio & No: posee y eso lo presume legitimado \\
Probar posesión anterior & Sí, obligatorio & No necesita probarla \\
Probar que su derecho es mejor & Sí, obligatorio & No \\
Presumido de buena fe & No automáticamente & Sí (presunción legal) \\
Caso especial (Estado reivindicante) & El Estado no prueba (su título emana de la ley) & Debe probar su derecho \\
\bottomrule
\end{tabularx}
\end{table}

\subsection{Presunciones del art.~2256}

En ayuda del reivindicante frente a la prueba diabólica existe una serie de presunciones, reguladas en el art.~2256, que distingue casos según el origen de los títulos que presentan las partes y la situación posesoria.

\begin{itemize}
    \item \textbf{Ambas partes presentan título del mismo antecesor.} Se presume que el derecho está en cabeza de quien primero haya obtenido la tradición. Esta situación es muy difícil que ocurra en la práctica, con los avances producidos en materia registral, por lo que se considera un caso hipotético.
    \item \textbf{Ambas partes presentan títulos de distintos antecesores.} También lo prevé el art.~2256.
    \item \textbf{El demandado no presenta título.} Uno de los supuestos es aquel en que el título del reivindicante es \textbf{posterior a la posesión}: ese título no es suficiente para que prospere la demanda, porque el demandado ya había entrado en posesión de la cosa.
\end{itemize}

% TODO: contenido incompleto o ambiguo en las notas originales (el apunte no desarrolla la solución del caso de títulos de distintos antecesores; la tabla lo resume como «se analizan los títulos y la posesión»)
% TODO: contenido incompleto o ambiguo en las notas originales (el apunte anuncia dos supuestos para el caso del demandado sin título pero desarrolla solo uno)

\begin{table}[H]
\centering
\small
\begin{tabularx}{\textwidth}{@{}>{\raggedright\arraybackslash}X>{\raggedright\arraybackslash}X@{}}
\toprule
\textbf{Situación} & \textbf{Presunción / efecto} \\
\midrule
Ambos con título del mismo antecesor & Prevalece quien primero obtuvo la tradición \\
Ambos con título de distintos antecesores & Se analizan los títulos y la posesión (art.~2256) \\
Demandado sin título; título del reivindicante posterior a la posesión & El título no es suficiente para ganar; prevalece la posesión anterior \\
Estado como reivindicante & La carga de la prueba se invierte: debe probarla el poseedor \\
\bottomrule
\end{tabularx}
\end{table}

\subsection{El Estado como reivindicante}

Cuando el reivindicante es el Estado, como su título emana de la ley, la carga de la prueba la tiene el poseedor. Salvo en ese caso, cuando el reivindicante es cualquier titular de dominio o cualquier particular, toda la carga probatoria recae sobre quien invoca su derecho.

\cornell{acciones reales}{
\cq{¿Qué es una acción real y qué defienden según el art.~2247?}{Es el medio procesal con que cuenta el titular de un derecho real para reclamar la intervención jurisdiccional ante su lesión; se diferencia de la acción personal. Defienden la existencia, plenitud y libertad de los derechos reales, sobre cosa propia o ajena.}
\cq{¿Cuáles son las cuatro acciones reales y cómo se identifica la que corresponde?}{Reivindicatoria (existencia, ante el despojo total), negatoria (libertad, ante la turbación), confesoria (plenitud, ante el impedimento) y de deslinde (incertidumbre de límites). Se identifica por su ámbito: derecho protegido, lesión, legitimados activos y pasivos.}
\cq{¿Cuándo procede la reivindicatoria y quiénes pueden iniciarla?}{Cuando el titular de un derecho real que se ejerce por la posesión ha sido despojado totalmente. La inician el titular y, \textit{iure propio}, el acreedor hipotecario; se dirige contra el poseedor o el tenedor, que debe nominar a quien corresponda.}
\cq{¿Qué cosas no son reivindicables y qué ordena la sentencia?}{Según el art.~2253: objetos inmateriales, cosas indeterminadas o fungibles, accesorios sin el principal y automotores inscriptos de buena fe (salvo hurto o robo). La sentencia es declarativa y de condena: restitución con accesorios (inmuebles libres de ocupantes) y daños si fueron solicitados.}
\cq{¿Qué es la turbación, qué acción habilita y qué ordena la sentencia?}{Conducta de quien actúa como si fuera el dueño o ejerce un derecho más allá de sus límites. Habilita la negatoria (libertad del derecho); la sentencia hace cesar la turbación y reduce la servidumbre a sus límites. Conviene actuar temprano: turbación, despojo y violencia.}
\cq{¿Cuál es la diferencia entre negatoria y confesoria?}{Negatoria: el tercero actúa como si fuera el dueño o excede los límites de su derecho. Confesoria: el tercero impide al titular ejercer plenamente sus derechos inherentes a la posesión o sus servidumbres activas; según el criterio atribuido al Dr.~Paños, es residual.}
\cq{¿Para qué sirve el deslinde?}{Para hacer cesar la incertidumbre sobre los límites de dos fondos contiguos, sin que exista agresor.}
\cq{¿Por qué es difícil ganar la reivindicación y cuándo se invierte la carga de la prueba?}{Es un juicio petitorio en el que el reivindicante debe probar su derecho, que es mejor que el del poseedor, y que estuvo en posesión y fue despojado (prueba diabólica), mientras el poseedor goza de presunciones. El art.~2256 prevé presunciones; solo si el reivindicante es el Estado la carga recae en el poseedor.}
}{Las acciones reales defienden el derecho real mismo. La reivindicatoria protege la existencia del derecho ante el despojo total; la negatoria, la libertad ante la turbación; la confesoria, de carácter residual, la plenitud ante el impedimento; el deslinde resuelve la incertidumbre sobre límites. En todas la legitimación activa corresponde al titular del derecho real y, en supuestos específicos, al acreedor hipotecario. La reivindicación de inmuebles es un juicio petitorio de gran dificultad probatoria.}

% ============================================================
% ============================================================


% ============================================================
\part{Publicidad registral y protección de la vivienda}
% ============================================================

\chapter{Publicidad registral inmobiliaria}\label{cap:registral}
% ============================================================

Los derechos reales inmobiliarios y su publicidad constituyen pilares fundamentales del tráfico jurídico en Argentina. Este capítulo explica por qué el derecho real exige ser conocido, cómo organiza esa publicidad la Ley 17.801, cuáles son los principios que rigen el sistema, qué documentos se inscriben, cómo opera el bloqueo registral y cuáles son los plazos críticos.

Su dominio permite al profesional proteger los derechos de los clientes en transacciones inmobiliarias, evitando fraudes y garantizando los títulos; comprender el sistema registral argentino; prevenir daños (el caso de la ``sombrilla registral'' muestra cómo una inscripción tardía puede costarle todo al cliente); aplicar en la práctica diaria los ocho principios registrales; asesorar sobre embargos, medidas cautelares y protección de patrimonios; y trabajar con escribanos y registros de la propiedad, comprendiendo plazos, requisitos y excepciones. La seguridad del tráfico inmobiliario depende de profesionales que entiendan cómo funciona la publicidad registral.

\section{Derecho real y publicidad}

Hablar de publicidad en materia de derechos reales es hablar propiamente del derecho real. Entre las diferencias esenciales entre derechos reales y derechos personales se encuentra que \textbf{el derecho real es absoluto y oponible \textit{erga omnes}}, mientras que el derecho personal es relativo y resulta oponible únicamente a las partes que integran el vínculo obligacional: el sujeto activo y el sujeto pasivo vinculado por la prestación.

El derecho real debe ser respetado por todos y, por lo tanto, necesariamente debe ser conocido por todos, porque \textbf{nadie puede ser obligado a respetar aquello que no conoce}. Por ello, la publicidad es una exigencia intrínseca y esencial de la inherencia del derecho real.

\begin{importante}
El derecho real es absoluto y oponible a todos; por eso exige publicidad. El derecho personal es relativo y solo es oponible a las partes del vínculo obligacional.
\end{importante}

\subsection{La posesión como forma de publicidad}

La publicidad por excelencia es la tradición, es decir, la posesión. Rige la máxima \textbf{``sin posesión no hay derecho real''}: la posesión es constitutiva en todos los derechos reales que se ejercen por la posesión. Esto es así porque el derecho real en materia inmobiliaria se adquiere, se transmite y se extingue \textbf{fuera del registro}.

% TODO: contenido incompleto o ambiguo en las notas originales (texto del artículo 2006 transcripto tal como figura en el apunte)
\begin{regla}{Artículo 2006 CCyCN}
El derecho real inmobiliario se constituye, transmite y extingue por actos jurídicos conforme a lo dispuesto en este Código. La posesión es la forma de ejercer el derecho real.
\end{regla}

\subsection{Efecto de la inscripción registral}

La inscripción registral es necesaria para la oponibilidad a terceros. Su efecto es \textbf{declarativo}, porque el derecho se constituye, se transmite y se extingue fuera del registro. Sin inscripción, el derecho será oponible solo a muy pocas personas: aquellas que tuvieron conocimiento del acto (la parte, el escribano, el abogado), ninguna de las cuales podrá prevalerse de la falta de inscripción. Para la oponibilidad a terceros, en cambio, es imprescindible que el acto jurídico de adquisición, transmisión, constitución o gravamen del derecho se haya inscripto en el registro.

\section{El sistema de la Ley 17.801}

Todo lo relativo a la inscripción registral conforma el sistema de publicidad inmobiliaria, regulado en Argentina por la \textbf{Ley 17.801}. Es una ley federal y, por lo tanto, aplicable en toda la Nación, sin perjuicio de que los ordenamientos provinciales tengan sus propias normas, en general decretos-leyes, para regular sus propios registros.

\begin{table}[H]
\centering
\begin{tabularx}{\textwidth}{@{}>{\raggedright\arraybackslash}p{0.22\textwidth}>{\raggedright\arraybackslash}X@{}}
\toprule
\textbf{Característica} & \textbf{Significado} \\
\midrule
Sistema real & El registro se realiza \textbf{por inmueble} y no por la persona de su titular. Cada inmueble está matriculado y tiene su número de orden; se le asigna una hoja denominada \textbf{folio real}, donde se asienta todo lo relativo a la adquisición, transmisión, gravamen y levantamiento de gravámenes. El sistema personal, en cambio, registra por titular. \\
Sistema de inscripción & Cada vez que se inscribe un acto de adquisición o transmisión, se asienta una síntesis del documento. En los sistemas de transcripción se transcribe el documento completo. \\
Sistema no convalidante & La inscripción no subsana ni convalida el título: si el título inscripto tiene defectos que pudieran acarrear su nulidad, seguirá teniéndolos y será atacable. Otros sistemas validan los títulos y prevén otro procedimiento para atacarlos. \\
\bottomrule
\end{tabularx}
\end{table}

\section{Principios registrales}

Poco tiempo después de la sanción de la Ley 17.801 se realizó un Congreso Internacional de Derecho Registral, en el que se aprobó la \textbf{Carta de Buenos Aires}. Allí se declaran los principios registrales que constituyen la base sobre la que se asienta el sistema. Se los explica en su esencia, sin pretender una definición concreta de cada uno.

\begin{table}[H]
\centering
\begin{tabularx}{\textwidth}{@{}>{\raggedright\arraybackslash}p{0.27\textwidth}>{\raggedright\arraybackslash}X@{}}
\toprule
\textbf{Principio} & \textbf{Explicación} \\
\midrule
Rogación & El registro procede a solicitud de parte, no de oficio. \\
Inscripción & Documentos inscribibles para su oponibilidad a terceros. \\
Presunción registral & Los asientos registrales se presumen veraces (fe pública). \\
Especialidad & Identificación perfecta e indubitativa del inmueble. \\
Legalidad & Examen de formas extrínsecas (calificación registral). \\
Prioridad & La inscripción primera tiene prioridad sobre las posteriores. \\
Tracto sucesivo & Encadenamiento perfecto entre todos los títulos. \\
Publicidad & Dar publicidad y seguridad al tráfico inmobiliario. \\
\bottomrule
\end{tabularx}
\end{table}

\paragraph{Rogación.} El registro procede a solicitud de parte y no de oficio. Su intervención es siempre rogada, ya sea para lograr un asiento inicial, para las mutaciones registrales, para anotar embargos y anotaciones de litis, o para pedir certificados o informes.

\paragraph{Inscripción.} Todo documento por el que se constituyen, transmiten, declaran, modifican o extinguen derechos reales sobre inmuebles debe ser inscripto para su oponibilidad a terceros (art.~2 de la Ley 17.801). Se vincula con el efecto declarativo de la inscripción y su efecto de oponibilidad.

\paragraph{Presunción registral.} Se infiere del art.~22 de la ley e implica que los asientos registrales se presumen veraces; se enlaza con la \textbf{fe pública registral}. Si se obtiene un informe expedido por el registro que indica que un inmueble pertenece a una persona determinada, esa información es el único canal para conocer el derecho real que resulta oponible. Aunque pueda haber existido algún caso de acciones judiciales contra el registro, lo que indica el registro se presume veraz.

\paragraph{Especialidad.} En todo título que se inscribe, plasmado en una escritura pública, el inmueble debe estar perfectamente individualizado: las medidas, la superficie, los linderos, la nomenclatura, y todo dato que indique indubitablemente respecto de qué inmueble se realiza la inscripción. Esto surge de todo el capítulo tercero de la ley, referido a la matriculación.

\paragraph{Legalidad o calificación.} Confiere al registro la facultad de examinar las \textbf{formas extrínsecas} del título cuya inscripción se pretende, no su contenido intrínseco. Si las formas reúnen todos los requisitos, el registro lo inscribe; si no, lo rechaza o lo inscribe \textbf{provisionalmente por un plazo de 180 días}, durante el cual debe subsanarse la observación, tanto si la inscripción fue solicitada por el escribano para una escritura pública como si fue requerida por el juzgado mediante una medida cautelar o una adjudicación. Los arts.~8 y 9 de la Ley 17.801 contemplan los recursos que pueden interponerse, administrativamente y, agotada la vía administrativa, judicialmente, si se considera que el registro se ha opuesto a la inscripción de manera infundada.

\paragraph{Prioridad.} Tiene prioridad el acto que ingresó primero al registro sobre otro que ingresó posteriormente, aun cuando el título posterior sea de fecha anterior. Lo que otorga la prioridad es la \textbf{fecha de inscripción y no la fecha del título}.

\paragraph{Tracto sucesivo.} Cada transmisión sucede a la otra mediante un \textbf{encadenamiento perfecto entre todos los títulos}: los unos derivan de los otros en forma regular, ordenada y encadenada. Las \textbf{excepciones} se dan cuando no existe transmisión de un derecho a otro; por ejemplo, en las inscripciones originarias y en los inmuebles adquiridos por prescripción adquisitiva (capítulo~\ref{cap:prescripcion}), que no responden al sistema de tradición traslativa de dominio.

\begin{regla}{Tracto abreviado (art.~16, Ley 17.801)}
El tracto abreviado permite que, en algunos casos regulados por la ley, puedan inscribirse simultáneamente dos transmisiones en un mismo asiento. No es una excepción al tracto sucesivo, sino una facilitación que mantiene el encadenamiento registral.
\end{regla}

\begin{ejemplo}{Tracto abreviado en la sucesión}
Juan compra un departamento y luego fallece. Sus hijos heredan el departamento y obtienen la declaratoria de herederos. El juzgado ordena inscribir el inmueble a nombre de los herederos. Si los herederos deciden vender el inmueble, resulta costoso realizar dos trámites: primero inscribir a nombre de los herederos y luego inscribir la venta. Aquí aparece el beneficio del art.~16 de la Ley 17.801: se solicita al juez del sucesorio que la inscripción a nombre de los herederos se realice directamente con la escritura traslativa de dominio por la venta. Es un procedimiento sumamente utilizado.
\end{ejemplo}

\paragraph{Publicidad.} Surge de todo lo expuesto. El objeto fundamental del registro y su finalidad específica es dar publicidad y conferir seguridad al tráfico inmobiliario, y abastece la buena fe de los derechos reales. La buena fe es registral: no puede alegarse buena fe si no se obtuvieron los informes que permiten conocer verazmente la información que consta en el registro.

\section{Documentos que se inscriben}

Lo que se inscribe siempre son \textbf{documentos}. Aunque en el lenguaje cotidiano se diga que ``se inscribe la transmisión del derecho'', ello siempre está plasmado en un instrumento público, porque se trata de materia inmobiliaria.

\begin{regla}{Documentos inscribibles (art.~2, Ley 17.801)}
Todos aquellos instrumentos que constituyen, transmiten o extinguen derechos reales, también los que dispongan embargos, inhibiciones y las otras providencias cautelares. Debe tratarse de instrumentos públicos, emanados del notario que realiza la escritura pública.
\end{regla}

Entre los \textbf{instrumentos notariales}, el notario es el profesional cuya incumbencia es la realización de la escritura pública, como fedatario, y tiene a su cargo la inscripción de la transmisión y adquisición del derecho real. También se inscriben los \textbf{instrumentos judiciales}, es decir, las resoluciones judiciales que ordenan o disponen la inscripción: por ejemplo, la sucesión por causa de muerte (el inmueble que pertenecía a una persona pasa a ser de los herederos) y la liquidación de la sociedad conyugal (cuando un inmueble se adjudica a uno de los cónyuges, el juez ordena la adjudicación y emite el documento que dispone la inscripción a nombre del adjudicatario).

\section{Bloqueo registral y reserva de prioridad}

El \textbf{bloqueo registral}, también llamado \textbf{reserva de prioridad}, es un procedimiento por el cual el registro \textbf{congela la información}. Todos los escribanos pueden pedir la información a través del \textbf{informe de dominio}, que da cuenta del titular, del porcentaje de titularidad, de los gravámenes y de las restricciones. Sin embargo, cuando un escribano solicita el certificado en función de la realización de un negocio inmobiliario, ese certificado tiene un efecto diferente: produce la \textbf{prioridad} del negocio jurídico para el cual fue solicitado y garantiza la inmutabilidad de todos los datos que el registro da a conocer durante un plazo de \textbf{15 días en la ciudad, 25 días en la provincia o 30 días fuera de la provincia}.

\begin{importante}
El informe de dominio solo brinda información. El certificado solicitado para un negocio inmobiliario, en cambio, produce el bloqueo registral: los datos del inmueble se mantienen inmutables durante el plazo de vigencia.
\end{importante}

\begin{nota}
\textbf{Metáfora de la sombrilla registral.} Durante los 15 días (o 25 a 30, según jurisdicción) que dura la validez del certificado, se abre una sombrilla protectora sobre el inmueble. Mientras está abierta, los nuevos gravámenes (como gotitas de lluvia) no pueden ingresar. Si la sombrilla se cierra antes de que el escribano haya inscripto la venta, los gravámenes caen sobre el nuevo propietario.
\end{nota}

\subsection{El caso de Juan y María}

\begin{ejemplo}{Compraventa entre Juan y María}
Juan, el enajenante, adquirió el inmueble por 50 mil dólares. Pacta con María, en un boleto (que es un derecho personal), que ella pagará 20 mil dólares al boleto y 30 mil dólares a la escritura. El boleto se celebra el 10 de junio y se acuerda escriturar el 10 de julio, a 30 días. Se lleva el título de propiedad original al escribano, quien solicita los certificados cuya expedición produce el bloqueo registral.
\end{ejemplo}

\begin{table}[H]
\centering
\begin{tabularx}{\textwidth}{@{}>{\raggedright\arraybackslash}p{0.27\textwidth}>{\raggedright\arraybackslash}X@{}}
\toprule
\textbf{Momento} & \textbf{Hecho} \\
\midrule
10 de junio & Celebración del boleto (derecho personal de María). \\
1.º al 15 de julio & Plazo de 15 días de vigencia de los certificados (``paraguas abierto''). \\
10 de julio & Otorgamiento de la escritura traslativa de dominio, dentro del período de vigencia de los certificados. \\
10 de julio al 15 de agosto & Plazo de 45 días para inscribir el título. \\
\bottomrule
\end{tabularx}
\end{table}
% TODO: contenido incompleto o ambiguo en las notas originales (las fechas del apunte no coinciden exactamente con los plazos de 15 y 45 días indicados)

La expedición del certificado abre un ``paraguas'' que protege el inmueble: durante ese tiempo, los datos que indican que el inmueble es de Juan, que no está gravado y que no tiene hipotecas se mantienen inalterables. Si en ese lapso el Juzgado Civil 5 ordena un embargo de 100 mil dólares sobre el inmueble matrícula 3267, titularidad de Juan, el gravamen \textbf{no puede ingresar} mientras el paraguas está abierto: como una gotita, el embargo queda suspendido.

Mientras tanto, es obligación del escribano otorgar la escritura traslativa de dominio, con el título de propiedad a la vista y los certificados vigentes. Hasta ese momento había solo un derecho personal de María; con la escritura pública celebrada dentro del período de vigencia (el 10 de julio) se cumple la tradición traslativa de dominio, pero \textbf{falta la inscripción}.

\subsection{Plazo de 45 días y retroactividad}

La ley concede al escribano un nuevo plazo de \textbf{45 días} para inscribir el título. Si lo inscribe dentro de ese plazo, se produce la \textbf{retroactividad}: la inscripción tiene efecto retroactivo al día de la celebración de la escritura, el 10 de julio. Cuando el paraguas se cierra, el embargo ordenado por el Juzgado Civil 5 cae como la gotita; pero como la inscripción tiene efecto retroactivo al día de la celebración, el inmueble ya no es de Juan sino de María. Por lo tanto, el embargo es rechazable.

Si el escribano no inscribe dentro de los 45 días (por ejemplo, inscribe el 20 de agosto en lugar de antes del 15 de agosto), pierde el privilegio de la retroactividad que otorga la ley.

\begin{importante}
Pueden inscribirse los títulos después de los 45 días, pero se pierde la retroactividad. En cambio, no puede otorgarse la escritura si no es dentro de la vigencia de los certificados (dentro de la protección del paraguas).
\end{importante}

En el caso planteado, la transmisión de Juan a María es oponible a partir del 20 de agosto. Cuando el paraguas se cerró, el embargo ingresó: María compra el inmueble con un embargo de 100 mil dólares. Se trata de un departamento de 50 mil dólares gravado por 100 mil, de modo que \textbf{lo ha perdido todo}. Si bien puede iniciarse un juicio contra el escribano y obtenerse una indemnización de daños, esta es escasa, y es probable que la deuda ni siquiera fuera de María. El verdadero desafío es \textbf{prevenir el daño}: repararlo ya tiene sabor a poco. Lo que María quería era una casa; lo necesario es estar presente para ayudar a las personas a seguir adelante.

\section{Importancia del sistema y caducidades}

El derecho real es de orden público, y el sistema registral está comprometido con la defensa de ese orden público. Permite tomar conocimiento de todo lo inscripto: al comprar un inmueble, al embargarlo, o cuando se es acreedor quirografario que ganó el juicio y quiere conocer si el demandado condenado tiene inmuebles a su nombre. Además, brinda \textbf{seguridad al tráfico inmobiliario}. Existen ciertos plazos de caducidad vinculados con la inscripción:

\begin{table}[H]
\centering
\begin{tabularx}{\textwidth}{@{}>{\raggedright\arraybackslash}p{0.38\textwidth}>{\raggedright\arraybackslash}X@{}}
\toprule
\textbf{Inscripción} & \textbf{Caducidad} \\
\midrule
Inscripción provisional & A los 180 días, si no se subsana la observación. \\
Medidas cautelares & A los 5 años. \\
Inscripciones hipotecarias & A los 35 años (véase el capítulo~\ref{cap:hip-prenda-antic}). \\
\bottomrule
\end{tabularx}
\end{table}

\cornell{publicidad registral inmobiliaria}{
\cq{¿Por qué el derecho real exige publicidad y qué efecto tiene la inscripción?}{Porque es absoluto y oponible \textit{erga omnes}, y nadie puede ser obligado a respetar lo que no conoce; el derecho personal es relativo. La inscripción tiene efecto declarativo: el derecho se constituye, transmite y extingue fuera del registro, y la inscripción lo hace oponible a terceros. Sin ella, es oponible solo a quienes conocieron el acto (partes, escribano, abogado).}
\cq{¿Por qué el sistema es real, de inscripción y no convalidante?}{Es real porque el registro se organiza por inmueble (matrícula y folio real) y no por persona; de inscripción porque se asienta una síntesis del documento; no convalidante porque la inscripción no subsana el título, que seguirá siendo atacable si tiene defectos.}
\cq{¿Cuáles son los ocho principios registrales?}{Rogación, inscripción, presunción registral, especialidad, legalidad o calificación, prioridad, tracto sucesivo y publicidad (Carta de Buenos Aires).}
\cq{¿Qué examina el registro en virtud de la legalidad y qué puede resolver?}{Las formas extrínsecas del título, no su contenido intrínseco. Si cumplen, inscribe; si no, rechaza o inscribe provisionalmente por 180 días. Los arts.~8 y 9 prevén recursos administrativos y judiciales.}
\cq{¿Qué otorga la prioridad y qué permite el tracto abreviado?}{La fecha de inscripción, no la del título. El tracto abreviado (art.~16) permite inscribir simultáneamente dos transmisiones en un mismo asiento (por ejemplo, herederos que venden), sin romper el encadenamiento.}
\cq{¿Cómo funciona el bloqueo registral?}{El certificado solicitado por el escribano para un negocio congela la información del inmueble por 15 días (ciudad), 25 (provincia) o 30 (fuera de la provincia); los nuevos gravámenes no pueden ingresar. Otorgada la escritura dentro de ese plazo, el escribano tiene 45 días para inscribir con efecto retroactivo.}
\cq{¿Qué ocurre si no se inscribe dentro de los 45 días?}{Se pierde la retroactividad: la inscripción es oponible desde la fecha real y el embargo que ingresó en el ínterin afecta al nuevo propietario (caso de María: departamento de 50 mil dólares con embargo de 100 mil). Reparar el daño es insuficiente; hay que prevenirlo.}
\cq{¿Cuáles son los plazos de caducidad?}{Inscripción provisional, 180 días; medidas cautelares, 5 años; inscripciones hipotecarias, 35 años.}
}{Los derechos reales inmobiliarios son absolutos y oponibles \textit{erga omnes}, lo que exige publicidad. La Ley 17.801 establece un sistema registral real, de inscripción y no convalidante, que descansa en ocho principios y cuya inscripción tiene efecto declarativo. El bloqueo registral protege el inmueble durante la vigencia del certificado y, si se inscribe dentro de los 45 días, la inscripción es retroactiva; fuera de ese plazo se pierde la retroactividad. Al profesional le corresponde vigilar estos plazos para prevenir daños irreparables.}

\chapter{Protección de la vivienda}\label{cap:vivienda}
% ============================================================

El régimen de protección de la vivienda, regulado en los artículos 244 a 256 del CCyCN, es la respuesta del derecho argentino a una tensión estructural: todo el patrimonio del deudor, incluida su casa, es garantía común de sus acreedores (art.~242), pero el acceso a una vivienda digna es, a la vez, un derecho humano reconocido internacionalmente. Este capítulo sigue el caso de \textbf{Juan} a través de cinco ejes: por qué su casa responde por sus deudas; cómo puede ``ponerla a salvo'' mediante la afectación; a quiénes puede proteger; cómo la protección ``viaja'' al dinero si vende; y cómo y cuándo la protección puede cancelarse. La pregunta de fondo es: \textit{¿hasta dónde puede el derecho proteger el techo de una persona sin desproteger por completo a quienes le prestaron confianza y dinero?}

El régimen es de uso cotidiano para quien ejerce la abogacía, la escribanía, la actividad notarial, contable o inmobiliaria: permite asesorar sobre la afectación de inmuebles, redactar escrituras de constitución, evaluar riesgos crediticios (si el inmueble puede servir de garantía o no), representar a acreedores o defender a una familia frente a una ejecución. En lo personal, permite proteger la vivienda propia frente a una quiebra, un despido o una demanda imprevista.

\section{Fundamentos del régimen}

\subsection{El patrimonio como garantía común}

El régimen de protección de la vivienda es de profundo interés social y está vinculado con el derecho de todo propietario de afectar el inmueble en el que vive a un régimen especial, para protegerlo de los avatares de la vida, que en su devenir genera en cada persona infinidad de obligaciones. La \textbf{obligación} es inherente a la vida: trabajar, ejercer la profesión, realizar una actividad comercial, formar una familia, contratar empleados, o incluso dañar a otro aun sin intención de dañar, son hechos susceptibles de generar obligaciones de dar, de hacer o de no hacer. Así como el riesgo es inherente a la vida, es inherente a la obligación: la finitud humana, los despidos y las pandemias ponen en riesgo el cumplimiento, generalmente por causas ajenas a la propia voluntad.

Desde que el mundo occidental decidió prohibir la tortura, los azotes, la esclavitud por deudas y la prisión del deudor incumplidor (véase el capítulo~\ref{cap:garantias-fund}), se desarrolló la \textbf{teoría del patrimonio}: el patrimonio del deudor responde a todas las obligaciones. Sin firmar nada en especial, solo por contraer obligaciones, todo el patrimonio queda afectado a su cumplimiento y es susceptible de ser \textbf{embargado} y \textbf{ejecutado} (subastado), incluso el inmueble en el que se vive, porque este es parte de la garantía de todos los acreedores. Ese escenario ha dado lugar a un sistema de protección especial de la vivienda.

\subsection{Origen histórico y antecedentes normativos}

El origen del sistema se remonta a la década de 1930 en los Estados Unidos. El estado de \textbf{Texas} fue el primero en desarrollar un régimen de protección, porque la debacle económica y los efectos de la crisis del 30 hicieron necesario proteger los inmuebles. Luego se desarrollaron regímenes similares en otros estados y, en la medida en que el derecho a la vivienda se asienta como derecho humano, el régimen adquiere raigambre internacional.

Antes del CCyCN existió el régimen de protección de la vivienda familiar regulado por la \textbf{Ley 14.394}, de mediados del siglo pasado (1954-1955), vigente hasta poco antes del 1 de agosto de 2015, que protegía la vivienda únicamente de aquel que tuviera familia. Frente a ello se plantean interrogantes sobre su justicia: ¿aquellos que no tenían familia, no la quisieron o no pudieron tenerla, y aquellos que la perdieron, acaso no tenían derecho a ver protegida su vivienda? Por ello el nuevo régimen incorpora la normativa en materia de tratados internacionales de derechos humanos, tal como lo expresan los arts.~1 y 2 del CCyCN.

\begin{definicion}{Ubicación normativa}
El régimen de protección de la vivienda se regula en los arts.~244 a 256 del CCyCN, dentro del Capítulo 3, dedicado a los bienes, y protege la vivienda del titular \textbf{tenga o no tenga familia}.
\end{definicion}

\begin{importante}
A diferencia del régimen anterior de la Ley 14.394, que solo amparaba a la familia constituida por matrimonio e hijos, el nuevo régimen protege la vivienda de toda persona, tenga o no tenga familia.
\end{importante}

\section{Constitución de la afectación}

\subsection{Concepto y efectos}

La \textbf{afectación} implica que el propietario puede afectar un inmueble que destine a vivienda, ya sea todo el inmueble o una parte, hasta un tope determinado de su valor. A partir de la constitución, el inmueble queda protegido: es \textbf{inembargable e inejecutable} por las deudas que el propietario contraiga posteriormente.

\begin{importante}
La protección alcanza a todas las deudas de \textbf{causa posterior} a la constitución, no de fecha posterior. Es la fecha de la causa de la obligación, no la fecha en que la deuda se torna exigible, la que determina si el inmueble está o no protegido frente a ese acreedor. A partir de la constitución, la afectación es \textbf{inoponible a los acreedores de causa anterior}.
\end{importante}

\begin{ejemplo}{Juan, titular de dominio y locatario comercial}
Juan es titular de dominio de un inmueble en el que vive. Además, alquila otro inmueble para instalar un comercio; esa locación la realiza en diciembre del año 2018. En diciembre del año 2019, Juan toma conocimiento de la existencia de este régimen y afecta su inmueble (el que habita) al régimen de protección de la vivienda. En ese momento no tenía grandes dificultades económicas: el comercio funcionaba bien y pagaba regularmente los alquileres. Posteriormente, los efectos económicos producidos por una pandemia reducen drásticamente sus ingresos y Juan ya no puede pagar los alquileres.

\textbf{Primera pregunta:} ¿puede el locador, acreedor, embargar y ejecutar el inmueble, siendo que la causa de la contratación es el contrato celebrado en 2018, aun cuando la deuda sea posterior a marzo de 2020? El locador puede ejecutar el inmueble: el sistema de protección le es \textbf{inoponible} a este acreedor, porque la causa de la obligación (el contrato de locación) radica en un tiempo anterior a la constitución de la afectación. La deuda es de causa anterior, aunque se haga exigible después.
% TODO: contenido incompleto o ambiguo en las notas originales (el apunte califica la respuesta como ``negativa para el locador'', en el sentido de que la protección no le es oponible)

Juan, con gran esfuerzo, comienza a restablecer su actividad y a regularizar sus obligaciones, entre ellas el alquiler. Pero hay otras deudas que todavía no ha podido cumplir, organizar o pagar. \textbf{Segunda pregunta:} ¿pueden los otros acreedores, de causa posterior, ejecutar este inmueble? No pueden, porque existe la protección de la vivienda, y ese régimen protege al inmueble de todas las obligaciones de causa posterior.
\end{ejemplo}

\begin{importante}
Incluso si un acreedor de causa anterior subasta el inmueble, el \textbf{remanente} de la subasta, cuando existe una afectación posterior, no ingresa libre al patrimonio del deudor, sino que conserva ese ``halo de protección'' y queda protegido frente a las deudas de causa posterior.
\end{importante}

La protección no opera frente a: (i) expensas, cargas y contribuciones; (ii) obligaciones ya garantizadas con derechos reales de garantía; (iii) obligaciones que tienen origen en construcciones o mejoras realizadas en la vivienda. Para todas las demás obligaciones de causa posterior, el inmueble está protegido a partir de la constitución.

\subsection{Legitimados y formas de constitución}

Están legitimados para afectar el inmueble el titular de dominio exclusivo, los condóminos en su conjunto y el titular del derecho de propiedad horizontal o de conjuntos inmobiliarios. Dentro de las amplias facultades del titular dominial se incluye esta afectación: a partir de ella el inmueble queda protegido, pero al mismo tiempo \textbf{sale, de alguna manera, de la vida negocial}. Mientras dure la afectación será inembargable e inejecutable, lo que tiene consecuencias prácticas en la vida negocial del propietario, en particular respecto de las fianzas y garantías (véase la sección sobre la tensión con los acreedores).

La afectación se constituye:
\begin{itemize}
    \item \textbf{por voluntad del propietario}, mediante una instrumentación con publicidad registral; si se realiza notarialmente, debe hacerse por \textbf{escritura pública} o por \textbf{testamento};
    \item \textbf{por vía judicial}, solo en el ámbito del derecho de familia, cuando esté comprometido el interés de menores o incapaces, o cuando existan personas con capacidad restringida.
\end{itemize}

\begin{nota}
En muchas ocasiones el propietario expresa que, mientras él esté, estará para trabajar y proteger a los suyos. Por ello, muchos clientes deciden la afectación como una voluntad \textit{post mortem}, mediante testamento.
\end{nota}

\section{Beneficiarios y habitación efectiva}

La protección del inmueble destinado a vivienda está vinculada con el derecho a la vivienda. A diferencia del régimen anterior, que solo protegía a la familia constituida (el matrimonio y los hijos), hoy, tras el paso hacia el derecho de las familias (el derecho a elegir entre distintos tipos de familias y a protegerse a uno mismo aunque no se tenga familia), es posible afectar el inmueble en el propio beneficio.

\begin{regla}{Art.~246 CCyCN}
La afectación puede realizarse en beneficio del propietario constituyente, de su cónyuge, de su conviviente, de sus descendientes o de sus ascendientes, aun cuando no convivan con él. También pueden ser beneficiarios los colaterales del propietario hasta el tercer grado.
\end{regla}

La inclusión del \textbf{conviviente} como posible beneficiario era impensada bajo la Ley 14.394, porque el conviviente no formaba parte de las familias reguladas. Pueden ser beneficiarios los descendientes y los ascendientes del constituyente, \textbf{aun cuando no convivan} con él, y los \textbf{colaterales hasta el tercer grado}, pero en su caso es imprescindible que vivan en el inmueble.

\begin{ejemplo}{Juan y la elección de beneficiarios}
Juan puede designarse a sí mismo como beneficiario, o designar como beneficiarios a su padre y a su nieto, por considerar que son quienes deben estar protegidos, porque tal vez su hijo ya no esté. La elección es discrecional: quien decide afectar o no el inmueble, y en beneficio de quiénes, es el propietario.
\end{ejemplo}

\begin{regla}{Art.~247 CCyCN}
Si la afectación es peticionada por el titular registral, se requiere que al menos uno de los beneficiarios habite el inmueble. En todos los casos, para que los efectos subsistan, basta que uno de ellos permanezca habitando allí.
\end{regla}

\begin{importante}
La convivencia se exige únicamente para el constituyente y para los colaterales; respecto de ascendientes o descendientes no es requisito. Mientras al menos uno de los beneficiarios continúe viviendo en el inmueble, la protección registral subsiste, aun cuando los demás beneficiarios no habiten allí.
\end{importante}

\section{Subrogación real, transmisión y asentimiento}

\begin{definicion}{Subrogación real}
Implica que una cosa, por decisión de la ley, ocupa el lugar jurídico que corresponde a otra. En caso de que el inmueble se venda, el precio ocupa su lugar y el precio conserva la protección.
\end{definicion}

Si Juan ha afectado el inmueble y lo enajena, necesita desafectarlo; en el momento en que lo desafecta, el embargo podría ingresar. Si lo desafecta y simultáneamente lo enajena, el precio percibido ingresaría a su patrimonio y, de no existir la subrogación real, sería inmediata e instantáneamente embargable y ejecutable.

\begin{importante}
La afectación de la vivienda que se enajena se transmite a la vivienda que se adquiere, o a los importes que la sustituyen, ya sea en concepto de indemnización o de precio.
\end{importante}

\begin{ejemplo}{El seguro de incendio}
Si el inmueble se incendia y se cobra un seguro de incendio, ese dinero ingresaría al patrimonio y sería garantía común de los acreedores de no existir el principio de subrogación real, que resulta sumamente importante para el desarrollo de esta protección a lo largo de la vida del propietario que la ha afectado.
\end{ejemplo}

La transmisión del inmueble afectado puede realizarse por acto \textbf{oneroso} o \textbf{gratuito}. Cuando ha sido afectado en beneficio del cónyuge o del conviviente (que viva en una unión convivencial registrada) se requiere el \textbf{asentimiento} de este. La distinción terminológica es central: se requiere \textbf{asentimiento, no consentimiento}, porque el inmueble es de Juan y, si es de Juan, él tiene la disposición material y jurídica del inmueble. En protección de otros sistemas de derecho, como el de la familia, existen otras normativas que protegen a la familia, aunque este no es el régimen de protección de la vivienda familiar.

\begin{ejemplo}{Juan, propietario exclusivo, y el asentimiento del cónyuge}
Si Juan pretende vender el inmueble del que es propietario exclusivo, perpetuo y en un cien por cien propio, y designó como beneficiarios a su cónyuge o conviviente, requerirá el asentimiento de este. Si el asentimiento no se otorga, lo que la ley protege es la protección de la vivienda, no el abuso del derecho de quien, por otras cuestiones, pretenda impedir la enajenación. En ese caso será necesaria la petición judicial, y el juez evaluará la razonabilidad o no de ese asentimiento remiso.
\end{ejemplo}

\subsection{Gravámenes, frutos y beneficios}

El inmueble puede gravarse con \textbf{derecho real de hipoteca}: se desafecta el régimen al solo efecto de realizar el gravamen hipotecario, al que, por supuesto, le es inoponible la protección. También para este gravamen se requiere el asentimiento. Los frutos civiles del inmueble afectado son embargables, salvo que sean indispensables para las necesidades de los beneficiarios.

\begin{table}[H]
\centering
\begin{tabularx}{\textwidth}{@{}>{\raggedright\arraybackslash}p{0.34\textwidth}>{\raggedright\arraybackslash}X@{}}
\toprule
\textbf{Aspecto} & \textbf{Beneficio o efecto} \\
\midrule
Impuesto a la transmisión gratuita de bienes & Exención luego del fallecimiento del causante \\
Honorarios profesionales & Reducidos, tanto para la constitución como para la sucesión \\
Concursos y quiebras & Se manejan siempre sobre valuaciones fiscales cuando el inmueble está afectado al régimen \\
Frutos civiles & Embargables, salvo que sean indispensables para las necesidades de los beneficiarios \\
\bottomrule
\end{tabularx}
\end{table}

\section{Desafectación y tensión con los acreedores}

\subsection{Causas de desafectación y cancelación de la inscripción}

La desafectación y la cancelación de la inscripción están reguladas en el art.~255. Por un lado se desafecta y, por otro, se cancela la inscripción en el registro, porque tanto la afectación como la cancelación deben inscribirse para su oponibilidad a terceros.

\begin{importante}
En estos casos la inscripción es \textbf{constitutiva}: solo cuando el tercero toma conocimiento de que el inmueble está afectado al régimen de protección, esa afectación le es oponible.
\end{importante}
% TODO: verificar consistencia entre fuentes originales: en este material se califica la inscripción de la afectación como constitutiva, mientras que el material sobre publicidad registral presenta como regla general que la inscripción tiene efecto declarativo.

\begin{regla}{Art.~255 CCyCN}
La desafectación procede únicamente en los siguientes casos:
\begin{enumerate}
    \item A solicitud del constituyente, con el asentimiento del cónyuge o conviviente de una unión convivencial inscripta, en el caso de que el propietario haya indicado como beneficiarios a estos.
    \item A solicitud de la mayoría de los herederos, cuando el régimen subsiste tras el fallecimiento del propietario porque este lo afectó en protección de sus hijos; mientras estos habiten el inmueble el régimen subsiste, salvo que exista una solicitud de la mayoría de los condóminos, computada según sus partes indivisas.
    \item A pedido de cualquiera de los interesados o de oficio, si fallece el constituyente y todos los beneficiarios ya no viven en el inmueble (en el caso de ascendientes, descendientes o colaterales, donde el requisito de habitar es diferente).
    \item En caso de expropiación, reivindicación o ejecución autorizada por la propia normativa (por ejemplo, el caso de las expensas), supuestos en los que el régimen de afectación jamás es oponible.
\end{enumerate}
\end{regla}
% TODO: contenido incompleto o ambiguo en las notas originales (inciso 2: la relación entre «mayoría de herederos» y «mayoría de condóminos» no está desarrollada; inciso 3: la referencia al requisito de habitar remite a lo explicado sobre habitación efectiva)

\begin{nota}
Si entre los beneficiarios existe un menor o un hijo con capacidad restringida, se trata de un caso concreto que deberá resolverse judicialmente.
\end{nota}

\begin{regla}{Art.~256 CCyCN}
Puede afectarse un inmueble rural cuando no exceda la unidad económica, según lo que sea previsible o de acuerdo con la normativa local, porque el destino del régimen es siempre habitacional.
\end{regla}

\subsection{Tensión con el derecho de los acreedores}

La protección retira el inmueble del conjunto de bienes afectados a la garantía de todos los acreedores quirografarios, lo que disminuye esa garantía. Surge así la pregunta: ¿no tienen estos acreedores derecho a cobrar? La respuesta es que \textbf{sí lo tienen}, pero el ordenamiento ha tenido que decidir: no ha sido posible proteger a todos. Se trata de un desafío del derecho, que impone decidir; y esta decisión, adoptada en todo el derecho occidental, se refleja en figuras semejantes en los distintos sistemas de derecho comparado.

La afectación tiene \textbf{publicidad registral} (capítulo~\ref{cap:registral}), de modo que el acreedor diligente no será sorprendido en su buena fe: antes de otorgar un crédito obtendrá el informe registral, que dará cuenta de que el inmueble está afectado y, por lo tanto, es inembargable e inejecutable.

\begin{ejemplo}{La vivienda afectada como garantía}
Cuando Juan quiera dar una garantía o una fianza y demuestre que tiene la propiedad de un inmueble, si ese inmueble es inembargable e inejecutable no será apto para afianzar ni una obligación propia ni una obligación de terceros. Algunos acreedores de fecha posterior pueden verse afectados al intentar el embargo y la ejecución y constatar que el inmueble está afectado al régimen.
\end{ejemplo}

\begin{importante}
El régimen de protección de la vivienda está en tensión con el derecho de los acreedores de causa posterior a la constitución. En esta disyuntiva se ha optado por la protección de la vivienda.
\end{importante}

La vivienda propia, el hogar del propio acreedor, también está protegida por este régimen frente a sus otros acreedores, porque la vivienda es el hábitat imprescindible para el desarrollo de la vida. Por eso la regulación no ha mirado solamente al propietario, sino al interés de toda la sociedad. Esta protección no es absoluta ni eterna: cede frente a deudas de causa anterior, frente a expensas y obligaciones garantizadas con derechos reales, y puede desafectarse por voluntad del constituyente, de sus herederos, o por disposición legal en casos de expropiación o ejecución autorizada.

\cornell{protección de la vivienda}{
\cq{¿Qué es la garantía común de los acreedores y cómo se originó la protección de la vivienda?}{Todo el patrimonio del deudor, incluida su vivienda, responde por sus obligaciones sin necesidad de pactarlo (art.~242). La protección nació en Texas en la década de 1930, como respuesta a la crisis económica, y luego adquirió raigambre internacional como derecho humano.}
\cq{¿Qué cambió respecto de la Ley 14.394?}{La Ley 14.394 (1954-1955) solo protegía la vivienda de quienes tenían familia constituida. El CCyCN (arts.~244 a 256) protege la vivienda de toda persona, tenga o no tenga familia.}
\cq{¿Qué efecto produce la afectación y frente a qué deudas no protege?}{El inmueble queda inembargable e inejecutable frente a deudas de causa posterior a la constitución. No protege frente a deudas de causa anterior (aunque se vuelvan exigibles después), expensas, cargas y contribuciones, obligaciones garantizadas con derechos reales de garantía ni obligaciones originadas en construcciones o mejoras. El remanente de una subasta conserva la protección frente a deudas de causa posterior.}
\cq{Caso de Juan: ¿puede el locador ejecutar su vivienda?}{Sí: la deuda de alquiler tiene causa en un contrato (2018) anterior a la afectación (2019); el régimen le es inoponible. Frente a deudas de causa posterior, el inmueble está protegido.}
\cq{¿Quién puede afectar el inmueble y de qué manera?}{El titular de dominio exclusivo, los condóminos en conjunto o el titular de propiedad horizontal o conjuntos inmobiliarios, por escritura pública o testamento; por vía judicial, solo en el ámbito del derecho de familia.}
\cq{¿Quiénes pueden ser beneficiarios y es obligatorio habitar el inmueble (arts.~246 y 247)?}{El constituyente, su cónyuge o conviviente, descendientes y ascendientes (aunque no convivan) y colaterales hasta el tercer grado (estos, solo si conviven). Para que subsista la protección basta que al menos un beneficiario siga habitando el inmueble.}
\cq{¿Qué es la subrogación real y qué se requiere para vender o hipotecar?}{El precio, la indemnización o el seguro ocupan el lugar jurídico del inmueble y conservan la protección. La transmisión, onerosa o gratuita, y la hipoteca requieren el asentimiento (no consentimiento) del cónyuge o conviviente beneficiario; si lo niega, el juez evalúa su razonabilidad.}
\cq{¿En qué casos procede la desafectación (art.~255) y por qué se inscribe?}{A solicitud del constituyente (con asentimiento si corresponde), de la mayoría de los herederos, de interesados o de oficio si fallecido el constituyente los beneficiarios ya no habitan el inmueble, y en caso de expropiación, reivindicación o ejecución autorizada. Debe inscribirse para ser oponible a terceros.}
\cq{¿Qué tensión genera el régimen y cómo se resolvió?}{Tensiona la protección del deudor con el derecho de cobro de los acreedores de causa posterior; se optó por priorizar la vivienda, con publicidad registral que resguarda la buena fe del acreedor diligente.}
}{El régimen de los arts.~244 a 256 permite al propietario sustraer su vivienda de la ejecución por deudas de causa posterior a la afectación, en beneficio propio o de su familia, mediante escritura pública, testamento o, excepcionalmente, vía judicial. La protección cede frente a deudas de causa anterior, expensas y obligaciones garantizadas con derechos reales, se traslada al precio por subrogación real, exige asentimiento del cónyuge o conviviente beneficiario para transmitir o gravar y se cancela solo por las causas del art.~255. Ante la tensión entre vivienda y acreedores, el ordenamiento optó por proteger la vivienda.}


% ============================================================
% APÉNDICE
% ============================================================
\appendix
\chapter{Ejercicios de evaluación sobre la teoría general}\label{ap:ejercicios}

Las siguientes consignas abarcan la Parte~I y los principios de la Parte~II, organizadas por niveles: básico, intermedio y avanzado.

\section{Nivel básico}

\begin{ejercicio}{1}
\textbf{Consigna.} Defina la obligación según el art.~724 e identifique sus tres elementos.

\textbf{Respuesta.} Relación jurídica en virtud de la cual el acreedor tiene el derecho de exigir del deudor una prestación destinada a satisfacer un interés lícito y, ante el incumplimiento, obtener forzadamente su satisfacción. Elementos: sujeto activo o acreedor, sujeto pasivo o deudor y prestación (dar, hacer o no hacer).
\end{ejercicio}

\begin{ejercicio}{2}
\textbf{Consigna.} ¿Cuál es la diferencia esencial entre un derecho personal y uno real? Explique con un ejemplo.

\textbf{Respuesta.} La naturaleza del vínculo: en el personal, entre dos sujetos y relativo; en el real, directo e inmediato entre titular y cosa, y absoluto (\textit{erga omnes}). Si se presta dinero a Juan (derecho personal), solo Juan debe; si se es dueño de un inmueble (derecho real), toda la comunidad debe respetar esa propiedad.
\end{ejercicio}

\begin{ejercicio}{3}
\textbf{Consigna.} ¿Por qué los derechos personales son «relativos» y los reales «absolutos»?

\textbf{Respuesta.} Los personales son relativos porque el vínculo existe solo entre las partes y el acreedor solo puede exigir a su deudor. Los reales son absolutos porque son oponibles \textit{erga omnes}: toda la comunidad debe respetar el derecho del titular.
\end{ejercicio}

\begin{ejercicio}{4}
\textbf{Consigna.} Defina el derecho real según el art.~1882 e identifique sus elementos.

\textbf{Respuesta.} Poder jurídico de estructura legal que se ejerce directamente sobre su objeto en forma autónoma y que atribuye a su titular las facultades de persecución y preferencia y las demás previstas en el Código. Elementos: el titular y la cosa.
\end{ejercicio}

\begin{ejercicio}{5}
\textbf{Consigna.} ¿Qué integra el patrimonio de una persona según el CCyC?

\textbf{Respuesta.} Los bienes, que incluyen tanto los derechos personales como los reales.
\end{ejercicio}

\begin{ejercicio}{6}
\textbf{Consigna.} ¿Qué establece el art.~16 y por qué es fundamental para los derechos reales?

\textbf{Respuesta.} Que los derechos pueden recaer sobre bienes susceptibles de valor económico y que, cuando son materiales, se llaman cosas (también se aplica a la energía y fuerzas naturales). Define el objeto de los derechos reales: la cosa.
\end{ejercicio}

\section{Nivel intermedio}

\begin{ejercicio}{7}
\textbf{Consigna.} ¿Qué significa el \textit{numerus clausus} y cuál es su consecuencia jurídica?

\textbf{Respuesta.} Solo la ley puede crear derechos reales; los particulares no pueden crear nuevos ni modificar el contenido de los existentes. La consecuencia es la nulidad (art.~1884).
\end{ejercicio}

\begin{ejercicio}{8}
\textbf{Consigna.} Explique la convalidación (art.~1885) y su relación con el \textit{nemo plus iuris}.

\textbf{Respuesta.} El \textit{nemo plus iuris} impide transmitir un derecho mejor o más extenso que el que se tiene, de modo que quien transmite un derecho que no tiene realiza un acto nulo. La convalidación establece que si el transmitente adquiere el derecho con posterioridad, el acto queda convalidado.
\end{ejercicio}

\begin{ejercicio}{9}
\textbf{Consigna.} ¿Qué son el \textit{ius persequendi} y el \textit{ius preferendi} y por qué son características del derecho real?

\textbf{Respuesta.} Persecución: perseguir la cosa en poder de quien se encuentre. Preferencia: ser preferido frente a titulares de derechos reales o personales de fecha posterior. Derivan de la relación directa entre titular y cosa y de su carácter absoluto.
\end{ejercicio}

\begin{ejercicio}{10}
\textbf{Consigna.} Enumere los catorce derechos reales del art.~1887.

\textbf{Respuesta.} Dominio, condominio, propiedad horizontal, conjuntos inmobiliarios, tiempo compartido, cementerio privado, superficie, usufructo, uso, habitación, servidumbre, hipoteca, anticresis y prenda.
\end{ejercicio}

\begin{ejercicio}{11}
\textbf{Consigna.} Diferencie derechos reales sobre cosa propia y sobre cosa ajena, con un ejemplo de cada uno.

\textbf{Respuesta.} Sobre cosa propia, el titular ejerce el derecho sobre su propia cosa (dominio: ser dueño de una casa). Sobre cosa ajena, sobre la cosa de otro, como carga o gravamen para el dueño (hipoteca: el banco tiene un derecho real sobre el inmueble en garantía de un crédito).
\end{ejercicio}

\begin{ejercicio}{12}
\textbf{Consigna.} Diferencie los derechos reales principales de los accesorios, con ejemplos.

\textbf{Respuesta.} Los principales son independientes (dominio, usufructo, servidumbre). Los accesorios son de garantía y existen en función de un crédito (hipoteca, prenda, anticresis); si se extingue el crédito principal, se extingue el derecho de garantía.
\end{ejercicio}

\section{Nivel avanzado}

\begin{ejercicio}{13}
\textbf{Consigna.} ¿Por qué la normativa de derechos reales es de orden público y qué consecuencias tiene?

\textbf{Respuesta.} Porque los derechos reales trascienden el interés particular: su ejercicio afecta a toda la comunidad. Consecuencias: los particulares no pueden crear derechos reales ni modificar su contenido; la violación acarrea nulidad; no se aplica la autonomía de la voluntad como en los contratos.
\end{ejercicio}

\begin{ejercicio}{14}
\textbf{Consigna.} ¿Qué establece el art.~17 y por qué las partes del cuerpo no son objeto de derechos reales?

\textbf{Respuesta.} Que los derechos sobre el cuerpo humano o sus partes no tienen valor comercial sino afectivo, científico, terapéutico, humanitario o social, y se rigen por leyes especiales. No son objeto de derechos reales porque el art.~16 exige valor económico para que un bien sea cosa.
\end{ejercicio}

\begin{ejercicio}{15}
\textbf{Consigna.} Pedro vende un inmueble a María en enero, pero no era propietario. En marzo adquiere la propiedad. ¿Qué ocurre con la venta de enero?

\textbf{Respuesta.} Por el \textit{nemo plus iuris}, Pedro no podía transmitir un derecho que no tenía: la venta nació nula en cuanto a la transmisión del derecho real. Pero por el principio de convalidación (art.~1885), al adquirir Pedro la propiedad en marzo, la transmisión queda convalidada y María adquiere el dominio retroactivamente.
\end{ejercicio}

\backmatter

\chapter*{Referencias mencionadas en las fuentes}
\addcontentsline{toc}{chapter}{Referencias mencionadas en las fuentes}
\markboth{Referencias}{Referencias}

Los apuntes originales no incluyen una bibliografía formal. Se conservan, sin completarlas con información externa, las referencias doctrinales, normativas y jurisprudenciales mencionadas en el texto.

\begin{itemize}
    \item Código Civil y Comercial de la Nación (Ley~26.994), en vigencia desde agosto de 2015; Código Civil de Vélez Sársfield.
    \item Constitución Nacional, arts.~14, 17, 18, 19 y 75, inc.~22; Pacto de San José de Costa Rica, art.~21.
    \item Ley~13.512 (propiedad horizontal, 1948); Ley~14.394; Ley~14.159 (art.~24); Ley~17.711; Ley~17.801 (registros de la propiedad inmueble); Ley~20.798 (art.~3574~bis del Código de Vélez); Ley~27.261 y Ley~27.271 (reformas de los arts.~2189 y 2210 del CCyC; las fuentes difieren sobre cuál de ellas sustituyó el art.~2189); Decreto-Ley~15.348/46 (prenda con registro).
    \item Ordenanza General~221 (1978) y Ley~9.000/94 de la Provincia (cementerios privados).
    \item Mariani de Vidal, Marina: estructura didáctica de la clasificación de los muros; concepción de las expensas como «savia» del consorcio; responsabilidad subsidiaria e ilimitada de los propietarios; caracterización de los countries bajo propiedad horizontal como «grandes castillos construidos sobre arena».
    \item Cámara Nacional de Apelaciones en lo Civil, plenario «Servicios Eficientes S.A. c/ Sabrá, Victor» (1999).
    \item Alterini, Jorge: referido como director de la revista La Ley, que solicitó un estudio sobre la doctrina del plenario diez años después (2009).
    \item Quinto Congreso de Derecho Privado: ponencia sobre conjuntos inmobiliarios.
    \item Carta de Buenos Aires, aprobada en un Congreso Internacional de Derecho Registral (principios registrales).
    \item Ashley (autor mencionado sobre el tiempo) y San Agustín (tiempo existencial); criterio práctico atribuido al Dr.~Paños sobre el ámbito de la acción confesoria.
\end{itemize}
% TODO: verificar referencia presente en las fuentes originales (las fuentes no ofrecen datos bibliográficos completos de estas referencias; no se completaron).

\end{document}
